\providecommand{\toplevelprefix}{../..}  %
\documentclass[../../book-main.tex]{subfiles}

\begin{document}

\chapter{Learning Representations for Real-World Data and Tasks}
\label{ch:applications}

\begin{quote}
    \hfill    ``{\em The best theory is inspired by practice, and the
    best practice is inspired by theory}.''

    $~$ \hfill --- Donald Knuth
\end{quote}
\vspace{5mm}

\section{Introduction}
The previous chapters have presented a systematic introduction to
mathematical problems, computational frameworks, and practical
algorithms associated with learning low-dimensional distributions
from high-dimensional data. Although most theoretical justifications
of these methods have been established for idealistic models of data
distributions such as (mixtures of) subspaces and/or Gaussians, the
principles and ideas behind these computational methods are
nevertheless powerful and general, and they are in fact meant to be
applicable to real-world datasets and tasks.

To help readers understand the material in this book better and learn
how to apply what you have learned so far to real-world data, here we
provide some demonstrations and vignettes of
several representative applications. Each application proposes a
solution to a real-world task with a real-world dataset (such as
visual data, motion data, and text data), using the methods we have
introduced in this book. The results presented in this chapter are
meant to serve the following two purposes:
\begin{itemize}
    \item firstly, to provide additional experimental details and
        empirical evidence which validate the methods presented
        earlier in the book, and demonstrate their significant
        potential in real-world contexts;
    \item secondly, to introduce the reader to certain modern
        empirical methods and tasks in deep learning which are not
        well-documented outside of research or production codebases.
\end{itemize}
However, in our honest opinion, the solutions and results featured
here are designed simply to verify that the methodology works. As
such, there is great room for future improvement, in both theoretical
understanding and practical
engineering, to further advance the
state-of-the-art performance. We will showcase some immediate and
obvious improvements through this chapter and especially in the last
Section \ref{sec:scalable}. We will leave the discussions of
potentially much more transformative ideas about representation
learning and remaining significant open problems about intelligence
in general for
future research in the final  \Cref{ch:future}.

\subsection{Outline of the Chapter}
\label{sec:application_outline}
For intelligence in nature, it is arguably true that the three most
important types of data with which our brain builds our memory and
knowledge about are {\em the visual data}, {\em our body motions},
and the {\em natural languages}. Hence, in this chapter, we will show
how to apply principles and methods introduced in this book to learn
good representations of these real-world data and show how such
representations facilitate important practical tasks associated with these data.

Firstly, from \Cref{sec:contrastive_learning} to
\ref{sec:CUPID}, we will start with the visual data and show how to
learn the distributions of 2D images to 3D objects since the visual
memory serves as the low-level but foundational model for us (and
animals) that stores knowledge of the physical world and helps
perceive and predict in a new environment. We will then, in Section
\ref{sec:egoallo}, show how to learn the distribution of body motions
as it is crucial for us (and animals) to act or interact swiftly and
accurately within the environments. Finally, in Section
\ref{sec:clm_text}, we will show how to model and learn the
distribution of natural languages.

Evidently, compared to the natural languages, our coverage of
physical data, visual and motion, will be much more extensive.
Although in the past few years, the AI industry has invested heavily
in modeling natural languages,\footnote{Including their derivatives
    such as programming languages, logical deductions, and mathematical
proofs.} people have started to realize the importance of modeling
knowledge about the physical world and our actions within. In most
recent years, there have been a fastly growing interest and shifting
attention to develop the so-called ``world model,'' ``spatial
intelligence,'' ``physical intelligence,'' or ``embodied
intelligence'', as they are crucial for AI systems or agents, say an
intelligent robot, which need to interact with an open physical
world. As we have learned in Chapter \ref{ch:intro}, that is
precisely the goal and scope of the {\em  Cybernetics} program
initiated by Norbert Wiener in the 1940s.

On the technical level, from  \Cref{sec:contrastive_learning} to
\Cref{sec:image_classification}, we will show how to use principles
introduced in this book to learn good representations of imagery data
in {\em unsupervised, weakly supervised}, and {\em supervised}
settings, respectively.\footnote{The goal in all three settings is to
    learn a good representation for the data distribution. They differ
    only in how much side information is given about the distribution
that may facilitate the learning.} This will serve as both an
introduction to imagery data processing, data augmentation
techniques, and a popular deep architecture Transformer, as a
precursor to the implementation of the CRATE architecture later.
These examples also give the
first demonstration of the drastic kinds of simplifications one can
already make using the principles introduced in the book. We will continue
with modifications to the \textit{transformer network architecture}
in \Cref{sec:image_classification} for images and \Cref{sec:clm_text}
or texts respectively, which demonstrate the capabilities of
simplified white-box architectures, including CRATE and its variants,
for \textit{encoding} within the
image and text domains. We will also demonstrate the capabilities of
the simplified architectures for \textit{autoencoding} in
\Cref{sec:image_completion}. %

Practically speaking, the primary reason why we want to learn a good
representation of a data distribution is to allow us to sample the
distribution (as a prior) to regenerate, estimate or predict the
state of the world, conditioned on current observation or information
given. We demonstrate the generative capabilities of the learned representations
\begin{itemize}
    \item for 2D images in \Cref{sec:auto-encode-img-gen} and
        \Cref{sec:cond_gen};
    \item for 3D objects in
        \Cref{sec:Michelangelo} and \Cref{sec:shape-and-pose};
    \item for human motions in \Cref{sec:egoallo};
    \item and for natural languages in \Cref{sec:clm_text}.
\end{itemize}

\subsection{Overall Technical Setup}\label{sec:experiment_setup}

In previous chapters, we alluded to different setups in which we used
representation-learning techniques to process real data at scale. In
this chapter, we will describe such setups in great detail. The
objective of this section is to get you, the reader, to be able to
reproduce any experiment discussed in this section (or indeed the
book) using just the description we will give in the book, the
principles introduced in previous chapters and expanded on in this
chapter, and hyperparameters taken from a smattering of papers whose
results are discussed in this chapter. To this end, we will
\textit{precisely} describe all procedures in a detailed language,
pseudocode, or mathematical notation that can be directly implemented
in code. Wherever possible, we will discuss how the concrete
implementations connect to the principles presented earlier in the book.

\begin{figure}[t!]
    \includegraphics[width=\textwidth]{\toplevelprefix/chapters/chapter8/figs/encoder_pipeline.pdf}
    \caption{\small\textbf{A diagram of the encoder pipeline.} Data
        \(\vX \in \cD\) is fed through the embedding
        \(f_{\theta}^{\emb}\) to get a sequence in \((\R^{d})^{*}\). The
        embedding is fed through a backbone \(f_{\theta}^{\mathrm{bb}}\)
        to get features \(\vZ_{\theta}(\vX)\) for each token. We can
        extract an aggregate feature \(\vz_{\theta}(\vX)\) using the
        extraction map \(f_{\theta}^{\mathrm{ext}}\). Finally, to use the
        aggregate feature in downstream tasks, we can use the
    task-specific head \(h_{\theta}\).}
    \label{fig:overall_encoder_pipeline}
\end{figure}

\begin{figure}[t!]
    \includegraphics[width=\textwidth]{\toplevelprefix/chapters/chapter8/figs/autoencoder_pipeline.pdf}
    \caption{\small\textbf{A diagram of the autoencoder pipeline.}
        Data \(\vX \in \cD\) is fed through the embedding
        \(f_{\theta}^{\emb}\) to get a sequence in \((\R^{d})^{*}\). The
        embedding is fed through an encoder backbone
        \(f_{\theta}^{\mathrm{bb}}\) to get features
        \(\vZ_{\theta}(\vX)\) for each token. To decode
        \(\vZ_{\theta}(\vX)\), we pass it through a decoder backbone
        \(g_{\eta}^{\mathrm{bb}}\). To map the decoder backbone output
        back to data space \(\cD\), we use an \textit{unembedding} layer
        \(g_{\eta}^{\mathrm{unemb}}\), overall obtaining a reconstruction
        \(\hat{\vX}_{\theta, \eta}(\vX)\) (here stylized to be a
    pixelated reconstruction of the input).}
    \label{fig:overall_autoencoder_pipeline}
\end{figure}

%\paragraph{Technical Setup.}
Let us define the set of possible data as \(\cD\) (eventually this
    will be the set of images \(\cI\), for example, or the set of text
\(\cT\)), and the set of finite sequences of tokens in \(\R^{d}\)
(i.e., the set of matrices with \(d\) rows) as \((\R^{d})^{*} \doteq
\bigcup_{T = 1}^{\infty}\R^{d \times T}\). In order to discuss the
wealth of applications we introduce in this chapter, we first recall
that in the rest of the book, we discuss two different types of model
architectures.
\begin{itemize}
    \item An \textit{encoder} architecture, parameterized by
        \(\theta\), which is composed of several components:
        \begin{itemize}
            \item An \textit{embedding} \(f_{\theta}^{\emb} \colon
                \cD \to (\R^{d})^{*}\), which converts the input data
                \(\cD\) into a series of \textit{tokens} which are
                mapped into, or \textit{embedded} in,
                \(d\)-dimensional space. \textit{In the rest of the
                    chapter, we will often identify tokens and embeddings
                with each other.}
            \item An \textit{encoder backbone}
                \(f_{\theta}^{\backbone} \colon (\R^{d})^{*} \to
                (\R^{d})^{*}\), which processes the series of
                embeddings using a sequence-to-sequence operation.
                This backbone is implemented by the network
                architectures discussed in the previous chapters, but
                we will give a more formal description as we go along.
            \item An \textit{aggregate feature extractor}
                \(f_{\theta}^{\ext} \colon (\R^{d})^{*} \to \R^{d}\),
                which extracts an aggregate representation of the
                whole sequence. This is used to define a single
                feature for the entire data sample.
            \item A \textit{task-specific head} \(h_{\theta} \colon
                \R^{d} \to \R^{m}\), which extracts an
                \(m\)-dimensional output for prediction.
        \end{itemize}
        We also define \(f_{\theta} \doteq f_{\theta}^{\backbone}
        \circ f_{\theta}^{\emb} \colon \cD \to (\R^{d})^{*}\). Given
        an input \(\vX\), we write \(\vZ_{\theta}(\vX) \doteq
        f_{\theta}(\vX)\) and \(\vz_{\theta}(\vX) \doteq
        f_{\theta}^{\ext}(\vZ_{\theta}(\vX))\). The overall pipeline
        is depicted in \Cref{fig:overall_encoder_pipeline}.
    \item An \textit{autoencoder} architecture, which is composed of
        several components:
        \begin{itemize}
            \item An \textit{embedding} \(f_{\theta}^{\emb} \colon
                \cD \to (\R^{d})^{*}\), which converts the input data
                \(\cD\) into a series of \textit{tokens} which are
                embedded in \(d\)-dimensional space.
            \item An \textit{encoder backbone}
                \(f_{\theta}^{\backbone} \colon (\R^{d})^{*} \to
                (\R^{d})^{*}\), which processes the series of
                embeddings using a sequence-to-sequence operation.
            \item A \textit{decoder backbone} \(g_{\eta}^{\backbone}
                \colon (\R^{d})^{*} \to (\R^{d})^{*}\), which
                conceptually undoes the operation of the encoder backbone.
            \item An \textit{unembedding} \(g_{\eta}^{\unemb} \colon
                (\R^{d})^{*} \to \cD\), which acts as an inverse of
                the embedding.
        \end{itemize}
        We also define \(f_{\theta} \doteq f_{\theta}^{\backbone}
        \circ f_{\theta}^{\emb} \colon \cD \to (\R^{d})^{*}\) and
        \(g_{\eta} \doteq g_{\eta}^{\unemb} \circ
        g_{\eta}^{\backbone} \colon (\R^{d})^{*} \to \cD\). Given an
        input \(\vX\), we write \(\vZ_{\theta}(\vX) \doteq
        f_{\theta}(\vX)\) and \(\hat{\vX}_{\theta, \eta}(\vX) \doteq
        g_{\eta}(\vZ_{\theta}(\vX))\). The overall pipeline is
        depicted in \Cref{fig:overall_autoencoder_pipeline}.
\end{itemize}

We will repeatedly use this notation many times in this chapter, so
please feel free to refer back to it if something doesn't make sense.
This decomposition of our networks also closely mirrors most code
implementations, and you can start your coding projects by defining
these networks.

\section{Unsupervised Learning of Image
Representations}\label{sec:contrastive_learning}

Learning high-quality and faithful representations of imagery data
distributions is a fundamental problem in machine intelligence, as we
have discussed in great length in Chapter \ref{ch:lossy-compression}.
In the literature, there have been many popular approaches proposed
for this task, many of which may appear not to use the techniques and
principles outlined in this manuscript. One such approach is called
\textit{contrastive learning}, so named because the learning
objective is (roughly speaking) about ensuring that features of
``similar'' data are similar, and features of ``dissimilar'' data are
far apart. Contrastive learning solutions are often
highly-engineered, empirically designed approaches. In this section,
we will introduce one such popular approach known as DINO
\citep{caron2021emerging}. Moreover, though, we will demonstrate how
to use the principles described in this book to drastically simplify
its design decisions while improving the learned representations.
This is known as the \textbf{SimDINO} \cite{wu2025simplifying}:
\\
\centerline{\url{https://robinwu218.github.io/SimDINO/}.}\\

\subsection{Data}\label{sub:contrastive_learning_data}
The data that we will use to explore and simplify the DINO
methodology are all two-dimensional (2D) image data.

For \textit{training},
we will use the ImageNet-1K and ImageNet-21K datasets. Each sample in
the dataset is an RGB image, of varying resolution, and a label
indicating the object or scene that the image contains (i.e., the
\textit{class} of the image). The ImageNet-1K dataset contains 1.28M
training images and 50K validation images partitioned into 1K
classes. The ImageNet-21K dataset contains 14.2M training images and
21.8K classes, but the classes are not disjoint (i.e., some classes
are subsets of others). Since we are doing self-supervised learning,
the labels will not be used during training, only during evaluation.
In order to increase the robustness of our model, we often apply
\textit{small data augmentations} to each image during processing,
such as flips, added small random noise, or random large crops; we do
not include this in our notation, as each augmentation of a natural
image is itself (very close to) a natural image in our dataset.

For \textit{evaluation}, we will use a litany of datasets. Of these,
the most common is CIFAR-10. CIFAR-10 is a dataset of 60K RGB 32
\(\times\) 32 natural images partitioned into 10 classes, with a
pre-established training set of 50K samples and a validation set of
10K samples. The purpose of using CIFAR-10 is to ensure that models
which train on one distribution of images (ImageNet) can generalize
to another distribution of images (CIFAR-10). We also refer to other
similar datasets, such as CIFAR-100 (disjoint from CIFAR-10), Oxford
Flowers, and Oxford Pets. Exemplars of ImageNet-1K and CIFAR-10 data
are shown in \Cref{fig:in1k_cifar10_examples}.

\begin{figure}[t!]
    \centering

    \hfill
    \begin{subfigure}{0.3\textwidth}
        \centering
        \includegraphics[width=\textwidth]{\toplevelprefix/chapters/chapter8/figs/imagenet.pdf}
        \caption{\small ImageNet-1K samples.}
    \end{subfigure}
    \hfill
    \begin{subfigure}{0.26\textwidth}
        \centering
        \includegraphics[width=\textwidth]{\toplevelprefix/chapters/chapter8/figs/cifar10.pdf}
        \caption{\small CIFAR10 samples.}
    \end{subfigure}
    \hfill
    \phantom{}

    \caption{\small\textbf{Images from ImageNet-1K \textit{(left)}
        and CIFAR-10 \textit{(right)}.} Notice that the CIFAR-10 images
        are much lower resolution, generally speaking, reducing the
    complexity of learning that distribution.}
    \label{fig:in1k_cifar10_examples}
\end{figure}

On a slightly more formal level, our data \(\vX\) will be images; we
let \(\cI\) be the set of all images. Since an image is a rectangular
array of pixels, and each pixel has a color given by RGB, CMYK, or
another color format, we say that an image is an element of \(\R^{c
\times h \times w}\) --- here \(c\) is the number of channels (i.e.,
\(3\) for RGB and \(4\) for CMYK), \(h\) is the image height, and
\(w\) is the image width. Consequently, the set of all images \(\cI
\doteq \bigcup_{c, h, w = 1}^{\infty}\R^{c \times h \times w}\) is
the set of all possible such data. Again, we will use this notation repeatedly.

\subsection{The Simplified DINO Task and Objective}
\label{sub:contrastive_learning_objective}

Our task is to learn a good representation of the data. Contrastive
learning, by and large, does this by defining what properties of the
input image we wish the features to reflect, constructing images
which share these properties but vary others, and setting up a loss
which promotes that the features of images with shared properties are
close and images with different properties are different. The
naturally optimal solution to this learning problem is that the
learned features preserve the desired properties of the input.
However, there are many practical and empirical complications that
arise in the course of training contrastive models.

In the case of DINO, the authors propose to use a methodology which
produces a single feature vector for the whole image and desires the
feature vector to contain ``global'' (i.e., image-level) information.
Accordingly, the loss will promote that images with similar global
information have similar features and images with different global
information have different features.

This seems intuitive, but as previously mentioned, there are several
empirical considerations, even while setting up the loss. First and
foremost, how should we promote similarities and differences? The
answer from DINO \citep{caron2021emerging} is\footnote{In the
author's view, inexplicably...} to convert the output features into
``logits'' corresponding to some probability distribution and take
their cross-entropy. More specifically, let \(\Delta_{m} \doteq \{\vx
        \in \R^{m} \colon x_{i} \geq 0\ \forall i \in [m], \sum_{i =
1}^{m}x_{i} = 1\}\) be the space of probability vectors in \(\R^{m}\)
and define the function \(d_{\CE} \colon \R^{m} \times \R^{m} \to \R\) by
\begin{equation}\label{eq:cross_entropy_difference}
    d_{\CE}(\vp, \vq) \doteq \CE(\vp, \vq), \quad \forall \vp, \vq
    \in \Delta_{m}
\end{equation}
where \(\CE \colon \Delta_{m} \times \Delta_{m} \to \R\) is the
cross-entropy, defined as
\begin{equation}\label{eq:expts_def_ce}
    \CE(\vp, \vq) \doteq -\sum_{i = 1}^{m}p_{i}\log q_{i}, \quad
    \forall \vp = (p_{1}, \dots, p_{m}), \vq = (q_{1}, \dots, q_{m})
    \in \Delta_{m}.
\end{equation}
Before we continue our discussion, let us build some intuition about
this distance function. We have, in particular,
\begin{align}
    \CE(\vp, \vq)
    &= -\sum_{i = 1}^{m}p_{i}\log q_{i} = \sum_{i = 1}^{m}p_{i}\log(p_{i}/q_{i})
    - \sum_{i = 1}^{m}p_{i}\log p_{i} \\
    &= \KL(\vp\mmid \vq) + H(\vp)
\end{align}
where \(\KL \colon \Delta_{m} \times \Delta_{m} \to \R\) is the KL
divergence, defined as
\begin{equation}
    \KL(\vp\mmid \vq) \doteq \sum_{i = 1}^{m}p_{i}\log(p_{i}/q_{i}),
\end{equation}
and \(H \colon \Delta_{m} \to \R\) is the entropy of a random variable. Note
that \(\KL(\vp\mmid \vq)\) is minimized if and only if \(\vp = \vq\).
So minimizing \(d_{\CE}\) does two things: it makes \(\vp = \vq\),
and it makes \(\vp\) and \(\vq\) have minimal entropy (i.e., vectors
    with \(1\) in one component and \(0\) elsewhere --- these are called
\textit{one-hot vectors}). Overall, the goal of this objective is not
\textit{just} to match \(\vp\) and \(\vq\) but also to shape them in
a certain way to make them low-entropy. Keep this in mind when we
discuss the formulation.

The next question is, how should we obtain samples with similar
global information? The answer from DINO (as well as nearly all
contrastive learning) is \textit{data augmentation} --- from each
sample, make several correlated samples which share the desired
properties. In the DINO case, we use different crops or
\textit{views} of the input image. Recall that we model an image as
an element of the set \(\cI\). In this notation, a view is a function
\(v \colon \cI \to \cI\). In the DINO case, the view is a
\textit{random resized crop}: it takes a randomly chosen rectangular
crop of the image (which has a fixed percentage \(p_{v} \in [0, 1]\)
of the total area of the image), resizes it proportionally so that
the \textit{shorter edge} is \(S_{\rsz}\) pixels long, then resizes
it to a fixed shape \((C, S_{v}, S_{v})\) where \(S_{v} \geq 1\) is
the size of the view and \(C\) is the number of channels in the original image.

\begin{figure}[t!]
    \centering
    \includegraphics[width=0.7\textwidth]{\toplevelprefix/chapters/chapter8/figs/global_local_views.pdf}
    \caption{\textbf{Local and global views in DINO.} Local views and
        global views take a rectangular crop of the input image and
        resize it to a square shape, which is then input into the network
    for processing.}
    \label{fig:dino_local_global_views}
\end{figure}

There are two types of views we want to use, depicted in
\Cref{fig:dino_local_global_views}:
\begin{itemize}
    \item \textit{global views}, which are random resized crops with
        area percentage parameter \(p_{\glo} \in [0, 1]\) and output
        shape \((C, S_{\glo}, S_{\glo})\);
    \item \textit{local views}, which are random resized crops with
        area percentage parameter \(p_{\loc} \in [0, 1]\) and output
        shape \((C, S_{\loc}, S_{\loc})\). Here \(p_{\loc} <
        p_{\glo}\) and \(S_{\loc} < S_{\glo}\).
\end{itemize}

DINO desires that the aggregate features \(\vz_{\theta}(\vX_{v})
\doteq (f_{\theta}^{\ext} \circ f_{\theta})(\vX_{v})\) of all views
\(\vX_{v} \doteq v(\vX)\) of an input image \(\vX\) be consistent
with each other. DINO does this by using a \textit{``DINO
head''}\footnote{Note that \(h_{\vW, \vmu}\) is the task-specific
    head, which in \Cref{sec:experiment_setup} is parameterized only by
    \(\theta\) as opposed to any specific parameters, but since we use
    two invocations of \(h\) with different values of the second
parameter, we keep the specified notation.} \(h_{\vW, \vmu}\),
parameterized by a matrix \(\vW \in \R^{s \times d}\) and a vector
\(\vmu \in \R^{s}\), to extract a probability vector \(\vp_{\theta,
\vW, \vmu}(\vX_{v}) \doteq h_{\vW, \vmu}(\vz_{\theta}(\vX_{v}))\)
from the aggregate feature \(\vz_{\theta}(\vX_{v})\), using the
following simple recipe:
\begin{equation}
    h_{\vW, \vmu}(\vz) \doteq \softmax([\vW\vz - \vmu]/\tau), \qquad
    \forall \vz \in \R^{d},
\end{equation}
where the \(\softmax \colon \R^{s} \to \Delta_{s}\) function is defined by
\begin{equation}
    \softmax\rp{\mat{x_{1} \\ \vdots \\ x_{s}}} \doteq
    \frac{1}{\sum_{i = 1}^{s}e^{x_{i}}}\mat{e^{x_{1}} \\ \vdots \\ e^{x_{s}}}
\end{equation}
and \(\tau > 0\) is a ``temperature'' parameter which controls the
entropy of the softmax's output.

In particular, DINO minimizes the difference between the probability
vector \(\vp_{\theta, \vW, \vmu}(\vX_{g}) \doteq h_{\vW,
\vmu}(\vz_{\theta}(\vX_{g}))\) for each global view \(\vX_{g} \doteq
v_{g}(\vX)\) and the probability vector \(\vp_{\theta, \vW}(\vX_{c})
\doteq h_{\vW, \vzero_{m}}(\vz_{\theta}(\vX_{c}))\) for each view
\(\vX_{c} \doteq v_{c}(\vX)\). Here, \(v_{c}\) can \textit{either} be
a local view or a global view. We will discuss the implementation of
\(f_{\theta}\) and \(f_{\theta}^{\ext}\) shortly in
\Cref{sub:contrastive_learning_architecture}. Overall, DINO solves the problem
\begin{eqnarray}
    &&\min_{\theta, \vW, \vmu}\cL_{\dino}(\theta, \vW, \vmu) \qquad
    \text{where} \nonumber \\
    &&\cL_{\dino}(\theta, \vW, \vmu) \doteq \Ex[d_{\CE}(\vp_{\theta,
    \vW, \vmu}(\vX_{g}), \vp_{\theta, \vW}(\vX_{c}))],\label{eq:dino_loss}
\end{eqnarray}
where the expectation is over data \(\vX\), global views \(v_{g}\),
and other views \(v_{c}\).

In this specific case, however, if you try to implement
\eqref{eq:dino_loss} and optimize it on a real network, it is very
likely that you will run into a problem: after running a few
iterations of the learning algorithm, the feature mapping
\(f_{\theta}^{\ext} \circ f_{\theta}\) will \textit{become the
constant function}! This certainly optimizes the above loss since it
minimizes the distance between features of different views of the
same image. But we obviously do not want to learn this solution.

Actually avoiding collapse is a very common consideration in
contrastive learning. So how do we do it in this case? The solution
from DINO, again, is empirically designed, and carefully tunes the
optimization of the parameter \(\vmu\) (which is updated using all
samples in the batch) and a ``temperature'' hyperparameter \(\tau\)
which is part of the implementation of \(h_{\vW, \vmu}\) and
discussed in \Cref{sub:contrastive_learning_architecture}. Given a
certain special set of hyperparameters that work well, this is indeed
enough to ensure non-collapse of the representation. However, outside
of this special configuration, training models to converge is
difficult, and the training is highly unstable.

To amend this state of affairs, let us discuss simplifications to the
formulation. First, instead of computing a probability vector using a
learned transformation \(h_{\vW, \vmu}\) of the aggregate features
\(\vz_{\theta}\), we can \textit{directly use the aggregate
representation}, ignoring the task-specific head (or equivalently,
setting it to the identity mapping). But now we need a way to compare
the vectors directly. Using our hypothesis from
\Cref{ch:deep-networks} that good representations should have
Euclidean (subspace) geometry, a much more natural measure of
difference is the \textit{squared \(\ell^{2}\) distance}
\(d_{\ell^{2}} \colon \R^{d} \times \R^{d} \to \R\), defined as
\begin{equation}\label{eq:cosine_similarity}
    d_{\ell^{2}}(\vx, \vy) \doteq \frac{1}{2}\norm{\vx -
    \vy}_{2}^{2}, \qquad  \forall \vx, \vy \in \R^{d}.
\end{equation}
This distance-based score is even more efficient to compute than the
cross-entropy score. Thus, \(d_{\ell^{2}}\) takes the place of
\(d_{\CE}\) in our simplification.

Before, collapse was avoided by using tricks to update \(\vmu\) and
\(\tau\). In our simplification, if we compare the features within
the representation space instead of converting them to probabilities,
we do not have either of these parameters and so must consider a
different way to avoid collapse. To do this, we return to the
fundamentals. The basic idea of avoiding collapse is that in order to
make sure that all samples do not return the same exact same
features, we need different samples to have different features. In
other words, we would like the \textit{covariance} of the features to
be \textit{large} in some sense. But from
\Cref{ch:lossy-compression,ch:deep-networks}, we already have a quantity
which measures the size of the covariance matrix. Namely, we use the
straightforward (population-level) \textit{Gaussian coding rate}
\(R\) to ensure that the features of global views of different
images, which have different global information, are well-separated
and not collapsed (hence expanded). The overall modified loss
\(\cL_{\simdino}\) becomes:
\begin{align}\label{eq:simdino_loss}
    \cL_{\simdino}(\theta) &\doteq \Ex[d_{\ell^{2}}(\vz_{\theta}(\vX_{g}),
    \vz_{\theta}(\vX_{c}))] \\
    &\qquad - \frac{\gamma}{2}\log\det\rp{\vI +
    \frac{d}{\eps^{2}}\Cov(\vz_{\theta}(\vX_{g}))},
\end{align}
where \(\eps > 0\) is fixed and the appropriate expectations are, as
before, taken over data \(\vX\), global view \(v_{g}\), and other
(local or global) view \(v_{c}\). The loss in \eqref{eq:simdino_loss}
is essentially the same loss introduced in Section \ref{sec:SimDINO}
of Chapter \ref{ch:lossy-compression} for the simplified DINO
(``SimDINO'' \cite{wu2025simplifying}). As we will see, when properly
implemented, it works at
least as well as the original DINO.

\subsection{Architecture: Vision
Transformer}\label{sub:contrastive_learning_architecture}

For the architecture, we use a standard vision transformer. Here is
how such an architecture works formally in the context of image data.
Recall from \Cref{sec:experiment_setup} that there are four
components to an encoder architecture, namely an embedding, a
backbone, a feature extractor, and a task-specific head. We discuss
these four parts presently.

\begin{figure}[t!]
    \centering
    \includegraphics[width=0.8\textwidth]{\toplevelprefix/chapters/chapter8/figs/patchify.pdf}
    \caption{\small\textbf{An example of an image turned into \(5
        \times 5\) square patches, which are placed in raster order.}
        Each patch is of the same size, and the grid of patches is of
        shape \((N_{H}, N_{W}) = (5, 5)\). The grid of patches is then
    unrolled into a sequence of length \(5 \times 5 = 25\) in raster order.}
    \label{fig:patchify_rasterize}
\end{figure}

\begin{figure}[t!]
    \centering
    \includegraphics[width=\textwidth]{\toplevelprefix/chapters/chapter8/figs/transformer_embedding.pdf}
    \caption{\small\textbf{The transformer embedding pipeline.} Given
        a sequence of unrolled patches in raster order \(\vX^{\patch}\),
        each unrolled patch is linearly projected into the feature space,
        and equipped with an (additive) positional encoding and an
        additional token known as the class token. The output is the
        first-layer-input feature \(\vZ_{\theta}^{1}(\vX) =
    f_{\theta}^{\emb}(\vX)\).}
    \label{fig:transformer_embedding}
\end{figure}

\paragraph{Embedding.} Given image data \(\vX \in \cI\), we embed it
as a sequence of tokens in \(\R^{d}\) using the map
\(f_{\theta}^{\emb}\), as follows. The first two steps are depicted
in \Cref{fig:patchify_rasterize}, and the latter two are depicted in
\Cref{fig:transformer_embedding}.
\begin{enumerate}
    \item First, we turn the image data \(\vX\) into a sequence of
        patches of shape \((C, P_{H}, P_{W})\) where \(P_{H}\) and
        \(P_{W}\) are the patch dimensions. We assume that \(P_{H}\)
        and \(P_{W}\) evenly divide the height and width of \(\vX\),
        respectively (in the notation of
            \Cref{sub:contrastive_learning_objective} we assume that
            \(P_{H}\) and \(P_{W}\) evenly divide \(S_{\loc}\) and
        \(S_{\glo}\)). Let the resulting grid of patches have
        \(N_{H}\) rows and \(N_{W}\) columns.
    \item We unroll each patch into a vector of length \(D \doteq
        CP_{H}P_{W}\). There are \(N \doteq N_{H}N_{W}\) patch
        vectors, which we place in ``raster order'' (top left \(\to\)
        top right \(\to\) bottom left \(\to\) bottom right) into a
        matrix \(\vX^{\patch} \in \R^{D \times N}\), where
        \(\vX^{\patch} \doteq f^{\patch}(\vX)\). Notice that \(D\)
        depends only on the patch size and number of channels. Since
        the latter quantity is normally constant among samples in the
        same dataset, \(D\) is the same for all images in the
        dataset, while \(N\) is different for larger and smaller images.
    \item We then perform the following operation on \(\vX^{\patch}
        \in \R^{D \times N}\) to project it to \(\R^{d \times n}\)
        where \(n \doteq N + 1\):
        \begin{equation}
            \vX^{\patch} \mapsto [\vz_{\cls}^{1}, \vW^{\emb}\vX] + \vE^{\pos}.
        \end{equation}
        Here we have three trainable parameters \(\vW^{\emb}\),
        \(\vz_{\cls}^{1}\), and \(\vE^{\pos}\) whose purpose is as follows:
        \begin{itemize}
            \item \(\vW^{\emb} \in \R^{d \times D}\) is a matrix
                which projects each patch vector to a token feature.
            \item \(\vz_{\cls}^{1} \in \R^{d}\) is a so-called
                \textit{class token} or \textit{register token}. The
                class token heuristically holds global information of
                the whole data and is used for downstream tasks. In
                the framework of compressive deep networks from
                \Cref{ch:deep-networks}, we expect that the class
                token is projected onto the same subspaces as the
                salient or semantically relevant tokens during the
                progression of the forward pass.
            \item \(\vE^{\pos} \in \R^{d \times N}\) is a so-called
                \textit{positional encoding} which distinguishes
                tokens of different patches from each other. That is,
                token features should have positional information, so
                that the overall map \(f^{\pre}\) is not invariant to
                permutations of the patches, and \(\vE^{\pos}\)
                inserts this positional information.
                \begin{itemize}
                    \item In this DINO case, where the transformer
                        receives differently-sized images, we learn a
                        positional encoding for the largest size
                        received during training, and interpolate to
                        get the positional encodings for smaller-sized inputs.
                \end{itemize}
        \end{itemize}
\end{enumerate}
Thus, in the end we have
\begin{equation}\label{eq:definition_of_embedding_module}
    f_{\theta}^{\emb}(\vX) \doteq \mat{\vz_{\cls}^{1},
    \vW^{\emb}f^{\patch}(\vX) + \vE^{\pos}}.
\end{equation}
All parameters \(\vz_{\cls}^{1}, \vW^{\emb}, \vE^{\pos}\) are
contained in the parameter set \(\theta\).

\begin{figure}[t!]
    \centering
    \includegraphics[width=\textwidth]{\toplevelprefix/chapters/chapter8/figs/transformer_backbone.pdf}
    \caption{\small\textbf{One layer \(f_{\theta}^{\ell}\) of the
        transformer backbone.} The input features go through
        layer-normalization, multi-head self-attention, and multi-layer
    perceptron blocks in sequence to form the output features of the layer.}
    \label{fig:transformer_backbone}
\end{figure}

\paragraph{Backbone.} Given a sequence of embeddings
\(\vZ_{\theta}^{1}(\vX) \doteq f_{\theta}^{\emb}(\vX) \in
(\R^{d})^{*}\), we process it using the backbone map
\(f_{\theta}^{\backbone}\) as follows and as depicted in
\Cref{fig:transformer_backbone}. The function
\(f_{\theta}^{\backbone}\) is composed of \(L\) \textit{layers}
\(f_{\theta}^{\ell}\), i.e.,
\begin{equation}
    f_{\theta}^{\backbone} = f_{\theta}^{L} \circ \cdots \circ f_{\theta}^{1}.
\end{equation}
The layer \(f_{\theta}^{\ell}\) has the following implementation:
\begin{align}\label{eq:vit-res-block}
    \vZ_{\theta}^{\ell + 1/2}(\vX)
    &= \vZ_{\theta}^{\ell}(\vX) +
    \MHSA_{\theta}^{\ell}(\LN_{\theta}^{1, \ell}(\vZ_{\theta}^{\ell}(\vX))) \\
    \vZ_{\theta}^{\ell + 1}(\vX)
    &= \vZ_{\theta}^{\ell + 1/2}(\vX) +
    \MLP_{\theta}^{\ell}(\LN_{\theta}^{2, \ell}(\vZ_{\theta}^{\ell + 1/2}(\vX)))
\end{align}
and \(f_{\theta}^{\ell}\) is defined such that
\(f_{\theta}^{\ell}(\vZ_{\theta}^{\ell}(\vX)) \doteq
\vZ_{\theta}^{\ell + 1}(\vX)\). Here we have used some operators,
such as \(\MHSA_{\theta}^{\ell}, \MLP_{\theta}^{\ell}\) and
\(\LN_{\theta}^{i, \ell}\) that are defined as follows:
\begin{itemize}
    \item The \(\MHSA_{\theta}^{\ell}\) operator is
        multi-head-self-attention, the predecessor of the multi-head
        subspace self-attention (cf \Cref{ch:deep-networks}). The
        formulation is as follows:
        \begin{align}\label{eq:mhsa}
            &\scalebox{0.95}{\(\displaystyle\MHSA_{\theta}^{\ell}(\vZ)
                    \doteq \vU_{\out}^{\ell}\mat{\SA([\vU_{\query}^{1,
                            \ell}]^{\top}\vZ, [\vU_{\attnkey}^{1,
                            \ell}]^{\top}\vZ,
                        [\vU_{\val}^{1, \ell}]^{\top}\vZ) \\ \vdots
                        \\ \SA([\vU_{\query}^{K, \ell}]^{\top}\vZ,
                            [\vU_{\attnkey}^{K, \ell}]^{\top}\vZ,
                            [\vU_{\val}^{K,
            \ell}]^{\top}\vZ)} + \vb_{\out}^{\ell}\vone_{n}^{\top},\)} \\
            \label{eq:self_attention}
            &\text{where} \qquad \SA(\vQ, \vK, \vV)
            \doteq \vV
            \underbrace{\softmax\rp{\frac{\vK^{\top}\vQ}{\sqrt{p}}}}_{\doteq
            \vA(\vQ, \vK)}
        \end{align}
        where \(p\) is a positive integer, \(\vU_{\query}^{k, \ell},
            \vU_{\attnkey}^{k, \ell}, \vU_{\val}^{k, \ell} \in \R^{d
        \times p}\), \(\vU_{\out}^{\ell} \in \R^{d \times Kp}\), and
        \(\vb_{\out}^{\ell} \in \R^{d}\) are trainable parameters
        contained in the parameter set \(\theta\), and the
        \(\softmax\) is defined column-wise as
        \begin{align}
            \softmax(\vM)
            &\doteq \mat{\softmax(\vm_{1}) & \cdots & \softmax(\vm_{p})}, \\
            \forall \vM
            &= \mat{\vm_{1}, \dots, \vm_{p}} \in \R^{n \times p}.
        \end{align}
        In practice, the dimensions are usually picked such that \(Kp
        = d\). The terms
        \begin{eqnarray}
            \vA_{\theta}^{k, \ell}(\vZ) &\doteq&
            \vA([\vU_{\query}^{k, \ell}]^{\top}\vZ,
            [\vU_{\attnkey}^{k, \ell}]^{\top}\vZ), \nonumber \\
            \SA_{\theta}^{k, \ell}(\vZ) &\doteq&
            \SA([\vU_{\query}^{k, \ell}]^{\top}\vZ,
                [\vU_{\attnkey}^{k, \ell}]^{\top}\vZ, [\vU_{\val}^{k,
            \ell}]^{\top}\vZ)
            \label{eq:attention_map}
        \end{eqnarray}
        are also known as the \textit{\(k\)th attention map} and
        \textit{\(k\)th attention head output} at layer \(\ell\),
        respectively. Furthermore, the operation \(\SA(\vQ, \vK,
        \vV)\) can be computed extremely efficiently using
        specialized software such as FlashAttention
        \citep{shah2025flashattention}.
    \item The \(\MLP_{\theta}^{\ell}\) is a two-layer perceptron, a
        regular nonlinearity used in deep networks, and has the form
        \begin{equation}
            \MLP_{\theta}^{\ell}(\vZ) \doteq
            \vW_{\down}^{\ell}\ReLU(\vW_{\up}^{\ell}\vZ +
            \vb_{\up}^{\ell}\vone_{n}^{\top}) +
            \vb_{\down}^{\ell}\vone_{n}^{\top}
        \end{equation}
        where \(\vW_{\up}^{\ell} \in \R^{q \times d},
            \vW_{\down}^{\ell} \in \R^{d \times q}, \vb_{\up}^{\ell} \in
        \R^{q}, \vb_{\down}^{\ell} \in \R^{d}\) are trainable
        parameters also contained in the parameter set \(\theta\),
        and \(\ReLU\) is the element-wise ReLU nonlinearity, i.e.,
        \(\ReLU(\vM)_{ij} = \max\{M_{ij}, 0\}\).
    \item Each layer-norm \(\LN_{\theta}^{i, \ell}\) for \(i \in \{1,
        2\}\) is a standard normalization, which applies column-wise
        to each token feature independently:
        \begin{equation}
            \LN_{\theta}^{i, \ell}(\vZ) = \LN_{\theta}^{i,
            \ell}(\mat{\vz_{1}, \dots, \vz_{n}}) =
            \mat{\LN_{\theta}^{i, \ell}(\vz_{1}), \dots,
            \LN_{\theta}^{i, \ell}(\vz_{n})}
        \end{equation}
        and has the form
        \begin{align}
            \LN_{\theta}^{i, \ell}(\vz)
            &= \frac{\vz
            - \operatorname{mean}(\vz)\vone_{d}}{\norm{\vz
            - \operatorname{mean}(\vz)\vone_{d}}_{2}} \hada
            \valpha^{i, \ell} + \vbeta^{i, \ell} \\
            \text{where}
            \qquad \operatorname{mean}(\vz) &= \frac{1}{d}\vone_{d}^{\top}\vz
        \end{align}
        where \(\hada\) denotes element-wise multiplication, and
        \(\valpha^{i, \ell}, \vbeta^{i, \ell} \in \R^{d}\) are
        trainable parameters contained in the parameter set
        \(\theta\). The layer-norm operator serves as a sort of
        normalization on each token, where the scale of each token
        afterwards is learnable and shared amongst all tokens.
\end{itemize}

The transformer is one of the most popular neural network
architectures in history, powering applications in almost all fields
of deep learning.

\paragraph{Feature extractor.} We use a post-processing step
\(f_{\theta}^{\ext}\) which extracts the \textit{class token
feature}, which (recall) is the feature meant to contain aggregate
information about the input image, and applies an MLP and
normalization to it. Namely, we have
\begin{align}
    \vz_{\theta}(\vX) &\doteq f_{\theta}^{\ext}(\vZ_{\theta}(\vX)) =
    f_{\theta}^{\ext}([\vz_{\theta}^{1}(\vX), \dots,
    \vz_{\theta}^{n}(\vX)]) \\
    &\doteq
    \frac{\MLP_{\theta}^{\ext}(\vz_{\theta}^{1}(\vX))}{\norm{\MLP_{\theta}^{\ext}(\vz_{\theta}^{1}(\vX))}_{2}}.
\end{align}

\paragraph{Task-specific (``DINO'') head.} For DINO, we use the
task-specific DINO head \(h_{\vW, \vmu}\). For SimDINO, we use
\textit{no} task-specific head \textit{at all}, as previously described.

\subsection{Optimization Strategy}\label{sub:contrastive_learning_optimization}

\begin{figure}[t!]
    \centering
    \includegraphics[width=\textwidth]{\toplevelprefix/chapters/chapter8/figs/dino_pipeline.pdf}
    \caption{\small \textbf{The DINO pipeline.} Student features and
        teacher features are computed for each input. The objective
        attempts to align the student features with the teacher features
        by projecting both sets of features into a high-dimensional
        probability simplex and computing a cross-entropy loss. Notably,
        because of the ``stop-grad'', the gradient is only computed
    w.r.t.~the \textit{student parameters' outputs}.}
    \label{fig:dino_pipeline}
\end{figure}

\paragraph{Optimizing DINO.} We have a loss function and an
architecture, so we now discuss the optimization strategy. The
optimization strategy for DINO uses \textit{two sets of weights for
the same architecture}: \textit{student} weights
\(\theta_{\student}\) and \textit{teacher} weights
\(\theta_{\teacher}\). These correspond to two different neural
networks, called the teacher network and student network, with the
same architecture. The teacher network encodes all global views,
while the student network encodes all ``other'' views. The goal of
the loss is to distill teacher outputs into the student model.
Namely, we train on the loss \(\cL_{\dino{}-\student\teacher}\):
\begin{equation}\label{eq:dino_loss_teacherstudent}
    \cL_{\dino{}-\student\teacher}(\theta_{\student},
    \theta_{\teacher}, \vW_{\student}, \vW_{\teacher}, \vmu) \doteq
    \Ex[d_{\CE}(\vp_{\theta_{\teacher}, \vW_{\teacher},
    \vmu}(\vX_{g}), \vp_{\theta_{\student}, \vW_{\student}}(\vX_{c}))].
\end{equation}
Now, we can fully describe the overall pipeline of DINO, depicted in
\Cref{fig:dino_pipeline}.

While it is easy to reason about \eqref{eq:dino_loss_teacherstudent},
it is impossible in practice to implement optimization algorithms
such as gradient descent with a loss given by
\(\cL_{\dino{}-\student\teacher}\). This is because the expectations
in the loss are impossible to evaluate, much less to take the
gradient of. In this extremely frequent case, we approximate the
expectation via finite samples. That is, at each timestep \(k\) we:
\begin{itemize}
    \item Subsample \(B\) data points from our dataset
        \(\{\vX_{1}^{(k)}, \dots, \vX_{B}^{(k)}\} \subset \cI\).
    \item For each data point \(\vX_{b}^{(k)}\), sample \(M_{\glo}\)
        global views \(v_{b, g}^{(k), i}\) and \(M_{\loc}\) local
        views \(v_{b, \ell}^{(k), i}\). Apply the views to
        \(\vX_{b}^{(k)}\) to obtain \(\vX_{b, g}^{(k), i} \doteq
        v_{b, g}^{(k), i}(\vX_{b}^{(k)})\) and \(\vX_{b, \ell}^{(k),
        i} \doteq v_{b, \ell}^{(k), i}(\vX_{b}^{(k)})\).
    \item For each \textit{local} view \(\vX_{b, \ell}^{(k), i}\),
        compute the following quantities:
        \begin{align}
            &\vz_{\theta_{\student}}(\vX_{b, \ell}^{(k), i}) \doteq
            (f_{\theta_{\student}}^{\ext} \circ
            f_{\theta_{\student}})(\vX_{b, \ell}^{(k), i}), \\
            &\vp_{\theta_{\student}, \vW_{\student}}(\vX_{b,
            \ell}^{(k), i}) \doteq h_{\vW_{\student},
            \vzero_{m}}(\vz_{\theta_{\student}}(\vX_{b, \ell}^{(k), i}(\theta)))
        \end{align}
        and for each \textit{global} view \(\vX_{b, g}^{(k), i}\),
        compute the following quantities (by an abuse of notation):
        \begin{align}
            &\vz_{\theta_{\student}}(\vX_{b, g}^{(k), i}) \doteq
            (f_{\theta_{\student}}^{\ext} \circ
            f_{\theta_{\student}})(\vX_{b, g}^{(k), i}), \\
            &\vp_{\theta_{\student}, \vW_{\student}}(\vX_{b, g}^{(k),
            i}) \doteq h_{\vW_{\student},
            \vzero_{m}}(\vz_{\theta_{\student}}(\vX_{b, g}^{(k), i})), \\
            &\vz_{\theta_{\teacher}}(\vX_{b, g}^{(k), i}) \doteq
            (f_{\theta_{\teacher}}^{\ext} \circ
            f_{\theta_{\teacher}})(\vX_{b, g}^{(k), i}), \\
            &\vp_{\theta_{\teacher}, \vW_{\teacher}, \vmu}(\vX_{b,
            g}^{(k), i}) \doteq h_{\vW_{\teacher},
            \vmu}(\vZ_{\theta_{\teacher}}(\vX_{b, g}^{(k), i})).
        \end{align}
    \item Compute the \textit{surrogate, approximate loss}
        \(\hat{\cL}_{\dino-\student\teacher}^{(k)}\), defined as follows:
        \begin{align}\label{eq:dino_loss_teacherstudent_empirical}
            &\hat{\cL}_{\dino{}-\student\teacher}^{(k)}(\theta_{\student},
            \theta_{\teacher}, \vW_{\student}, \vW_{\teacher}, \vmu) \doteq
            \frac{1}{BM_{\glo}(M_{\glo} + M_{\loc} - 1)}\sum_{b =
            1}^{B}\sum_{i = 1}^{M_{\glo}}\\
            &\Bigg[\sum_{j =
                1}^{M_{\loc}}d_{\CE}(\vp_{\theta_{\teacher},
                    \vW_{\teacher}, \vmu}(\vX_{b, g}^{(k), i}),
                    \vp_{\theta_{\student}, \vW_{\student}}(\vX_{b,
                \ell}^{(k), j})) \\
                &\qquad \qquad \qquad \qquad + \sum_{\substack{j = 1 \\ j \neq
                i}}^{M_{\glo}}d_{\CE}(\vp_{\theta_{\teacher},
                    \vW_{\teacher}, \vmu}(\vX_{b, g}^{(k), i}),
                    \vp_{\theta_{\student}, \vW_{\student}}(\vX_{b, g}^{(k),
            j}))\Bigg]
        \end{align}
        as well as its gradients with respect to
        \(\theta_{\student}\) and \(\vW_{\student}\), which should be
        computed under the assumption that \(\theta_{\teacher}\),
        \(\vW_{\teacher}\), and \(\vmu\) are constants --- namely
        that they are \textit{detached from the computational graph}
        and not dependent on \(\theta_{\student}\) and \(\vW_{\student}\).
    \item Update the student parameters \(\theta_{\student}\) and
        \(\vW_{\student}\) via an iterative gradient-based
        optimization algorithm, and update \(\theta_{\teacher}\),
        \(\vW_{\teacher}\), and \(\vmu\) via exponential moving
        averages with decay parameters \(\nu^{(k)}\), \(\nu^{(k)}\),
        and \(\rho^{(k)}\) respectively, i.e.,
        \begin{align}
            &(\theta_{\student}^{(k + 1)}, \vW_{\student}^{(k + 1)})
            = \textsc{OptUpdate}^{(k)}(\theta_{\student}^{(k)},
                \vW_{\student}^{(k)}; \nabla_{(\theta_{\student},
            \vW_{\student})}\hat{\cL}_{\dino-\student\teacher}^{(k)}) \\
            &\theta_{\teacher}^{(k + 1)}
            = \nu^{(k)}\theta_{\teacher}^{(k)} + (1 -
            \nu^{(k)})\theta_{\student}^{(k + 1)} \\
            &\vW_{\teacher}^{(k + 1)}
            = \nu^{(k)}\vW_{\teacher}^{(k)} + (1 -
            \nu^{(k)})\vW_{\student}^{(k + 1)} \\
            &\vmu^{(k + 1)}
            = \rho^{(k)}\vmu^{(k)} + (1 -
            \rho^{(k)})\cdot\frac{1}{BM_{\glo}}\sum_{b =
            1}^{B}\sum_{i =
            1}^{M_{\glo}}\vW^{(k)}\vz_{\theta_{\teacher}}(\vX_{b, g}^{(k), i}),
        \end{align}
        For example, if the chosen optimization algorithm were
        stochastic gradient descent, we would have the update
        \(\theta_{\student}^{(k + 1)} \doteq \theta_{\student}^{(k)}
            -
        \delta^{(k)}\nabla_{\theta_{\student}}\hat{\cL}_{\dino{}-\student\teacher}^{(k)}\),
        and so on.
\end{itemize}
Notice that the optimization procedure is rather irregular: although
all four parameters change at each iteration, only two of them are
directly updated from a gradient-based method. The other two are
updated from exponential moving averages, and indeed treated as
constants when computing any gradients. After training, we discard
the student weights and use the teacher weights for our trained
network \(f\), as this exponentially moving average has been
empirically shown to stabilize the resulting model (this idea is
known as Polyak averaging or iterate averaging).

The way that \(\nu\) and \(\rho\) change over the optimization
trajectory (i.e., the functions \(k \mapsto \nu^{(k)}\) and \(k
\mapsto \rho^{(k)}\)) are hyperparameters or design decisions, with
\(\nu^{(1)} < 1\) and \(\lim_{k \to \infty}\nu^{(k)} = 1\) usually,
and similar for \(\rho\). The temperature hyperparameter \(\tau\),
used in the DINO head \(h_{\vW, \vmu}\), also changes over the
optimization trajectory (though this dependence is not explicitly notated).

Using the surrogate (``empirical'') loss transforms our intractable
optimization problem, as in optimizing the loss in
\eqref{eq:dino_loss_teacherstudent}, into a tractable stochastic
optimization problem which is run to train essentially every deep
learning model in the world. This conversion is extremely natural
once you have seen some examples, and we will hopefully give these
examples throughout the chapter.

\begin{figure}[t!]
    \centering
    \includegraphics[width=0.7\textwidth]{\toplevelprefix/chapters/chapter8/figs/simdino_pipeline.pdf}
    \caption{\small\textbf{The SimDINO pipeline.} Here, in contrast
        to the DINO pipeline in \Cref{fig:dino_pipeline}, the loss is
        computed directly on the features without need of further
        manipulation. This shaves off several large matrices' worth of
        parameters and simplifies the pipeline, simultaneously making it
    more stable to train.}\label{fig:simdino_pipeline}
\end{figure}

\paragraph{Optimizing SimDINO.} The simplified DINO population-level
objective is very similar in spirit but much simpler in execution, i.e.,
\begin{align}\label{eq:simdino_loss_teacherstudent}
    \cL_{\simdino-\student\teacher}(\theta_{\student},
    \theta_{\teacher}) &\doteq
    \Ex\rs{d_{\ell^{2}}(\vz_{\theta_{\teacher}}(\vX_{g}),
    \vz_{\theta_{\student}}(\vX_{c}))} \\
    &\qquad \qquad - \frac{\gamma}{2}\log\det\rp{\vI
+ \frac{d}{\eps^{2}}\Cov(\vz_{\theta_{\student}}(\vX_{g})))}.
\end{align}
Thus, as elaborated in \Cref{fig:simdino_pipeline}, the SimDINO
pipeline is strictly simpler than the DINO pipeline. We can use a
simpler version of the DINO training pipeline to optimize SimDINO. At
each timestep \(k\), we:
\begin{itemize}
\item Subsample \(B\) data points from our dataset \(\{\vX_{1}^{(k)},
    \dots, \vX_{B}^{(k)}\} \subset \cI\).
\item For each data point \(\vX_{b}^{(k)}\), sample \(M_{\glo}\)
    global views \(v_{b, g}^{(k), i}\) and \(M_{\loc}\) local views
    \(v_{b, \ell}^{(k), i}\). Apply the views to \(\vX_{b}^{(k)}\) to
    obtain \(\vX_{b, g}^{(k), i} \doteq v_{b, g}^{(k),
    i}(\vX_{b}^{(k)})\) and \(\vX_{b, \ell}^{(k), i} \doteq v_{b,
    \ell}^{(k), i}(\vX_{b}^{(k)})\).
\item For each \textit{local} view \(\vX_{b, \ell}^{(k), i}\) compute
    \(\vz_{\theta_{\student}}(\vX_{b, \ell}^{(k), i}) \doteq
        (f_{\theta_{\student}}^{\ext} \circ
    f_{\theta_{\student}})(\vX_{b, \ell}^{(k), i})\). For each
    \textit{global} view \(\vX_{b, g}^{(k), i}\) compute
    \(\vz_{\theta_{\student}}(\vX_{b, g}^{(k), i}) \doteq
        (f_{\theta_{\student}}^{\ext} \circ
    f_{\theta_{\student}})(\vX_{b, g}^{(k), i})\) and
    \(\vz_{\theta_{\teacher}}(\vX_{b, g}^{(k), i}) \doteq
        (f_{\theta_{\teacher}}^{\ext} \circ
    f_{\theta_{\teacher}})(\vX_{b, g}^{(k), i})\).
\item Compute the \textit{surrogate, approximate loss}
    \(\hat{\cL}_{\simdino-\student\teacher}^{(k)}\), defined as follows:
    \begin{align}\label{eq:simdino_loss_teacherstudent_empirical}
        &\hat{\cL}_{\simdino{}-\student\teacher}^{(k)}(\theta_{\student},
        \theta_{\teacher})  \\
        \doteq &
        \scalebox{0.975}{\(\displaystyle\frac{1}{BM_{\glo}(M_{\glo} +
                M_{\loc} - 1)}\sum_{b =
                1}^{B}\sum_{i = 1}^{M_{\glo}}\Bigg[\sum_{j =
                    1}^{M_{\loc}}d_{\ell^{2}}(\vz_{\theta_{\teacher}}(\vX_{b,
                        g}^{(k), i}), \vz_{\theta_{\student}}(\vX_{b,
                            \ell}^{(k),
            j}))\)} \\
            & + \sum_{j =
            1}^{M_{\glo}}d_{\ell^{2}}(\vz_{\theta_{\teacher}}(\vX_{b,
                g}^{(k), i}), \vz_{\theta_{\student}}(\vX_{b, g}^{(k),
        j}))\Bigg] \\
        & - \frac{\gamma}{M_{\glo}}\sum_{i =
        1}^{M_{\glo}}R_{\eps}([\vz_{\theta_{\student}}(\vX_{1,
                g}^{(k), i}), \dots, \vz_{\theta_{\student}}(\vX_{B, g}^{(k),
        i})])
    \end{align}
    where \(R_{\eps}\) is the Gaussian coding rate estimated on
    finite samples, described in \Cref{ch:deep-networks}. The
    gradient of \(\hat{\cL}_{\simdino-\student\teacher}^{(k)}\) with
    respect to \(\theta_{\student}\) should (again) be computed,
    under the assumption that \(\theta_{\teacher}\) is constant.
\item Update the student parameters \(\theta_{\student}\) via an
    iterative gradient-based optimization algorithm, and update
    \(\theta_{\teacher}\) via an exponential moving average with
    decay parameter \(\nu^{(k)}\), i.e.,
    \begin{align}
        \theta_{\student}^{(k + 1)}
        &= \textsc{OptUpdate}^{(k)}(\theta_{\student}^{(k)};
        \nabla_{\theta_{\student}}\hat{\cL}_{\simdino-\student\teacher}^{(k)})
        \\
        \theta_{\teacher}^{(k + 1)}
        &= \nu^{(k)}\theta_{\teacher}^{(k)} + (1 -
        \nu^{(k)})\theta_{\student}^{(k + 1)}.
    \end{align}
\end{itemize}
Again, we re-iterate that the gradient is only taken
w.r.t.~\(\theta_{\student}\), treating \(\theta_{\teacher}\) as a
constant. Here, note that while the choice of \(\nu\) is still a
design decision, the hyperparameters \(\rho\) and \(\tau\) are removed.

\subsection{Evaluation Methodology}\label{sub:contrastive_learning_evals}
There are several ways to evaluate a trained transformer model. We
highlight two in this section. Let us define the \textit{center crop}
view \(v_{\cc} \colon \cI \to \cI\) which is a \textit{deterministic
resized crop}:
\begin{itemize}
\item it resizes the image so that the shortest edge is of size
    \(S_{\rsz}\) (similar to random resized crops with area
    percentage parameter \(1\));
\item then it takes the \textit{center} \(S_{\cc} \times S_{\cc}\) crop;
\end{itemize}
so that the final shape is \((C, S_{\cc}, S_{\cc})\). Notice that the
view \(v_{\cc}\) is completely deterministic given an input. For an
input \(\vX\), we write \(\vX_{\cc} \doteq v_{\cc}(\vX)\). Here
\(S_{\cc} \leq S_{\rsz}\).

\paragraph{Linear Probing.}

The first, and most architecture-agnostic, way to evaluate an encoder
model \(\vX \mapsto \vz_{\theta}(\vX)\) is to employ \textit{linear
probing}. Linear probing is, in a sentence, running logistic
regression on the aggregate features computed by the encoder. This
tells us how much semantic information exists in the representations,
as well as how easily this information can be extracted. (That is: to
what extent do the features of images with different semantics live
on different subspaces of the feature space?)

More formally, let us suppose that we want to evaluate the quality
and faithfulness of the features of the encoder on image-label data
\((\vX, \vy)\), where there are \(N_{\cls}\) classes and \(\vy \in
\{0, 1\}^{N_{\cls}}\) is a ``one-hot encoding'' (namely, zeros in all
positions except a \(1\) in the \(i\)th position if \(\vX\) is in the
\(i\)th class). One way to do this is to solve the logistic regression problem
\begin{equation}\label{eq:linear_probing}
\min_{\vW \in \R^{N_{\cls} \times d}}\Ex[\CE(\vy, \vW \vz_{\theta}(\vX_{\cc}))].
\end{equation}
More practically, if we have labeled data \(\{(\vX_{b}, \vy_{b})\}_{b
= 1}^{B}\), we can solve the \textit{empirical} logistic regression
problem (akin to \eqref{eq:dino_loss_teacherstudent}
vs.~\eqref{eq:dino_loss_teacherstudent_empirical}) given by
\begin{equation}\label{eq:linear_probing_empirical}
\min_{\vW \in \R^{N_{\cls} \times d}}\frac{1}{B}\sum_{b =
1}^{B}\CE(\vy_{b}, \vW \vz_{\theta}(\vX_{b, \cc})).
\end{equation}
This problem is a convex optimization problem in \(\vW\), and thus
can be solved efficiently via (stochastic) gradient descent or a
litany of other algorithms. This linear probe, together with the
encoder, may be used as a classifier, and we can evaluate the
classification accuracy. The usual practice is to train the model
first on a large dataset (such as ImageNet-1K), then train the linear
probe on a dataset (such as the training dataset of CIFAR-10), and
evaluate it on a third (``holdout'') dataset which is drawn from the
same distribution as the second one (such as the evaluation dataset
of CIFAR-10).

\paragraph{\(k\)-nearest Neighbors.}  We can also evaluate the
performance of the features on classification tasks \textit{without
needing to explicitly train a classifier} by using the \(k\)-nearest
neighbor algorithm to get an average predicted label. Namely, given a
dataset \(\{\vz_{b}\}_{b = 1}^{B} \subseteq \R^{d}\), define the
\(k\)-nearest neighbors of another point \(\vz \in \R^{d}\) as
\(\operatorname{NN}_{k}(\vz, \{\vz_{b}\}_{b = 1}^{B})\). Using this
notation, we can compute the predicted label \(\hat{\vy}_{\theta}(\vX
\mid \{(\vX_{b}, \vy_{b})\}_{b = 1}^{B})\) as
\begin{eqnarray}
&&\hat{\vy}_{\theta}(\vX \mid \{(\vX_{b}, \vy_{b})\}_{b = 1}^{B}) =
\vone(i^{\star}) \quad \text{where} \nonumber \\
&& i^{\star} \doteq \argmax_{i \in [Q]}\sum_{b =
1}^{B}\vy_{b}\indvar[\vz_{\theta}(\vX_{\cc, b}) \in
\operatorname{NN}_{k}(\vz_{\theta}(\vX_{\cc}))].
\end{eqnarray}
Here, \(\vone(i) \in \Delta_{N_{\cls}}\) is (by an abuse of notation,
cf.~indicator variables) the one-hot probability vector supported at
\(i\), i.e., \(1\) in the \(i\)th coordinate and \(0\) elsewhere.
That is, this procedure takes the most common label among the \(k\)
nearest points in feature space. The \(k\)-nearest neighbor
classification accuracy is just the accuracy of this predicted label, namely,
\begin{equation}
\Ex_{\vX, \vy}[\indvar(\hat{\vy}_{\theta}(\vX \mid \{(\vX_{b},
\vy_{b})\}_{b = 1}^{B}) = \vy)]
\end{equation}
or more commonly its corresponding empirical version, where \((\vX,
\vy)\) ranges over a finite dataset (\textit{not} the existing
samples \((\vX_{b}, \vy_{b})\) which are used for the \(k\) neighbors).

\paragraph{Fidelity of the Attention Maps.}

Another way to check the performance of the representations, for a
transformer-based encoder, is to examine the fidelity of the
attention maps \(\vA^{L, k} \in \R^{n \times n}\) as defined in
\Cref{eq:attention_map}, at the last layer \(L\), and given by the
following pipeline:
\begin{equation}
\scalebox{0.925}{\(\displaystyle\vX \mapsto \cdots \mapsto \vZ^{L -
        1} = [\underbrace{\vz_{1}^{L -
            1}}_{\text{class token}}, \underbrace{\vz_{2}^{L - 1} \dots,
        \vz_{n}^{L - 1}}_{\text{patch tokens}}] \mapsto \vA^{k, L} =
        \mat{\vA_{1, 1}^{k, L} & \vA_{1, 2:}^{k, L} \\ \vA_{2:, 1}^{k, L} &
\vA_{2:, 2:}^{k, L}}.\)}
\end{equation}
In particular, we examine what the attention maps for a given input
reveal about the salient objects in the input image, i.e., which
parts of the image provide the most globally-relevant information to
the class token. One particular way to do this is to examine the
component of the attention map where the class token is extracted as
the query and removed from the value matrix, i.e., \(\vA_{2:, 1}^{k,
L} \in \R^{1 \times (n - 1)} = \R^{1 \times N}\) or its transpose
\(\va^{k, L} = (\vA_{2:, 1}^{k, L})^{\top} \in \R^{N}\). Notice that
this vector \(\va^{k, L}\), which we label as the ``\textit{saliency
vector} at the \(k\)th attention head at layer \(L\),'' has a value
for every patch, \(1, \dots, N\), and we use this value to describe
how relevant each patch is toward the global information. In
particular for visualization's sake we create a new image where each
patch is replaced by its corresponding value in the saliency vector,
showcasing the contribution of each patch; we call this image the
``\textit{saliency map} at the \(k\)th attention head at layer
\(L\)''. To visualize the total relevance of each patch toward the
global information across all heads, we can average the saliency
vector, i.e., \(\tilde{\va}^{L} \doteq \frac{1}{K}\sum_{k =
1}^{K}\va^{k, L}\) and expand into the \textit{average saliency map}.
The average saliency maps should highlight the relevant parts of the
input image.

\paragraph{Object Detection and Segmentation.}

We can evaluate how the representations capture the fine-grained
(i.e., smaller or more detailed) properties of the input by using
them for \textit{semantic segmentation}. Roughly, this means that we
use the features to construct bounding boxes for all objects in the
input. There are several ways to do this, and several ways to score
the resulting bounding boxes compared to ground truth. Each
combination of methods corresponds to a particular segmentation
metric. We do not formally describe them here as they are not
particularly insightful, but the DINO paper \citep{caron2021emerging}
and DINOv2 paper \citep{oquab2023dinov2} contain references to all
metrics that are used in practice.

\subsection{Experimental Setup and Results}
\label{sub:contrastive_learning_experiment_results}

Since SimDINO is directly built upon DINO, we compare the optimal
settings for DINO as given by their original paper
\citep{caron2021emerging} with the same settings applied to SimDINO
\cite{wu2025simplifying}
for a fair comparison.

\paragraph{Objective Function.} We use \(10\) local views (i.e.,
\(M_{\loc} = 10\)) of resolution \(96 \times 96\) (i.e., \(S_{\loc} =
96\)) and \(2\) global views (i.e., \(M_{\glo} = 2\)) of resolution
\(224 \times 224\) (i.e., \(S_{\glo} = 224\)) for all experiments.
The corresponding portions of the original images cropped for local
and global views are \(p_{\loc} \in [\frac{1}{20}, \frac{3}{10}]\)
and \(p_{\glo} \in [\frac{3}{10}, 1]\) (chosen randomly per-view).
The smaller edge size within the resized crops is \(S_{\rsz} = 256\),
and the center crop (evaluation) view edge size is \(S_{\cc} = 224\).
All of these settings apply to both DINO and SimDINO.

\paragraph{Model Architecture.} For all inputs, we set the patch size
to be \(16 \times 16\) (namely, \(P_{H} = P_{W} = 16\)). We use the
small, base, and large models of the ViT \citep{dosovitskiy2020image}
architecture as the embedding and backbone. The feature extractor is
a three-layer MLP with a hidden size of \(2048\) and an output
dimension of \(256\), followed by an \(\ell^{2}\)-normalization, as
specified in \Cref{sub:contrastive_learning_architecture}. For DINO
architectures (i.e., not SimDINO architectures), the DINO head
\(\vW\) is a matrix in \(\R^{65536 \times 256}\), and the parameter
\(\vmu\) is a vector in \(\R^{65536}\).

\paragraph{Datasets and Optimization.} For pre-training, both our
DINO reproduction and SimDINO use the ImageNet-1K dataset across all
methods. We use AdamW \citep{Loshchilov2017DecoupledWD} as the
optimizer, which is a very standard choice. We follow the following
hyperparameter recommendations:
\begin{itemize}
\item The batch size is \(B = 1024\).
\item The learning rate (for AdamW and the student model) has
    ``base'' value \(2 \times 10^{-3}\). In the first \(10\) epochs
    the learning rate linearly increases from \(0\) to the base value
    (i.e., at the \(i\)th epoch the learning rate is \((i/10) \cdot 2
    \times 10^{-3}\), for \(1 \leq i \leq 10\)). Then over the next
    \(90\) epochs the learning rate decays via a so-called
    \textit{cosine schedule} back down to \(0\). The definition of a
    cosine schedule is given in many places, including
    \href{https://pytorch.org/docs/stable/generated/torch.optim.lr_scheduler.CosineAnnealingLR.html}{PyTorch
    documentation}, and it is commonly used when training deep vision models.
\item The weight decay (the W in AdamW) follows a cosine schedule
    from 0.04 to \(0.4\) over training.
\item The EMA rate \(\nu\) follows a cosine schedule from \(0.996\)
    to \(1.0\) over training. Specifically for DINO, the centering
    EMA rate \(\rho\) is fixed at \(0.9\).
\item Specifically for DINO, the teacher temperature
    \(\tau_{\teacher}\) is fixed at \(0.1\), while the student
    temperature \(\tau_{\student}\) linearly increases from \(0.04\)
    to \(0.07\) during the first \(30\) epochs and is fixed at
    \(0.07\) thereafter.
\end{itemize}
We use some (essentially information-preserving) data augmentations,
such as flips, color jittering, Gaussian blur, and solarization, for
each seen image during training, before taking the local and global
views. The exact hyperparameters governing these are not listed here,
but are referenced in the DINO paper \citep{caron2021emerging}.

For linear probing, the linear probe is usually trained using the
AdamW optimizer with learning rate \(2 \times 10^{-4}\), weight decay
\(0.01\), and batch size \(512\), but these are often modified on a
case-by-case basis to minimize the loss.

\begin{table}
\caption{\small\textbf{Classification performance} on hold-out test
    data for DINO and SimDINO, using both \(k\)-nearest neighbor accuracy
    (\(k = 20\)) and linear probing. At the same number of iterations
    (\(100\)), SimDINO is clearly better in terms of performance, and is
    more stable (the DINO training running on ViT-L backbone with the
        provided settings has very unstable optimization and obtains NaN loss
    in short order). We also compare to other standout methods, namely
SwAV and MoCov3, which DINO was built on.}

\centering
\begin{tabular}{@{}lcccc@{}} %
    \toprule
    Method & Model & Epochs & 20-NN & Linear Probing
    \\
    \midrule
    DINO & ViT-B & 100 & 72.9 & 76.3 \\
    SimDINO & ViT-B & 100 & \bf 74.9 & \bf 77.3 \\
    DINO & ViT-L & 100 & -- & -- \\
    SimDINO & ViT-L & 100 & \bf 75.6 & \bf 77.4 \\
    \midrule
    \color{gray} SwAV & \color{gray} ViT-S & \color{gray} 800 &
    \color{gray} 66.3 & \color{gray} 73.5 \\
    \color{gray} MoCov3 & \color{gray} ViT-B & \color{gray} 300  &
    \color{gray} -- & \color{gray} 76.7 \\
    \bottomrule
\end{tabular}
\label{tab:dino_imagenet_linear_probing}
\end{table}

\begin{figure}[t!]
\centering
\includegraphics[width=\textwidth]{\toplevelprefix/chapters/chapter8/figs/dino_attention_maps.png}
\caption{\small\textbf{A qualitative comparison of saliency maps}
    generated by DINO \textit{(middle row)} and by SimDINO
    \textit{(bottom row)}. For each image, we compute and display the
    average saliency map in the last layer \(L\). The saliency maps are
    similar across models, meaning that all models converge to a similar
    notion of what objects are important. Note that although
    \(X_{\evaluation}\) is a square image, it is interpolated back into
rectangular shape to make this visualization.}
\label{fig:dino_attention_maps_saliency}
\end{figure}

\begin{table}
\caption{\small\textbf{Segmentation performance} of pre-trained DINO
    and SimDINO models on COCO val2017 \citep{lin2014microsoft}, a
    segmentation dataset which contains object location metadata. We do
    not train on COCO, merely using the pre-trained embedding and
    backbone, and the bounding boxes are extracted from the features via
    a method called MaskCut \citep{wang2023cut}. Nevertheless, SimDINO
    surpasses DINO at object detection and segmentation under fair
    comparison, and even surpasses DINO with smaller patch size (side
    length \(8\) instead of \(16\)). Smaller patch sizes are known to
    help performance, especially with detection and segmentation tasks,
so this result is quite surprising and encouraging.}
\centering
\begin{tabular}{@{}llcccccccc@{}}
    \toprule
    &  & \multicolumn{3}{c}{Detection $\uparrow$} &
    \multicolumn{3}{c}{Segmentation $\uparrow$} \\
    Method & Model & AP$_{50}$  & AP$_{75}$ & AP & AP$_{50}$ &
    AP$_{75}$ & AP  \\
    \midrule
    SimDINO &ViT-L/16 &\bf 5.4 &1.9 &2.4 &4.5 &1.4 &1.9 \\
    SimDINO &ViT-B/16 &5.2 & \bf 2.0 & \bf 2.5 & \bf4.7 & \bf 1.5 & \bf 2.0 \\
    DINO &ViT-B/16 &3.9 &1.5 &1.8 &3.1 &1.0 &1.4 \\
    \midrule
    \color{gray} DINO & \color{gray} ViT-B/8 & \color{gray}5.1 &
    \color{gray}2.3 & \color{gray}2.5 & \color{gray}4.1 &
    \color{gray}1.3 & \color{gray}1.8 \\
    \bottomrule
\end{tabular}
\label{tab:dino_segmentation}
\end{table}

\paragraph{Evaluation Results.} In terms of downstream classification
performance, we obtain the performance in
\Cref{tab:dino_imagenet_linear_probing}. We observe that the
performance of SimDINO is much higher than that of DINO under fair
comparison. Also, it is much more stable: the prescribed settings of
DINO cannot train a ViT-L(arge) model. On the other hand,
\Cref{fig:dino_attention_maps_saliency} shows visualizations of the
average saliency maps in DINO and our simplified DINO, observing that
the saliency maps look quite similar across models, indicating that
the models learn features which are at least as good at capturing
fine-grained details. The segmentation and object detection
performances in \Cref{tab:dino_segmentation} confirm this claim
quantitatively, where SimDINO features show substantive improvement
over those of DINO.

\section{Weakly Supervised Image Representation via Text
Binding}\label{sec:CLIP}
%\yima{Robin, please provide a vignette about CLIP or its simplified
% version. Sam will give a general description of concepts and
% methods for joint image/text embedding and related work in Chapter 7. }

Another influential contrastive learning approach departs from purely
visual comparisons and instead leverages the natural alignment
between images and language. Rather than defining similarity solely
through multiple views of the same image, this line of work exploits
the fact that many images in the wild are accompanied by textual
descriptions that provide weak semantic supervision. A prominent
example of this paradigm is CLIP (Contrastive Language–Image
Pretraining) \citep{Radford2021-ir}. In CLIP, images and their
corresponding captions are treated as positive pairs, while
mismatched image–text pairs are treated as negatives. The learning
objective encourages the representation of an image to be close to
the representation of its associated caption, and far from captions
of other images. By grounding visual representations in natural
language, CLIP learns semantically rich features that capture
high-level concepts beyond appearance alone.

\subsection{Data}

The data that we will use to explore the CLIP methodology are {\em
text-image pairs}, rather than images alone. The original CLIP work
constructed a large web-scale corpus of approximately 400 million
text-image pairs collected from publicly available sources on the
Internet. Each sample consists of an image and an associated
natural-language string (e.g., a caption, title, or short
description) that co-occurs with the image on the web, providing a
weak and noisy form of supervision. To encourage broad coverage of
visual concepts, the CLIP work describes building this dataset by
searching for pairs whose text matches one of roughly 500{,}000
queries, and approximately balancing the results by including up to
20{,}000 pairs per query. This dataset is not publicly released, so
most community replications and extensions of CLIP rely instead on
open web-scale datasets of text-image pairs. The most common of which
is the LAION~\citep{schuhmann2022laion5bopenlargescaledataset}
family: LAION-400M provides 400 million CLIP-filtered pairs, and
LAION-5B scales this recipe to 5.85 billion pairs.

Because raw web-scraped pairs are often extremely noisy (e.g., images
that do not match their associated text, boilerplate alt-text, spam,
duplicates, or unsafe content), CLIP-style datasets are typically
constructed with a post-processing pipeline that filters and cleans
candidates before training. Typical methods include basic quality
filters (e.g., dropping pairs with extremely short text or very small
or corrupted images), language identification to retain captions in a
target language, de-duplication, and the removal of potentially
malicious content.

\subsection{The CLIP Task and Objective}
Our goal is to learn a good representation of the data. In the CLIP
setting, however, the notion of “similarity” is not induced by two
augmented views of the same image, but by the weak semantic
supervision provided by natural-language descriptions. Concretely, we
would like an image representation to be close to the representation
of its associated caption, and far from captions of other images.
Equivalently, if two images admit similar textual descriptions, their
learned visual features should be similar, whereas images described
by unrelated text should be separated in latent space.

We formalize the notion of ``similarity'' between representations
using \textit{cosine similarity}. Concretely, let \(f_{\theta}\) be
an image encoder and \(g_{\plainphi}\) be a text encoder (whose
architecture will be elaborated in the next section), both mapping
into a common latent space \(\vZ \in \mathbb{R}^{d}\). Given an image
\(\vX \in \cI\) and a text string \(\vT \in \cT\), we form normalized embeddings
\begin{equation}
\vz^{I}_{\theta}(\vX) \doteq \frac{f_{\theta}(\vX)}{\|f_{\theta}(\vX)\|_2},
\qquad
\vz^{T}_{\plainphi}(\vT) \doteq
\frac{g_{\plainphi}(\vT)}{\|g_{\plainphi}(\vT)\|_2}.
\end{equation}
Then the cosine similarity between embeddings is defined as:
\begin{equation}
s(\vX,\vT)
=
\left\langle \vz^{I}_{\theta}(\vX), \vz^{T}_{\plainphi}(\vT)\right\rangle
\in [-1,1].
\end{equation}
Therefore, our objective is to maximize the cosine similarity between
matched pairs and minimize it between mismatched pairs. To achieve
that, we can use a simple symmetric cross-entropy loss as discussed
in \Cref{sub:text-cond}. Concretely, given a mini-batch of \(n\)
text-image pairs \((\vX_i, \vT_i)_{i=1}^{n}\), we can define the loss as:
\begin{equation}\label{eq:practical_clip_loss}
\scalebox{0.925}{\(\displaystyle\cL_{\text{CLIP}} = -\frac{1}{n}\left(
            \sum_{i=1}^{n}
            \frac{\exp(s(\vX_i,\vT_i)/\tau)}{\sum_{k=1}^{n}\exp(s(\vX_i,\vT_k)/\tau)}
            +
            \sum_{i=1}^{n}
            \frac{\exp(s(\vT_i,\vX_i)/\tau)}{\sum_{k=1}^{n}\exp(s(\vT_i,\vX_k)/\tau)}
\right),\)}
\end{equation}
where \(\tau > 0\) is a temperature parameter that controls the
sharpness of the softmax function.

\subsection{Architecture: Vision Tower}\label{sub:clip_vision_tower}
The objective of CLIP is to learn both image and text representations
in the same latent space. Accordingly, its architecture trains a
vision encoder \(f_{\theta}: \cI \to \mathbb{R}^{d} \) and a text
encoder \(g_{\plainphi}: \cT \to \mathbb{R}^{d} \) , also known as
the ``dual-tower architecture''.

In terms of architecture, most vision backbones can be used: for
instance, convolutional networks such as ResNet~\citep{He2016-lc}, as
well as the Vision Transformer~\citep{dosovitskiy2020image} (detailed
in \Cref{sub:contrastive_learning_architecture}) are all valid
choices for CLIP. In practice, the Vision Transformer is more
commonly used due to better performance and versatility. Similarly to
DINO, we take the class token feature \(\vz_{\cls}\) as the
representation for each image: \(f_{\theta}(\vX) = \vz_{\cls}\). This
is because the class token is meant to aggregate global information
about the entire image and is commonly used for downstream prediction tasks.

% \paragraph{Feature Extractor.}
% Following the backbone, CLIP requires an additional feature
% extractor \(f^{\text{ext}}_{\theta}: \mathbb{R}^{t \times d^\prime}
% \to \mathbb{R}^{1 \times d^\prime} \), where \(t\) is the sequence
% length of the backbone's output and \(d^\prime\) is its hidden
% dimension. There are two mainstream designs of CLIP’s
% \(f^{\text{ext}}_{\theta}\), differing only in how they aggregate
% the backbone’s sequence output into a single \(d\)-dimensional
% image representation. The first is a \emph{mean extractor}: assume
% the backbone produces a sequence of latent features
% \([\vz_1,\dots,\vz_t] \in \mathbb{R}^{t\times d'}\), we take the
% pooled feature to be the average of token embeddings,
% \[
% \vz^{ext} =
% f^{\text{ext}}_{\theta}([\vz_1,\dots,\vz_t]) =
% \frac{1}{t}\sum_{i=1}^{t}\vz_i \in \mathbb{R}^{d'}.
% \]
% The second is a \emph{class-token extractor}, as is standard for
% Vision Transformers: assume the backbone output includes a
% designated class token along other latent features \([\vz_{\cls},
% \vz_1,\dots,\vz_t] \in \mathbb{R}^{(t+1)\times d'}\), then we have
% \[
% \vz^{ext} = \vz_{\cls}.
% \]

% \paragraph{Task-specific Head.}
% The feature extractor is typically followed by a projection head
% \(f^{\text{emb}}_{\theta} : \mathbb{R}^{d'} \to \mathbb{R}^{d}\)
% (often a learned linear map) to project \( \vz^{ext} \) into the
% \(d\)-dimension latent space shared of by the final image and text
% representation
% \[
% \vz^{ext} =
% \vW^{proj}\vz^{pool}
% \qquad
% \vW^{proj} \in \mathbb{R}^{d^\prime \times d}\,
% \]
% which is the final output of \(f^{I}_{\theta}\) and will later be
% used in the objective.

\subsection{Architecture: Text Tower}
The CLIP text tower \(g_{\plainphi}\) typically adopts a
Transformer-based architecture similar to the Vision Transformer
described earlier. The main difference lies in the embedding layer,
which is used to convert the text sequence into a sequence of latent
representations in the latent space \(\vZ \in \mathbb{R}^{d}\).

\paragraph{Embedding.} Given a text sequence \(\vT \in \cT\) (e.g., a
sentence or document), we embed it as a sequence of tokens in
\(\R^d\) using the map \(g_{\plainphi}^{\text{emb}}\), as follows.
\begin{enumerate}
\item First, we tokenize \(\vT\) into a sequence of discrete
    symbols (subwords/characters) from a vocabulary \(\cV\) of
    size \(|\cV|\). Concretely, a tokenizer
    \(t^{\tok} : \cT \to \cV^{N}\) produces a length-\(N\) token
    sequence\footnote{In practice one often augments this
        sequence with special tokens (e.g., \(\langle \text{bos}\rangle\) and
            \(\langle \text{eos}\rangle\)
            tokens that are used to indicate the beginning and end of the
        sequence) and may pad shorter sequences, but we keep the notation
    generic and treat these as elements of \(\cV\).}
    \((\tau_{1},\dots,\tau_{N})\), where
    \(\tau_i \in \cV\).
\item We convert tokens to integer indices by a vocabulary map
    \(\cV\to\{1,\dots,|\cV|\}\), yielding an index
    sequence \((s_{1},\dots,s_{N})\). We place these indices into
    a matrix \(\vS \in \{0,1\}^{|\cV|\times N}\), where the
    \(i\)-th column is the one-hot vector \(\ve_{s_i}\).
\item We then perform the following operation on \(\vS\) to project
    it to \(\R^{d \times (N+1)}\):
    \begin{equation}
        \vS \mapsto [\vz_{\cls}^{1}, \vW^{\tok}\vS] + \vE^{\pos}.
    \end{equation}
    This is the final output of \(g_{\plainphi}^{\text{emb}}\). Here
    we have three trainable parameters \(\vW^{\tok}\),
    \(\vz_{\cls}^{1}\), and \(\vE^{\pos}\) whose purpose is as follows:
    \begin{itemize}
        \item \(\vW^{\tok} \in \R^{d \times |\cV|}\) is a matrix
            whose \(j\)-th column is the learned embedding vector
            for the \(j\)-th vocabulary item. Thus \(\vW^{\tok}\vS
            \in \R^{d\times N}\) maps each discrete token to a
            continuous feature in \(\R^{d}\).
        \item \(\vz_{\cls}^{1} \in \R^{d}\) is a special
            \textit{classification} (or \textit{summary}) token.
            Heuristically, it is meant to aggregate global
            information about the entire sequence and is commonly
            used for downstream prediction tasks.
        \item \(\vE^{\pos} \in \R^{d \times n}\) is a
            \textit{positional encoding} which distinguishes
            tokens at different positions. Since self-attention
            is permutation-equivariant in the token dimension,
            \(\vE^{\pos}\) injects order information so that the
            overall map is not invariant to permutations of the
            tokens.
    \end{itemize}
\end{enumerate}
All parameters \(\vz_{\cls}^{1}, \vW^{\tok}, \vE^{\pos}\) are
contained in the parameter set \(\theta\). The vocabulary \(\cV\) and
tokenizer \(t^{\tok}\) are fixed.

Other components of the text tower, such as the backbone and the
task-specific head, follow similar design choices as the image tower
and are omitted for brevity. It is also worth noting that follow-up
works have demonstrated substantial flexibility in the choice of
CLIP's text encoder: fixed text representations contained in large
language models, and even simple bag-of-words features, can also be
effective. We discuss these variants in \Cref{sub:follow-up}.

% \paragraph{Backbone.} Now, our text input (a string) has been
% transformed into a sequence of latent representations
% \([\vz_{\cls}, \vz_{1}, \ldots, \vz_{n}]\) in the latent space
% \(\R^{d}\). Therefore, its subsequent processing follows the same
% procedure as in \cref{sub:contrastive_learning_architecture} for
% image latent representations.

% \paragraph{Feature Extractor.} The text transformer adopts a
% class-token extractor similar to the one in \cref{sub:clip_vision_tower}.

% \paragraph{Task-specific Head.} The text transformer adopts a
% linear head similar to the one in \cref{sub:clip_vision_tower}.

\subsection{Optimization Strategy}
We train CLIP using a standard end-to-end stochastic optimization
procedure (e.g. SGD with momentum, AdamW). Below we present a generic process.

At each timestep \(k\), we:
\begin{itemize}
\item \textbf{Subsample paired data.} Sample a minibatch of \(n\)
    paired examples
    \[
        \{(\vX_i^{(k)}, \vT_i^{(k)})\}_{i=1}^{n} \subseteq \cX \times \cT,
    \]
    where \(\vX_i^{(k)}\) is an image and \(\vT_i^{(k)}\) is its
    associated text.

\item \textbf{Compute latent representations.} For each pair
    \((\vX_i^{(k)}, \vT_i^{(k)})\), compute normalized embeddings
    using the vision and text encoders \(f_{\theta}\) and \(g_{\plainphi}\):
    \begin{equation}
        z_{\theta}^{I}(\vX_i^{(k)}) \doteq
        \frac{f_{\theta}(\vX_i^{(k)})}{\|f_{\theta}(\vX_i^{(k)})\|_2},
        \qquad
        z_{\plainphi}^{T}(\vT_i^{(k)}) \doteq
        \frac{g_{\plainphi}(\vT_i^{(k)})}{\|g_{\plainphi}(\vT_i^{(k)})\|_2}.
    \end{equation}

\item \textbf{Compute CLIP loss.} Compute the CLIP loss as defined in
    \eqref{eq:practical_clip_loss}:
    \begin{equation}
        \scalebox{0.875}{\(\displaystyle\hat{\cL}_{\text{CLIP}}^{(k)}
                \doteq
                -\frac{1}{n}\left(
                    \sum_{i=1}^{n}
                    \frac{\exp(s(\vX_i^{(k)},\vT_i^{(k)})/\tau)}{\sum_{k=1}^{n}\exp(s(\vX_i^{(k)},\vT_k^{(k)})/\tau)}
                    +
                    \sum_{i=1}^{n}
                    \frac{\exp(s(\vT_i^{(k)},\vX_i^{(k)})/\tau)}{\sum_{k=1}^{n}\exp(s(\vT_i^{(k)},\vX_k^{(k)})/\tau)}
        \right).\)}
    \end{equation}

\item \textbf{Compute gradients.} Compute the gradients of the
    stochastic loss with respect to both parameter sets:
    \begin{equation}
        \nabla_{\theta}\hat{\cL}_{\text{CLIP}}^{(k)}(\theta,\plainphi),
        \qquad
        \nabla_{\plainphi}\hat{\cL}_{\text{CLIP}}^{(k)}(\theta,\plainphi).
    \end{equation}

\item \textbf{Update parameters.} Apply one step of an
    optimization algorithm to \((\theta,\plainphi)\), yielding the
    iteration
    \begin{equation}
        (\theta^{(k+1)}, \plainphi^{(k+1)})
        \doteq
        \textsc{OptUpdate}^{(k)}\!\left(
            (\theta^{(k)}, \plainphi^{(k)});
            \nabla_{\theta}\hat{\cL}_{\text{CLIP}}^{(k)},
            \nabla_{\plainphi}\hat{\cL}_{\text{CLIP}}^{(k)}
        \right).
    \end{equation}
\end{itemize}

\subsection{Evaluation Methodology}\label{sub:clip_eval}
Since CLIP is architecture-agnostic, most standard evaluation
protocols for vision and language encoders can be applied. For
example, one can perform linear probing by training a lightweight
classification head on top of the vision encoder to assess the
quality of the learned representations in image classification tasks.
Further, one can perform zero-shot classification without any
additional training thanks to the dual-tower architecture of CLIP.

Recall that in a standard \(n\)-way image classification setup, we
are given a dataset of images
\(\{\vX_i\}_{i=1}^{N}\subseteq\cI\), each annotated with one of \(n\)
class labels from the set
\(\cL \doteq \{\ell_1,\ldots,\ell_n\}\) (e.g., \(\ell_j \in
\{\text{cat},\text{dog},\ldots\}\)).
Our goal is to build a prediction pipeline \(P\), leveraging a
pre-trained model, such that
\(P:\cI\to\{1,\ldots,n\}\) maps each image \(\vX_i\) to a predicted class index
\(\hat{y}_i \doteq P(\vX_i)\). Writing the ground-truth class label
as its index \(y_i \in \{1,\ldots,n\}\),
the (top-1) classification accuracy of \(P\) on this dataset is then computed as
\[
\frac{1}{N}\sum_{i=1}^{N}\boldsymbol{1}\{\hat{y}_i = y_i\}.
\]

In CLIP, we can use the text encoder \(g_{\plainphi}\) to build a
text dictionary \(\vD\in\R^{n\times d}\), whose \(j\)-th row is the
normalized latent representation of the \(j\)-th class label.
Concretely, for each class \(\ell_j \in \cL\), we construct a textual
description (prompt) \(\vT_j\) (e.g.,
``a photo of a \(\ell_j\)''), and compute its normalized embedding
\(\vz_j^T \doteq
\frac{g_{\plainphi}(\vT_j)}{\|g_{\plainphi}(\vT_j)\|_2} \in \R^{d}\).
Stacking these row-wise yields
\[
\vD \doteq
\mat{(\vz_1^T)^{\top}\\ \vdots \\ (\vz_n^T)^{\top}}
\in \R^{n\times d},
\]

Given an image \(\vX_i \in \cI\), we compute its image representation
via the vision encoder and normalize it as
\[
\vz_i^I \doteq \frac{f_{\theta}(\vX_i)}{\|f_{\theta}(\vX_i)\|_2} \in \R^{d}.
\]
We can then compute a similarity vector
\[
\vd_i \doteq \vD\,\vz_i^I \in \R^{n},
\]
whose \(j\)-th entry \(\vd_i[j]\) is the cosine similarity between
the \(i\)-th image and the \(j\)-th class label
in the shared latent space. By the CLIP training objective, images
are encouraged to be more similar to their
matching texts than to mismatched ones; hence a larger cosine
similarity indicates closer alignment between
\(\vX_i\) and class \(\ell_j\). Therefore, we predict the class index by
\[
\hat{y}_i \doteq \arg\max_{j \in \{1,\ldots,n\}} \vd_i[j].
\]
Finally, we can calculate the classification accuracy as mentioned above.

This zero-shot decision protocol is not limited to classification,
but also applies naturally to \emph{retrieval} tasks,
where the goal is to return the most relevant items from a large
database given a query (e.g., retrieve the most
relevant images for a text query, or the most relevant texts for an
image query). For instance, in text-to-image
retrieval, instead of building a dictionary from class-name prompts,
we build an \emph{image dictionary} by encoding
and normalizing all database images \(\{\vX_m\}_{m=1}^{M}\), stacking
their embeddings into
\(\vD\in\R^{M\times d}\). Given a query text \(\vT\), we compute its
normalized embedding
\(\vz^T \doteq g_{\plainphi}(\vT)/\|g_{\plainphi}(\vT)\|_2\), form
similarity scores
\(\vd \doteq \vD\vz^T\in\R^{M}\), and then rank images by these
scores (equivalently, return the top-\(K\)
indices).
% More broadly, this simple embedding-and-dot-product pipeline is
% arguably the most distinctive and
% practically important feature of CLIP.

\subsection{Experimental Setup and Results}

Next, we present results from the original CLIP paper comparing CLIP
to a supervised ResNet baseline on several classic classification
benchmarks. Across all benchmarks, CLIP is evaluated using the
zero-shot classification protocol described in \Cref{sub:clip_eval}.
Full details of the ResNet experimental setup are provided
in~\citep{He2016-lc}; here, we briefly introduce the evaluation
datasets and summarize CLIP’s training and evaluation settings
relevant to these comparisons.

\paragraph{Evaluation Data.}
\begin{itemize}
\item \textbf{CIFAR-10~\citep{krizhevsky2009learning}.}
    CIFAR-10 is a widely used small-scale image classification
    benchmark consisting of 60{,}000 RGB images of size \(32\times
    32\) across 10 classes (6{,}000 images per class).
    It is typically split into 50{,}000 training images and 10{,}000
    test images.

\item \textbf{CIFAR-100~\citep{krizhevsky2009learning}.}
    CIFAR-100 has the same overall size and image resolution as
    CIFAR-10 (60{,}000 RGB \(32\times 32\) images with a
    50{,}000/10{,}000 train/test split), but it is more challenging
    because it contains 100 fine-grained classes (600 images per class).
    The 100 classes are additionally grouped into 20 superclasses.

\item \textbf{PASCAL VOC 2007 (VOC2007)~\citep{everingham2010pascal}.}
    PASCAL VOC 2007 is a classic computer vision benchmark for object
    recognition, supporting tasks such as image classification,
    object detection, and segmentation.
    It focuses on 20 object categories (e.g., person, car, dog).

\item \textbf{ImageNet-1K~\citep{imagenet_cvpr09}.}
    ImageNet-1K is a subset of ILSVRC 2012 containing 1{,}000 object classes.
    This subset includes 1{,}281{,}167 training images, 50{,}000
    validation images, and 100{,}000 test images, and has played a
    central role in large-scale visual recognition and pretraining.
\end{itemize}

\paragraph{Model Setup.}
In the CLIP models demonstrated here, Vision Transformer is adopted
as  the vision encoder \(f_{\theta}\) and is instantiated as one of
three ViT variants: ViT-B/32, ViT-B/16, or ViT-L/14. Here,
\(\text{B}\) (``Base'') and \(\text{L}\) (``Large'') specify the
transformer width/depth (e.g., ViT-B typically uses 12 transformer
blocks with embedding dimension \(d=768\) and 12 attention heads,
whereas ViT-L uses 24 blocks with \(d=1024\) and 16 heads), while the
suffix \(/32\), \(/16\), and \(/14\) denotes the \emph{patch size} in
pixels used to patchify the input image . Smaller patch sizes (e.g.,
16 or 14) yield more tokens per image at a fixed resolution,
increasing compute but often improving accuracy due to finer spatial
granularity.

\paragraph{Optimization Setup.}
All ViT-based CLIP models are trained for 32 epochs using the Adam
optimizer with decoupled weight decay (i.e., AdamW-style
regularization), where weight decay is applied to all parameters
except gains and biases. The learning rate is decayed using a cosine
schedule. The learnable temperature parameter \(\tau\) (used to scale
the contrastive logits) is initialized to the equivalent of \(0.07\),
and is clipped to prevent scaling the logits by more than 100, which
is found necessary for training stability. Finally, training uses a
very large mini-batch size of \(32{,}768\).

\paragraph{Results.}\Cref{tab:clip_vs_resnet} reports zero-shot top-1
accuracy of three CLIP ViT variants (B/32, B/16, L/14) and compares
them to supervised ResNet-50 and ResNet-101 baselines.
Across all four datasets, CLIP consistently outperforms the supervised ResNets.
Scaling the CLIP backbone further yields additional gains, with
ViT-L/14 reaching \(98.0\) on CIFAR-10, \(87.5\) on CIFAR-100,
\(89.6\) on VOC2007, and \(83.9\) on ImageNet.
Nevertheless, this table should be interpreted as a comparison
between two \emph{systems} rather than a controlled backbone
ablation: CLIP benefits from large-scale image--text pretraining and
is evaluated via zero-shot prompting, whereas the ResNet baselines
rely on supervised pretraining and require labeled data to fit a
probe on each benchmark.
Thus, the consistent margins primarily demonstrate the strength of
CLIP's transferable representations under the zero-shot protocol,
rather than establishing a purely architecture- or data-matched
superiority of ViT over ResNet.

\begin{table}[t]
\centering
\begin{tabular}{@{}llcccc@{}}
    \toprule
    Backbone & Variant & CIFAR-10 & CIFAR-100 & VOC2007 & ImageNet \\
    \midrule
    CLIP-ViT & B/32 & 95.1 & 80.5 & 87.7 & 76.1 \\
    CLIP-ViT & B/16 & 96.2 & 83.1 & 89.2 & 80.2 \\
    CLIP-ViT & L/14 & 98.0 & 87.5 & 89.6 & 83.9 \\
    \midrule
    ResNet & 50  & 91.8 & 74.5 & 83.8 & 74.3 \\
    ResNet & 101 & 93.0 & 77.2 & 84.4 & 75.8 \\
    \bottomrule
\end{tabular}
\caption{Performance on downstream datasets for various CLIP
backbones and ResNet baselines.}
\label{tab:clip_vs_resnet}
\end{table}

\subsection{Simplified Extension: LIFT}\label{sub:follow-up}

While effective, CLIP exhibits several limitations. First, jointly
training both a text encoder and an image encoder from scratch is
computationally expensive, often requiring extremely large batches
and massive datasets to reach strong alignment and downstream
transfer. Second, models trained in this way can struggle with
compositional understanding—capturing word order in text, spatial
layout in images, and object–attribute or object–object
relations—partly because retrieval-style contrastive supervision can
reward shortcut solutions that downweight fine-grained compositional
features. LIFT (Language–Image alignment with a Fixed Text
encoder)~\citep{yang2025languageimagealignmentfixedtext} revisits a
core assumption underlying these pipelines: that optimal alignment
demands joint end-to-end training of both encoders:\\
\centerline{\url{https://jingfeng0705.github.io/LIFT/}.}\\
Instead, LIFT leverages the observation that modern large language
models (LLMs) already produce highly informative text embeddings.
Concretely, LIFT fixes a strong pretrained text encoder (e.g., one
derived from or fine-tuned on an LLM), computes text embeddings
offline, and trains only the image encoder to align to these fixed targets.

\paragraph{Architecture.}
LIFT retains the same dual-tower CLIP formulation in terms of a
vision encoder and a text encoder that map into a shared latent
space, but crucially differs in that the text tower is \emph{fixed}
and can be instantiated by any strong LLM-based text encoder
\(g_{\plainphi}:\cT\to\R^{d}\) to produce text embeddings, where
\(\plainphi\) is not updated during training. In our implementation,
we adopt the NV-Embed-V2 text encoder, for which the embedding
dimension is \(d=4096\). Since the text tower \(g_{\plainphi}\) is
fixed, we can pre-compute all text embeddings
\(\{g_{\plainphi}(\vT)\}\) offline and, during training, only compute
gradients and update the image encoder \(f_{\theta}\) online. The
image tower \(f_{\theta}:\cI\to\R^{d}\) follows the same structure as
in \Cref{sub:clip_vision_tower}, with the projection head output
dimension set to \(d\) to match the dimension of the fixed text
embedding space. Apart from fixing \(g_{\plainphi}\) and matching the
projection dimension, the remaining components---including the
contrastive loss and the optimization procedure---follow the same
design principles as CLIP.

\paragraph{Evaluation Methodology.}
Compositional understanding is a known limitation of CLIP. We
evaluate LIFT and CLIP on seven
SugarCrepe~\citep{hsieh2023sugarcrepe} tasks, where each image
\(\vX\) is paired with a \emph{positive} (correct) caption
\(\vT_{\text{pos}}\) and a \emph{negative} caption
\(\vT_{\text{neg}}\) constructed by adding, replacing, or swapping an
object, attribute, or relation in \(\vT_{\text{pos}}\). See Figure
\ref{fig:SugarCrepe-examples} for some examples. Models are asked to
identify the correct caption by comparing caption--image cosine
similarities. Formally, for each sample \(i\in\{1,\dots,N\}\), we
compute the normalized embeddings
\[
\vz_i^{I} \doteq \frac{f_{\theta}(\vX_i)}{\|f_{\theta}(\vX_i)\|_2},\qquad
\vz_{i,\text{pos}}^{T} \doteq
\frac{g_{\plainphi}(\vT_{i,\text{pos}})}{\|g_{\plainphi}(\vT_{i,\text{pos}})\|_2},\qquad
\vz_{i,\text{neg}}^{T} \doteq
\frac{g_{\plainphi}(\vT_{i,\text{neg}})}{\|g_{\plainphi}(\vT_{i,\text{neg}})\|_2}.
\]
We then compute similarity scores via cosine similarities:
\[
s_{i,\text{pos}} \doteq \langle \vz_i^{I}, \vz_{i,\text{pos}}^{T} \rangle,
\qquad
s_{i,\text{neg}} \doteq \langle \vz_i^{I}, \vz_{i,\text{neg}}^{T} \rangle.
\]
The model is considered correct on sample \(i\) if
\(s_{i,\text{pos}} > s_{i,\text{neg}}\), and the overall
compositional accuracy is
\[
\frac{1}{N}\sum_{i=1}^{N}\boldsymbol{1}\{s_{i,\text{pos}} > s_{i,\text{neg}}\}.
\]

\begin{figure}[t!]
\centering
\includegraphics[width=\textwidth]{\toplevelprefix/chapters/chapter8/figs/negative_caption_demo.pdf}
\caption{The original captions (top) and their negative counterparts
    (bottom) from two SugarCrepe tasks: \texttt{replace relation} (left)
and \texttt{swap attribute (right)}.}
\label{fig:SugarCrepe-examples}
\end{figure}

\paragraph{Experimental Results.}
As shown in \Cref{tab:Compositionality}, when trained on DataComp-1B,
LIFT outperforms CLIP on all seven tasks with a \textbf{6.8\%}
average accuracy gain. LIFT achieves significant gains on \texttt{add
attribute}, \texttt{replace attribute}, and \texttt{replace relation}
tasks. These improvements are strong evidence that LLM-based text
encoders \(g_{\plainphi}\) capture more informative text
representations that enable more accurate modeling of
object–attribute associations and object–object relations. On the
other hand, we can see LIFT shows relatively low accuracy on
\texttt{swap object} and \texttt{swap attribute} compared to other
SugarCrepe tasks. We attribute this limitation to the contrastive
learning objective, which primarily focuses on aligning lower-order
statistics. Addressing this challenge requires exploring more refined
information-theoretic measures for language-image alignment, a key
direction for future work.

\begin{table}[t]
\centering
\resizebox{\textwidth}{!}{%
    \begin{tabular}{@{}lllccccccc@{}}
        \toprule
        & & & \multicolumn{2}{c}{Add} & \multicolumn{3}{c}{Replace} &
        \multicolumn{2}{c}{Swap} \\
        Method & Dataset & Sample Seen & Obj & Att & Obj & Att & Rel
        & Obj & Att \\
        \midrule
        OpenCLIP & DataComp & 1.28B & 82.3 & 73.7 & 91.7 & 79.4 &
        61.2 & 59.6 & 56.9 \\
        LIFT     & DataComp & 1.28B & \textbf{89.0} & \textbf{86.1} &
        \textbf{93.2} & \textbf{86.0} & \textbf{70.6} & \textbf{64.1}
        & \textbf{63.4} \\
        \bottomrule
    \end{tabular}%
}
\caption{The performance of LIFT and CLIP on seven SugarCrepe tasks.
    We use an open-source CLIP implementation,
    OpenCLIP~\citep{ilharco_gabriel_2021_5143773}, and train both
    OpenCLIP and LIFT on the open-source text--image pair dataset
    DataComp~\citep{gadre2023datacompsearchgenerationmultimodal}. We
report accuracy for each task, with the best results bolded.}
\label{tab:Compositionality}
\end{table}

\section{Supervised Image Representation via
Classification}\label{sec:image_classification}

In the previous sections, we have shown how to simplify some overly
complex, often heuristically designed, learning objectives based on
the principles of compression introduced in this book. However,
objectives associated with some of the most popular learning tasks
are already rather simple. In these cases, it is difficult to further
simplify the objective. Thus, in this and future sections, we will
focus on principled ways to modify the \textit{deep network
architectures} for a variety of tasks.

Let us first start with arguably the most classical task in machine
learning: \textit{image classification}, which is often used as a
standard task to evaluate pattern recognition algorithms or deep
network architectures. From our discussion of white-box architectures
in \Cref{ch:deep-networks}, we only need a semantically meaningful
task to learn good representations with white-box architectures. We
will validate this idea in this section.

First, the dataset stays largely the same as
\Cref{sub:contrastive_learning_data}. Both the training and test data
consist of labeled images, i.e., image-label pairs \((\vX, \vy) \in
\R^{C \times H \times W} \times \{0, 1\}^{N_{\cls}}\). We still apply
various data augmentations (e.g., flips, Gaussian blurring,
solarization, etc.) to each sample in each new batch.

\subsection{Task and Objective} \label{sub:image_classification_objective}

Unlike before, our task is not just to learn a good representation of
the data, but also to simultaneously build a classifier. Formally, we
have labeled data pairs \((\vX, \vy)\), where \(\vy \in \{0,
1\}^{N_{\cls}}\) is a one-hot vector denoting the class membership of
\(\vX\). We consider a deterministic \textit{center crop view}
\(v_{\cc}\) of the input data \(\vX\) (cf
\Cref{sub:contrastive_learning_objective}). We want to jointly train
a feature mapping \((f_{\theta}, f_{\theta}^{\ext})\) and a
\textit{classification head} \(h_{\theta}\), defined as follows:
\begin{equation}
h_{\theta}(\vz) \doteq \softmax(\vW^{\head}\vz + \vb^{\head}), \qquad
\forall \vz \in \R^{d}
\end{equation}
where \((\vW^{\head}, \vb^{\head}) \in \R^{N_{\cls} \times d} \times
\R^{N_{\cls}}\) are trainable parameters in the parameter set
\(\theta\), such that the map \(\vX_{\cc} \mapsto
\vp_{\theta}(\vX_{\cc}) \doteq h_{\theta}(\vz_{\theta}(\vX_{\cc}))\)
predicts a smoothed label for the view \(\vX_{\cc} = v_{\cc}(\vX)\)
of the input \(\vX\). The learning problem attempts to minimize the
distance between \(\vp_{\theta}\) and \(\vy\) measured through cross-entropy:
\begin{equation}\label{eq:classification_ce_loss}
\min_{\theta}\bc{\cL_{\CE}(\theta) \doteq \Ex[\CE(\vy,
\vp_{\theta}(\vX_{\cc}))]}.
\end{equation}

\subsection{The CRATE Architecture}\label{sub:image_classification_architecture}

The architecture that we use is the CRATE architecture, described in
some detail in \Cref{ch:deep-networks}. The overall setup is similar
to that of the regular transformer in
\Cref{sub:contrastive_learning_architecture}, with a few changes.
While the embedding step is the same as both DINO and SimDINO in
\Cref{sub:contrastive_learning_architecture}, the feature extraction
step is the same as SimDINO in
\Cref{sub:contrastive_learning_architecture} as it just extracts the
feature corresponding to the class token, and the classification head
is described in \Cref{sub:image_classification_objective}, the
backbone architecture is different. Each layer takes the form
\begin{align}\label{eq:CARTE updates}
\vZ_{\theta}^{\ell + 1/2}(\vX)
&= \vZ_{\theta}^{\ell}(\vX) + \MSSA_{\theta}^{\ell}(\LN_{\theta}^{1,
\ell}(\vZ_{\theta}^{\ell}(\vX))), \\
\vZ_{\theta}^{\ell + 1}(\vX)
&= \ISTA_{\theta}^{\ell}(\LN_{\theta}^{2, \ell}(\vZ_{\theta}^{\ell +
1/2}(\vX))),
\end{align}
where the \(\MSSA_{\theta}^{\ell}\) and \(\ISTA_{\theta}^{\ell}\)
blocks are as described in \Cref{ch:deep-networks}, namely:
\begin{itemize}
\item The \(\MSSA\) operator is multi-head-subspace-self-attention,
    defined as follows:
    \begin{equation}
        \MSSA_{\theta}^{\ell}(\vZ) \doteq
        \vU_{\out}^{\ell}\mat{\SA([\vU^{1, \ell}]^{\top}\vZ, [\vU^{1,
            \ell}]^{\top}\vZ, [\vU^{1, \ell}]^{\top}\vZ)\\ \vdots
            \\ \SA([\vU^{K, \ell}]^{\top}\vZ, [\vU^{K, \ell}]^{\top}\vZ,
        [\vU^{1, \ell}]^{\top}\vZ)} + \vb_{\out}^{\ell}\vone_{n}^{\top}
    \end{equation}
    where \(\vU^{k, \ell} \in \R^{d \times p}\), \(\vU_{\out}^{\ell}
    \in \R^{d \times Kp}\), and \(\vb_{\out}^{\ell} \in \R^{d}\) are
    trainable parameters belonging to the parameter set \(\theta\),
    and (recall) the self-attention operator \(\SA\) is defined in
    \eqref{eq:self_attention}.
\item The \(\ISTA\) operator is the
    iterative-shrinkage-thresholding-algorithm operator, defined as follows:
    \begin{equation}
        \ISTA_{\theta}^{\ell}(\vZ) \doteq \ReLU(\vZ - \beta
            (\vD^{\ell})^{\top}(\vD^{\ell}\vZ - \vZ) + \beta\lambda
        \vone_{d}\vone_{n}^{\top}),
    \end{equation}
    so named because the map \(\vX \mapsto \ReLU(\vX - \beta
            \vD^{\top}(\vD\vX - \vZ) + \beta  \lambda
    \vone_{d}\vone_{n}^{\top})\) is one step of the well-established
    ISTA algorithm to find an element-wise non-negative sparse
    representation for \(\vZ\) with respect to the complete
    dictionary \(\vD\) (cf \Cref{sec:dictionary_learning}).
\end{itemize}

We call this architecture CRATE, and a layer of the backbone is
depicted in \Cref{fig:crate_backbone}. CRATE models, on top of being
interpretable, are generally also highly performant and parameter-efficient.

\subsection{Optimization} \label{sub:image_classification_optimization}

We train our classifier using a simple end-to-end stochastic
optimization procedure, where we subsample data and views, compute
the average loss and its gradient over these samples, and use an
optimization algorithm to change the parameters. At each timestep \(k\), we:
\begin{itemize}
\item Subsample \(B\) different labeled samples \(\{(\vX_{b}^{(k)},
    \vy_{b}^{(k)})\}_{b = 1}^{B} \subseteq \cI \times \{0, 1\}^{N_{\cls}}\).
\item For each labeled sample \((\vX_{b}^{(k)}, \vy_{b}^{(k)})\),
    compute the central crop view \(v_{b, \cc}^{(k)}\) and apply it
    to \(\vX_{b}^{(k)}\) to get \(\vX_{b, \cc}^{(k)} \doteq v_{b,
    \cc}^{(k)}(\vX_{b}^{(k)})\).
\item Compute the predictions \(\vp_{\theta}(\vX_{b, \cc}^{(k)})
        \doteq (h_{\theta} \circ f_{\theta}^{\ext} \circ
    f_{\theta})(\vX_{b, \cc}^{(k)})\).
\item Form the surrogate stochastic loss
    \begin{equation}
        \hat{\cL}_{\CE}^{(k)}(\theta) \doteq \frac{1}{B}\sum_{b =
        1}^{B}\CE(\vy_{b}^{(k)}, \vp_{\theta}(\vX_{b, \cc}^{(k)})).
    \end{equation}
\item Compute one step of an optimization algorithm on \(\theta\),
    giving the following iteration:
    \begin{equation}
        \theta^{(k + 1)} \doteq
        \textsc{OptUpdate}^{(k)}(\theta^{(k)};
        \nabla_{\theta}\hat{\cL}_{\CE}^{(k)}).
    \end{equation}
\end{itemize}

\subsection{Evaluation Methodology} \label{sub:image_classification_evals}

We use the same evaluation procedure as
\Cref{sub:contrastive_learning_evals}. To summarize, for all
evaluations (as well as training) we use a center crop view
\(v_{\cc}\) which reshapes the input image and takes a large central
crop of size \((C, S_{\cc}, S_{\cc})\) where \(C\) is the number of
channels in the input image. We can then do linear probing, attention
map visualization, and detection/segmentation benchmarks, given the
output of this view.

\subsection{Experimental Setup and
Results}\label{sub:image_classification_experiments}

Since CRATE is directly based on the transformer, we compare the
optimal settings for ViT as given by
\cite{dosovitskiy2020image,touvron2020training} with the same
settings applied to CRATE for a fair comparison.

\paragraph{Model Architecture.} The center crop resizes the whole
image so that the shorter edge is of size \(256\) (i.e., \(S_{\rsz} =
256\)) before taking a center crop of size \(224 \times 224\) (i.e.,
\(S_{\cc} = 224\)), both in evaluation and training. We take patch
size \(16\) (i.e., \(P_{H} = P_{W} = 16\)). We use the tiny, small,
base, and large models of the ViT \cite{dosovitskiy2020image}
architecture as the embedding and backbone, swapping out the MHSA and
MLP components for MSSA and ISTA, respectively, using the same number
of heads and head dimension in the case of MSSA, and therefore
reducing the number of training parameters drastically. For CRATE, we
set \((\beta, \lambda) = (1, 0.1)\).

\paragraph{Datasets and Optimization.} For pre-training, we use the
ImageNet-1K dataset. We use the LION optimizer
\citep{chen2024symbolic} to pre-train both our ViT replication as
well as CRATE. We set the base learning rate as \(2.4 \times
10^{-4}\), the weight decay as \(0.5\), and batch size as \(B =
2048\). Our learning rate schedule increases the learning rate
linearly to the base learning rate over the first \(5\) epochs, and
decreases to \(0\) using a cosine schedule over the next \(145\)
epochs (training all models for \(150\) epochs each). For
pre-training, we apply a usual regime of data augmentations (flips,
Gaussian blurs, solarization, etc.) to the image data, and also add
small noise to the labels (this is called \textit{label smoothing}
\citep{muller2019does}).

For linear probing, we use several evaluation datasets such as
CIFAR10, Oxford-Flowers, and Oxford-IIT-Pets. We use the AdamW
optimizer to train the linear probe, using learning rate \(5 \times
10^{-5}\), weight decay \(0.01\), and batch size \(B = 256\). We also
apply the aforementioned data augmentations to the image data.

\begin{table}
\caption{\small \textbf{Linear probing classification accuracy of
    CRATE and ViT} on various datasets with different model sizes when
    the backbone is pre-trained for classification on ImageNet-1K. We
    observe that given the same model configuration, CRATE has comparable
    classification performance with a simpler, more principled, and more
parameter-efficient design.}
\centering
\resizebox{1\columnwidth}{!}
{
    \begin{tabular}{@{}lcccc|cc@{}}
        \toprule
        \textbf{Model} & CRATE-T  &  CRATE-S & CRATE-B & CRATE-L & {
        \color{gray} ViT-T} &  { \color{gray}ViT-S } \\
        \midrule
        \midrule
        \# parameters & 6.09M & 13.12M & 22.80M & 77.64M & {
        \color{gray} 5.72M} & { \color{gray} 22.05M} \\
        \midrule
        ImageNet-1K & 66.7 & 69.2 & 70.8 & 71.3 & { \color{gray}
        71.5} & { \color{gray} 72.4} \\
        ImageNet-1K ReaL & 74.0 & 76.0 & 76.5 & 77.4 & { \color{gray}
        78.3 } & { \color{gray} 78.4} \\
        CIFAR10 & 95.5 & 96.0 & 96.8 & 97.2 & { \color{gray} 96.6} &
        { \color{gray} 97.2} \\
        CIFAR100 & 78.9 & 81.0 & 82.7 & 83.6 & { \color{gray} 81.8} &
        { \color{gray} 83.2}\\
        Oxford Flowers-102 & 84.6 & 87.1 & 88.7 & 88.3 & {
        \color{gray} 85.1} & { \color{gray} 88.5}\\
        Oxford-IIIT-Pets & 81.4 & 84.9 & 85.3 & 87.4 & { \color{gray}
        88.5} & { \color{gray} 88.6} \\
        \bottomrule
    \end{tabular}
}
\label{tab:crate_classification_linear_probing}
\end{table}

\begin{figure}[t!]
\centering
\includegraphics[width=\textwidth]{\toplevelprefix/chapters/chapter8/figs/crate_semantic_heads.pdf}
\caption{\small\textbf{Interpretable saliency maps in CRATE} with
    patch size \(8\). When given images with similar properties (perhaps
    but not necessarily from the same class), the saliency maps
    corresponding to different attention heads in the last layer each
    highlight a specific property. One can observe that the average
    saliency map (not included) then highlights all relevant objects in
    the image, showing that it uses all fine-grained details of the input
    image for classification. This is the \textit{first} machine learning
    system to do this, to the authors' knowledge, much less automatically
without training on any segmentation data.}
\label{fig:crate_semantic_heads}
\end{figure}

\begin{table}
\caption{\small \textbf{Object detection and fine-grained
    segmentation via MaskCut on COCO {val2017}~\citep{lin2014microsoft}}.
    Here all models are trained with patch size \(8\) instead of \(16\).
    CRATE conclusively performs better than the ViT at detection and
segmentation metrics when both are trained using supervised classification.}
\centering
\begin{tabular}{@{}lcccccccc@{}}
    \toprule
    &  \multicolumn{3}{c}{Detection (\(\uparrow\))} &
    \multicolumn{3}{c}{Segmentation (\(\uparrow\))} \\
    Model & AP$_{50}$ & AP$_{75}$ & AP & AP$_{50}$ & AP$_{75}$ & AP  \\
    \midrule
    CRATE-S/8 & \textbf{2.9} & \textbf{1.0} & 1.1 & 1.8 & \textbf{0.7} & 0.8 \\
    CRATE-B/8 & \textbf{2.9} & \textbf{1.0} & \textbf{1.3} &
    \textbf{2.2} & \textbf{0.7} & \textbf{1.0} \\
    ViT-S/8 & 0.1& 0.1 & 0.0 & 0.0 & 0.0 & 0.0 \\
    ViT-B/8 & 0.8 & 0.2 & 0.4 & 0.7 & 0.5 & 0.4 \\
    \bottomrule
\end{tabular}
\label{tab:crate_detection_segmentation}
\end{table}

\paragraph{Experiment Results.}

\Cref{tab:crate_classification_linear_probing} demonstrates that
CRATE models achieve parity or improvement compared to the popular
Vision Transformer (ViT) architecture at similar parameter counts, at
least in terms of the linear separability of their features with
respect to different classes. In terms of attention map fidelity,
\Cref{fig:crate_semantic_heads} demonstrates a truly extraordinary
result: without needing to train on any segmentation or object
detection data, \textit{not only} do the saliency maps effectively
capture all relevant parts of the input image, the saliency maps
\textit{self-organize} to each correspond to a discrete set of
concepts, even across samples and classes! This is the first system
to do this, to the authors' knowledge, and it can do this without
using any extra data except for the image classification data.
\Cref{tab:crate_detection_segmentation} confirms these qualitative
insights quantitatively, showing significant improvement over ViTs
trained in the same supervised classification setup.

\section{Image Representation via Masked
Autoencoding}\label{sec:image_completion}

In this section, we show how to apply the CRATE architecture to the
\textit{image completion} problem, also known as \textit{masked
autoencoding} (MAE), which can be viewed as a generalization of the
low-rank matrix completion problem discussed in \Cref{ch:classic-models}.
Masked autoencoding, since its introduction in the deep learning
context by \cite{he2022masked}, has been a staple and simple
self-supervised representation learning method, which aims to endow
each patch feature within \(\vZ_{\theta}\) with aggregate information
as well as information about its neighbors, such that both the patch
feature and aggregate features are rich sources of information for
the whole sample.

The data we use for this task is the same as the image datasets
discussed in \Cref{sub:contrastive_learning_data}. As usual, we still
apply data augmentations to each sample in each new batch.

\subsection{The MAE Task and Objective}\label{sub:image_completion_objective}

As the name suggests, masked autoencoding involves a view \(v_{m}\)
which, given an input, performs a random resized crop (cf
\Cref{sub:contrastive_learning_objective}) to turn the input image
into a square image of size \((C, S_{\mask}, S_{\mask})\), then
\textit{masks} (i.e., sets to zero) a fixed percentage \(p_{\mask}
\in [0, 1]\) of pixels in the input. For efficiency
reasons\footnote{The original implementation of MAE by
\cite{he2022masked} embeds the whole image, \textit{removes} the
tokens that would be masked, feeds the resulting token set through
the encoder, adds back learned placeholder tokens in the masked spots
and adds back the appropriate positional encoding, and feeds the
resulting token set through the decoder to get the autoencoding
prediction. This is more efficient since the encoder has fewer tokens
to go through, but conceptually is the same as the method discussed
in the text, and the resulting models' performance in the masked
autoencoding task and downstream evaluations is very similar.}, the
masking is done patch-wise, i.e.,  after embedding the whole image,
\(p_{\mask}\) percentage of \textit{patches} are set to zero. The
goal of MAE is to train an encoder \(f_{\theta} \colon \cI \to
(\R^{d})^{*}\) and a decoder \(g_{\eta} \colon (\R^{d})^{*} \to \cI\)
that can reconstruct an input from its masking, i.e., writing
\(\hat{\vX}_{\theta, \eta} \doteq g_{\eta} \circ f_{\theta}\), we have
\begin{equation}
\min_{\theta, \eta}\bc{\cL_{\mathrm{MAE}}(\theta, \eta) \doteq
\Ex\norm{\hat{\vX}_{\theta, \eta}(\vX_{m}) - \vX}_{F}^{2}}
\end{equation}
Essentially this means that the features \(\vZ_{\theta}(\vX_{m})\) of
the view \(\vX_{m} \doteq v_{m}(\vX)\) must contain information about
the \textit{masked patches} as well as the existing patches. From the
perspective of the compression-based white-box models in
\Cref{ch:deep-networks}, if a white-box autoencoder  \((f_{\theta},
g_{\eta})\) succeeds at this task, it means that the learned
subspaces and dictionaries perform a \textit{redundant} encoding of
the data such that it can reconstruct missing parts of the data from
encoded other parts of the data. This means that information about
each patch is stored in other patches. Therefore, each patch feature
should contain both information about the patch and information about
the statistics of the whole image. Thus, again, we expect that the
representations should contain both local and global semantically
relevant information, and therefore representations of different
patches with similar local and global information should be related
(i.e., on the same subspace or encoded together by a dictionary).

\subsection{CRATE-Based MAE
Architecture}\label{sub:image_completion_architecture}

\begin{figure}[t!]
\centering
\includegraphics[width=\textwidth]{\toplevelprefix/chapters/chapter8/figs/crate_ae_backbone.pdf}
\caption{\small\textbf{One layer of the encoder and decoder in a
    CRATE autoencoder backbone.} The encoder and decoder layers both feed
    their inputs through multi-head subspace self-attention and a
    dictionary learning or dictionary encoding step. Note that the
    encoder and decoder layers are symmetrically designed; the conceptual
    goal of each decoder layer is to invert an encoder layer, so this
    symmetry is very much by design (see e.g.,
\Cref{ch:consistent-representations}).}
\end{figure}

We use a CRATE encoder and decoder, depicted in
\Cref{fig:transformer_backbone}, though of course it is possible to
use a regular transformer encoder and decoder. Details follow.

\paragraph{The Encoder.} The encoder is the same as the CRATE encoder
in \Cref{sub:image_classification_architecture}, with the caveat that
there is no feature extractor \(f_{\theta}^{\ext}\). However, both
the embedding \(f_{\theta}^{\emb}\) and the backbone
\(f_{\theta}^{\backbone}\) are the same.

\paragraph{The Decoder Backbone.} The decoder backbone is the CRATE
decoder described in \Cref{ch:consistent-representations}. For completeness, we
describe it now. Given a feature sequence \(\vZ_{\theta}(\vX) \doteq
f_{\theta}(\vX) \in (\R^{d})^{*}\), we can process it using the
decoder backbone \(g_{\eta}^{\backbone}\) as follows. The function
\(g_{\eta}^{\backbone}\) is composed of \(L\) layers \(g_{\eta}^{\ell}\), i.e.,
\begin{equation}
g_{\eta}^{\backbone} = g_{\eta}^{L} \circ \cdots \circ g_{\eta}^{1}.
\end{equation}
The layer \(g_{\eta}^{\ell}\) has the following implementation.
First, define \(\tilde{\vZ}_{\theta, \eta}^{1}(\vX) \doteq
\vZ_{\theta}(\vX)\). Then, we obtain
\begin{align}
\tilde{\vZ}_{\theta, \eta}^{\ell + 1/2}(\vX)
&= [\tilde{\vD}^{\ell}]^{\top}\LN_{\eta}^{1,
\ell}(\tilde{\vZ}_{\theta, \eta}^{\ell}(\vX)) \\
\tilde{\vZ}_{\theta, \eta}^{\ell + 1}(\vX)
&= \tilde{\vZ}_{\theta, \eta}^{\ell + 1/2}(\vX) -
\MSSA_{\eta}^{\ell}(\LN_{\eta}^{2, \ell}(\tilde{\vZ}_{\theta,
\eta}^{\ell + 1/2}))
\end{align}
and \(g_{\eta}^{\ell}\) is defined such that
\(g_{\eta}^{\ell}(\tilde{\vZ}_{\theta, \eta}^{\ell}) \doteq
\tilde{\vZ}_{\theta, \eta}^{\ell + 1}(\vX)\). Here, the relevant
concept is that \(g_{\eta}^{\ell}\) should learn an approximate
inverse of \(f_{\theta}^{L + 1 - \ell}\), as discretizations of a
forward- and reverse-time diffusion process, respectively. In
particular, \(\tilde{\vD}^{\ell}\) should approximate \(\vD^{L + 1 -
\ell}\), and similarly, the \(\MSSA_{\eta}^{\ell}\) parameters should
be similar to the parameters of \(\MSSA_{\theta}^{L + 1 - \ell}\).
The output is \(\tilde{\vZ}_{\theta, \eta} \doteq
\tilde{\vZ}_{\theta, \eta}^{L + 1}\).

\paragraph{The Un-embedding Module.} To transform
\(\tilde{\vZ}_{\theta, \eta}(\vX)\) back into an estimate for
\(\vX\), we need to undo the effect of the embedding module
\(f_{\theta}^{\emb}\) using the unembedding module
\(g_{\eta}^{\unemb}\). As such, harkening back to the functional form
of the embedding module in \eqref{eq:definition_of_embedding_module}, i.e.,
\begin{equation}
f_{\theta}^{\emb}(\vX) \doteq \mat{\vz_{\cls}^{1},
\vW^{\emb}f^{\patch}(\vX) + \vE^{\pos}}
\end{equation}
it implies that our inverse operation \(g_{\eta}^{\unemb}\) looks
like the following:
\begin{align}
g_{\eta}^{\unemb}(\tilde{\vZ}) &\doteq
g_{\eta}^{\unemb}(\mat{\tilde{\vz}^{1}, \dots, \tilde{\vz}^{n}}) \\
&=
g^{\unpatch}(\vW^{\unemb}([\tilde{\vz}^{2}, \dots, \tilde{\vz}^{n}] -
\tilde{\vE}^{\pos})),
\end{align}
where \(g^{\unpatch}\) does the inverse operation of the unrolling
and flattening operation that \(f^{\patch}\) does.\footnote{Again,
the ``inverse positional encoding'' \(\tilde{\vE}^{\pos}\) is learned
for a large input, and for smaller inputs may be interpolated. It is
even possible to directly set \(\tilde{\vE}^{\pos}\) equal to the
positional encoding \(\vE^{\pos}\) and use the same interpolated
positional encodings for each input in both the encoder and decoder.}

This architecture is a white-box autoencoder \((f_{\theta},
g_{\eta})\) where (recall) \(f_{\theta} = f_{\theta}^{\backbone}
\circ f_{\theta}^{\emb}\) and \(g_{\eta} = g_{\eta}^{\unemb} \circ
g_{\eta}^{\backbone}\). In particular, we can use it to compute an
estimate for a masked view \(\hat{\vX}_{\theta, \eta}(\vX_{m}) =
(g_{\eta} \circ f_{\eta})(\vX_{m})\) which should approximately equal
\(\vX\) itself.

\subsection{Optimization}\label{sub:image_completion_optimization}

As in \Cref{sub:image_classification_optimization}, we use a simple
optimization setup: we sample images and masks, compute the loss on
those samples and the gradients of this loss, and update the
parameters using a generic optimization algorithm and the
aforementioned gradients. For each timestep \(k\), we:
\begin{itemize}
\item Subsample \(B\) different samples \(\{\vX_{b}^{(k)}\}_{b =
    1}^{B} \subseteq \cI\).
\item For each sample \(\vX_{b}^{(k)}\), compute a different
    randomized resized crop and mask \(v_{b, m}^{(k)}\) and apply it
    to \(\vX_{b}^{(k)}\) to get \(\vX_{b, m}^{(k)} \doteq v_{b,
    m}^{t}(\X_{b}^{(k)})\).
\item Compute the estimated autoencoding \(\hat{\vX}_{\theta,
        \eta}(\vX_{b, r}^{(k)}) \doteq (g_{\eta} \circ
    f_{\theta})(\vX_{b, r}^{(k)})\).
\item Form the surrogate stochastic loss
    \begin{equation}
        \hat{\cL}_{\mathrm{MAE}}^{(k)}(\theta, \eta) \doteq
        \frac{1}{B}\sum_{b = 1}^{B}\norm{\hat{\vX}_{\theta,
        \eta}(\vX_{b, r}^{(k)}) - \vX_{b}^{(k)}}_{F}^{2}.
    \end{equation}
\item Compute one step of an optimization algorithm on \((\theta,
    \eta)\), giving the following iteration:
    \begin{equation}
        (\theta^{(k + 1)}, \eta^{(k + 1)}) \doteq
        \textsc{OptUpdate}^{(k)}(\theta^{(k)}, \eta^{(k)};
        \nabla_{(\theta, \eta)}\hat{\cL}_{\mathrm{MAE}}^{(k)}).
    \end{equation}
\end{itemize}

\subsection{Evaluation} \label{sub:image_completion_optimization_1}

This is the first autoencoder network we discuss in this chapter. We
use the same center crop view \(v_{\cc}\) as in
\Cref{sub:contrastive_learning_evals,sub:image_classification_evals},
resizing the final image to a square with side length \(S_{\cc} =
S_{\mask}\) pixels to match the shapes of the input images seen during training.

In addition to evaluating the masked autoencoding loss itself, it is
also possible to evaluate the features \(\vZ_{\theta}(\vX_{\cc})\) of
the view \(\vX_{\cc} \doteq v_{\cc}(\vX)\) of the data \(\vX\)
directly. For attention map fidelity evaluation, obtaining
\(\vZ_{\theta}(\vX_{\cc})\) is sufficient, but for linear probing we
need to extract a summarized or aggregate feature from
\(\vZ_{\theta}\). To do this, we can use a (parameter-free) feature
extraction map that returns the feature corresponding to the class token, i.e.,
\begin{equation}
f_{\theta}^{\ext}(\vZ) \doteq f_{\theta}^{\ext}([\vz^{1}, \dots,
\vz^{n}]) = \vz^{1},
\end{equation}
as in (for example)
\Cref{sub:image_classification_objective,sub:image_classification_architecture}.
With this, we have a way to obtain aggregate features
\(\vz_{\theta}(\vX_{\cc}) \doteq (f_{\theta}^{\ext} \circ
f_{\theta})(\vX_{\cc})\), at which point we can perform linear
probing, segmentation evaluations, and so on.

\subsection{Experiments}\label{sub:image_completion_experiments}

Since CRATE-MAE is directly based on ViT-MAE, we compare the optimal
settings for ViT-MAE as given by \citep{he2022masked} with the same
settings applied to CRATE-MAE for a fair comparison.

\paragraph{Model Architecture.} During training, the masked crop
\(v_{m}\) resizes the whole image so that the shorter edge is of size
\(256\) (i.e., \(S_{\rsz} = 256\)) before taking a random crop of
size \(224 \times 224\) (i.e., \(S_{\mathrm{mask}} = 224\)), and
masking \(p_{\mathrm{mask}} = \frac{3}{4}\) of the patches. We take
patch size \(16\) (i.e., \(P_{H} = P_{W} = 16\)). We use the small
and base variants of the ViT-MAE architecture as the embedding and
backbone for both the encoder and decoder, swapping out the MHSA and
MLP components for MSSA, ISTA, and linear layers, respectively. We
use the same number of heads and head dimension in the case of MSSA.
However, the original ViT-MAE uses an encoder which uses nearly all
the total layers and a decoder which only uses a few layers; we
allocate half the total number of layers (which stays the same from
ViT-MAE to CRATE-MAE) to our encoder and decoder, as suggested by the
conceptual and theoretical framework in
\Cref{ch:consistent-representations}. For
CRATE-MAE we set \((\beta, \lambda) = (1, 0.1)\).

\paragraph{Datasets and Optimization.} For pre-training, we use the
ImageNet-1K dataset. We use the AdamW optimizer to pre-train both our
ViT-MAE replication as well as CRATE-MAE. We set the base learning
rate as \(3 \times 10^{-5}\), the weight decay as \(0.1\), and batch
size as \(B = 4096\). Our learning rate schedule increases the
learning rate linearly to the base learning rate over the first
\(40\) epochs, and decreases to \(0\) using a cosine schedule over
the next \(760\) epochs (training all models for \(800\) epochs
each). For pre-training, we apply the usual regime of data
augmentations (flips, Gaussian blurs, solarization, etc) to the image data.

For linear probing, we use several evaluation datasets such as
CIFAR10, CIFAR100, Oxford-Flowers, and Oxford-IIT-Pets. For linear
probing, we precompute the features of all samples in the target
dataset and apply a fast linear regression solver, e.g., from a
standard package such as Scikit-Learn.

\paragraph{Experiment Results.} \Cref{tab:crate_mae_linear_probing}
demonstrates that CRATE-MAE models achieve, roughly speaking, parity
with the popular ViT-MAE architecture at similar parameter counts,
and also that the feature learning performance (as measured by
performance on downstream classification tasks) increases with scale.
Meanwhile, \Cref{fig:crate_mae_semantic_heads} demonstrates that the
encoder saliency maps (and therefore the fine-grained features
learned by the encoder) indeed isolate and highlight the key parts of
the input image.

\begin{table}
\caption{\small\textbf{Linear probing classification accuracy of
    CRATE-MAE and ViT-MAE} on various datasets with different model sizes
    when the backbone is pre-trained for masked autoencoding on
    ImageNet-1K. Given the same parameter count, CRATE-MAE achieves
    roughly similar performance while simultaneously enjoying a simpler
and more principled architecture design.}
\centering
\resizebox{1\columnwidth}{!}
{
    \begin{tabular}{@{}lcc|cc@{}}
        \toprule
        \textbf{Model} & CRATE-MAE-S(mall) & CRATE-MAE-B(ase) &
        {\color{gray} ViT-MAE-S} & {\color{gray} ViT-MAE-B} \\
        \midrule
        \midrule
        \# parameters & 25.4M & 44.6M & 47.6M & {\color{gray}143.8M} \\
        \midrule
        CIFAR10 & 79.4 & 80.9 & {\color{gray} 79.9} & {\color{gray} 87.9} \\
        CIFAR100 & 56.6 & 60.1 & {\color{gray} 62.3} & {\color{gray} 68.0} \\
        Oxford Flowers-102 & 57.7 & 61.8 & {\color{gray} 66.8} &
        {\color{gray} 66.4} \\
        Oxford-IIIT-Pets & 40.6 & 46.2 & {\color{gray} 51.8} &
        {\color{gray} 80.1} \\
        \bottomrule
    \end{tabular}
}

\label{tab:crate_mae_linear_probing}
\end{table}

\begin{figure}[t!]
\centering
\includegraphics[width=\textwidth]{\toplevelprefix/chapters/chapter8/figs/crate_mae_semantic_heads.pdf}
\caption{\small\textbf{Saliency maps of CRATE-MAE.} Each pair of
    images consists of the original image (left) and a selected saliency
    map (right) corresponding to an attention head in the last layer. As
    is usual for CRATE models, but unusual for general transformer-like
models, the saliency maps correspond to the objects in the input image.}
\label{fig:crate_mae_semantic_heads}
\end{figure}

\section{Image Generation via Auto-Encoding and Sampling}
\label{sec:auto-encode-img-gen}

In the previous section, we have introduced masked autoencoding (MAE)
as one approach to learn an autoencoding of natural images $\x \sim p(\x)$:
\begin{equation}
\x\xrightarrow{\hspace{1mm} \mathcal{E} = f \hspace{1mm}} \z
\xrightarrow{\hspace{1mm} \mathcal{D} = g \hspace{1mm}} \hat{\x}.
\end{equation}
where $\z = f(\x)$ is certain feature space. Such an autoencoding is
designed to learn correlation, or co-occurrence, among  patches of a
natural image.

However, MAE does not  enforce the (distribution of) learned features
$\z$ to have any explicit structures or properties, which can be very
important if we want to use the so-learned features for certain
subsequent tasks such as recognition or controlled generation. In
addition, the simple MMSE loss used for image completion does not
necessarily respect the possible nonlinear structures in the image
manifold and typically generate an expected value, hence usually
blurred version, for the missing image patches, as we have shown in
Figure \ref{fig:mae_autoencoding-small} in the previous Chapter
\ref{ch:conditional-inference}.

Hence, in this section, we will introduce some other methods to learn
representations for (the distribution of) natural images via
auto-encoding that aim, to large extent, to remedy the above two
issues. More precisely, we want the distribution of the encoded
features to have certain desired structures and images decoded from
such features are much more faithful to the distribution of natural
images -- that is, of much better visual quality.

\subsection{Variational Auto-Encoding of Imagery Data}
% As we have introduced in Section \ref{sec:vae} of Chapter
% \ref{ch:consistent-representations}, variational auto-encoding is a
% method that aims to promote the distribution of learned features
% $\z$ and the conditional distributions between $\x$ and $\z$ to be
% more Gaussian like.\footnote{One may view such prior is to promote
% the learned the features to be diverse, in a similar spirit to
% maximizing the coding rate of all features.}

% Here we introduce a popular implementation of a VAE for imagery data...

% \subsection{Variational Auto-Encoding of Imagery Data}
\label{sec:VAE-implement}

%refer back to Chapter 6 VAE here.
As we have introduced in Section~\ref{sec:vae} of
Chapter~\ref{ch:consistent-representations}, variational
auto-encoding is a method that aims to promote the distribution of
learned features $\z$ and the conditional distributions between $\x$
and $\z$ to be more Gaussian like.\footnote{One may view such prior
is to promote the learned features to be diverse, in a similar spirit
to maximizing the coding rate of all features.}

Here we introduce a popular implementation of a VAE for imagery data,
following the perceptual compression model leveraged in Stable
Diffusion~\citep{rombach2022high}. The key idea is to train an
autoencoder that maps high-resolution images into a low-dimensional
latent space, while imposing a KL-regularization on the latent
distribution to encourage it to be close to a standard Gaussian
$\mathcal{N}(\boldsymbol{0}, \boldsymbol{I})$. This latent space can
subsequently serve as an efficient arena for downstream generative
models such as latent diffusion. Both the encoder and decoder are
jointly trained end-to-end. We outline the details as follows.

\paragraph{The Encoder ($\mathcal{E}$).}
The encoder $\mathcal{E}$ maps an input RGB image $\x \in
\mathbb{R}^{H \times W \times 3}$ into a low-dimensional latent
space. Its architecture, adopted from~\citep{esser2021taming}, is a
fully convolutional network. It begins with an input convolution that
projects the 3-channel RGB input into a feature map with a base
channel count of $128$. This is followed by four stages, each
consisting of two ResNet blocks~\citep{he2016deep}. The channel
dimension is scaled by multipliers $(1, 2, 4, 4)$ across stages,
yielding feature maps with $128, 256, 512, 512$ channels,
respectively. Spatial downsampling by a factor of $2\times$ is
performed between consecutive stages via strided convolutions,
resulting in three downsamplings and an overall downsampling factor
of $f = 2^3 = 8$. After the four stages, a middle block consisting of
two additional ResNet blocks further refines the features. Finally,
an output convolution projects the features to $2c$ channels, which
are split into the two quantities that parameterize the latent
distribution (described below).

Crucially, unlike a standard autoencoder, the encoder does not output
a single deterministic latent vector. Instead, it parameterizes a
diagonal Gaussian posterior $q(\z | \x) =
\mathcal{N}(\boldsymbol{\mu}, \boldsymbol{\sigma}^2 \boldsymbol{I})$
by outputting two quantities:
\begin{itemize}
\item A mean $\boldsymbol{\mu} = \mathcal{E}_\mu(\x) \in
    \mathbb{R}^{h \times w \times c}$, and
\item A log-variance $\log \boldsymbol{\sigma}^2 =
    \mathcal{E}_\sigma(\x) \in \mathbb{R}^{h \times w \times c}$.
\end{itemize}
The latent representation is then obtained via a so called
\emph{reparameterization trick}~\citep{kingma2013auto}:
\begin{equation}
\z = \boldsymbol{\mu} + \boldsymbol{\sigma} \hada
\boldsymbol{\epsilon}, \quad \boldsymbol{\epsilon} \sim
\mathcal{N}(\boldsymbol{0}, \boldsymbol{I}),
\end{equation}
where $\hada$ denotes element-wise multiplication. This formulation
enables gradient-based optimization through the stochastic sampling step.

\paragraph{The Decoder ($\mathcal{D}$).}
The decoder $\mathcal{D}$ learns the reverse mapping from the latent
space back to the pixel space, producing the reconstructed image
$\tilde{\x} = \mathcal{D}(\z) \in \mathbb{R}^{H \times W \times 3}$.
Its architecture is approximately symmetric to the encoder: an input
convolution first projects the $c$-channel latent into the feature
space, followed by a middle block of two ResNet blocks. Then, four
stages---each with two ResNet blocks---progressively upsample the
spatial resolution, with the channel dimensions mirroring those of
the encoder in reverse order ($512, 512, 256, 128$). A final output
convolution maps the features back to 3 RGB channels. Both the
encoder and decoder architectures follow the convolutional design
introduced in~\citep{esser2021taming}.

\subsubsection{Data}
We assume access to a dataset $\mathcal{D} = \{\x_i\}_{i=1}^N$ of
samples from the original data distribution. Each sample is an image
of the form $\x \in \mathbb{R}^{H \times W \times 3}$, where $H$ and
$W$ are the height and width, respectively.
In~\citep{rombach2022high}, the autoencoder is trained on the
OpenImages dataset with images cropped to $256 \times 256$ resolution.

Given the downsampling factor $f = 8$ and latent channel count $c =
4$, the encoder maps each input image to a latent tensor $\z \in
\mathbb{R}^{H/f \times W/f \times c}$. For the default training
resolution, this yields a latent of size $32 \times 32 \times 4$. At
inference time, the autoencoder generalizes to other resolutions; for
example, a $512 \times 512 \times 3$ input is encoded into a $64
\times 64 \times 4$ latent.

\subsubsection{Objective and Evaluation Metrics}
To train the VAE, our goal is to jointly optimize the encoder
$\mathcal{E}$ and decoder $\mathcal{D}$ such that the reconstructed
data $\tilde{\x} = \mathcal{D}(\mathcal{E}(\x))$ is as close as
possible to the original data $\x$, while regularizing the latent
distribution to be close to a standard Gaussian prior. An empirically
effective choice of the reconstruction loss is a weighted combination
of L1, LPIPS~\citep{zhang2018unreasonable}, and adversarial losses as
in GAN~\citep{goodfellow2014generative}. Additionally, a KL
divergence term is included to regularize the latent space. Formally,
the training objective can be expressed as:

\begin{align}
\label{eq:vae_loss}
\min_{\theta_{\mathcal{E}}, \theta_{\mathcal{D}}} \;
\mathcal{L}_{\mathrm{VAE}}(\theta_{\mathcal{E}}, \theta_{\mathcal{D}})
&= \mathbb{E}_{\x \sim \mathcal{D}} \Big[
    \|\tilde{\x} - \x\|_1 + \lambda_1 \, \mathrm{LPIPS}(\tilde{\x}, \x) \\
    &\qquad\quad + \lambda_2 \,
    \mathcal{L}_{\mathrm{GAN}}(\tilde{\x}, \x) + \; \lambda_3 \,
    D_{KL}\!\big(q(\z | \x) \,\big\|\, p(\z)\big)
\Big],
\end{align}
where $\theta_{\mathcal{E}}$ and $\theta_{\mathcal{D}}$ are the
parameters of the encoder and decoder, respectively, and $\lambda_1,
\lambda_2, \lambda_3$ are hyperparameters that balance the
contributions of each loss term. The first three terms enforce
sample-wise consistency in the same manner as the reconstruction
objectives in Section~\ref{sec:RAE-implement}. The fourth term---the
KL divergence between the encoder posterior $q(\z | \x) =
\mathcal{N}(\boldsymbol{\mu}, \boldsymbol{\sigma}^2 \boldsymbol{I})$
and the standard Gaussian prior $p(\z) = \mathcal{N}(\boldsymbol{0},
\boldsymbol{I})$---regularizes the latent space. For a diagonal
Gaussian posterior, this admits a closed-form expression:
\begin{equation}
D_{KL}\!\big(q(\z | \x) \,\big\|\, p(\z)\big) = \frac{1}{2}
\sum_{j=1}^{d} \left( \mu_j^2 + \sigma_j^2 - \log \sigma_j^2 - 1 \right),
\end{equation}
where $d = h \times w \times c$ is the total dimensionality of the latent space.

The implementation detail from~\citep{rombach2022high} is that the KL
weight $\lambda_3$ is set to a very small value ($\approx 10^{-6}$).
This is because the primary purpose of the KL term here is not to
enable unconditional sampling from the latent space (as in a
traditional VAE), but rather to prevent the latent space from
developing arbitrarily high variance, which would degrade the
signal-to-noise ratio for downstream models operating in this space.

For evaluation, we measure the reconstruction quality by computing
the FID score~\citep{heusel2017gans} on the reconstructed images from
the ImageNet-1k validation set, denoted as rFID. A lower rFID
indicates better reconstruction quality.

\subsection{Representation Auto-Encoding of Imagery Data}
\label{sec:RAE-implement}
% \yima{Conceptual connections and necessity... Practical
% implementations details... (Robin and Peter)}

% rae stuff
\definecolor{deemph}{gray}{0.6}
\newcommand{\gc}[1]{\textcolor{deemph}{#1}}
\newcommand{\tablestyle}[2]{\setlength{\tabcolsep}{#1}\renewcommand{\arraystretch}{#2}\centering\scriptsize}
% \newcommand{\round}[1]{\num[round-mode=places, round-precision=2]{#1}}

As discussed in \Cref{ch:lossy-compression}, learning structured
representations of the data is crucial for many downstream tasks. In
visual representation learning, this is often achieved by
self-supervised learning methods such as DINO
\cite{caron2021emerging} and SimDINO (studied in Section
\ref{sec:contrastive_learning}). In
\Cref{ch:consistent-representations}, we posit
that a good representation of data should be effectively decoded back
into the original data space. This requires a consistent
auto-encoding framework that can accurately reconstruct the original
data from its learned representation. It is then natural to ask: is
the representation learned by methods like DINO consistent and can it
be effectively decoded in an auto-encoding framework?

To find out, we can train a \textit{representation auto-encoder}
(RAE) that consists of an pretrained representation encoder (e.g.
DINO) and a trainable decoder, following the work:\\
\centerline{\url{https://rae-dit.github.io}.}\\
The encoder $f$ maps an input sample
$\x$ to its learned representation $\z = f(\x)$, while the decoder
$g$ maps the representation back to the original data space
$\x'=g(\z)$. As detailed in \Cref{ch:consistent-representations}, one simple
objective is to impose sample-wise consistency by training the
decoder to minimize the reconstruction loss between the original data
$\x$ and the reconstructed data $\x'$. We outline the details as follows.

\paragraph{The Encoder.} We use a pretrained representation encoder
$f \colon \cX \to \cZ$ that maps an input data sample $\x \in \cX$ to
its learned representation $\z = f(\x) \in \cZ$. The encoder is kept
frozen during training. Here, we assume a ViT-based encoder, as
introduced in \Cref{sub:contrastive_learning_architecture}.

\paragraph{The Decoder.} For architectural simplicity, we use a
trainable ViT-based decoder $g \colon \cZ \to \cX$ that maps the
learned representation $\z$ back to the original data space,
producing the reconstructed data $\x' = g(\z)$.

\subsubsection{Data}
We assume access to a dataset $\cD = \{\x_i\}_{i=1}^N$ of samples
from the original data distribution $\cX$. Each sample is an image of
the form $\x \in \R^{3 \times H \times W}$, where $H$ and $W$ are the
height and width, respectively. The encoder $f$ is set to a patch
size of $p_e$ and feature dimension of $d$. Consequently, for each
input image $\x$, we would obtain $N = HW/p_e^2$ tokens of dimension
$d$. The decoder $g$ with patch size $p_d$ then maps these visual
tokens back to reconstruct $\x' \in \R^{3 \times H\frac{p_d}{p_e}
\times W\frac{p_d}{p_e}}$ in the original pixel space. By default, we
set $p_d = p_e$ to ensure that the reconstructed image has the same
resolution as the original image.

\subsubsection{Objective and Evaluation Metrics} To train RAE, our goal
is to train a decoder $g$ such that the reconstructed data $\x' =
g(f(\x))$ is as close as possible to the original data $\x$. This can
be achieved by minimizing a sample-wise consistency loss over the
dataset $\cD$. An empirically effective choice of the loss function
is a weighted combination of L1, LPIPS \cite{zhang2018unreasonable}
\footnote{This is a perceptual loss based on computing distances
between features across different layers of a pretrained network. It
empirically captures human perceptual similarity better than
pixel-wise losses. Refer to the original paper for details.}, and
adversarial losses as in GAN \cite{goodfellow2014generative}.
Formally, the training objective can be expressed as:

\begin{equation}
\min_{\theta} \cL_{\mathrm{RAE}}(\theta) = \Ex_{\x \sim \cD} \left[
    \norm{\x' - \x}_1 + \lambda_1 \mathrm{LPIPS}(\x', \x) + \lambda_2
\cL_{\mathrm{GAN}}(\x', \x) \right],
\end{equation}
where $\theta$ are the parameters of the decoder $g$, and $\lambda_1,
\lambda_2$ are hyperparameters that balance the contributions of each loss term.

For evaluation, we measure the reconstruction quality by computing
the FID score \cite{heusel2017gans} on the reconstructed images from
the ImageNet-1k validation set, denoted as rFID. A lower rFID
indicates better reconstruction quality.

\subsubsection{Experimental Results}
We can compare RAEs trained on different pretrained encoders,
including DINOv2, MAE, and SigLIP2. As a baseline, we also consider a
vanilla variational autoencoder (VAE) trained directly on the image
pixels without any pretrained encoder. One popular choice is SD-VAE.
RAE decoders are trained on the ImageNet-1K training set and
evaluated on the ImageNet-1K validation set. The results are
summarized in \Cref{tab:rae_recon_results}. From
\Cref{tab:decoder_scaling_no_noise}, we can see that RAEs with
pretrained encoders outperform SD-VAE in reconstruction quality.
Furthermore, larger decoder sizes lead to improved reconstruction
quality, as shown in \Cref{tab:decoder_scaling_no_noise}, while
remaining more efficient than VAEs. \Cref{tab:encoder_scaling_dinov2}
shows that the reconstruction quality is stable across different
sizes of the DINO encoders. Finally, \Cref{tab:linear_probing}
demonstrates that RAEs also retain high-quality representations,
achieving much higher linear probing accuracy compared to VAEs.

\begin{table*}[t]
\centering
{\captionsetup[subfloat]{labelfont=scriptsize}
    \caption{
    \small\textbf{Reconstruction and linear probing results of RAEs.} }
    \begin{tabular}{@{}cc@{}}
        \subfloat[
            \scriptsize \textbf{Encoder choice.} All encoders outperform SD-VAE.
            \label{tab:recon_deterministic}
        ]
        {
            \begin{minipage}[t]{0.45\linewidth}{
                    \vspace{0pt}
                    \begin{center}
                        \tablestyle{4pt}{1.1}
                        \begin{tabular}{lc}
                            \toprule
                            Model & rFID \\
                            \midrule
                            \rowcolor{gray!20}
                            DINOv2-B & 0.49 \\
                            \rowcolor{white}
                            SigLIP2-B & 0.53 \\
                            MAE-B & 0.16 \\
                            \midrule
                            SD-VAE & \gc{0.62} \\
                            \bottomrule
                        \end{tabular}
                    \end{center}
                }
            \end{minipage}
        }
        &
        \subfloat[
            \scriptsize \textbf{Larger decoders} improve rFID while
            remaining much more efficient than VAEs.
            \label{tab:decoder_scaling_no_noise}
        ]
        {
            \begin{minipage}[t]{0.45\linewidth}{
                    \vspace{0pt}
                    \begin{center}
                        \tablestyle{4pt}{1.1}
                        \begin{tabular}{lcc}
                            \toprule
                            Decoder & rFID & GFLOPs \\
                            \midrule
                            ViT-B & 0.58 & 22.2 \\
                            ViT-L & 0.50 & 78.1 \\
                            \rowcolor{gray!20}
                            ViT-XL & 0.49 & 106.7 \\
                            \midrule
                            \rowcolor{white}
                            SD-VAE & \gc{0.62} & \gc{310.4} \\
                            \bottomrule
                        \end{tabular}
                \end{center}}
            \end{minipage}
        }
        \\[1em]
        \subfloat[
            \scriptsize \textbf{Encoder scaling.}  rFID is stable
            across RAE sizes.
            \label{tab:encoder_scaling_dinov2}
        ]{
            \begin{minipage}[t]{0.45\linewidth}{
                    \vspace{0pt}
                    \begin{center}
                        \tablestyle{4pt}{1.1}
                        \begin{tabular}{lc}
                            \toprule
                            Encoder & rFID \\
                            \midrule
                            DINOv2-S & 0.52 \\
                            \rowcolor{gray!20}
                            DINOv2-B & 0.49 \\
                            \rowcolor{white}
                            DINOv2-L & 0.52 \\
                            \bottomrule
                            \multicolumn{2}{c}{~} \\
                            \multicolumn{2}{c}{~} \\
                        \end{tabular}
                        \vspace{-0.1cm}
                    \end{center}
                }
            \end{minipage}
        }
        &
        \subfloat[
            \scriptsize \textbf{Representation quality.} RAEs have
            much higher linear probing accuracy than VAEs.
            \label{tab:linear_probing}
        ]
        {
            \begin{minipage}[t]{0.45\linewidth}{
                    \vspace{0pt}
                    \begin{center}
                        \tablestyle{4pt}{1.1}
                        \begin{tabular}{lc}
                            \toprule
                            Model & Top-1 Acc. \\
                            \midrule

                            \rowcolor{gray!20} DINOv2-B & 84.5 \\
                            \rowcolor{white}
                            SigLIP2-B & 79.1 \\
                            MAE-B & 68.0 \\
                            \midrule
                            SD-VAE & \gc{8.0} \\
                            \bottomrule
                        \end{tabular}
                    \end{center}
                }
            \end{minipage}
        }
    \end{tabular}
}
% \vspace{-4pt}
\label{tab:rae_recon_results}
\end{table*}
To complement the quantitative metrics in \Cref{tab:rae_recon_results},
\Cref{fig:rae_recon_qualitative} provides a visual comparison of
reconstructions. RAE reconstructions preserve fine-grained textures,
object boundaries, and semantic details as faithfully as the SD-VAE
baseline, consistent with its better rFID scores.

\begin{figure}[t!]
\centering
\includegraphics[width=\textwidth]{chapters/chapter8/figs/recon_grid.pdf}
\caption{\small\textbf{RAE Reconstruction examples.} From left to
    right: input image, RAE (DINOv2-B), RAE (Siglp-B), RAE (MAE-B),
SD-VAE. Zoom in for details.}
\label{fig:rae_recon_qualitative}
\end{figure}

% \yima{Besides the quantitative results, please have some
% illustrations to show qualitative improvements... }

\subsection{Sampling from Learned Representations via Denoising}

Once we have learned a consistent and structured representation of
the distribution of natural images, say via VAE or
RAE, we typically like to efficiently generate\textit{ samples} $\z
\in \cZ$ from the learned representation distribution so that we can
regenerate the natural images associated with these samples via the
learned decoder $\hat{\x} = g(\z) \in \cX$. This auto-encoding
process is illustrated by the top half of Figure
\ref{fig:mental-picture} in Chapter \ref{ch:intro}.

To learn hence sample the distribution of the features $\z$ in the
latent space $\cZ$, we can leverage the denoising and diffusion
method introduced in Chapter \ref{ch:denoising}. Concretely, suppose
$p(\z)$ is the distribution of the learned features $\z$ in the
feature space. We may learn a denosing process that, starting from an
initial standard Gaussian distribution $\vg \sim \cN(0, \vI)$,
converges to this distribution $p(\z)$. Following the method
introduced in Chapter \ref{ch:denoising}, we want to learn a denoiser
$\bar{\z}$ such that it iteratively
denoises a sample from the initial distribution $\cN(0, \vI)$ to the
target distribution of the learned representations $p(\z)$. This
process is illustrated by the bottom half of Figure
\ref{fig:mental-picture} in Chapter \ref{ch:intro}.

Here to learn the denoising,  let us take the simple case of flow
matching introduced in Chapter \ref{ch:denoising}. We can define the
denoiser as a time-dependent vector field $\bar{\z}_{\theta}(t,
\z_t)$ parameterized by $\theta$ that transforms a sample from the
initial distribution to the target distribution over time $t \in [0,
1]$. Here, $\z_t$ is an interpolation between $\z_0$ and $\z_1$ at time $t$:
\begin{equation}
\z_t = (1 - t)\z_0 + t\z_1,
\end{equation}
where $\z_0 \sim p(\z)$ is a sample from the learned representation
distribution and $\z_1 \sim \cN(0, \vI)$.
The training objective for the denoiser can then be expressed as:
\begin{equation}
\min_{\theta} \cL_{\mathrm{FM}}(\theta) = \Ex_{\z_0 \sim p(\z), \z_1
\sim \cN(0, I), t \sim \mathrm{Unif}(0, 1)} \left[
\norm{\bar{\z}_{\theta}(t,\z_t) - (\z_1 - \z_0)}_2^2 \right].
\end{equation}

Typically, the denoising vector field $\bar{\z}_{\theta}(t,
\z_t)$ are modeled by a deep network. Popular choices of the deep
network can be either a U-Net or a diffusion image transformer (DiT).
Below we provide more details about how they can be implemented in practice.

% %more detail about the UNet, and refer back to PCA. Add the references.
% \subsubsection{Realizing Denoising with a U-Net.}
% A common architecture for implementing the parametric denoiser
% $\bar{\z}_{\theta}(t, \z_t)$ is the \textit{U-Net}, an encoder-decoder
% network with skip connections. The encoder progressively downsamples the
% noisy input to extract multi-scale features, while the decoder upsamples
% and refines predictions. Skip connections concatenate encoder features to
% corresponding decoder layers, preserving spatial detail. To capture
% long-range dependencies, self-attention layers are often inserted at
% lower resolutions. The convolutional structure provides strong spatial
% inductive biases, which can be beneficial when training data is limited.

% The timestep $t$ is typically embedded via
% sinusoidal functions (analogous to positional encodings in transformers)
% and injected through \textit{adaptive group normalization}: the scale and
% shift parameters of normalization layers are regressed from the timestep
% embedding. For class-conditional generation, a class embedding is simply
% added to the timestep embedding before injection. For text conditioning,
% a pretrained text encoder (e.g., CLIP) produces a sequence of token
% embeddings, which are incorporated via \textit{cross-attention}: at
% selected layers, the image features serve as queries while the text
% embeddings serve as keys and values. This allows the network to attend to
% relevant parts of the text prompt when denoising each spatial location.

\subsection{Realizing Denoising with a U-Net}
\label{sec:UNet-implement}

The dominant architecture for implementing the denoising network
$\boldsymbol{\epsilon}_\theta$ in diffusion models has been the
\textit{U-Net}~\citep{ronneberger2015unet}, an encoder-decoder
convolutional network with skip connections originally developed for
biomedical image segmentation. Ho et al.~\citep{ho2020denoising}
first adopted the U-Net as the backbone for Denoising Diffusion
Probabilistic Models (DDPM), establishing it as the default
architecture for the field. Subsequently, Dhariwal and
Nichol~\citep{dhariwal2021diffusion} conducted a systematic series of
architectural ablations that substantially improved the U-Net design,
producing what they call the \emph{Ablated Diffusion Model} (ADM).
Their improved model was the first diffusion model to surpass GANs in
sample quality on ImageNet. Building on these foundations, Rombach et
al.~\citep{rombach2022high} moved the U-Net from pixel space into the
latent space of a pretrained VAE (Section~\ref{sec:VAE-implement})
and introduced cross-attention layers for flexible multi-modal
conditioning, giving rise to the \emph{Latent Diffusion Model} (LDM)
that underpins Stable Diffusion. In this section, we trace the
evolution of the U-Net backbone through these three stages, before
discussing the architectural limitations that have motivated the
recent transition to transformer-based alternatives
(Section~\ref{sec:DiT-implement}).

\paragraph{Encoder-Decoder Structure with Skip Connections.}
The U-Net follows an encoder-decoder structure. The encoder
progressively downsamples the input through a series of resolution
stages, extracting features at multiple spatial scales; the decoder
then upsamples these features back to the original resolution to
produce the output prediction. The distinguishing feature of the
U-Net is its \emph{skip connections}: at each resolution level, the
encoder feature map is concatenated channel-wise to the corresponding
decoder feature map. This allows the decoder to combine coarse,
semantically rich features from the bottleneck with fine-grained
spatial details from earlier encoder layers, which is essential for
producing spatially precise predictions.

In the DDPM U-Net~\citep{ho2020denoising}, the architecture begins
with an input convolution that projects the noised input into a
feature space with a base channel count (e.g., 128). The encoder then
proceeds through multiple resolution stages---four for $32 \times 32$
inputs, or six for $256 \times 256$ inputs---where each stage
consists of two convolutional residual blocks. Spatial downsampling
by a factor of $2\times$ is performed between consecutive stages via
strided convolutions. Self-attention layers are inserted at the $16
\times 16$ resolution level between the convolutional blocks,
providing limited long-range modeling at low resolution. At the
bottleneck (the lowest resolution), a middle block with additional
residual blocks and self-attention further refines the features. The
decoder mirrors the encoder: residual blocks at each resolution level
followed by upsampling via transposed convolutions, with skip
connections concatenating the corresponding encoder activations. The
architecture uses group normalization~\citep{wu2018group} throughout,
replacing the weight normalization from which the design originates.

When the U-Net operates in pixel space, its input and output are the
noised image $\x_t \in \mathbb{R}^{H \times W \times 3}$ and the
predicted noise $\hat{\boldsymbol{\epsilon}} \in \mathbb{R}^{H \times
W \times 3}$, respectively. When operating in the latent space of a
pretrained VAE (as in LDM), the input becomes $\z_t \in \mathbb{R}^{h
\times w \times c}$ and the output is $\hat{\boldsymbol{\epsilon}}
\in \mathbb{R}^{h \times w \times c}$, with $h = H/f$, $w = W/f$ as
defined in Section~\ref{sec:VAE-implement}.

\paragraph{Timestep Conditioning via Adaptive Group Normalization.}
The denoising network must be conditioned on the diffusion timestep
$t$ so that it can adapt its behavior across different noise levels.
Following the positional encoding scheme from the
Transformer~\citep{vaswani2017attention}, the timestep $t$ is first
mapped to a sinusoidal frequency embedding:
\begin{equation}
\mathrm{Embed}(t)_{2i} = \sin\!\left(\frac{t}{10000^{2i/d_e}}\right), \quad
\mathrm{Embed}(t)_{2i+1} = \cos\!\left(\frac{t}{10000^{2i/d_e}}\right),
\end{equation}
where $d_e$ is the embedding dimension and $i$ indexes its entries.
This embedding is then passed through a small MLP (typically two
linear layers with a nonlinearity) to produce a timestep vector
$\boldsymbol{e}_t \in \mathbb{R}^{d_e}$.

Rather than simply concatenating $\boldsymbol{e}_t$ to the input, the
U-Net injects the timestep information into every residual block via
\emph{Adaptive Group Normalization}
(AdaGN)~\citep{dhariwal2021diffusion}. After the first convolution
within each residual block, a group normalization operation is
applied, followed by an affine modulation whose scale and shift
parameters are regressed from the timestep embedding:
\begin{equation}
\mathrm{AdaGN}(\boldsymbol{h}, \boldsymbol{e}_t) = \boldsymbol{y}_s
\hada \mathrm{GroupNorm}(\boldsymbol{h}) + \boldsymbol{y}_b,
\label{eq:adagn}
\end{equation}
where $\boldsymbol{h}$ denotes the intermediate feature map,
$[\boldsymbol{y}_s, \boldsymbol{y}_b] =
\mathrm{Linear}(\boldsymbol{e}_t)$ are the scale and shift vectors
obtained from a linear projection of the timestep embedding, and
$\hada$ denotes element-wise multiplication. This mechanism is
analogous to adaptive instance
normalization~\citep{dumoulin2017learned} and FiLM
conditioning~\citep{perez2018film}, and allows each layer to modulate
its activations based on the current noise level. Ablation
experiments in~\citep{dhariwal2021diffusion} show that removing AdaGN
degrades FID by over 2 points, confirming its importance.

For class-conditional generation, a class embedding
$\mathrm{Embed}(c)$ (a learned lookup table) is simply added to the
timestep embedding before injection, so that $\boldsymbol{e}_t$ is
replaced by $\boldsymbol{e}_t + \mathrm{Embed}(c)$ in
Eq.~\eqref{eq:adagn}. This mechanism is lightweight but limited to
categorical labels; richer conditioning modalities such as text
require the cross-attention mechanism introduced in Section~\ref{sec:unet-ldm}.

\subsubsection{Architecture Improvements: From DDPM to ADM}
\label{sec:unet-adm}

Dhariwal and Nichol~\citep{dhariwal2021diffusion} conducted a
systematic ablation study to identify which architectural
modifications most improve sample quality. Starting from the baseline
DDPM U-Net~\citep{ho2020denoising}, they explored the following
changes on ImageNet $128 \times 128$: (a) expanding self-attention
from the $16 \times 16$ resolution to three resolutions ($32 \times
32$, $16 \times 16$, $8 \times 8$); (b) increasing the number of
attention heads, using 64 channels per head to better match the
Transformer convention; (c) adopting the
BigGAN~\citep{brock2019large} residual block for both upsampling and
downsampling, which improves gradient flow; and (d) including the
AdaGN conditioning layer (described above) in every residual block by
default. All of these changes improve FID, and their effects compound
positively. The resulting model, which they name ADM, serves as the
new default U-Net design for subsequent work.

In addition to architectural improvements, Dhariwal and Nichol
introduced \emph{classifier guidance} as a mechanism for trading off
sample diversity against fidelity. A classifier $p_\phi(c|\z_t, t)$
is trained on noisy images, and its gradients are used to steer the
sampling trajectory toward high-probability regions of the desired class:
\begin{equation}
\hat{\boldsymbol{\mu}}_\theta(\z_t, t) =
\boldsymbol{\mu}_\theta(\z_t, t) + s \,
\boldsymbol{\Sigma}_\theta(\z_t, t) \, \nabla_{\z_t} \log p_\phi(c|\z_t),
\end{equation}
where $s > 1$ is a guidance scale that sharpens the effective class
distribution. With this technique, ADM achieved a state-of-the-art
FID of 4.59 on class-conditional ImageNet $256 \times 256$,
surpassing the long-standing GAN benchmark of BigGAN-deep (FID 6.95)
for the first time. The guided model (ADM-G) further improved FID to
3.94 when combined with an upsampling stack (ADM-U).

The main drawback of classifier guidance is that it requires training
a separate classifier on noisy images. This limitation was
subsequently addressed by \emph{classifier-free
guidance}~\citep{ho2022classifier}, which avoids the external
classifier entirely by jointly learning conditional and unconditional
denoising within a single model---a technique we describe in detail
in Section~\ref{sec:DiT-implement}.

\subsubsection{The U-Net in Latent Diffusion (Stable Diffusion)}
\label{sec:unet-ldm}

While ADM operates directly in pixel space, Rombach et
al.~\citep{rombach2022high} observed that much of the computational
burden is spent modeling perceptually irrelevant high-frequency
details. Their key idea is to separate \emph{perceptual compression}
from \emph{generative modeling}: a pretrained autoencoder
(Section~\ref{sec:VAE-implement}) first compresses the image into a
low-dimensional latent space, and the diffusion model then operates
entirely within this latent space. For a $256 \times 256$ image with
downsampling factor $f = 8$, the U-Net processes a $32 \times 32
\times 4$ latent tensor rather than a $256 \times 256 \times 3$ pixel
tensor---a reduction of roughly $48\times$ in spatial elements. This
yields a speedup of at least $2.7\times$ in both training and
sampling throughput while improving FID.

Beyond computational savings, the latent diffusion framework
introduces a general-purpose \emph{cross-attention conditioning
mechanism} that enables the U-Net to condition on arbitrary
modalities such as text prompts, semantic maps, or bounding boxes. In
brief, a domain-specific encoder $\tau_\theta$ maps the conditioning
input $\boldsymbol{y}$ into an intermediate representation
$\tau_\theta(\boldsymbol{y}) \in \mathbb{R}^{M \times d_\tau}$, which
is then injected into the U-Net at multiple resolution levels via
cross-attention layers: the image features serve as queries while the
conditioning embeddings serve as keys and values, allowing each
spatial location to attend to relevant parts of the conditioning
signal. We defer the full mathematical description of this
cross-attention mechanism and the conditional training objective to
Section~\ref{subsec:basic_realization}.

In Stable Diffusion~\citep{rombach2022high}, the text encoder
$\tau_\theta$ is instantiated as the pretrained CLIP ViT-L/14
model~\citep{radford2021learning}, which maps text prompts to a
sequence of token embeddings. The U-Net itself retains the ADM
improvements---BigGAN residual blocks, attention at three resolution
levels, AdaGN for timestep injection---and additionally interleaves
cross-attention layers at these same resolution levels for text
conditioning. The full model contains approximately 860M parameters
in the U-Net and 123M in the text encoder.

\subsubsection{Limitations of the U-Net Architecture}
\label{sec:unet-limitations}

Despite its success, the U-Net has several limitations that have
motivated alternative backbones. First, \textbf{scaling} a U-Net is
cumbersome: its encoder-decoder structure involves many coupled
design choices (channel multipliers, attention placement,
up/downsampling block types), and it is unclear how to allocate
additional compute optimally. In contrast, the plain Transformer
scales along two simple axes, depth and width, and exhibits
well-characterized scaling laws based on empirical
studies~\citep{kaplan2020scaling}. Second, the U-Net's convolutional
layers encode a strong \textbf{spatial inductive bias} (locality,
translation equivariance) that helps in low-data regimes but limits
long-range modeling on large datasets. Self-attention partially
addresses this, yet in standard U-Net designs it is applied only at
low-resolution stages due to its quadratic cost, leaving
high-resolution features entirely local. Third, the U-Net is a
\textbf{image-specific} architecture that does not readily generalize
so well to other modalities. The Transformer, having proven effective
across language, vision, audio, and multimodal
settings~\citep{vaswani2017attention}, offers a unified backbone that
can leverage cross-domain scaling insights.

These considerations motivated the Diffusion Transformer
(DiT)~\citep{peebles2023scalable}, which replaces the U-Net entirely
with a plain Transformer operating on latent patches. As we describe
in Section~\ref{sec:DiT-implement}, DiT matches and surpasses the
best U-Net-based models while exhibiting cleaner scaling behavior
(Gflops--FID correlation of $-0.93$) and a simpler architecture.

% % refer the transformer arch back to previous Chapters. Add the references.
% \subsubsection{Realizing Denoising with a DiT.}
% An alternative to the U-Net is the \textit{Diffusion Transformer (DiT)},
% which replaces convolutions entirely with transformer blocks. Following
% the Vision Transformer, the input is divided into patches and linearly
% embedded as token sequences. Unlike the U-Net, the DiT has no explicit
% spatial inductive bias---all structure is learned from data. This
% generality enables the same architecture to handle images, video, and
% other modalities with minimal modification. Perhaps most importantly, DiT
% follows favorable scaling laws: larger models and more compute yield
% predictable improvements in sample quality.

% Timestep and class conditioning in DiT is
% achieved through \textit{adaptive layer normalization} (adaLN-Zero): the
% scale and shift parameters of each layer normalization are regressed from
% the conditioning embedding (formed by summing timestep and class
% embeddings), with these parameters initialized near zero for training
% stability. For text conditioning, two approaches are common: (1)
% cross-attention, as in the U-Net, where text token embeddings from a
% frozen encoder are attended to by the image tokens; or (2) prepending the
% text tokens directly to the image token sequence, allowing the
% transformer's self-attention to jointly process both modalities.

% Once we have trained the denoiser, we can use it to generate new
% samples from the learned representation distribution via the denoising
% process. The generated feature samples can then be decoded back to
% the image space with the corresponding decoder $g$.

\subsection{Realizing Denoising with a Diffusion Transformer}
\label{sec:DiT-implement}

An alternative to the U-Net for implementing the denoising network
$\boldsymbol{\epsilon}_\theta$ is the \textit{Diffusion Transformer
(DiT)}~\citep{peebles2023scalable}, which replaces convolutions
entirely with transformer blocks. The key insight is that the U-Net's
spatial inductive bias is \emph{not crucial} to the performance of
diffusion models: a standard transformer operating on sequences of
latent patches can match and surpass the U-Net while inheriting the
favorable scaling properties of the transformer architecture
class~\citep{vaswani2017attention}. Unlike the U-Net, the DiT has no
explicit spatial inductive bias---all spatial structure is learned
from data. This generality enables the same architecture to handle
images, video, and other modalities with minimal modification.

DiT operates within the latent space of a pretrained VAE
(Section~\ref{sec:VAE-implement}), making the complete image
generation pipeline a hybrid architecture: a convolutional VAE for
perceptual compression and a transformer-based DDPM for generative
modeling. Following the Vision Transformer
(ViT)~\citep{dosovitskiy2020image}, the input latent is divided into
patches and linearly embedded as a token sequence. We describe the
architecture in detail below.

\paragraph{Patchify and Positional Encoding.}
The input to DiT is a noised latent $\z_t \in \mathbb{R}^{h \times w
\times c}$, where $h = H/f$ and $w = W/f$ are the spatial dimensions
of the latent space defined by the VAE encoder (e.g., $32 \times 32
\times 4$ for $256 \times 256$ images with $f = 8$, $c = 4$). The
first layer, called \emph{patchify}, partitions $\z_t$ into a grid of
non-overlapping patches of size $p \times p$ and linearly embeds each
patch into a $d$-dimensional token representation. This produces a sequence of
\begin{equation}
T = \left(\frac{h}{p}\right) \times \left(\frac{w}{p}\right)
\end{equation}
tokens in $\mathbb{R}^d$. Standard sinusoidal positional embeddings
(the sine-cosine variant from ViT) are added to each token to encode
spatial position. The patch size $p$ is a key design parameter:
halving $p$ quadruples $T$ and thus at least quadruples the total
Gflops of the transformer, while leaving the parameter count
essentially unchanged. In~\citep{peebles2023scalable}, the design
space includes $p \in \{2, 4, 8\}$; the finest setting $p = 2$ yields
the best generation quality.

\paragraph{DiT Block with adaLN-Zero.}
After patchification, the token sequence is processed by a stack of
$N$ transformer blocks. Each block follows the standard transformer
structure~\citep{vaswani2017attention}: layer normalization,
multi-head self-attention (MHSA), a second layer normalization, and a
pointwise feedforward network (MLP with GELU activations).

The central design question is how to inject the conditioning
information---the diffusion timestep $t$ and the class label
$c$---into each block. Peebles and Xie~\citep{peebles2023scalable}
compare four strategies: (1) \emph{in-context conditioning}, which
appends $t$ and $c$ as extra tokens in the input sequence; (2)
\emph{cross-attention}, which introduces an additional multi-head
cross-attention layer that attends to the conditioning embeddings;
(3) \emph{adaptive layer normalization (adaLN)}, which regresses the
scale and shift parameters of each layer normalization from the
conditioning signal; and (4) \emph{adaLN-Zero}, which extends adaLN
with an additional gating mechanism initialized to zero. Among these,
adaLN-Zero achieves the lowest FID while being the most
compute-efficient, and is therefore adopted as the default.

In the adaLN-Zero block, the conditioning information is first formed
by summing the embeddings of the timestep and class label:
\begin{equation}
\boldsymbol{e} = \mathrm{Embed}(t) + \mathrm{Embed}(c),
\end{equation}
where $\mathrm{Embed}(t)$ is obtained by passing a sinusoidal
frequency embedding of $t$ through a two-layer MLP with SiLU
activations, and $\mathrm{Embed}(c)$ is a learned class embedding.
From $\boldsymbol{e}$, a linear layer regresses six sets of
dimension-wise parameters: $\boldsymbol{\gamma}_1,
\boldsymbol{\beta}_1, \boldsymbol{\alpha}_1$ for the self-attention
sub-block and $\boldsymbol{\gamma}_2, \boldsymbol{\beta}_2,
\boldsymbol{\alpha}_2$ for the feedforward sub-block. The adaptive
layer normalization applies as:
\begin{equation}
\mathrm{adaLN}(\boldsymbol{h}, \boldsymbol{\gamma},
\boldsymbol{\beta}) = \boldsymbol{\gamma} \hada
\mathrm{LayerNorm}(\boldsymbol{h}) + \boldsymbol{\beta},
\end{equation}
where $\hada$ denotes element-wise multiplication. The parameters
$\boldsymbol{\alpha}_1$ and $\boldsymbol{\alpha}_2$ serve as
dimension-wise \emph{gating factors} applied immediately before each
residual connection:
\begin{equation}
\begin{aligned}
    \boldsymbol{h}' &= \boldsymbol{h} + \boldsymbol{\alpha}_1 \hada
    \mathrm{MHSA}\!\big(\mathrm{adaLN}(\boldsymbol{h},
    \boldsymbol{\gamma}_1, \boldsymbol{\beta}_1)\big), \\
    \boldsymbol{h}'' &= \boldsymbol{h}' + \boldsymbol{\alpha}_2 \hada
    \mathrm{MLP}\!\big(\mathrm{adaLN}(\boldsymbol{h}',
    \boldsymbol{\gamma}_2, \boldsymbol{\beta}_2)\big).
\end{aligned}
\end{equation}
The ``Zero'' in adaLN-Zero refers to the initialization: all
$\boldsymbol{\alpha}$ parameters are initialized to the zero vector
by the MLP that regresses them, so that each DiT block acts as the
\emph{identity function} at the start of training. This
zero-initialization strategy, inspired by similar practices in
ResNets~\citep{goyal2017accurate} and diffusion U-Nets, is important
for training stability and yields significantly lower FID than vanilla adaLN.

We note that adaLN applies the \emph{same} conditioning function to
all tokens uniformly. For text-to-image generation where the
conditioning signal is a variable-length sequence of text tokens, two
alternative approaches are common: (1) cross-attention, where text
token embeddings from a frozen encoder are attended to by the image
tokens; or (2) prepending the text tokens directly to the image token
sequence, allowing the transformer's self-attention to jointly
process both modalities.

\paragraph{Output Projection.}
After the final DiT block, the model must decode the sequence of
token representations back into a spatial prediction. A final
adaptive layer normalization (conditioned on $\boldsymbol{e}$) is
applied, followed by a linear layer that projects each token into a
tensor of shape $p \times p \times 2c$. The factor of $2c$ accounts
for the two quantities the model predicts: the noise estimate
$\hat{\boldsymbol{\epsilon}} \in \mathbb{R}^{h \times w \times c}$
and the diagonal covariance $\boldsymbol{\Sigma}_\theta \in
\mathbb{R}^{h \times w \times c}$. The decoded tokens are then
rearranged (un-patchified) back into their original spatial layout of
$\mathbb{R}^{h \times w \times c}$ for each output.

\subsubsection{Training}
The DiT is trained as a class-conditional latent diffusion model on
the ImageNet dataset~\citep{deng2009imagenet} at $256 \times 256$ and
$512 \times 512$ image resolution. The diffusion process operates on
the latent representations produced by the pretrained VAE from
Section~\ref{sec:VAE-implement}: for $256 \times 256$ images, the
input to DiT is $\z_t \in \mathbb{R}^{32 \times 32 \times 4}$; for
$512 \times 512$ images, it is $\z_t \in \mathbb{R}^{64 \times 64 \times 4}$.

The training objective combines two terms. The noise prediction loss
trains $\boldsymbol{\epsilon}_\theta$ via the simple mean-squared error:
\begin{equation}
\mathcal{L}_{\mathrm{simple}}(\theta) = \mathbb{E}_{\z_0,
\boldsymbol{\epsilon}, t}\left[\|\boldsymbol{\epsilon}_\theta(\z_t,
t) - \boldsymbol{\epsilon}\|_2^2\right],
\end{equation}
where $\z_0$ is the clean latent, $\boldsymbol{\epsilon} \sim
\mathcal{N}(\boldsymbol{0}, \boldsymbol{I})$ is the sampled noise,
and $\z_t = \sqrt{\bar{\alpha}_t}\,\z_0 + \sqrt{1 -
\bar{\alpha}_t}\,\boldsymbol{\epsilon}$ is the noised latent at
timestep $t$. The learned covariance $\boldsymbol{\Sigma}_\theta$ is
additionally trained with the full variational lower bound $D_{KL}$
term, following the approach of Nichol and
Dhariwal~\citep{nichol2021improved}. A linear variance schedule is
used with $t_{\max} = 1000$ steps, ranging from $1 \times 10^{-4}$ to
$2 \times 10^{-2}$.

All models are trained with AdamW~\citep{loshchilov2019decoupled}
using a constant learning rate of $1 \times 10^{-4}$, no weight
decay, and a batch size of $256$. Data augmentation consists solely
of random horizontal flips. An exponential moving average (EMA) of
the model weights is maintained with a decay of $0.9999$; all
reported results use the EMA model. Notably, unlike supervised ViT
training, DiT requires \emph{no} learning rate warmup and \emph{no}
regularization---training is highly stable across all model
configurations, with no observed loss spikes.

\subsubsection{Classifier-Free Guidance}
To produce class-controlled sampling, DiT can easily employ
\emph{classifier-free guidance}~\citep{ho2022classifier}. During
training, the class label $c$ is randomly dropped with some
probability and replaced by a learned null embedding $\varnothing$,
so that the model learns both the conditional
$\boldsymbol{\epsilon}_\theta(\z_t, c)$ and unconditional
$\boldsymbol{\epsilon}_\theta(\z_t, \varnothing)$ noise predictions.
At inference time, the guided noise prediction is computed as:
\begin{equation}
\hat{\boldsymbol{\epsilon}}_\theta(\z_t, c) =
\boldsymbol{\epsilon}_\theta(\z_t, \varnothing) + s \cdot
\big(\boldsymbol{\epsilon}_\theta(\z_t, c) -
\boldsymbol{\epsilon}_\theta(\z_t, \varnothing)\big),
\label{eq:cfg}
\end{equation}
where $s > 1$ is the guidance scale ($s = 1$ recovers standard
unguided sampling). The motivation follows from Bayes' rule: since
$\nabla_{\z_t} \log p(c | \z_t) \propto \nabla_{\z_t} \log p(\z_t |
c) - \nabla_{\z_t} \log p(\z_t)$, the guided estimate in
Eq.~\eqref{eq:cfg} effectively steers the sampling trajectory toward
regions of high conditional likelihood $p(c | \z_t)$. This technique
has a dramatic effect on generation quality: DiT-XL/2 achieves an FID
of $9.62$ without guidance, which improves to $2.27$ with a guidance
scale of $s = 1.5$---surpassing all prior diffusion models on the
class-conditional ImageNet $256 \times 256$ benchmark.

\subsubsection{Generation}
Once the denoising network $\boldsymbol{\epsilon}_\theta$ (i.e., the
DiT) has been trained, new images can be generated by running the
learned reverse process. Starting from pure noise $\z_{t_{\max}} \sim
\mathcal{N}(\boldsymbol{0}, \boldsymbol{I})$, the model iteratively
denoises by sampling $\z_{t-1} \sim p_\theta(\z_{t-1} | \z_t)$ at
each step, using the predicted noise
$\hat{\boldsymbol{\epsilon}}_\theta$ and covariance
$\boldsymbol{\Sigma}_\theta$ to parameterize the reverse transition.
After $t_{\max}$ denoising steps, the resulting clean latent $\z_0$
is decoded back to the image space via the VAE decoder: $\tilde{\x} =
\mathcal{D}(\z_0) \in \mathbb{R}^{H \times W \times 3}$.

\section{Conditioned Image  Representation and Generation}
\label{sec:cond_gen}
%\yima{From the above latent diffusion model for images, it is
% natural to transition to conditioned image generation. Text
% condition via CLIP or SimCLIP... and DPS implementation,
% classifier-free, or cross-attention etc. (Robin, Jingfeng, Daniel,
% and Peter, decide what to feature...)}

% --- draft by Daniel Yeh ---
In the previous sections, we established the fundamental machinery
for generative modeling: an autoencoder architecture to encode the
data $\x$ into a (compressed and more structured) latent $\z$:
\begin{equation}
\x\xrightarrow{\hspace{1mm} \mathcal{E} \hspace{1mm}} \z
\xrightarrow{\hspace{1mm} \mathcal{D} \hspace{1mm}} \hat{\x}.
\end{equation}

Then, a generative process is learned to model the distribution of
these latents $p(\z)$. In this section, we shift our focus from
\textit{unconditional generation} to \textit{control}. We assume the
representation $\z$ is already learned (via methods detailed in
Chapter \ref{ch:conditional-inference}); our specific task here is to
steer the sampling process using a user-provided ``control'' signal
$\vc$. That is, we aim to sample from the conditional distribution
$p(\z | \vc)$. In practice, $\vc$ could represent partial attributes
of the desired sample image: masked image, a semantic mask, class
information, or a text caption describing the desired content.

\subsection{Task Formulation}
\label{subsec:cond_problem}

Formally, the goal of conditioned generation is to learn a mapping
that transforms a simple source distribution (typically Gaussian
noise) into a target conditional distribution. The generated feature
$\hat{\z}$ is subsequently decoded to a corresponding image
$\hat{\x}$. As illustrated as follows, the system is defined by its
inputs and outputs:
\begin{enumerate}
\item \textbf{Stochastic Input ($\epsilon$):} A noise variable
    sampled from a standard normal prior $\epsilon \sim
    \mathcal{N}(\bm{0}, \I)$. This provides the ``seed'' for
    diversity, ensuring that we can generate multiple variations of
    the same concept.
\item \textbf{Control Signal ($\vc$):} A structured input defining
    the desired semantic or spatial content. This can be a text
    caption (e.g., ``a red car''), a semantic segmentation mask, or a layout.
\item \textbf{Output ($\hat{\x}$):} A high-fidelity data sample
    $\hat{\x} = \mathcal{D}(\hat{\z})$ that aligns with the user's
    intent $\vc$ while remaining on the natural image manifold.
\end{enumerate}
We employ a domain-specific \textbf{Condition Encoder}, denoted as
$\tau_\phi$, to project the control signal $\vc$ into a useful
embedding space. Hence, the overall controlled generation process can
be formally described as a composition of a generator $G_\theta$ and
the decoder $\mathcal{D}$:
\begin{equation}
\underbrace{\epsilon \sim \mathcal{N}(\bm{0}, \I)}_{\text{Source}}, \vc
\xrightarrow{\quad G_\theta(\epsilon, \tau_\phi(\vc)) \quad}
\hat{\z} \sim p(\z\mid \vc)
\xrightarrow{\quad \mathcal{D}(\hat{\z}) \quad}
\hat{\x} \approx \x.
\end{equation}

\paragraph{The Modality Gap.}
While the task is conceptually straightforward—mapping $(\epsilon,
\vc) \to \x$—it requires bridging disjoint topological spaces,
particularly in text-to-image generation. The visual data $\x$ lies
in a continuous, high-dimensional pixel space $\mathbb{R}^{H \times W
\times 3}$, whereas the condition $\vc$ typically resides in a
discrete symbolic space (language). There is no natural Euclidean
metric to measure the distance $\| \x - \vc \|$ directly.
Therefore, the critical component $\tau_\phi(\vc)$ mentioned above
must align the modalities. As discussed in Chapter
\ref{ch:conditional-inference}, modern approaches rely on pre-trained
contrastive encoders (like CLIP) or large language models (like
T5~\citep{raffel2020exploring}) to map discrete language tokens into
a continuous semantic embedding space. The generative backbone
$G_\theta$ (i.e., Denoising U-Net) must then learn to map the noise
distribution to the data distribution, guided by these semantic
embeddings via a specific fusion mechanism, such as Cross-Attention.

\paragraph{The One-to-Many Inverse Problem.}
Conditioning is inherently ill-posed because the control signal is
usually \textit{under-determined}. A text prompt $\vc=$ ``a dog''
provides semantic category information but omits the dog's pose,
lighting, background, and fur pattern. Consequently, the inverse
mapping $\vc \to \x$ is not a single point, but a complex multi-modal
distribution. If we were to train a simple deterministic regression
model to minimize the Euclidean error $\|\x - G_\theta(\vc)\|^2$, the
model would converge to the conditional mean $\mathbb{E}[\x|\vc]$—an
average of all possible dogs—resulting in a blurry, unrealistic
output. To generate high-fidelity results, we can utilize generative
frameworks (such as U-Net or DiT mentioned in
Section~\ref{sec:auto-encode-img-gen}) to model the full conditional
distribution $p(\z|\vc)$, allowing us to sample diverse, sharp
instances that satisfy the condition.

% To generate high-fidelity results, we
% must utilize the Flow Matching framework (Eq. 8.7 \yima{what is
% this?}) to model the full
% conditional distribution $p(\z|\vc)$, allowing us to sample diverse,
% sharp instances that satisfy the condition.

% \subsection{The Basic Realization: Text-to-Image Generation}
% \label{subsec:basic_realization}

% To realize the conditional framework defined in Section
% \ref{subsec:cond_problem}, we require a concrete architecture capable
% of bridging the modality gap between discrete text tokens and
% continuous pixel arrays. As illustrated in Figure
% \ref{fig:sd_pipeline}, modern state-of-the-art models (such as Stable
% Diffusion~\citep{rombach2022high}, DALL-E
% 2~\citep{ramesh2022hierarchical}, and
% Imagen~\citep{saharia2022photorealistic}) universally follow a
% pipeline consisting of three decoupled components: (1) a Text Encoder
% to interpret the prompt, (2) a Generation Model to synthesize a
% compressed ``latent representation'' of the image, and (3) a Decoder to map
% this representation back to a high-fidelity image in the pixel
% domain. We detail the
% implementation of each component below, using the Latent Diffusion
% Model (LDM)~\citep{rombach2022high} as the primary case study.

% \begin{figure}[h]
% \centering
% \includegraphics[width=\textwidth]{\toplevelprefix/chapters/chapter8/figs/sd_pipeline.pdf}
% \caption{\textbf{The three-stage inference pipeline of Stable Diffusion.}
%     (1) \textbf{Text Encoder} (e.g., CLIP) converts the input prompt
%     $\vc$ into semantic embeddings.
%     (2) \textbf{Generation Model} (a time-conditional U-Net)
%     iteratively denoises a random latent tensor $\z_T$ to produce a
%     clean intermediate representation $\hat{\z}$, conditioned on the
%     text embeddings via cross-attention.
%     (3) \textbf{Decoder} (from a VAE) maps the denoised latent
%     $\hat{\z}_0$ back to the pixel space to synthesize the final
% high-resolution image $\hat{\x}$.}
% \label{fig:sd_pipeline}
% \end{figure}

% \paragraph{The Text Encoder.}
% The first component converts the user's text prompt $\vc$ (e.g., ``A
% fluffy ginger cat on the couch'') into a machine-readable vector
% sequence. Its role is semantic understanding: because the generative
% model's knowledge is limited to its training data, a powerful
% pre-trained encoder is crucial for generalizing to novel concepts or
% rare words. In practice, Stable Diffusion typically utilizes a frozen
% CLIP text encoder~\citep{radford2021learning}, producing a sequence
% of semantically rich embeddings $\tau_\theta(\vc) \in
% \mathbb{R}^{d_\tau \times M}$ (where $d_\tau$ is typically 768 or
% 1024). \yima{Please refer to Section \ref{sec:CLIP}, and make sure
% notation and dimension match!} Recent
% works like Imagen have demonstrated that scaling up this encoder
% (e.g., using T5-XXL \yima{citation?}) yields significantly larger
% improvements in
% image fidelity than scaling the generative backbone itself.

% \paragraph{The Generation Model (Latent Representation).}
% The core generative process does not produce the final image
% directly. Instead, to avoid the prohibitive computational cost of
% modeling high-resolution pixels (e.g., $1024 \times 1024$), it
% generates a compressed ``intermediate product''—a latent representation $\z$.
% \begin{itemize}
% \item \textbf{Data Format:} The input image $\x \in \mathbb{R}^{H
%     \times W \times 3}$ is compressed by a factor of $f=8$ into a
%     latent tensor $\z_0 \in \mathbb{R}^{h \times w \times
% 4}$.\yima{How is it ``compressed''? Via a VAE? Do you follow some
% work here? citation?} For a
%     standard $512^2$ image, this results in a manageable $64 \times
%     64 \times 4$ latent space.
% \item \textbf{Conditioning Mechanism:} The generation model is
%     implemented as a Time-Conditional U-Net. It predicts the noise
%     $\epsilon$ added to the latent $\z_t$, conditioned on the text
%     embeddings via {\em Cross-Attention}. To enable this interaction,
%     the spatial latent features $\phi_i(\z_t)$ \yima{what
% ``spatial'' mean here? it was not introduced. what is
% $\phi_i(\cdot)$? how it is specified?} are flattened into a
%     sequence $\mathbb{R}^{N \times d}$ to act as the Query ($Q$),
%     while the text embeddings serve as Key ($K$) and Value ($V$):
%     \begin{equation}
%         Q = W_Q^{(i)} \cdot \phi_i(\z_t), \quad K = W_K^{(i)} \cdot
%         \tau_\theta(\vc), \quad V = W_V^{(i)} \cdot \tau_\theta(\vc).
%     \end{equation}
%     \yima{is $i$ here the index for layer?}
%     This mechanism allows the model to spatially align specific
%     visual regions (e.g., ``cat's eye'') with specific text tokens.
% \end{itemize}

% Training involves a noise prediction task. We sample a random
% timestep $t$, add Gaussian noise $\epsilon$ to the clean latent
% $\z_0$ to create a noisy latent $\z_t$, and train the U-Net backbone
% $\epsilon_\theta$ to predict the added noise. The model minimizes the
% standard reweighted diffusion objective:
% \begin{equation}
% \mathcal{L}_{\text{LDM}} = \mathbb{E}_{\z_0, \vc, \epsilon \sim
% \mathcal{N}(0,1), t} \left[ \| \epsilon - \epsilon_\theta(\z_t, t,
% \tau_\theta(\vc)) \|^2 \right].
% \end{equation}
% During inference (sampling), we start with pure Gaussian noise $\z_T
% \sim \mathcal{N}(\bm{0}, \I)$. In each step, the U-Net predicts the
% noise component, which is then subtracted from the current state to
% gradually refine the latent representation $\z_t \to \z_{t-1}$ until
% a clean latent $\hat{\z}_0$ is formed.

% \paragraph{The Generation Model (Latent Representation).}
% The core generative process does not produce the final image
% directly. Instead, to avoid the prohibitive computational cost of
% modeling high-resolution pixels (e.g., $1024 \times 1024$), it
% generates a compressed "intermediate product"—a latent representation $\z$.
% \begin{itemize}
% \item \textbf{Data Format:} The input image $\x \in \mathbb{R}^{H
%     \times W \times 3}$ is compressed by a factor of $f=8$ into a
%     latent tensor $\z_0 \in \mathbb{R}^{h \times w \times 4}$.
%     This compression is performed by a pre-trained Variational
% Autoencoder (VAE)~\citep{rombach2022high} that learns a
% perceptually equivalent but computationally efficient lower-dimensional space.
%     For a standard RGB input image of size $512 \times 512 \times
% 3$ (i.e., $H=W=512$), this results in a manageable $64 \times
%     64 \times 4$ latent space.

% \item \textbf{Conditioning Mechanism:} The generation model is
%     implemented as a Time-Conditional U-Net. It predicts the noise
%     $\epsilon$ added to the latent $\z_t$, conditioned on the text
%     embeddings via \textbf{Cross-Attention}.
%     To enable this interaction, we treat the intermediate visual
% features from the U-Net layer $i$ as the Query source, and the text
% embeddings as the Key/Value source.
%     Let $\phi_i(\z_t) \in \mathbb{R}^{d_\epsilon^{(i)} \times N}$
% denote the flattened spatial features at layer $i$ (where
% $d_\epsilon^{(i)}$ varies by layer, e.g., 320, 640, or 1280
% channels), and let $\tau_\theta(\vc) \in \mathbb{R}^{d_\tau \times
% M}$ denote the sequence of $M$ text embeddings (typically $d_\tau=768$).
%     Since these dimensions differ, we project both into a common
% subspace of dimension $d$ (head dimension) using learnable weights.

%     We first define the Query ($\boldsymbol{Q}$), Key
% ($\boldsymbol{K}$), and Value ($\boldsymbol{V}$) matrices using
% projection weights $\boldsymbol{U}^{(i)}_{Q} \in \mathbb{R}^{d
% \times d_\epsilon^{(i)}}$ and $\boldsymbol{U}^{(i)}_{K},
% \boldsymbol{U}^{(i)}_{V} \in \mathbb{R}^{d \times d_\tau}$:
%     \begin{equation}
%         \boldsymbol{Q} = \boldsymbol{U}^{(i)}_{Q} \phi_i(\z_t), \quad
%         \boldsymbol{K} = \boldsymbol{U}^{(i)}_{K} \tau_\theta(\vc), \quad
%         \boldsymbol{V} = \boldsymbol{U}^{(i)}_{V} \tau_\theta(\vc).
%     \end{equation}
%     Adapting the column-wise formulation from [placeholder], the
% cross-attention operation is then computed as:
%     \begin{equation}
%         \text{CrossAttn}(\boldsymbol{Q}, \boldsymbol{K},
% \boldsymbol{V}) \doteq \boldsymbol{V} \cdot
% \operatorname{softmax}\left( \frac{\boldsymbol{K}^\top
% \boldsymbol{Q}}{\sqrt{d}} \right),
%     \end{equation}
%     where the $\operatorname{softmax}$ is defined column-wise as
%     \begin{align}
%         \operatorname{softmax}(\boldsymbol{A})
%         &\doteq \begin{bmatrix}
% \operatorname{softmax}(\boldsymbol{a}_{1}) & \cdots &
% \operatorname{softmax}(\boldsymbol{a}_{N}) \end{bmatrix}, \\
%         \forall \boldsymbol{A}
%         &= \begin{bmatrix} \boldsymbol{a}_{1}, \dots,
% \boldsymbol{a}_{N} \end{bmatrix} \in \mathbb{R}^{M \times N}.
%     \end{align}
%     In this context, the matrix $\boldsymbol{A} =
% \frac{\boldsymbol{K}^\top \boldsymbol{Q}}{\sqrt{d}}$ represents the
% alignment scores, and the column-wise softmax normalizes the
% attention weights over the $M$ text tokens for each of the $N$
% visual regions. This mechanism allows the model to spatially align
% specific visual regions (e.g., ``cat's eye'') with specific text tokens.
% \end{itemize}

% Training involves a noise prediction task. We uniformly sample a random
% timestep $t \in \{1, \dots, T\}$, add Gaussian noise $\epsilon$ to
% the clean latent
% $\z_0$ to create a noisy latent $\z_t$. To inform the model of the
% noise level, the scalar timestep $t$ is first mapped to a
% sinusoidal embedding (encoding low and high frequency components),
% followed by a Multi-Layer Perceptron (MLP) composed of `Linear`
% $\to$ `SiLU` $\to$ `Linear` layers. This produces a rich time
% embedding vector (e.g., 1280-dim) which is added to the feature
% maps in each ResNet block of the U-Net backbone $\epsilon_\theta$.
% The model then minimizes the standard reweighted diffusion objective:
% \begin{equation}
% \mathcal{L}_{\text{LDM}} = \mathbb{E}_{\z_0, \vc, \epsilon \sim
% \mathcal{N}(0,1), t} \left[ \| \epsilon - \epsilon_\theta(\z_t, t,
% \tau_\theta(\vc)) \|^2 \right].
% \end{equation}
% During inference (sampling), we start with pure Gaussian noise $\z_T
% \sim \mathcal{N}(\bm{0}, \I)$. In each step, the U-Net predicts the
% noise component, which is then subtracted from the current state to
% gradually refine the latent representation $\z_t \to \z_{t-1}$ until
% a clean latent $\hat{\z}_0$ is formed.

% \paragraph{The Decoder (Image Reconstruction).}
% The final component is a neural network trained to translate the
% sample latent vector $\hat{\z}$ back into a natural image
% $\hat{\x}$. Unlike the generation model which requires image-text
% pairs, this Decoder is part of a VAE (Variational Autoencoder) that
% can be trained separately on large-scale unlabeled image data.
% During the training phase, the VAE minimizes a reconstruction error
% (combining L1 loss and Perceptual Loss) to ensure the decoder
% $\mathcal{D}$ can accurately map compressed latents back to pixels
% ($\mathcal{D}(\mathcal{E}(\x)) \approx \x$). Once trained, the
% Encoder is discarded from the inference pipeline, and the Decoder is
% frozen. It simply takes the denoised latent $\hat{\z}_0$ produced by
% the Generation Model and upsamples it to the final pixel result $\hat{\x}$.

% \textit{Note on Visualization:} When visualizing the generation for
% every time-step using Stable Diffusion, users often observe the
% image evolving from a
% "blurry blob" to a sharp result. This distinctive "blur-to-sharp"
% transition occurs because, at every step of the sampling process, the
% system passes the current noisy latent $\z_t$ through the Decoder to
% visualize progress. Since the Decoder is trained on clean data, it
% interprets the noise in the latent space as high-frequency texture
% information, rendering the noisy latent as a blurry, low-detail image
% rather than static noise.

\subsection{The Basic Realization: Text-to-Image Generation}
\label{subsec:basic_realization}

To realize the conditional framework defined in
\ref{subsec:cond_problem}, we require a concrete architecture capable
of bridging the modality gap between discrete text tokens and
continuous pixel arrays. As illustrated in
Figure~\ref{fig:sd_pipeline}, modern state-of-the-art models (such as
Stable Diffusion~\citep{rombach2022high}) decouple this complex task
into three independent components: (1) a \textbf{Perceptual
Compression} module (VAE) to translate between pixels and latents,
(2) a \textbf{Semantic Conditioning} interface (CLIP Text Encoder) to
interpret prompts, and (3) a \textbf{Latent Generation} backbone
(U-Net) to synthesize content. We detail the implementation of each
component below.

\begin{figure}[t!]
\centering
\includegraphics[width=\textwidth]{\toplevelprefix/chapters/chapter8/figs/sd_pipeline.pdf}
\caption{\textbf{The three-stage inference pipeline of Stable
    Diffusion.} (1) \textbf{Text Encoder} (e.g., CLIP) converts the input
    prompt $\boldsymbol{c}$ into semantic embeddings. (2)
    \textbf{Generation Model} (a time-conditional U-Net
    $\epsilon_\theta$) iteratively denoises a random latent tensor
    $\boldsymbol{z}_T$ to produce a clean intermediate representation
    $\hat{\boldsymbol{z}}$, conditioned on the text embeddings via
    cross-attention. (3) \textbf{Decoder} (from a VAE) maps the denoised
    latent $\hat{\boldsymbol{z}}$ back to the pixel space to synthesize
    the final high-resolution image $\hat{\boldsymbol{x}}$. \textit{Note:
        The VAE Encoder is utilized only during training and is omitted from
this inference pipeline.}}
\label{fig:sd_pipeline}
\end{figure}

\paragraph{Component 1: Perceptual Compression (The VAE).}
The first pillar of the architecture is a Variational Autoencoder
(VAE) trained to abstract away high-frequency details that are
perceptually insignificant but computationally expensive to model. As
we introduced in Section~\ref{sec:VAE-implement}, this component
defines the \textbf{Latent Space} in which the generation occurs.
\begin{itemize}
\item \textbf{The Encoder ($\mathcal{E}$):} The encoder maps an input
    RGB image $\boldsymbol{x} \in \mathbb{R}^{H \times W \times 3}$
    into a low-dimensional latent representation $\boldsymbol{z} =
    \mathcal{E}(\boldsymbol{x}) \in \mathbb{R}^{h \times w \times
    c}$. Following \cite{rombach2022high}, it utilizes a
    KL-regularized encoder with a downsampling factor of
    $f=H/h=W/w=8$. For a standard input of $512 \times 512 \times 3$,
    this yields a compact latent tensor of dimensions $64 \times 64
    \times 4$. In the context of the diffusion process (Component 3),
    this encoded representation serves as the ``clean'' ground-truth
    latent $\boldsymbol{z}_0$.

\item \textbf{The Decoder ($\mathcal{D}$):} The decoder learns the
    reverse mapping to reconstruct the image from the latent space,
    denoted as $\tilde{\boldsymbol{x}} = \mathcal{D}(\boldsymbol{z})
    = \mathcal{D}(\mathcal{E}(\boldsymbol{x}))$. It is trained to
    minimize the perceptual difference between $\boldsymbol{x}$ and
    $\tilde{\boldsymbol{x}}$. During the inference phase of the
    diffusion model, this decoder is frozen and used to map the
    generated (denoised) latent vectors back to the high-resolution pixel space.
\end{itemize}
Note that the VAE is pre-trained with a combination of perceptual
loss and patch-based adversarial loss to ensure the quality of
reconstructions. Once trained, its weights are frozen, and it serves
as a static interface for the diffusion model.

\paragraph{Component 2: Semantic Conditioning (The Text Encoder).}
To guide the generation process, we require a robust semantic
representation of the user's prompt $\vc$ (e.g., ``A fluffy ginger
cat''). Instead of training a text encoder from scratch, Latent
Diffusion Models typically leverage the frozen text encoder from
CLIP~\citep{radford2021learning}. This pre-trained transformer maps
the input text tokens to a sequence of semantic embeddings
$\tau_\theta(\vc) \in \mathbb{R}^{d_\tau \times M}$, where $M$ is the
sequence length (e.g., 77 tokens) and $d_\tau$ is the embedding
dimension (typically 768 or 1024). These embeddings serve as the
semantic ``Key'' and ``Value'' sources for the generation backbone.

\paragraph{Component 3: Latent Generation (The Time-Conditional U-Net).}
The core generation engine is implemented as a Time-Conditional U-Net
$\epsilon_\theta$. Unlike pixel-space models, this network operates
entirely within the compressed latent space. Its architecture is
composed of a series of ResNet blocks for feature extraction and
Cross-Attention layers for conditioning.
\begin{itemize}
\item \textbf{Input Streams:} The model processes three distinct signals:
    \begin{enumerate}
        \item \textbf{Visual Stream:} The noisy latent $\z_t \in
            \mathbb{R}^{h \times w \times 4}$, which represents the
            image at the current diffusion step.
        \item \textbf{Timestep ($t$):} To inform the model of the
            current noise level, the scalar timestep $t \in \{1,
            \dots, T\}$ is mapped to a sinusoidal embedding (encoding
            low and high frequency components), followed by a
            Multi-Layer Perceptron (MLP) composed of `Linear` $\to$
            `SiLU` $\to$ `Linear` layers. This produces a global time
            embedding vector (e.g., 1280-dim) which is injected into
            every ResNet block in the U-Net.
        \item \textbf{Context Stream:} The text embeddings
            $\tau_\theta(\vc)$ from Component 2.
    \end{enumerate}

\item \textbf{Conditioning Mechanism (Cross-Attention):} The
    interaction between the visual content and the text prompt is
    handled via Cross-Attention layers interleaved throughout the
    U-Net. Let $\phi_i(\z_t) \in \mathbb{R}^{d_\epsilon^{(i)} \times
    N}$ denote the flattened spatial features at U-Net layer $i$
    (where $N = h_i \times w_i$ is the number of visual tokens and
        $d_\epsilon^{(i)}$ is the channel depth, e.g., 320, 640, or
    1280), and let $\tau_\theta(\vc) \in \mathbb{R}^{d_\tau \times
    M}$ denote the text embeddings.

    Since the visual dimension $d_\epsilon^{(i)}$ and text dimension
    $d_\tau$ differ, we project both into a common subspace of
    dimension $d$ (the attention head dimension) using learnable
    weights. We first define the Query ($\boldsymbol{Q}$), Key
    ($\boldsymbol{K}$), and Value ($\boldsymbol{V}$) matrices:
    \begin{equation}
        \boldsymbol{Q} = \boldsymbol{U}^{(i)}_{Q} \phi_i(\z_t), \quad
        \boldsymbol{K} = \boldsymbol{U}^{(i)}_{K} \tau_\theta(\vc), \quad
        \boldsymbol{V} = \boldsymbol{U}^{(i)}_{V} \tau_\theta(\vc),
    \end{equation}

    where $\boldsymbol{U}^{(i)}_{Q} \in \mathbb{R}^{d \times
    d_\epsilon^{(i)}}$ and $\boldsymbol{U}^{(i)}_{K},
    \boldsymbol{U}^{(i)}_{V} \in \mathbb{R}^{d \times d_\tau}$.
    Adapting the column-wise formulation from
    \eqref{sub:image_classification_architecture}, the
    cross-attention operation is computed as:
    \begin{equation}
        \text{CrossAttn}(\boldsymbol{Q}, \boldsymbol{K},
        \boldsymbol{V}) \doteq \boldsymbol{V} \cdot
        \operatorname{softmax}\left( \frac{\boldsymbol{K}^\top
        \boldsymbol{Q}}{\sqrt{d}} \right),
    \end{equation}
    where the $\operatorname{softmax}$ is defined column-wise as
    \begin{align}
        \operatorname{softmax}(\boldsymbol{A})
        &\doteq
        \begin{bmatrix} \operatorname{softmax}(\boldsymbol{a}_{1}) &
            \cdots & \operatorname{softmax}(\boldsymbol{a}_{N})
        \end{bmatrix}, \\
        \forall \boldsymbol{A}
        &=
        \begin{bmatrix} \boldsymbol{a}_{1}, \dots, \boldsymbol{a}_{N}
        \end{bmatrix} \in \mathbb{R}^{M \times N}.
    \end{align}
    In this context, the matrix $\boldsymbol{A} =
    \frac{\boldsymbol{K}^\top \boldsymbol{Q}}{\sqrt{d}}$ represents
    the alignment scores, and the column-wise softmax normalizes the
    attention weights over the $M$ text tokens for each of the $N$
    visual regions. This mechanism allows the model to spatially
    align specific visual regions (e.g., ``cat's eye'') with specific
    text tokens.
\end{itemize}

\paragraph{Training Objective.}
The training process unifies these components. We sample a random
image $\x$, compress it using the frozen Encoder $\z_0 =
\mathcal{E}(\x)$, and add Gaussian noise $\epsilon$ to produce
$\z_t$. The U-Net $\epsilon_\theta$ is then trained to predict this
noise, conditioned on the text $\vc$ and timestep $t$:
\begin{equation}
\mathcal{L}_{\text{LDM}} = \mathbb{E}_{\mathcal{E}(\x), \vc, \epsilon
\sim \mathcal{N}(0,1), t} \left[ \| \epsilon - \epsilon_\theta(\z_t,
t, \tau_\theta(\vc)) \|^2 \right].
\end{equation}
During inference, the Encoder is discarded. The process starts from
pure noise $\z_T \sim \mathcal{N}(\bm{0}, \I)$, iteratively denoises
it using the U-Net, and finally passes the resulting $\hat{\z}_0$
through the Decoder $\mathcal{D}$ to synthesize the image.

\paragraph{Evaluation Metrics.}
Since conditional generation lacks a single ``correct'' ground truth,
we rely on statistical metrics to evaluate performance. We denote the
distribution of real images as $p_r$ and the distribution of
generated images as $p_g$.

\begin{itemize}
\item \textbf{FID (Fréchet Inception Distance):} This metric
    quantifies the realism and diversity of generated samples by
    measuring the distance between $p_r$ and $p_g$ in the feature
    space of a pre-trained Inception-v3 network. We approximate these
    feature distributions as multidimensional Gaussians
    $\mathcal{N}(\bm{\mu}_r, \bm{\Sigma}_r)$ and
    $\mathcal{N}(\bm{\mu}_g, \bm{\Sigma}_g)$. The FID score is then
    defined as the Fréchet distance between them:
    \begin{equation}
        \text{FID}(p_r, p_g) = \| \bm{\mu}_r - \bm{\mu}_g \|^2 +
        \text{Tr}\left( \bm{\Sigma}_r + \bm{\Sigma}_g -
        2(\bm{\Sigma}_r \bm{\Sigma}_g)^{1/2} \right).
    \end{equation}
    A lower FID indicates that the generated distribution is
    statistically closer to the real distribution, implying higher
    image quality and diversity.

\item \textbf{CLIP Score:} This metric measures the semantic
    alignment between a generated image $\hat{\x}$ and its
    conditioning prompt $\vc$. Let $E_{\text{img}}$ and
    $E_{\text{txt}}$ be the image and text encoders of a pre-trained
    CLIP model as discussed in Chapter
    \ref{ch:conditional-inference}. The score is calculated as the
    cosine similarity between their embeddings:
    \begin{equation}
        \text{CLIP-Score}(\hat{\x}, \vc) = \max \left( 0,
            \frac{E_{\text{img}}(\hat{\x}) \cdot
            E_{\text{txt}}(\vc)}{\|E_{\text{img}}(\hat{\x})\|
        \|E_{\text{txt}}(\vc)\|} \right) \cdot w,
    \end{equation}
    where $w$ is a scaling factor (typically 2.5). A higher score
    indicates that the visual content faithfully reflects the textual
    description.
\end{itemize}

\paragraph{Advanced Text Conditioned Architectures (SDXL and Flux).}
To address limitations in resolution and prompt adherence, advanced
architectures have evolved beyond the basic LDM.
SDXL~\citep{podell2023sdxl} introduces a dual-encoder strategy to
capture richer textual nuances, computing embeddings from both CLIP
ViT-L and OpenCLIP ViT-bigG and concatenating them along the channel
axis. Regarding data handling, SDXL employs \textit{Aspect Ratio
Bucketing}, grouping training images of similar aspect ratios into
batches to prevent the model from learning to crop subjects
arbitrarily. Furthermore, it implements \textit{Micro-Conditioning},
where metadata such as image resolution and crop coordinates are
embedded and concatenated with the text condition $\vc$, allowing the
model to explicitly distinguish between centered and cropped subjects.

Meanwhile, Flux~\citep{flux2024} represents a paradigm shift from the
U-Net architecture to a Multimodal Diffusion Transformer (MMDiT). It
processes text and visual tokens in a unified stream using Joint
Attention blocks, enabled by modern transformer components like
Rotary Positional Embeddings (RoPE) and RMSNorm. Mathematically, Flux
replaces standard diffusion with \textbf{Rectified Flow Matching}.
Unlike standard diffusion which predicts noise $\epsilon$, Flow
Matching simplifies the training target to predicting a straight-line
velocity field $v_\theta$ between noise and data. We define the
process over a continuous time $t \in [0, 1]$, where $t=0$ represents
the noise distribution ($\z_0$) and $t=1$ represents the data
distribution ($\z_1$).\footnote{Note that this time convention ($0
\to \text{Noise}, 1 \to \text{Data}$) is the inverse of the standard
diffusion notation used in previous sections.} The objective is to minimize:
\begin{equation}
\mathcal{L}_{\text{RFM}} = \mathbb{E}_{t, \z_1, \z_0} \left[ \|
v_\theta(\z_t, t, \vc) - (\z_1 - \z_0) \|^2 \right],
\end{equation}
where $\z_t = t \z_1 + (1-t) \z_0$ is the interpolated latent state.
The model $v_\theta$ learns to predict the constant velocity $(\z_1 -
\z_0)$ required to transport the noise $\z_0$ directly to the data
$\z_1$ along a straight trajectory. This linearity significantly
improves sampling efficiency and prompt adherence compared to the
curved stochastic paths of traditional diffusion models.

\subsection{Advanced Conditioning Mechanisms}
\label{subsec:advanced_conditioning}

\begin{figure}[t!]
\centering
\includegraphics[width=\textwidth]{\toplevelprefix/chapters/chapter8/figs/cnet_pipeline.pdf}
\caption{\textbf{ControlNet architecture pipeline.} A trainable copy
    of the Stable Diffusion (SD) U-Net encoder and middle blocks,
    conditioned on a hint image via the Hint Encoder
    $\mathcal{E}_{\text{hint}}$, injects feature maps into the frozen SD
    UNet decoder through zero convolutions to enable spatial control.
    (Image example from
\url{https://github.com/huggingface/diffusers/blob/main/docs/source/en/using-diffusers/controlnet.md}).}
\label{fig:controlnet_pipeline}
\end{figure}

While text conditioning provides semantic control, it struggles with
fine-grained spatial constraints (e.g., specific poses or edge maps).
To address this, ControlNet~\citep{zhang2023adding} introduces a
neural architecture that adds spatial control to large, pre-trained
text-to-image diffusion models as shown in Figure~\ref{fig:controlnet_pipeline}.

\textbf{Architecture and Connectivity.}
ControlNet operates by cloning the trainable parameters of a
pre-trained diffusion model. Specifically, it creates a trainable
copy of the 12 encoder blocks and 1 middle block of the Stable
Diffusion U-Net, while keeping the original parameters $\Theta$
completely locked (executed under \texttt{torch.no\_grad()}). This
preserves the generation quality learned from billions of images
while allowing the model to learn conditional tasks on small datasets.

To match the spatial resolution of the latent feature maps $\z \in
\mathbb{R}^{h \times w \times c}$ (typically $64 \times 64$ as
defined in Sec.~\ref{subsec:basic_realization}), the spatial
condition $\vc_s$ (e.g., a $512 \times 512$ Canny edge map) is first
processed by a lightweight \textbf{Hint Encoder}
$\mathcal{E}_{\text{hint}}(\cdot)$. Unlike the VAE, this encoder is
deeper but narrower, consisting of \textbf{eight} convolution layers
with $3 \times 3$ kernels and SiLU activations. It uses three
stride-2 layers to downsample the resolution, gradually increasing
channels ($16 \to 32 \to 96 \to 256$) before finally projecting to
the model dimension (320 channels). Crucially, the final layer of
this encoder is zero-initialized.

The critical innovation allowing these two streams to merge is the
\textbf{Zero Convolution}—a $1 \times 1$ convolutional layer
$\mathcal{Z}(\cdot; \theta_z)$ initialized with both weights $W$ and
bias $B$ set strictly to zero. The connectivity follows a specific pattern:
\begin{enumerate}
\item \textbf{Input Injection:} The encoded hint
    $\mathcal{E}_{\text{hint}}(\vc_s)$ is added \textbf{only} to the
    first input block of the ControlNet copy.
\item \textbf{Output Injection:} The output of each trainable block
    in the ControlNet copy is passed through a Zero Convolution and
    then added to the corresponding \textbf{skip connection} of the
    frozen Stable Diffusion Decoder.
\end{enumerate}

Formally, consider the $i$-th block of the U-Net. Let
$\boldsymbol{h}_i$ denote the input feature map to this layer. We
abstract the operation of this frozen block as a function
$\mathcal{F}_i(\cdot; \Theta_i)$. Accordingly, let
$\mathcal{F}_{\text{trainable}}^{(i)}(\cdot; \Theta_{c,i})$ denote
its trainable copy in the ControlNet branch. The output
$\boldsymbol{y}_i$ is computed as:
\begin{equation}
\boldsymbol{y}_i = \mathcal{F}_i(\boldsymbol{h}_i; \Theta_i) +
\mathcal{Z}\left(
    \mathcal{F}_{\text{trainable}}^{(i)}(\boldsymbol{h}_i +
\mathcal{E}_{\text{hint}}(\vc_s); \Theta_{c,i}); \theta_z \right).
\label{eq:controlnet}
\end{equation}

Because the Zero Convolution is initialized to zero, the control term
vanishes at the start of training ($\mathcal{Z}(\cdot) = 0$),
ensuring the model behaves exactly like the original Stable
Diffusion. However, the gradients for the weights $W$ are non-zero:
\begin{equation}
\frac{\partial \mathcal{Z}(I; \{W, B\})}{\partial W_{i,j}} = I_{p,i} \neq 0,
\end{equation}
where $I$ is the input feature map. This ensures that although the
initial influence is zero, learning commences immediately in the
first backward pass.

\textbf{Training Dynamics and Experiments.}
ControlNet is trained to minimize the standard diffusion objective,
but with the additional spatial condition $\vc_s$ injected into the
noise predictor. The specific learning objective is:
\begin{equation}
\mathcal{L} = \mathbb{E}_{\z_0, t, \vc_t, \vc_s, \epsilon \sim
\mathcal{N}(0,1)} \left[ \| \epsilon - \epsilon_\theta(\z_t, t,
\vc_t, \vc_s) \|^2 \right],
\end{equation}
where $\vc_t$ represents the text prompt conditioning and $\z_t$ is
the noisy latent. A unique phenomenon observed is ``sudden
convergence'': the model initially ignores the spatial condition,
then suddenly snaps to alignment as the error drops sharply. To force
the model to rely on structural signals rather than just the text
(which might contradict the edge map), a crucial strategy involves
\textbf{CFG-Dropout}: randomly replacing 50\% of the text prompts
$\vc_t$ with empty strings during training. This compels the network
to extract semantic meaning directly from the spatial condition
$\vc_s$. Empirically, ControlNet is highly efficient, requiring only
$\sim 23\%$ more GPU memory compared to standard fine-tuning, yet
effectively learning robust control from datasets as small as 1,000 samples.

\paragraph{Subject-Driven Personalization.}
While foundational models possess a vast semantic prior, they lack
the ability to synthesize specific subject instances (e.g., a user's
specific cat) consistently across different contexts. To address
this, DreamBooth \citep{ruiz2023dreambooth} introduces a
``personalization'' fine-tuning protocol. The core challenge is
implanting a new subject into the model's output domain without
``catastrophic forgetting''—where the model loses its understanding
of the general class (e.g., forgetting what a generic cat looks like)
or suffers from ``language drift'' (overwriting the word ``cat'' with
the specific cat's features).

\begin{figure}[t!]
\centering
\includegraphics[width=\textwidth]{\toplevelprefix/chapters/chapter8/figs/db_pipeline.pdf}
\caption{\textbf{The training and inference pipeline of DreamBooth.}
    During \textit{subject-driven finetuning}, the text-to-image model is
    updated using a combined loss: a reconstruction loss to learn the
    specific subject (e.g., ``[V] cat'') and a prior preservation loss to
    maintain generic class knowledge (e.g., ``a cat''). At
    \textit{inference}, the personalized model can generate the specific
    subject in novel contexts (e.g., ``laying on a wooden table'') using
the unique identifier.}
\label{fig:dreambooth_pipeline}
\end{figure}

\textbf{Identifier and Data Strategy.}
Unlike standard training which requires millions of pairs, DreamBooth
operates on a ``few-shot'' dataset provided by the user, typically
containing just 3--5 images of the subject. To bind this subject to
the model's text space, a unique identifier $[V]$ is required. The
authors propose a rare-token selection strategy: rather than using
random characters (which might be tokenized into common letters with
strong priors), the method searches the tokenizer's vocabulary for
rare token sequences (e.g., in the T5-XXL tokenizer range $\{5000,
\dots, 10000\}$) that have weak semantic associations. This
identifier is always paired with a coarse Class Noun (e.g., ``a $[V]$
cat'') in the prompt. This ``tethering'' is crucial: it leverages the
model's strong prior for the class (knowing how to render a ``cat''
in various poses) while binding the specific identity features to the
identifier $[V]$.

\textbf{Prior Preservation Training.}
To prevent the model from overfitting to the few-shot images and
losing diversity, DreamBooth employs an autogenous class-specific
prior preservation loss. As illustrated in Figure
\ref{fig:dreambooth_pipeline}, the training process minimizes a
combined objective:
\begin{equation}
\mathcal{L}_{\text{DB}} = \mathbb{E}_{\z_t, \vc, \epsilon, t} \left[
    \underbrace{\| \epsilon - \epsilon_\theta(\z_t, t, \vc_{\text{subj}})
    \|^2}_{\text{Reconstruction Loss}} + \lambda \underbrace{\| \epsilon
    - \epsilon_\theta(\z'_t, t, \vc_{\text{prior}}) \|^2}_{\text{Prior
Preservation Loss}} \right].
\label{eq:dreambooth}
\end{equation}
The first term fine-tunes the model weights (typically the entire
U-Net and text encoder) to reconstruct the subject images conditioned
on the unique prompt $\vc_{\text{subj}}$ (e.g., ``a $[V]$ cat''). The
second term acts as a regularizer: it supervises the model with its
own generated samples of the generic class $\vc_{\text{prior}}$
(e.g., ``a cat''). This forces the model to retain its original
understanding of the class distribution, ensuring that the identifier
$[V]$ learns the specific subject nuances while the class noun
``cat'' remains generic.

\textbf{Multi-Concept Personalization.}
While DreamBooth is effective, fine-tuning the entire model is
computationally expensive and storage-heavy. Custom Diffusion
\cite{kumari2023multi} improves efficiency by updating only the key
($W_K$) and value ($W_V$) projection matrices in the cross-attention
layers to demonstrate that these specific weights are sufficient to
encode new concepts.

\begin{figure}[t!]
\centering
\includegraphics[width=\textwidth]{\toplevelprefix/chapters/chapter8/figs/gen4gen_teaser.pdf}
\caption{\textbf{The Gen4Gen pipeline for multi-concept
    personalization.} To overcome the limitations of existing datasets in
    multi-subject scenarios, Gen4Gen \cite{yeh2024gen4gen} automates the
    creation of a high-quality benchmark (MyCanvas). Training on this
    data significantly boosts the model's ability to generate complex
    scenes with multiple specific subjects (e.g., ``a $[V1]$ cat and a
$[V2]$ dog'').}
\label{fig:gen4gen_pipeline}
\end{figure}

However, accurately synthesizing ``multiple'' personalized concepts
(e.g., ``a $[V1]$ cat and a $[V2]$ dog'') within a single image
remains difficult due to the scarcity of high-quality, text and image
aligned training data for complex compositions during the pretraining
stage of Stable Diffusion . To addressing this, Gen4Gen
\cite{yeh2024gen4gen} proposes a data-centric solution shown in
Figure \ref{fig:gen4gen_pipeline}. Instead of modifying the model
architecture, it introduces a generative data pipeline that leverages
foundation models to construct ``MyCanvas''—a benchmark dataset
featuring complex multi-object scenes paired with detailed, dense
captions. By fine-tuning on this synthesized data, where object
identities and spatial relationships are explicitly curated, models
can significantly improve their ability to disentangle and preserve
multiple personalized subjects simultaneously (e.g., a specific cat
and dog interacting in a complex scene) without suffering from
identity blending or omission.

More recently, the field has shifted toward Tuning-Free architectures
like PhotoMaker \cite{li2024photomaker} and InstantID
\cite{wang2024instantid}. Instead of iterative optimization, these
methods employ a ``stacking'' strategy: an external ID-Encoder (often
based on face recognition backbones) extracts a high-fidelity
identity embedding from the reference image, which is then injected
directly into the diffusion process via decoupled cross-attention or
prompt stacking. This enables zero-shot personalization of arbitrary
subjects in a single forward pass without any gradient updates.

\subsection{Extension to Video Generation}
\label{subsec:video_gen}

The principles of conditioned 2D generation extend naturally to video
generation by treating video as a sequence of temporally correlated
frames $\boldsymbol{x} \in \mathbb{R}^{T \times H \times W \times
3}$. However, directly modeling this high-dimensional space is
computationally prohibitive. Consequently, modern state-of-the-art
approaches typically operate within the compressed latent space
defined by the VAE (Sec.~\ref{subsec:basic_realization}), where the
video is represented as a tensor $\boldsymbol{z} \in \mathbb{R}^{T
\times h \times w \times c}$.

\paragraph{Spatiotemporal Inflation.}
To leverage the robust visual priors established by 2D image
generators, architectures such as VideoLDM~\citep{blattmann2023align}
and AnimateDiff~\citep{guo2023animatediff} introduce the concept of
``inflating'' 2D layers into pseudo-3D layers. This is achieved by
reshaping the input tensor $\boldsymbol{z}$ to alternate between
spatial and temporal processing:
\begin{itemize}
\item \textbf{Spatial Layers (Frozen):} The 2D backbone (ResNet and
    Spatial Attention) processes frames independently by reshaping
    the input to $(B \cdot T, h, w, c)$. This preserves the
    pre-trained texture and object generation capabilities.
\item \textbf{Temporal Layers (Trainable):} To model motion, the
    tensor is reshaped to isolate the temporal axis for each spatial
    location. We treat the spatial grid points ($B \cdot h \cdot w$)
    as the batch dimension and time ($T$) as the sequence length. A
    \textbf{Temporal Self-Attention} layer is then applied to
    exchange information across frames.
\end{itemize}

Mathematically, let $\boldsymbol{u} \in \mathbb{R}^{c \times T}$
denote the feature vector at a specific spatial pixel location across
time (analogous to the spatial features $\phi_i(\z_t)$ in
Sec.~\ref{subsec:basic_realization}). We define the temporal Query
($\boldsymbol{Q}$), Key ($\boldsymbol{K}$), and Value
($\boldsymbol{V}$) matrices using learnable projection weights
$\boldsymbol{U}^{(temp)}$:
\begin{equation}
\boldsymbol{Q} = \boldsymbol{U}^{(temp)}_Q \boldsymbol{u}, \quad
\boldsymbol{K} = \boldsymbol{U}^{(temp)}_K \boldsymbol{u}, \quad
\boldsymbol{V} = \boldsymbol{U}^{(temp)}_V \boldsymbol{u},
\end{equation}
where $\boldsymbol{U}^{(temp)}_{Q}, \boldsymbol{U}^{(temp)}_{K},
\boldsymbol{U}^{(temp)}_{V} \in \mathbb{R}^{c \times c}$. Following
the column-wise formulation established in
Sec.~\ref{subsec:basic_realization}, the temporal self-attention is computed as:
\begin{equation}
\text{TempAttn}(\boldsymbol{Q}, \boldsymbol{K}, \boldsymbol{V})
\doteq \boldsymbol{V} \cdot \operatorname{softmax}\left(
\frac{\boldsymbol{K}^\top \boldsymbol{Q}}{\sqrt{c}} \right) + \boldsymbol{u}.
\end{equation}
Here, the term $\boldsymbol{K}^\top \boldsymbol{Q} \in \mathbb{R}^{T
\times T}$ represents the temporal attention map, relating every
frame to every other frame. Crucially, the output projection layer of
this attention block is typically zero-initialized. This ensures that
at the start of training, the contribution of the temporal layers
vanishes, allowing the model to behave exactly like the pre-trained
2D generator and gradually learn temporal dynamics.

\paragraph{Interactive World Models.}
Recent advances have moved beyond passive text-to-video synthesis
toward interactive world simulation, effectively treating the video
generator as a predictive model of physical environments. Formally,
the goal is to model the latent transition probability conditioned on
an agent's action $\boldsymbol{a}_t$:
\begin{equation}
p(\boldsymbol{z}_{t+1} \mid \boldsymbol{z}_{1:t}, \boldsymbol{a}_t),
\end{equation}
where $\boldsymbol{a}_t$ represents the control signal (e.g.,
keyboard input or motor command) at time $t$. In frameworks like
Genie~\citep{bruce2024genie}, the action $\boldsymbol{a}_t$ is
discretized into a token and injected into the network using the same
conditioning mechanisms used for text. For instance,
$\boldsymbol{a}_t$ can be mapped to an embedding vector and added to
the timestep embedding or utilized as a key/value pair in
Cross-Attention layers. This allows the generative model to function
as a data-driven physics engine, predicting the visual consequences
of actions (e.g., a character jumping) by minimizing the
reconstruction or flow matching objectives. Large-scale
implementations like Sora~\citep{openai2024sora} further suggest that
scaling this autoregressive or diffusion-based prediction leads to
emergent 3D consistency and object permanence.

\section{Image and Text Conditioned 3D Object Generation}
\label{sec:Michelangelo}\label{sec:object-generation}
%\yima{Zibo and Sam, please move the image-text conditioned 3D shape
% generation material (Michelangelo) to this section. Try to be
% self-contained and consistent with the rest of the book for
% notation and terminologies, especially with the next Section for
% generation-based 3D reconstruction. }

So far, in previous sections, we have mainly shown how to learn
representations of the distribution of {\em two-dimensional} (2D)
images, as well as many image-based applications or tasks such
representations facilitate: such as image based object recognition,
image segmentation or completion, and image
generation.\footnote{Based on techniques introduced earlier: the
vision transformer (ViT)~\cite{NIPS2017_3f5ee243}, the language-image
contrastive learning (CLIP)~\cite{pmlr-v139-radford21a}, and latent
diffusion models (VAE-SD)~\cite{rombach2022high}.}

However, our world is three-dimensional (3D).\footnote{or
four-dimensional (4D) if we consider time as an extra dimension as
the world is dynamical.} Although our eyes perceive only  2D
projections of the world -- stereo pairs through our two eyes, we are
able to develop a visual memory that represents a full 3D world model
for our living environment and objects within.

This section demonstrates that, using the methods introduced in this
book, it is possible to learn the distribution of 3D shapes of
natural objects and then generate new 3D object shapes from this
distribution, conditioned on either an input 2D image or a text
description of the object. In particular, the work featured here is
based on the Michelangelo project~\cite{NEURIPS2023_ea1a7f7b}, which
very much generalizes the techniques that we have introduced for 2D
image representation and generation to 3D shape
data\footnote{Generally speaking, polygon meshes and video data can
all be referred to as ``3D data''; however, in this section, the term
3D data refers to polygon meshes or their equivalent representations,
explicit or implicit.}.

We first give a brief introduction to commonly used data forms that
represent 3D shapes and corresponding datasets. Then, following the
open-sourced project named \textbf{Michelangelo}  \cite{NEURIPS2023_ea1a7f7b}:\\
\centerline{\url{https://github.com/NeuralCarver/Michelangelo},}\\
we present a computational framework, objectives, and associated
network architectures that attempt to solve the task of image and
text conditioned 3D object generation. Finally, we present
experimental results and analysis to verify the proposed solution, as
well as discussion for further improvements.

% Introduce basic 3D representations and common used 3D dataset.
\subsection{3D Shape Representations, Datasets, and Tokenization}

\paragraph{3D Shape Representations}
\label{sec:3D representation}
Unlike 3D images, 3D shapes and data have not been formally
introduced before in this book. Therefore, this subsection provides a
brief introduction to 3D data.
Readers wishing to learn more details are advised to consult the
sources via the cited references.

3D data to which we refer here mainly mean continuous real physical
3D shapes that are typically digitized into discrete, computable
forms in a certain representation space. Finding good representations
of 3D shapes stands as one of the central problems in the fields of
computer vision and computer graphics. The four most commonly used,
and largely mathematically equivalent, types of 3D shape
representations are: point clouds, voxels, polygon meshes, and
implicit  functions,\footnote{such as the distance-sign function
(DSF).} as shown in
Figure~\ref{fig:3d_representations}~\footnote{Images in the first
three columns came from \cite{shi2022deep}, and the last column image
came from
https://medium.com/@sunil7545/implicit-vs-parametric-3d-shape-representation-9d4c01c8c60c}.

\begin{figure}[t!]
\centering
\includegraphics[width=\textwidth]{\toplevelprefix/chapters/chapter8/figs/3drep.pdf}
\caption{\textbf{Illustrations of different 3D representations for
    Stanford Bunny.} (a) the point cloud; (b) the voxel; (c) the polygon
mesh; (d) the implicit function.}
\label{fig:3d_representations}
\end{figure}

\textbf{Point clouds} precisely characterize the 3D shape of an
object through discretely sampled points on the object's surface.
They directly record 3D spatial coordinates $(x,y,z)$ of points on
the surface (and maybe some additional photometric information at
those points), as illustrated in
Figure~\ref{fig:3d_representations}(a) Point Cloud.
A point cloud is defined as a set of discrete coordinates in 3D Euclidean space:
\begin{equation}
\label{eq:pointcloud}
P = \{ \vp_i \in \mathbb{R}^d\ |\ i=1,2,\dots,N \}
\end{equation}
where $d \geq 3$ denotes the feature dimensionality, potentially including
(concatenations of) coordinates $(x_i,y_i,z_i)$, colors
$(r_i,g_i,b_i)$, and surface normals $(n_{xi},n_{yi},n_{zi})$.

\textbf{Voxels}
represent volumetric units in 3D space, constituting a natural
extension of the concept of 2D pixels to three dimension. Such a
representation is obtained by discretizing a 3D space into a structured grid.
Typical voxel models partition 3D space into uniform cubic cells,
where each voxel is indexed spatially and stores region-specific
attributes such as occupancy, density, material, or color, as shown
in Fig.~\ref{fig:3d_representations}(b) Voxel.

\textbf{Polygon meshes}
serve as arguably the most commonly used representation for 3D shapes
in computer graphics and physical simulation domains. It approximates
continuous surfaces through discrete  vertices, edges, and polygonal
faces, with triangle meshes being the predominant variant, as
depicted in Fig.~\ref{fig:3d_representations}(c) Polygon Mesh.
Geometrically, polygon meshes consist of vertex coordinates in $\mathbb{R}^3$,
complemented by their topological connectivity. Formally, a mesh $M$ is a tuple
$(V,F)$, where vertex set $V=\{\vv_i\in\mathbb{R}^3\}_{i=1}^n$
denotes the vertex coordinates and face set $F=\{f_j\}_{j=1}^m$
defines topology.

\textbf{Implicit functions}
represent surfaces via level sets of a scalar function. This approach
provides a unified framework for the mathematical description of
complex surfaces (Fig.~\ref{fig:3d_representations}(d) Implicit Function).
The standard implicit formulation is:
\begin{equation}
\label{eq:implicit_function}
F: \mathbb{R}^3 \rightarrow \mathbb{R}, \quad F(x,y,z) = d,
\end{equation}
where $(x,y,z)$ denote the coordinates of a query point, and $d$
denotes the distance between the query point and a surface of
interest.\footnote{Notice that we have used the same concept to
describe the low-dimensional support of a distribution in a
high-dimensional space, in \Cref{sec:manifold-distance}.} Hence the
surface of interest is implicitly defined as:
\begin{equation}
\label{eq:implicit_function_surface}
\mathcal{S} = \{ (x,y,z) \in \mathbb{R}^3 \mid F(x,y,z)=0. \}
\end{equation}
The conversion from an implicit function representation of 3D shape
to a polygon mesh can generally be achieved through iso-surface
extraction methods such as marching cubes~\cite{10.1145/280811.281026}.

Mathematically speaking, all four shape representations are largely
equivalent in terms of their representation capabilities.
Nevertheless, each data form may come with different computational
advantages or limitations. Hence, some representation is often
preferred for different tasks. In this section, we will adopt the
traditional point clouds representation and show how to learn the
distribution of 3D shapes based on this form of data. In the next
section, we will show how to use a (generalized) voxel representation
for learning the distribution of 3D shapes and poses.

\paragraph{3D Datasets.}
There have been many publicly available datasets for each of the
above representations mentioned above. They come with varying
quantities and qualities. In this section, we will use the
ShapeNet~\cite{chang2015shapenet} to study a joint shape-image-text
representation and its generation performance. ShapeNet provides
approximately 55,000 manufactured meshes across 55 categories.
Each mesh has a category tag and corresponding texts, such as
fine-grained categories or brief descriptions given by the creator.
To prepare the triplet data (3D shape, image, text), the provided
texts are augmented in two ways.
First, the shape tag and corresponding description are combined in
the format ``a 3D model of (shape tag), in the style of
(description)'' or ``a 3D model of (shape tag), (description).''
Then, inspired by ULIP~\cite{10203465}, multiple templates containing
65 predefined phrases are leveraged to provide more text information
during training.
As for the image data, each mesh is rendered under four camera poses,
augmenting and improving the rendering diversity via the
depth-conditioned ControlNet~\cite{10377881}.

\paragraph{3D Shape Data Tokenization.}
\label{sec:3d_tokenization}
As discussed in above, 3D data
manifests in various forms, including point clouds, voxels, polygon
meshes, and implicit neural representations. Fundamentally, point
clouds, voxels, and meshes serve as discrete approximations of
continuous 3D manifolds. In computer graphics and engineering,
polygon meshes are particularly ubiquitous because they explicitly
preserve surface topology, offering an efficient balance between
visual fidelity and geometric structure.

However, directly applying modern deep learning
architectures—specifically Transformers—to these raw representations
presents significant challenges. Unlike 2D images, which possess a
regular grid structure easily split into patches, 3D data is often
sparse, unstructured, and irregular. To leverage the power of
sequence modeling and self-attention mechanisms in 3D, we must bridge
the gap between raw geometric data and the discrete, ordered
sequences required by Transformers. This process is known as
\textit{3D Tokenization}.

\begin{itemize}
\item \textbf{Point Cloud Tokenization.}
    Raw point clouds are typically tokenized into local groups to capture
    geometric context. This process usually begins with Farthest
    Point Sampling (FPS)~\cite{eldar1997farthest},
    an iterative algorithm that selects a subset of points $\{x_1, \dots, x_m\}$
    from the original set $\mathcal{P}$, such that each selected point maximizes
    the distance to the set of previously selected points:
    $x_{i} = \arg\max_{x \in \mathcal{P}} \min_{j=1}^{i-1} \|x - x_j\|$.
    Following FPS, points are grouped via k-Nearest Neighbors (k-NN)
    to form local clusters.
    Each cluster serves as a token to handle the permutation
    invariance of points
    within a cluster, a lightweight network such as
    PointNet~\cite{qi2017pointnet} is applied.
    Mathematically, for a cluster of points $\{p_1, \dots, p_k\}$,
    PointNet approximates a symmetric function $f$ by applying a shared
    Multi-Layer Perceptron (MLP) $h$ followed by a symmetric
    aggregation operator (typically max-pooling):
    \begin{equation}
        \boldsymbol{t} = \gamma \left( \max_{i=1}^k \{ h(p_i) \} \right)
    \end{equation}
    where $\boldsymbol{t}$ is the resulting feature token and
    $\gamma$ is a post-processing MLP.
    This converts an unstructured cloud into a sequence of feature vectors.

    % {\color{blue} Update: We have added the implement details of
    % Farthest Point
    %     Sampling, and use math formula to describe the process of
    % PointNet, and the output token format. The following is the
    % original version, and is replaced by the above paragraph.}

    % Point clouds are typically
    %     tokenized using grouping techniques such as Farthest Point
    %     Sampling (FPS) \yima{what is this?} followed by k-Nearest
    % Neighbors (k-NN). Each
    %     cluster of points is treated as a token, often processed by a
    %     permutation-invariant networkwork (such as PointNet
    % \yima{reference? what does it do exactly, in mathematical
    % terms?} ) to extract a
    %     feature vector.

\item \textbf{Voxel Tokenization.}
    Analogous to the patch embedding strategy in Vision Transformers
    (ViT) for 2D images (see Chapter
    \ref{sub:contrastive_learning_architecture}), voxel grids can be
    partitioned into non-overlapping 3D sub-grids. Given a voxel grid
    $V \in \mathbb{R}^{D \times H \times W \times C}$, we extract
    patches of size $P \times P \times P$. Each patch is flattened
    into a vector $v_p \in \mathbb{R}^{P^3 \cdot C}$ and mapped via a
    learnable linear projection to the token dimension $d_{model}$.
    The result is a sequence of tokens $Z \in \mathbb{R}^{L \times
    d_{model}}$ where $L = \frac{DHW}{P^3}$. While straightforward,
    this method inherits the cubic complexity of voxel resolution
    ($O(N^3)$), often creating memory bottlenecks for high-resolution shapes.

    % Similar to 2D image patches, voxel
    %     grids can be partitioned into 3D sub-grids or ``cublets." While
    %     straightforward, this method often suffers from the cubic
    %     complexity of voxel resolution (the $O(N^3)$ memory bottleneck).

    %     Extending the 2D patch embedding strategy to 3D,
    % we partition the input voxel grid $V \in \mathbb{R}^{D \times H
    % \times W \times C}$ into
    % non-overlapping 3D patches or ``cubelets'' of size $p \times p \times p$.
    % These cubelets are flattened into 1D vectors and mapped via a
    % learnable linear projection
    % to the token dimension $d_{model}$.
    % The resulting output is a sequence of tokens with length $L =
    % \frac{DHW}{p^3}$.
    % However, this approach suffers from the cubic complexity of
    % voxel resolution,
    % as the sequence length $L$ scales with $O(N^3)$, causing
    % significant memory bottlenecks.
    % {\color{blue} Update: We have added the implement details of
    % the output tokens and the reference section about the patch
    % embedding strategy in Vision Transformers for 2D images.}

    %     \yima{Not clear at all, what are the output tokens? how are
    % they generated in what format? If this process is similar to
    % tokenizing 2D image, please refer to an earlier section that
    % describes tokenizing 2D image, if that is what you mean here.}

\item \textbf{Mesh Tokenization.}
    Tokenizing meshes is challenging due to their graph-like
    structure involving both vertices and faces.
    Unlike images, meshes lack a canonical order. To enable sequence
    modeling, common approaches involve
    explicit serialization, where raw geometric and topological
    elements are directly quantized and ordered into sequences.

    % Current methods typically serialize the mesh elements into a
    % sequence of discrete tokens to enable autoregressive modeling
    % (where the model predicts the next vertex/face token based on
    % previous ones to generate valid topology).

    % Prominent approaches like MeshGPT~\cite{siddiqui2024meshgpt}
    % serialize the mesh into two interleaved sequences: (1) a
    % **Vertex Sequence**, where vertex coordinates are quantized and
    % ordered (often via z-order curve sorting), and (2) a **Face
    % Sequence**, where faces are represented as indices pointing to
    % the vertex list. Alternatively, graph-based methods may treat
    % local mesh patches (geodesic neighborhoods) as tokens, encoding
    % them via Graph Convolutional Networks (GCNs) into feature vectors.

    Taking PolyGen~\cite{nash2020polygen} as an example, the mesh is
    represented as a set of
    vertices $V = \{v_1, \dots, v_n\}$ (where each $v_i \in
    \mathbb{R}^3$) and faces
    $F = \{f_1, \dots, f_m\}$ (where each $f_j$ is a tuple of vertex indices).
    The model learns the joint distribution $P(V, F)$ by factoring it
    into two sequential stages:
    a Vertex Model that predicts quantized coordinates, and a Face
    Model that predicts vertex indices
    \begin{equation}
        P(V, F) = P(V) \cdot P(F \mid V) = \prod_{i=1}^{n} P(v_i \mid
        v_{<i}) \cdot \prod_{j=1}^{m} P(f_j \mid f_{<j}, V).
    \end{equation}
    Crucially, this factorization dictates two distinct token formats
    corresponding to each stage.
    For the Vertex Model ($P(V)$), the tokens are quantized coordinate values,
    continuous coordinates are discretized (e.g., into 8-bit integers),
    creating a sequence of tokens from a fixed vocabulary of size 256.
    Conversely, for the Face Model ($P(F \mid V)$), the tokens are
    vertex indices (pointers).
    While effective, operating on these raw primitives often results
    in extremely long sequence lengths, limiting scalability.

    % {\color{blue} Update: We take the PolyGen as a representative
    % work and add the implement details of the mesh tokenization
    % process. The following is the original version, and is replaced
    % by the above paragraph.}
    % Tokenizing meshes presents unique
    %     challenges due to their irregular connectivity and graph-like
    %     structure. Common approaches involve traversing the surface to
    %     serialize faces or vertices into sequences, or decomposing the
    %     mesh into discrete primitives. These methods are particularly
    %     prevalent in autoregressive mesh generation \yima{what is
    % this, how it is related to the tokenization we are discussing
    % here}, where the objective
    %     is to generate not just the 3D geometry, but also a valid,
    % clean topology.\yima{This paragraph does NOT describe exactly
    % how to tokenize meshes at all, not even mention any existing method! }

\end{itemize}

While these explicit tokenization strategies effectively discretize
continuous manifolds, they operate directly in the high-dimensional
observation space, often resulting in excessively long sequence lengths
and computational redundancy. To address these scalability constraints
and achieve higher compression ratios, the field has gravitated towards
learning-based discrete representations. Instead of tokenizing raw
geometric primitives directly, these following methods first project the 3D data
into a compact, low-dimensional latent space where semantic information
is aggregated, effectively decoupling geometric reconstruction from
structural modeling.

\textbf{Standard Latent Vector Quantization.} This approach learns a discrete
codebook to represent the shape via Vector Quantization (VQ). A
representative method is  AutoSDF~\cite{mittal2022autosdf} In this
framework, a high-dimensional
input (e.g., a voxel grid $X$ or  implicit function field $F_\theta$)
is compressed by an encoder $E$ into
lower-dimensional spatial grid of latent vectors $Z \in \mathbb{R}^{h
\times w \times d \times c}$.
Each vector $z_i \in Z$ is then replaced by its nearest neighbor from
a learnable codebook $\mathcal{C}=\{c_k\}_{k=1}^K$:
\begin{equation}
q(z_i) = \arg\min_{c_k \in \mathcal{C}} | z_i - c_k |_2^2.
\end{equation}
This reduces the shape to a sequence of discrete codebook indices that preserve
spatial structure while drastically reducing dimensionality. These
discrete indices
serve as the tokens for subsequent Transformer processing. However,
this regular grid-based approach typically treats occupied and empty
space uniformly, leading to computational redundancy.

% {\color{blue} Update: We add more details of latent vector
% quantization process and wew specify that the 3D shape is converted
% to a sequence of tokens.}
% \yima{Then each 3D shape is converted to a single vector, instead
% of a sequence of tokens?}

\textbf{From Structural Optimization to Semantic Alignment.}
To address the redundancy of regular grids, approaches such as
3DILG~\cite{zhang20223dilg} proposes Irregular Latent Grids.
Instead of a dense voxel grid where tokens are wastefully
allocated to empty space, 3DILG defines the shape as a set
of latent codes floating at arbitrary, adaptive positions
in continuous space, typically concentrated on the shape's surface.
Specifically, given an input shape represented as a point cloud, the
method first encodes it into a set of feature-position pairs
$\mathcal{P} = \{(x_i, f_i)\}_{i=1}^M$, where $x_i \in \mathbb{R}^3$
are key points selected via Farthest Point Sampling (FPS) and $f_i
\in \mathbb{R}^C$ are their associated continuous feature vectors
extracted by a PointNet-based~\cite{qi2017pointnet} encoder. The
conversion of this unstructured set into a token sequence involves
three key steps:
\begin{itemize}
\item \textbf{Coordinate Quantization.} First, for a set of $M$ key
    points selected via Farthest Point Sampling (see Point Cloud
    Tokenization), the continuous spatial coordinates $x_i \in
    \mathbb{R}^3$ are discretized into 8-bit integers:
    \begin{equation}
        x_i \rightarrow (x_{i,1}, x_{i,2}, x_{i,3}), \quad
        \text{where } x_{i,k} \in {0, 1, \dots, 255}.
    \end{equation}
    This step effectively converts continuous positions into
    discrete, tokenizable units.

\item \textbf{Feature Quantization.} Concurrently, the local feature
    vector associated with each point is compressed into a discrete
    codebook index $z_i$ via Vector Quantization (VQ):
    \begin{equation}
        z_i \in \{0, 1, \dots, D-1\},
    \end{equation}
    where $D$ represents the size of the codebook.

\item \textbf{Sequence Construction.} To facilitate autoregressive
    generation with Transformers, the unstructured set of tuples must
    be serialized. 3DILG employs lexicographical sorting based on the
    quantized coordinates. The final shape is represented as a
    structured sequence $\mathcal{S}$ that explicitly couples
    positional coordinate and local geometric feature information:
    \begin{equation}
        \mathcal{S} = \big\{ (x_{0,1}, x_{0,2}, x_{0,3}, z_0), \dots,
        (x_{M-1,1}, x_{M-1,2}, x_{M-1,3}, z_{M-1}) \big\}.
    \end{equation}
\end{itemize}
This hybrid formulation drastically reduces sequence length compared
to dense grids while preserving high-fidelity surface details.

While 3DILG optimizes the spatial structure of tokens,
it largely remains a geometric compression method.
In the following, we provide a detailed exposition
of Michelangelo, a work that advances this paradigm
by addressing the semantic gap. It encodes 3D shapes
into a sequence of latent tokens that are not only
geometrically compact but also explicitly aligned
with the feature space of 2D vision-language models
(e.g., CLIP)~\cite{pmlr-v139-radford21a}), thereby
facilitating high-quality conditional generation.

% Introduce task and challenges
\subsection{Task}
Our goal here is to learn a conditional  3D shape generative model
capable of sampling 3D shapes from the learned distribution, subject
to a given condition, such as a given image or some texts, as
illustrated in Figure \ref{fig:michelangelo_network}.

Following the Latent Diffusion Model~\cite{rombach2022high}, to fully
leverage the capabilities of the generative model, a VAE is first
trained to compress 3D data samples into a latent space.
Generally, a good latent representation should possess two key
characteristics: it must be both highly compressible yet faithfully
reconstructible for 3D data samples, and conducive to generative model training.
However, compared to 2D image data, 3D data remains
scarce~\footnote{The situation persists even after the releases of
Objaverse-XL~\cite{10.5555/3666122.3667676}.}.
Consequently, learning an effective latent representation for 3D data
is highly challenging. Figure \ref{fig:michelangelo_network} shows
the overall diagram of the framework proposed by the Michelangelo
project. We will explain implementation details for each component below.
\begin{figure}[t!]
\centering
\includegraphics[trim=0cm 0cm 0cm 0cm,
clip=true,width=.99\linewidth]{\toplevelprefix/chapters/chapter8/figs/michelangelo_network.pdf}
\caption{
\textbf{Alignment-before-generation pipeline}. The approach
contains two models: the Shape-Image-Text-Aligned Variational
Auto-Encoder (SITA-VAE) and the Aligned Shape Latent Diffusion
Model (ASLDM). The SITA-VAE consists of four modules: an image
encoder, a text encoder, a 3D shape encoder, and a 3D shape
decoder. Encoders encode input pairs into an aligned space, and
the 3D shape decoder reconstructs 3D shapes given embeddings from
the aligned space. The ASLDM maps the image or text condition to
the aligned shape latent space for sampling a high-quality 3D
shape embedding, which is later reconstructed to high-fidelity 3D
shapes by the 3D shape decoder.
}
\label{fig:michelangelo_network}
\end{figure}

% Introduce network and loss func
\subsection{Architecture and Objective}

\paragraph{Shape-Image-Text Aligned Variational Auto-Encoder.}

Inspired by CLIP~\cite{pmlr-v139-radford21a}, the approach adopted by
the Michelangelo project injects semantic information into the latent
space via a contrastive loss to alleviate problems observed when
training VAEs exclusively from 3D data, which forms the
Shape-Image-Text-Aligned Variational Auto-Encoder (SITA-VAE).

The SITA-VAE contains four components: a pre-trained and fixed CLIP image
encoder $\mathcal{E}_{\mathrm{i}}$ and CLIP text encoder
$\mathcal{E}_{\mathrm{t}}$, a trainable 3D shape encoder
$\mathcal{E}_{\mathrm{s}}$ and neural field decoder $\mathcal{D}_{\mathrm{s}}$.
The CLIP image encoder and text encoder take 2D images $\vI \in \R^{H \times
W \times 3}$ and tokenized texts $\vT \in \R^{d_{{\mathrm{t}}} \times
L_{\mathrm{t}}}$ as input, and generate image tokens
$\vE_{\mathrm{i}} \in \R^{d \times (1 + L_{\mathrm{i}})}$ and text
tokens $\vE_{\mathrm{t}} \in \R^{d \times L_{\mathrm{t}}}$, where
$(1+L_{\mathrm{i}})$ and $L_{\mathrm{t}}$ are the sequence length of
image tokens $\vE_{\mathrm{i}}$ and text tokens $\vE_{\mathrm{t}}$.
% Advantage
The approach takes advantage of the pre-trained image encoder and
text encoder from CLIP. These two encoders are trained on large-scale
image-text pairs and are robust enough to capture a well-aligned
vision-language space, which enriches the semantics of the 3D shape
representation after multi-modal alignment via contrastive learning.

% 3D shape encoder
\textbf{3D shape encoder} aims to extract powerful feature
representations to characterize each 3D shape effectively. To achieve
this, point clouds $\vP \in \R^{(3 + C) \times N}$ are first sampled
from the surface of 3D shapes, where $N$ represents the number of
points, and $C$ denotes additional point features such as normal or
color.

To better capture high-frequency geometric details,
we utilize Fourier positional encoding to map the
low-dimensional spatial coordinates
$\boldsymbol{v} \in \mathbb{R}^3$
of each point into a higher-dimensional frequency space.
This encoding $\gamma(\boldsymbol{v})$ is defined as a
sequence of sinusoidal functions with exponentially
increasing frequencies:$$\gamma(\boldsymbol{v}) = [\sin(2^0\pi
\boldsymbol{v}), \cos(2^0\pi \boldsymbol{v}), \dots, \sin(2^{L-1}\pi
\boldsymbol{v}), \cos(2^{L-1}\pi \boldsymbol{v})],$$
where $L$ is the number of frequency bands.
To construct the final input representation, we explicitly
decompose the input point cloud $\vP \in \mathbb{R}^{(3+C)
\times N}$ into spatial coordinates $\boldsymbol{v}$ and
auxiliary features $\boldsymbol{f} \in \mathbb{R}^{C \times
N}$. The high-dimensional positional embeddings
$\gamma(\boldsymbol{v})$ are then concatenated with the
preserved features $\boldsymbol{f}$ and projected via a
learnable linear layer. Formally, the 3D shape encoder input
$\boldsymbol{X} \in \mathbb{R}^{d \times N}$ is obtained
as:$$\boldsymbol{X} = \text{Linear}([\gamma(\boldsymbol{v});
\boldsymbol{f}])$$where $[\cdot ; \cdot]$ denotes the
concatenation operation along the feature dimension.

% Next, a linear layer is used to project the concatenation of
% the Fourier positional encoded point clouds $\vP$ to the 3D shape
% encoder input $\vX \in \R^{d \times N}$.
% \yima{What does ``Fourier positional encoded'' mean? How is $\X$
% related to $\vP$, exactly?} {\color{blue} Update: We have added the
% formula for Fourier positional encoding and provided a mathematical
% expression describing the process of transforming from $\vP$ to $\X$.}

To process this dense feature sequence $\vX$ efficiently,
% \yima{Which data, $\vP$ or $\X$?} {\color{blue} Update: We have
% changed the expression to avoid the ambiguity.}
the encoder adopts a
Perceiver-based architecture~\cite{pmlr-v139-jaegle21a}. Unlike
standard Transformers that scale quadratically with input length,
the Perceiver employs a latent bottleneck mechanism to decouple
the processing complexity from the input size. Instead of tokenizing the
input points directly for self-attention, the model utilizes a fixed set
of small, learnable latent queries $\vQ \in \R^{(1+L_{\mathrm{s}}) \times
d}$, where $1 + L_{\mathrm{s}}$ is the number of query tokens.
Mechanistically, a Cross-Attention layer projects the high-dimensional
input space into this compact latent space. Formally, following the
formulation in Chapter ~\ref{subsec:basic_realization}, let
$\boldsymbol{U}^{(temp)}_{Q}, \boldsymbol{U}^{(temp)}_{K},
\boldsymbol{U}^{(temp)}_{V}$ be the learnable projection matrices.
The initial latent
representation $\vZ_0$ is computed as:
$$\vZ_0 = \text{Cross-Attention}
(\vQ, \vX) = (\boldsymbol{U}^{(temp)}_{V}\vX) \cdot
\text{softmax}\left(\frac{(\boldsymbol{U}^{(temp)}_{K}\vX)^\top(\boldsymbol{U}^{(temp)}_{Q}\vQ)}
{\sqrt{d_k}}\right).$$
Here, the input point cloud $\vX$ serves
as the Keys and Values, while the learnable queries $\vQ$ act as the
Queries. This operation compresses the variable-length input $N$ into a
fixed-length latent sequence $\vZ_0 \in \R^{(1+L_{\mathrm{s}}) \times
d}$.

In the context of Michelangelo, these queries are structured to
capture hierarchical semantics: they consist of one global head token
$\vQ_g \in \R^{1 \times d}$ with high-level semantics and
$L_{\mathrm{s}}$ local tokens $\vQ_l \in \R^{L_{\mathrm{s}} \times d}$
containing low-level geometric structure information. Then, several self-
attention blocks are used to iteratively improve the feature
representation and obtain the final shape embeddings, $\vE_{\mathrm{s}}
\in \R^{(1+L_{\mathrm{s}}) \times d}$.

% Drawing inspiration from
% previous transformer-based architectures for point cloud
% understanding~\cite{pmlr-v139-jaegle21a}, the 3D shape encoder is
% built on a perceiver-based transformer. Specifically, a
% cross-attention layer is used to inject the 3D shape information from
% the input $\vX$ into a series of learnable query tokens $\vQ \in
% \R^{d \times (1+L_{\mathrm{s}})}$, where $1 + L_{\mathrm{s}}$ is the
% length of query tokens $\vQ$, consisting of one global head token
% $\vQ_g \in \R^{d}$ with high-level semantics and $L_{\mathrm{s}}$
% local tokens $\vQ_l \in \R^{d \times L_{\mathrm{s}}}$ containing
% low-level geometric structure information. Then, several
% self-attention blocks are used to iteratively improve the feature
% representation and obtain the final shape embeddings,
% $\vE_{\mathrm{s}} \in \R^{d \times (1+L_{\mathrm{s}})}$.

% \yima{Probably more detail or illustration for the above
% description about how to tokenize 3D shape data, say point clouds,
% through a perceiver. How is it trained, separately or jointly, with
% what architecture and data etc.? How is it related to the ``3D
% tokenization'' methods introduced earlier?} {\color{blue} Update:
% We have supplemented the description of how the input is tokenized
% through the perceiver-based architecture, and the entire process is
% trained jointly. }

% Learn aligned latent
\textbf{Alignment among 3D shapes, images, and texts} plays a crucial
role in SITA-VAE and the conditional generative models. Since 3D data
is an order of magnitude smaller than images and text data, to learn
a better-aligned representation among 3D shapes, images, and texts,
the 3D shape encoder is constrained to be close to a pre-aligned
vision-language space, which is pre-trained on large-scale image-text
pairs with rich image and text representations by leveraging the
contrastive learning strategy. Consider an input triplet of 3D shapes
$\vX$, images $\vI$ and tokenized texts $\vT$. The triplet encoders
generate the corresponding shape embedding $\ve_{\mathrm{s}}$, image
embedding $\ve_{\mathrm{i}}$ and text-embedding $\ve_{\mathrm{t}}$ by
projecting the extracted shape tokens $\vE_{\mathrm{s}}$, image
tokens $\vE_{\mathrm{i}}$ and text tokens $\vE_{\mathrm{t}}$ as three
vectors with the same dimension, which is expressed as:
$\ve_{\mathrm{s}} = \mathcal{F}_{\mathrm{s}}(\vE_{\mathrm{s}}),
\ve_{\mathrm{i}} = \mathcal{F}_{\mathrm{i}}(\vE_{\mathrm{i}})$, and
$\ve_{\mathrm{t}} = \mathcal{F}_{\mathrm{t}}(\vE_{\mathrm{t}})$,
where $\mathcal{F}_{\mathrm{s}}$ is a learnable shape embedding
projector, image embedding projector $\mathcal{F}_{\mathrm{i}}$ and
text embedding projector $\mathcal{F}_{\mathrm{t}}$ are pre-trained
and frozen during training and inference. The contrastive loss is:

\begin{align}
&\mathcal{L}_{(\mathrm{s},\mathrm{i})} = - \frac{1}{2} \sum\limits_{(j,k)}
\left(
\log \frac{\exp (\ve_{\mathrm{s}}^j \ve_{\mathrm{i}}^k) }{
\sum\limits_l \exp (\ve_{\mathrm{s}}^j \ve_{\mathrm{i}}^l) } +
\log \frac{\exp (\ve_{\mathrm{s}}^j \ve_{\mathrm{i}}^k) }{
\sum\limits_l \exp (\ve_{\mathrm{s}}^l \ve_{\mathrm{i}}^k) }
\right), \\
&\mathcal{L}_{(\mathrm{s},\mathrm{t})}  = - \frac{1}{2} \sum\limits_{(j,k)}
\left(
\log \frac{\exp (\ve_{\mathrm{s}}^j \ve_{\mathrm{t}}^k) }{
\sum\limits_l \exp (\ve_{\mathrm{s}}^j \ve_{\mathrm{t}}^l) } +
\log \frac{\exp (\ve_{\mathrm{s}}^j \ve_{\mathrm{t}}^k) }{
\sum\limits_l \exp (\ve_{\mathrm{s}}^l \ve_{\mathrm{t}}^k) }
\right),
\end{align}
where $(j,k)$ indicates the positive pair in training batches, and
since pre-trained encoders from CLIP are utilized, the model is free
from the constraint $\mathcal{L}_{(\mathrm{i}, \mathrm{t})}$.

\textbf{3D shape decoder}, $\mathcal{D}_{\mathrm{s}}$, takes the
shape embeddings $\vE_{\mathrm{s}}$ as inputs to reconstruct the 3D
neural field with high quality. The KL divergence loss
$\mathcal{L}_{KL}$ is used to facilitate the generative process to
maintain the latent space as a continuous distribution.
% pre-fc, post-fc
Besides, a projection layer is leveraged to compress the latent from
dimension $d$ to lower dimensions $d_0$ for a compact representation.
Then, another projection layer is used to transform the sampled
latent from dimension $d_0$ back to high dimension $d$ for
reconstructing neural fields of 3D shapes.
% decode neural field
Like the encoder, the decoder model also builds on a transformer with the
cross-attention mechanism. Given a query 3D point $\vx \in \R^{3}$ in the field
and its corresponding shape latent embeddings $\vE_{\mathrm{s}}$, the decoder
computes cross attention iteratively for predicting the occupancy of the query
point $\mathcal{O} (\vx)$. The training loss is expressed as:

\begin{equation}
\mathcal{L}_r = \mathbb{E}_{\vx \in \R^3} [\text{BCE}(
\mathcal{D}(\vx \mid \vE_{\mathrm{s}} ), \mathcal{O}(\vx) )],
\end{equation}
where BCE is binary cross-entropy loss, and the total loss for
training Shape-Image-Text Aligned Variational Auto-Encoder (SITA) is written as:

\begin{equation}
\mathcal{L}_{SITA} = \lambda_c (\mathcal{L}_{(\mathrm{s},\mathrm{i})}
+ \mathcal{L}_{(\mathrm{s},\mathrm{t})}) + \mathcal{L}_r +
\lambda_{KL} \mathcal{L}_{KL}.
\label{eq:sita_vae}
\end{equation}

\paragraph{Aligned Shape Latent Diffusion Model.}
After training the SITA-VAE, an alignment space among 3D shapes,
images, and texts is obtained, as well as a 3D shape encoder and
decoder that compress the 3D shape into low-dimensional shape latent
embeddings and reconstruct shape latent embeddings to a neural field
with high quality. Building on the success of the Latent Diffusion
Model (LDM)~\cite{rombach2022high} in text-to-image generation, which
strikes a balance between computational overhead and generation
quality, a shape latent diffusion model is proposed on the aligned
space to learn a better probabilistic mapping from 2D images or texts
to 3D shape latent embeddings. By leveraging the alignment space and
the shape latent diffusion model, it is possible to generate
high-quality 3D shapes that better conform to the visual or textual
conditional inputs.

The Aligned Shape Latent Diffusion Model (ASLDM) builds on a
UNet-like transformer~\cite{10203178}, aiming to fit a distribution
of the shape latent embeddings, accompanied by an auto-encoder for
encoding data samples into the latent space and reconstructing the
data samples given the sampled latent. By learning in the latent
space, the latent diffusion model is computationally efficient, and
leveraging such a compact representation enables the model to fit the
target distribution faster. Specifically, the model
$\epsilon_{\theta}$ focuses on generating shape latent embeddings
$\vE_{\mathrm{s}}$ conditioned on $\vC$, which is represented by the
CLIP image or text encoder. Following LDM~\cite{rombach2022high}, the
objective is
\begin{equation}
\mathcal{L} = \mathbb{E}_{\vE_{\mathrm{s}}, \boldsymbol{\epsilon} \sim
\mathcal{N}(\Zero,\vI),t} \left[ \Vert \boldsymbol{\epsilon}
- \epsilon_\theta(\vE_{\mathrm{s}}^{(t)}, \vC, t) \Vert^2_2 \right],
\end{equation}
where $t$ is uniformly sampled from $\{1, ..., T\}$ and
$\vE_{\mathrm{s}}^{(t)}$ is a noisy version of
$\vE_{\mathrm{s}}^{(0)}$. During inference, after sampling Gaussian
noise, the model gradually denoises the signal until reaching
$\vE_{\mathrm{s}}^{(0)}$. Following classifier-free guidance
(CFG)~\cite{ho2021classifierfree}, the conditional latent diffusion
model is trained with classifier-free guidance. In the training
phase, the condition $\vC$ randomly converts to an empty set
$\emptyset$ with a fixed probability of $10\%$. Then, the sampling is
performed with the linear combination of conditional and unconditional samples:
\begin{equation}
\epsilon_\theta ( \vE_{\mathrm{s}}^{(t)}, \vC, t ) = \epsilon_\theta(
\vE_{\mathrm{s}}^{(t)}, \emptyset, t ) + \lambda
( \epsilon_\theta( \vE_{\mathrm{s}}^{(t)}, \vC, t ) -\epsilon_\theta(
\vE_{\mathrm{s}}^{(t)}, \emptyset, t ) ),
\label{eq:cfg-michelangelo}
\end{equation}
where $\lambda$ is the guidance scale for trading off the sampling
fidelity and diversity.

\subsection{Experiments and Analysis}

\textbf{Metrics.}
The Intersection of Union (IoU) is used to reflect the accuracy of
reconstructions. Two new metrics are proposed for evaluating 3D shape
generation methods. The first is a shape-image score (SI-S), which
uses a 3D shape encoder and image encoder to extract corresponding
shape and image embeddings and compute the Cosine Similarity of these
two modalities. Another is a shape-text score (ST-S), which computes
the similarity between the generated 3D shape and the conditional
text input in the aligned shape embedding and text embedding space.
Both metrics evaluate the similarity between results and their
corresponding conditions. Moreover, both the pre-trained
ULIP~\cite{10203465} and SITA are used to compute SI-S and ST-S, in
terms of SI-S (ULIP), ST-S (ULIP), SI-S (SITA) and ST-S (SITA),
respectively. Besides, the metrics of P-IS and P-FID as introduced in
Point-E~\cite{nichol2022point} are followed, using a pre-trained
PointNet++~\cite{NIPS2017_d8bf84be} to compute the point cloud
analogous Inception Score~\cite{NIPS2016_8a3363ab} and
FID~\cite{NIPS2017_8a1d6947} to evaluate the diversity and quality of
the generated 3D shapes.

\textbf{Baselines.}
We compare the Michelangelo approach to other methods, including
OccNet~\cite{8953655}, ConvOcc~\cite{10.1007/978-3-030-58580-8_31},
IF-Net~\cite{9156292}, 3DILG~\cite{NEURIPS2022_894ca1c4}, and
Learnable Query Version of 3DS2V~\cite{10.1145/3592442},  on
reconstruction tasks to validate the ability of the model to recover
a neural field given shape embeddings on the ShapeNet
dataset~\cite{chang2015shapenet}. For the conditional generation
stage, two recent powerful 3D generation methods, 3DILG and 3DS2V,
are chosen as baselines. Their shape representation modules are first
fine-tuned on a mixture dataset of ShapeNet and the 3D Cartoon
Monster. Then, the text and image conditional generative models of
3DILG and 3DS2V are retrained with the same protocols.

\textbf{Numerical comparison}.
The numerical results are reported in
Table~\ref{tab:michelangelo_recon} and Table~\ref{tab:michelangelo_gen}.
Table~\ref{tab:michelangelo_recon} shows that Michelangelo has the
best reconstruction performance on 55 overall categories, surpassing
the rest. Results of the selected categories further demonstrate that
the model can faithfully reconstruct 3D shapes in each of the 55 categories.
Table~\ref{tab:michelangelo_gen} reports the numerical results for
conditional 3D shape generation.
Michelangelo achieves the best performance on all the SI-S and ST-S
metrics, indicating that it can map the information from the image or
text to its corresponding 3D shape information for generating
high-fidelity results. Moreover, the P-FID demonstrates that the
model can produce high-quality shape-tokens for generating realistic
3D shapes, and P-IS indicates the diversity of the samples.
Specifically, the four left columns show that Michelangelo surpasses
the baselines on image-conditioned generation, demonstrating that it
can better map visual information to 3D shapes.
The four right columns validate the generative quality of
text-conditioned generation. Since natural language, compared to 2D
images, usually provides limited and abstract information, learning a
model to map text information to 3D shapes is challenging. However,
benefiting from training on the aligned latent space, Michelangelo
significantly improves text-conditioned generation, as shown in the
right columns of Table~\ref{tab:michelangelo_gen}, which reflects
that the model effectively maps natural language information to 3D
shapes and generates diverse and high-quality results.

\begin{table}[th]
\centering
\resizebox{1\columnwidth}{!}
{
\begin{tabular}{rcccccccc}
    \hline
    & Overall & Selected & Table & Chair & Airplane & Car & Rifle
    & Lamp   \\ \hline
    OccNet~\cite{8953655}
    & 0.825 & 0.81 & 0.823 & 0.803 & 0.835 & 0.911 & 0.755 & 0.735  \\
    ConvOccNet~\cite{10.1007/978-3-030-58580-8_31}
    & 0.888 & 0.873 & 0.847 & 0.856 & 0.881 & 0.921 & 0.871 & 0.859  \\
    IF-Net~\cite{9156292}
    & 0.934 & 0.924 & 0.901 & 0.927 & 0.937 & 0.952 & 0.914 & 0.914 \\
    3DILG~\cite{NEURIPS2022_894ca1c4}
    & 0.950 & 0.948 & 0.963 & 0.95 & 0.952 & 0.961 & 0.938 & 0.926 \\
    3DS2V(LQ)~\cite{10.1145/3592442}
    & 0.955 & 0.955 & \textbf{0.965} & 0.957 & 0.962 & 0.966 &
    0.947 & 0.931 \\
    Ours
    & \textbf{0.966} & \textbf{0.964} & \textbf{0.965} &
    \textbf{0.966} & \textbf{0.966} & \textbf{0.969} &
    \textbf{0.967} & \textbf{0.95} \\ \hline
\end{tabular}
}
\caption{\textbf{Numerical results for reconstruction comparison on
IoU($\uparrow$, a larger value is better)}. The results show that
Michelangelo performs best in 55 categories. The results of selected
categories further demonstrate that the model can reconstruct each
category faithfully.}
\label{tab:michelangelo_recon}
\end{table}

\begin{table}[th]
\resizebox{1\columnwidth}{!}{%
\centering
\begin{tabular}{rcccccccc}
    \hline
    & \multicolumn{4}{c}{Image-Conditioned} &
    \multicolumn{4}{c}{Text-Conditioned} \\
    \cmidrule(lr){2-5}\cmidrule(lr){6-9}
    & SI-S (ULIP)$\uparrow$ & SI-S (SITA)$\uparrow$   &
    P-FID$\downarrow$ & P-IS$\uparrow$ & ST-S (ULIP)$\uparrow$ &
    ST-S (SITA)$\uparrow$  & P-FID$\downarrow$ & P-IS$\uparrow$ \\ \hline
    3DILG
    & 9.134  & 11.703 & 4.592 & 12.247 & 10.293 & 6.878  & 10.283 & 12.921\\
    3DS2V
    & 13.289 & 15.156  & 2.921 & 12.92 & 12.934 & 9.833  & 5.704 & 13.149\\
    Ours
    & \textbf{13.818} & \textbf{15.206}  & \textbf{1.586} &
    \textbf{13.233} & \textbf{16.647} & \textbf{13.128}  &
    \textbf{2.075} & \textbf{13.558}\\ \hline
\end{tabular}
}
\caption{\textbf{Numerical results for conditional generation
comparison}. The results show that Michelangelo achieves the best
generative performance. The SI-S and ST-S indicate that the model
generates high-fidelity results by effectively mapping the condition
information to its related 3D shapes. Moreover, P-FID reflects that
the model generates the most realistic 3D shapes, and P-IS indicates
that the generated samples are diverse. $\uparrow$ means a larger
value is better, and $\downarrow$ otherwise.}
\label{tab:michelangelo_gen}
\end{table}

\textbf{Visual comparison}.
The visual comparisons of the image/text-conditioned 3D shape
generations are illustrated in
Figure~\ref{fig:michelangelo_img_cond_gen} and
Figure~\ref{fig:michelangelo_text_cond_gen}.
Figure~\ref{fig:michelangelo_img_cond_gen} shows that
3DILG~\cite{NEURIPS2022_894ca1c4} pays more attention to the global
shape in the auto-regressive generation process, where its results
lack depictions of details of 3D shapes. While
3DS2V~\cite{10.1145/3592442} generates more details of 3D shapes, the
results have discontinuous or noisy surfaces. Besides, both methods
struggle to generate a complete shape when the given conditions map
to a complex object, fine machine, or rare monster.
Figure~\ref{fig:michelangelo_text_cond_gen} shows the visual
comparison of text-conditional generation. In the upper-half rows,
the results given simple and abstract concepts are shown, while in
the lower-half rows, the results given detailed texts like
descriptions for deterministic parts of the target shape are shown.
Similar to the observation above, 3DILG~\cite{NEURIPS2022_894ca1c4}
generates over-smooth shape surfaces with fewer details, and
3DS2V~\cite{10.1145/3592442} produces fewer details on discontinuous
object surfaces. In contrast, Michelangelo produces correct shapes
that conform to the given concepts or detailed descriptions with
delicate details on smooth surfaces.

\begin{figure}[t!]
\centering
\includegraphics[width=.98\linewidth]{\toplevelprefix/chapters/chapter8/figs/michelangelo_img_cond_gen.pdf}
\caption{
\textbf{Visual results for image-conditioned generation
comparison}. The figure shows that
3DILG~\cite{NEURIPS2022_894ca1c4} generates over-smooth surfaces
and lacks details of shapes, whereas 3DS2V~\cite{10.1145/3592442}
generates few details with noisy and discontinuous surfaces of
shapes. In contrast to baselines, Michelangelo produces smooth
surfaces and portrays shape details. Please zoom in for more visual details.
}
\label{fig:michelangelo_img_cond_gen}
\end{figure}

\begin{figure}[t!]
\centering
\includegraphics[width=.98\linewidth]{\toplevelprefix/chapters/chapter8/figs/michelangelo_text_cond_gen.pdf}
\caption{
\textbf{Visual results for text-conditioned generation comparison}.
In the first two rows, the models are tested with abstract texts, and
the result shows that only Michelangelo generates a 3D shape that
conforms to the target text with a smooth surface and fine details.
The last two rows show the results given texts containing detailed
descriptions, which further demonstrates that Michelangelo can
capture both the global conditional information and the local
information for generating high-fidelity 3D shapes. Keywords are
highlighted in red; please zoom in for more visual details.}
\label{fig:michelangelo_text_cond_gen}
\end{figure}

\textbf{The effectiveness of training generative model in the aligned
space}. A visual comparison is performed to ablate the effectiveness
of training the generative model in the aligned space, as illustrated
in Figure~\ref{fig:michelangelo_ablate_align}. The upper samples are
generated from the generative model trained in the aligned space,
while the lower samples are generated from the generative model
trained without the aligned space. The results demonstrate that the
upper samples conform to the given text while the lower samples do
not, which indicates that training the generative model in the
aligned space leads to high-fidelity samples.

\begin{figure}[t!]
\centering
\includegraphics[width=.98\linewidth]{\toplevelprefix/chapters/chapter8/figs/michelangelo_1aligned_space_for_generation.pdf}
\caption{
\textbf{Ablation study the effectiveness of training generative
model in the aligned space}. This figure illustrates visual
comparisons for ablation studies on the effectiveness of training
the generative model in the aligned space. Compared with the
lower samples based on the conditional texts, the upper samples
are closer to the conditions semantically, which indicates the
effectiveness of the training generative model in the aligned space.
}
\label{fig:michelangelo_ablate_align}
\end{figure}

\textbf{The effectiveness of vision-language models}. In addition to
the well-known vision-language model (VLM)
CLIP~\cite{pmlr-v139-radford21a}, another vision-language model (VLM)
SLIP~\cite{10.1007/978-3-031-19809-0_30} is introduced for training
the SITA-VAE for a comprehensive comparison. First, the impact of the
vision-language model on SITA-VAE's reconstruction ability is
evaluated, and the results are shown in
Figure~\ref{fig:michelangelo_ablate_vlm}. The results show that the
model composed with CLIP achieves the best performance. Then, the
vision-language model's impact on the ability to align multi-modal
space is evaluated. Standard and zero-shot classification tasks are
selected to reflect the impact of the vision-language model. Note
that the classification is performed by a feature matching operation,
where multiple 3D shapes and phrases are provided to the SITA-VAE; it
returns the similarity between 3D shapes and each phrase as
classification results, which indicates that the more aligned the
multi-modal space is, the higher the classification accuracy. The
results show that the model composed with CLIP achieves the best performance.

\textbf{The impact of the learnable query embeddings}. An ablation
study of learnable query embeddings is performed with the same
experiments as above, and the results show that using 512 learnable
query embeddings leads to the best performance on reconstructions and
classifications.

\begin{figure}[t!]
\centering
\includegraphics[width=1.\linewidth]{\toplevelprefix/chapters/chapter8/figs/michelangelo_2aligned_space_for_generation.pdf}
\caption{
\textbf{Ablation study of the effectiveness of vision-language
models and the impact of learnable query embeddings}. This figure
shows the ablation study on the effectiveness of the
vision-language model and the impact of learnable query
embeddings. According to the table, the model with CLIP and 512
learnable query embeddings achieves the best reconstruction and
classification performance, indicating its ability to recover 3D
shapes and align multi-modal space.
}
\label{fig:michelangelo_ablate_vlm}
\end{figure}

\paragraph{Nearest Neighbor Analysis.}
To verify whether the model learns to generate results based on the
given conditions or memorizes the training split of the dataset, the
training set is traversed to retrieve the nearest neighbors of the
generated samples, as illustrated in
Figure~\ref{fig:michelangelo_nearestneighbor}. Specifically, three
nearest neighbors for each sample are exhibited for comparison.
According to the visualization, the generated 3D shapes differ from
each retrieved nearest neighbor, demonstrating that the model
hallucinates new 3D shapes rather than overfitting the training sets.

\begin{figure}[t!]
\centering
\includegraphics[width=1.\linewidth]{\toplevelprefix/chapters/chapter8/figs/michelangelo_nearestneighbor.pdf}
\caption{
\textbf{Nearest Neighbor Analysis}. We traverse the whole
training set to find three nearest neighbors for the generated 3D
shapes, and the results reveal that our model could produce novel
3D shapes based on given images instead of memorizing specific ones.
}
\label{fig:michelangelo_nearestneighbor}
\end{figure}

\paragraph{Discussion.}
Though the Michelangelo project has shown that we can achieve
excellent results in generating 3D objects, it still has some
limitations. First, it requires a large number of samples of ground
truth 3D shapes for training and learning the distribution, whereas
such 3D shape data are much more expensive to obtain and available
datasets are typically orders of magnitude smaller than those for 2D
images. As we have alluded to in \Cref{sec:3D-visual-model} at the
end of the previous chapter, our brain is able to learn a full 3D (or
4D) visual model by having only 2D images as such data are much
natural and economic to acquire. Hence, ultimately, to truly scale up
the generative 3D model, we have to develop more advanced methods
that would enable us to learn the shape representation (within a 3D
shape-image-text aligned space) from only (multi-view) 2D images, say
via differentiable rendering.
Furthermore, since we represent each 3D shape as an occupancy field,
it needs to convert the 3D mesh into a watertight one, which will
inevitably degrade the original quality of the 3D mesh.

\section{Generation-Based 3D Shape and Pose
Reconstruction}\label{sec:CUPID}\label{sec:shape-and-pose}

% refactor
\newcommand{\pose}{\boldsymbol{\theta}}
\newcommand{\object}{\mathcal{O}}
\newcommand{\Image}{\boldsymbol{I}}
\newcommand{\image}{\vI}
% \newcommand{\mask}{\boldsymbol{M}^{\textrm{occ}}}
\newcommand{\render}{\boldsymbol{\mathcal{P}}}
\newcommand{\latent}{\boldsymbol{z}}
\newcommand{\timestep}{t}
\newcommand{\latentt}{\boldsymbol{z}_{\timestep}}
\newcommand{\latentzero}{\boldsymbol{z}_{0}}
\newcommand{\noise}{\boldsymbol{\epsilon}}
\newcommand{\velocity}{\boldsymbol{v}}
\newcommand{\parameter}{\phi}
\newcommand{\encoder}{\varphi}
\newcommand{\point}{\boldsymbol{x}}
\newcommand{\pixcoord}{\boldsymbol{u}}
\newcommand{\feat}{\boldsymbol{f}}
\newcommand{\sflow}{\boldsymbol{\mathcal{G}_{S}}}
\newcommand{\lflow}{\boldsymbol{\mathcal{G}_{L}}}
\newcommand{\transpose}{\operatorname{T}}
\newcommand{\Grid}{\mG}
\newcommand{\sgrid}{\mS}
\newcommand{\numpts}{L}
% \newcommand{\res}{N}
\newcommand{\sres}{D}
\newcommand{\vaeencoder}{\psi}
\newcommand{\trellis}{{\textsc{TRELLIS}}\xspace}
% \newcommand{\dino}{\textsc{DINO}\xspace}
\newcommand{\interp}{\operatorname{BilinearInterp}}
\newcommand{\slatenc}{\operatorname{SlatEncoder}}
\newcommand{\twotothree}{$2$D $\to$ $3$D\xspace}
\newcommand{\threetotwo}{$2$D $\to$ $3$D\xspace}
\newcommand{\low}{\mathrm{l}}
\newcommand{\high}{\mathrm{h}}

% \paragraph{Generation-based 3D Reconstruction}
The fundamental reason why we develop the sense of vision is to learn
to perceive our three-dimensional (3D) environment. Our brain has
dedicated a significant portion of it, {\em the visual cortex}, to
perform this task. It enables us to accurately and efficiently
reconstruct shapes of 3D objects and their spatial
relationships\footnote{say relative poses relative to our viewpoint
or between the objects themselves.} from the 2D images perceived
through our eyes. There has been a rich and long history of
understanding mathematically how, in principle, we can fully
reconstruct 3D geometry from (multiple) 2D images. Interested readers
may refer to the textbook \cite{MaY2003} for a detailed account of
this important topic.

\begin{figure}[t!]
\centering
\includegraphics[width=0.8\linewidth]{chapters/chapter8/figs/cupid_definition.pdf}
\caption{\textbf{Task: to infer the 3D shape of an object with a
canonical pose relative to the camera pose from a single image.}
Left: Input image. Middle: Output recovers a canonical 3D model
of the object (with intrinsic properties such as shape and
texture) along with the exact camera pose (extrinsic property)
that generates the input.
Right: Re-rendering the 3D model from the estimated pose faithfully
reproduces the input.}
\label{fig:cupid_teaser}
\end{figure}

Nevertheless, as we have alluded to in the Section
\ref{sec:3D-visual-model} in the previous Chapter
\ref{ch:conditional-inference}, ideally, 3D reconstruction should be
done in a Bayesian manner as our brain does so too. Our perception of
the 3D environment seems effortless is mainly because we reconstruct
a 3D scene and recognize 3D objects based on a rich ``visual
memory.'' More precisely, we have learned through our years of
accumulated experience the distribution of out 3D environment and
objects within. The distribution is rather low-dimensional and
structured hence we can correctly and efficiently complete 3D
geometry of a scene and objects in it even if we have seen only small
part of it, say from a single view, as illustrated by the example in
Figure \ref{fig:cupid_teaser}. Conceptually, this ability to
``complete'' missing information is very similar to how natural
images can be completed from a small fraction of patches, as we have
seen in Section \ref{sec:image_completion}.

As a distribution, our 3D visual memory is not only low-dimensional
but also highly structured. For instance, it seems that our brain
stores semantics (identity of objects) of a scene in the {\em
inferior temporal} cortex (IT cortex) and the 3D spatial information
of the scene in the {\em hippocampus}. Our brain  not only
reconstructs the 3D geometry of the scene,\footnote{Say 3D geometry
that is represented in different but largely equivalent forms: point
clouds, meshes, signed distance functions, neural radiance fields, or
Gaussian splatters.} but also parses its content into individual
objects and their spatial relationships with the viewer and among the
objects themselves. Such a structured ``what-is-where''
representation enables us to conduct spatial inference within the
scene from either an {\em egocentric, object-centric, or allocentric}
perspective. These abilities are necessary and crucial for tasks such
as navigating in a scene, interacting and manipulating objects within.

Hence, if the goal of learning a model of the 3D scene is to support
physical interaction with an  environment and objects
within,\footnote{which is the objective of the so-called ``Embodied
AI'' for intelligent agents or robots.} we must learn a
representation of the 3D environment that shares the same properties
and characteristics as our visual memory mentioned above.

\subsection{Task and Objective}
\label{sub:cupid_objective}

\paragraph{Object 3D Shape and Poses from a Single 2D View}
In this section, we takes a first step to show how this can be
achieved by illustrating with a challenging and seemingly ill-posed
problem: {\em inferring the complete 3D shape and and pose, including
both ego-centric and object-centric pose, of an object from a single
2D image.} This problem, by its own, lies at the heart of computer
vision, aiming to reverse the physics of image formation. An image
$\image$ is the result of a rendering process,
\[
\image = \render(\object, \pose),
\]
where $\object$ represents the 3D object in a canonical,
view-agnostic frame and $\pose$ denotes the camera viewpoint. Given
only $\image$, the task of disentangling the object's intrinsic
properties (e.g., shape, texture) from the extrinsic camera pose is
fundamentally ambiguous.

Previous approaches typically oversimplify this ambiguity. Generative
models often prioritize the canonical object $\object$ at the expense
of the camera pose $\pose$, yielding reconstructions that fail to
align faithfully with the input view. Conversely, traditional
view-centric methods fix the pose, which limits their ability to
generate or ``in-paint'' occluded regions and seamlessly paste the
reconstructed object back into the original image.

Here we present an approach known as \textbf{Cupid}~\citep{huang2025cupid}:
\\
\centerline{\url{https://cupid3d.github.io},}\\
an open-source framework that jointly models the distribution over
both the 3D object and camera pose (Figure~\ref{fig:cupid_teaser}).
By simultaneously generating the 3D model and reasoning about the
viewpoint, Cupid ensures the output can be accurately \textit{pasted}
into the original image. This approach recasts single-view
reconstruction as a conditional sampling problem, achieving geometric
fidelity by combining observational data with learned generative
priors, thereby enabling efficient object- and scene-level
reconstruction (Figure~\ref{fig:cupid_results}).

\begin{figure}[t!]
\centering
\includegraphics[width=0.95\textwidth]{chapters/chapter8/figs/cupid/figures/results.pdf}
\caption{\textbf{Results for \textit{generative 3D reconstruction}
from a single test image.} Given an input image (top left), Cupid
estimates camera pose (bottom left) and reconstructs 3D model (bottom
right), re-rendering the input (top right). It is robust to changes
in scale, placement, and lighting while preserving fine details, and
supports component-aligned scene reconstruction (bottom row). All
results are produced in seconds via feed-forward sampling of the
learned model. }
\label{fig:cupid_results}
\vspace{-5mm}
\end{figure}

\paragraph{Task Formulation and Objective.}
We formulate generative 3D reconstruction as the task of estimating
the joint posterior distribution
\[
p(\object, \pose \mid \image)
\]
under the constraint that the observation $\image$ is a rendering of
the object $\object$ from pose $\pose$. To model this conditional
distribution, we employ a flow-based generative model. Specifically,
we first map the 3D object and camera pose into a unified volumetric
latent representation
\(\latent = \encoder(\object, \pose)\).
We then use a Rectified Flow~\citep{lipmanflow} model to learn the conditional
generation of this latent variable. The model defines a linear
interpolation between a sample from the data distribution,
$\latentzero$, and a sample from a noise distribution, $\noise$, over
a normalized time interval $\timestep \in [0,1]$:
\begin{equation}
\latentt = (1 - \timestep) \latentzero + \timestep \noise.
\end{equation}
The generative process involves reversing this trajectory by learning
a time-dependent velocity field
\[
\velocity(\latentt, \image, \timestep) = \nabla_{\timestep} \latentt
\]
that transports noisy samples towards the data manifold, conditioned
on the input image $\image$. We parameterize this velocity field with
a neural network $\velocity_{\parameter}$ and train it using the
Conditional Flow Matching (CFM) objective~\cite{lipmanflow}:
\begin{equation}
\mathcal{L}_{\mathrm{CFM}}(\parameter) \doteq \mathbb{E}_{\timestep,
\latentzero, \noise} \left\| \velocity_{\parameter}(\latentt, \image,
\timestep) - (\noise - \latentzero) \right\|_{2}^{2}.
\end{equation}
By sampling from this learned model, we can generate a latent code
$\latent$ which can then be decoded into the final 3D object and camera pose.

\subsection{3D Shape Representation and Tokenization}

A central challenge in 3D shape generation is finding a tokenization
that is both compact and expressive, capable of high-fidelity
reconstruction. While
\textbf{Michelangelo}~\citep{zhao2023michelangelo} utilizes a learned
latent vector representation—encoding shapes into a sequence of 1024
tokens—these vectors lack explicit 3D spatial coordinates. In this
project, we seek a latent representation that retains explicit 3D
positional information, allowing us to align shape tokens directly
with image pixels. By establishing this spatial correspondence, we
can leverage local pixel information to enhance generation details.
Consequently, \textbf{Cupid} adopts a 3D voxel-based latent
representation and tokenization approach as described in
Section\ref{sec:3d_tokenization}. Here, we represent a shape with
sparse voxels accompanied by features, which are usually defined by
RGB color or more advanced DINO features~\citep{DINOv2} from 2D
foundation models. We then train a decoder to map such a
representation back to the shape, parameterized by SDF, and to the
apperance parameterized by the Gaussian splats~\citep{kerbl20233d}.

Specifically, given a set of sparse voxel features $\boldsymbol{F} =
\{\boldsymbol{f}_i\}_{i=1}^N$ located at coordinates
$\boldsymbol{x}_i$, we employ two distinct decoders,
$\mathcal{D}_{geo}$ and $\mathcal{D}_{app}$, to disentangle geometry
from appearance. The decoding process is formulated as:
\begin{equation}
s_i = \mathcal{D}_{geo}(\boldsymbol{f}_i), \quad \mathcal{G}i =
\mathcal{D}_{app}(\boldsymbol{f}_i),
\end{equation}
where $s_i \in \mathbb{R}$ represents the Signed Distance Field (SDF)
value used for surface reconstruction, and $\mathcal{G}_i =
\{\boldsymbol{c}_i, \alpha_i, \boldsymbol{s}_i, \boldsymbol{q}_i\}$
denotes the set of 3D Gaussian parameters—comprising color, opacity,
scaling, and rotation—required for volumetric rendering. We will then
use the sparse voxel features as the latent representation for latent
flow matching~\citep{rombach2022high}.

\subsection{Architecture and Implementation}
\label{sub:3d_reconstruction_architecture}
Since we utilize 3D voxels for object representation, applying
diffusion naively results in excessive memory consumption, as a large
portion of the grid consists of unnecessary empty space, as
Section~\ref{sec:3d_tokenization} describes. To mitigate this, the
\textbf{Cupid} architecture is designed as a cascaded, two-stage flow
model that progressively refines the 3D
reconstruction~\citep{xiang2025structured}. The first stage generates
a coarse geometric shape and estimates the camera pose, while the
second stage generates high-fidelity geometry and appearance. A key
innovation of this framework lies in the explicit representation of
the camera pose and its integration into the generative process. The
overall pipeline is depicted in \Cref{fig:cupid_architecture}.

\paragraph{Decoupled Object and Pose Representation.}

As discussed in Section~\ref{sec:3D representation}, the 3D object
$\object$ is tensorized into a sparse voxel-based representation:
\begin{equation}
\object \triangleq \{\point_i, \feat_i\}_{i=1}^{\numpts},
\end{equation}
where $\point_i \in \mathbb{R}^3$ denotes the coordinate of an active
voxel, and $\feat_i$ is a feature vector encoding local geometry and
appearance. These features are derived by compressing a dense 3D grid
of multi-view DINO~\citep{oquab2023dinov2} features via a 3D VAE. The
VAE decoder is subsequently trained to reconstruct $\object$ into a
Signed Distance Field (SDF) for mesh extraction. This 3D latent
representation is illustrated in the upper section of the central
blue block in Figure~\ref{fig:cupid_architecture}.

The camera pose $\pose$, represented by its projection matrix
$\boldsymbol{P} = \boldsymbol{K}[\boldsymbol{R}|\boldsymbol{t}]$, is
reparameterized
in a novel way to facilitate joint generation. Instead of a compact
vector, we represent the pose as a dense field of 3D-to-2D
correspondences, which we term an \textit{in-cube pixel distribution}:
\[
\pose \triangleq \{\point_i, \pixcoord_i\}_{i=1}^{\numpts},
\]
where $\pixcoord_i = (u_i, v_i) \in [0,1]^2$ are the normalized 2D
pixel coordinates corresponding to the 3D voxel center $\point_i$.
These correspondences are obtained via perspective projection:
\[
\pixcoord_i = \pi(\boldsymbol{P}, \point_i).
\]
This set of correspondences can be visualized as a 3D UV cube, where
the $(u,v)$ values act like view-dependent colors on a 3D grid. This
is illustrated at the bottom of the blue block in the middle of
Figure \ref{fig:cupid_architecture}.

Given a generated set of correspondences, the global camera matrix
$\boldsymbol{P}^{*}$ is recovered by solving a {\em Perspective-n-Point}
(PnP) problem using a least-squares solver~\citep{abdel2015direct}:
\begin{equation}
\boldsymbol{P}^{*} = \argmin_{\boldsymbol{P}} \sum_{i=1}^{\numpts}
\big\|\pi(\boldsymbol{P},\point_i) - \pixcoord_i\big\|^2.
\label{eq:pnp}
\end{equation}
This formulation elegantly transforms the joint object-pose
generation problem into one of generating an object-centered 3D
shape/appearance with an associated observer/camera centered pose of
the object (via the UV coordinates). This is illustrated at the
bottom left corner of the the red block on the right of Figure
\ref{fig:cupid_architecture}.

\begin{remark}[Structured 3D Modeling]
Notice that the voxel-type representations chosen here for shape and
pose are different from the point cloud type representations adopted
in the conditioned shape generation task studied in the preceding
\Cref{sec:Michelangelo}. In that task, the goal is to generate an
object that only needs to be strongly correlated to the given image
(or text). It does not require an accurate pose alignment between the
object (in its canonical pose) and the given image (viewpoint). For
the task considered here, however, we want to ``reconstruct'' an
object whose shape, appearance, and pose faithfully agrees with the
given image. Therefore, it is necessary for the internal
representation to explicitly model the object's canonical shape and
relative pose simultaneously. Note that such a ``decoupling'' of an
object (canonical) shape/appearance and a pose from which it being
observed closely imitates the aforementioned roles of the IT cortex
and Hippocampus in our brain, respectively. As we will see later,
such a decoupled representation will enable us to obtain a structured
and consistent 3D model of a scene that may contain multiple objects
from multiple views. We contend that a so-structured 3D model is
necessary for us to conduct view-centric or object-centric
reasoning,\footnote{For example, answering questions such as ``which
object is in front of, behind, to the left or right of, or on top of
which object.} generation, and interaction with the scene.
\end{remark}

\begin{figure*}[t]
\centering
\includegraphics[width=\textwidth]{chapters/chapter8/figs/cupid_method.pdf}
\caption{\textbf{The Cupid Two-Stage Generative Reconstruction
Pipeline.} Given an input image, the first stage ($\sflow$) generates
a coarse occupancy cube and a UV cube, which encodes 3D-to-2D
correspondences. A PnP solver recovers the camera pose
$\boldsymbol{P}^{*}$ from these correspondences. The second stage
($\lflow$) is conditioned on this recovered pose. It injects
pixel-aligned features (sampled from DINOv2 and low-level feature
maps) into the generative process to synthesize high-fidelity
geometry and appearance for the final 3D object.}
\label{fig:cupid_architecture}
\end{figure*}

\subsection{Cascaded Flow Modeling}
\label{sec:flow_model}

We employ a two-stage cascaded flow model to jointly sample a 3D
object and its corresponding camera pose, denoted as $\latent = \{
(\point_i, \feat_i), (\point_i,\pixcoord_i) \}_{i=1}^{\numpts}$. In
the first stage, the occupancy and pose generation model ($\sflow$)
produces two key outputs: (i) an occupancy cube identifying active
voxels, and (ii) a UV cube encoding the object-centric camera pose.
In the second stage, conditioned on these predictions, the
pose-aligned geometry and appearance model ($\lflow$) synthesizes
DINO~\citep{oquab2023dinov2} features $\feat_i$ for each active voxel
to yield the final latent $\latent$.

\paragraph{Occupancy and Pose Generation.}
Given a conditioning image \(\image\), this stage generates an
occupancy cube and a UV cube, both at resolution $r$. The occupancy
cube, \(\boldsymbol{G}_o \in \{0,1\}^{r \times r \times r \times
1}\), contains binary values distinguishing active from inactive
voxels. The UV cube, \(\boldsymbol{G}_{uv} \in [0,1]^{r \times r
\times r \times 2}\), stores normalized pixel coordinates $(u,v)$ for
every voxel; notably, voxels sharing identical pixel coordinates lie
along the same camera ray.\footnote{We do not explicitly model
occlusion, as recent studies suggest transformers effectively model
light transport implicitly~\citep{jin2024lvsm, zeng2025renderformer}.}

To improve computational efficiency, we train a 3D VAE to compress
the UV cube into a low-resolution feature grid \(\boldsymbol{S}_{uv}
\in \mathbb{R}^{s\times s\times s \times C}\), achieving
near-lossless pose recovery (mean RRE/RTE $< 0.5^\circ$). To
fine-tune TRELLIS, we concatenate the original feature grid
\(\boldsymbol{S}_o\) with \(\boldsymbol{S}_{uv}\) and introduce
linear projection layers at both the input and output of the flow
network $\sflow$. Once the occupancy and UV cubes are generated, we
extract the set $\{ (\point_i, \pixcoord_i(\pose))
\}_{i=1}^{\numpts}$ by collecting active voxels and solving for the
camera pose via Equation~\ref{eq:pnp}.

% {\color{blue} Updated: We have included 3D tokenization in section
% \ref{sec:3d_tokenization}.}
% \yima{How these voxel-type 3D data are tokenized. Discuss
% relationship or difference from the Perceiver for tokenizing point
% clouds in the Michelangelo project. Also, describe basic deep
% architectures used for the two flow-matching based generative
% models learned. }

\paragraph{Pose-Aligned Geometry and Appearance Generation.}
In the second stage, we generate detailed 3D latent features $\{
\feat_i \}_{i=1}^{\numpts}$ exclusively at active voxel locations.
Experiments indicate that the standard $\lflow$ model from
TRELLIS~\citep{xiang2025structured}, which relies on globally
attended image information, is prone to color drift and the loss of
fine-grained details. We mitigate this by utilizing the camera pose
derived in the first stage to inject locally attended pixel
information into each voxel.

Specifically, we compute the features for the $i$-th voxel using the
estimated pose as follows:
\begin{equation}
\begin{aligned}
\feat^\dino_{i} &= \interp(\pixcoord_i, \dino(\image)) \in \mathbb{R}^{1024}, \\
\{\feat^{\high}_{i}\}_{i=1}^{\numpts} &= \slatenc\big(\{\point_i,
\feat^\dino_i\}_{i=1}^{\numpts}\big), \quad \feat^{\high}_{i} \in
\mathbb{R}^{8},
\end{aligned}
\label{eq:latent}
\end{equation}
where $\pixcoord_i$ represents the projection of the $i$-th 3D voxel
center onto the image plane, $\interp$ denotes bilinear
interpolation, and $\slatenc$ is the 3D VAE encoder. Although
DINO~\citep{oquab2023dinov2} captures high-level semantics, it often
lacks the low-level cues necessary for precise reconstruction. To
compensate, we extract complementary low-level features
$\feat^{\low}$ from $\image$ using a lightweight convolutional head,
sampling them at $\pixcoord_i$ via $\interp$. Finally, at each time
step $t$, we fuse the noisy voxel feature $\feat_i^{t}$ with the
pixel-aligned features via a linear layer before inputting them into
the flow transformer:
\begin{equation}
l_t = \mathrm{Linear}\big([\feat_i^{t}\,\oplus\, \feat^{\high}_{i}
\,\oplus\, \feat^{\low}_{i}]\big).
\end{equation}
This pose-aligned fusion strategy significantly enhances both 3D
geometric accuracy and appearance fidelity relative to the input image.

\begin{remark}[Decoupled Generative Processes]
Notice that the above two-stage framework decouples the distribution
of 3D objects into several related components, shape, pose, and
appearance and learn them separately. Although this means we have to
learn two generative denoising models together, the resulting learned
representations are much better structured and enable much more
precise and flexible spatial reasoning and generation.
\end{remark}

\subsection{Experiments}
\label{sec:experiments}
\label{sec:Evaluation}

To validate the architecture, we evaluated the system across three
critical dimensions: monocular geometry prediction, input-view
consistency, and full single-image-to-3D reconstruction. The
benchmarking process focused on learning-based systems, excluding
per-scene optimization methods to ensure a fair comparison of
inference capabilities. We ablate pose-conditioning.

\paragraph{The Comparative Landscape.}
We benchmarked against three distinct paradigms currently dominating the field:
\begin{enumerate}[label=(\roman*)]
\item \textit{Point-map Regression}, such as
VGGT~\citep{wang2025vggt} and MoGe~\citep{wang2025moge}, which
simultaneously predict per-pixel 3D points and camera poses but often
lack robust priors for occluded geometry;
\item \textit{View-centric Reconstruction}, including
LRM~\citep{hong2023lrm} and LaRa~\citep{LaRa}, which generate 3D
models directly in view space to bypass test-time pose estimation; and
\item \textit{Decoupled Generation}, exemplified by
OnePoseGen~\citep{geng2025one}, which separates canonical 3D
generation from pose estimation, relying on post-hoc alignment.
\end{enumerate}

\newcommand{\cmark}{\checkmark} % Checkmark
\newcommand{\xmark}{\times} % Cross

\begin{table*}[t]
\renewcommand{\arraystretch}{1.0}
\centering
\caption{
\textbf{Monocular geometry accuracy.} Cupid outperforms all 3D reconstruction and generation baselines and matches point-map regression methods that predict only partial geometry. Note that VGGT uses a ground-truth object mask, which may overestimate accuracy.
}
\label{tab:geometry}
\resizebox{0.99\textwidth}{!}{
\begin{tabular}{l c c c c c c | c c c c c}
\toprule
\multirow{2}{*}{\textbf{Method}} & \multirow{2}{*}{\textbf{3D}} &
\multicolumn{5}{c|}{\textbf{Toys4k}} &
\multicolumn{5}{c}{\textbf{GSO}} \\
& &
\textbf{mIOU} & \textbf{CD} & \textbf{CD} & \textbf{F-score} & \textbf{F-score} &
\textbf{mIOU} & \textbf{CD} & \textbf{CD} & \textbf{F-score} & \textbf{F-score} \\
& &
(avg)$\uparrow$ & (avg)$\downarrow$ & (med)$\downarrow$ & (0.01)$\uparrow$ & (0.05)$\uparrow$ &
(avg)$\uparrow$ & (avg)$\downarrow$ & (med)$\downarrow$ & (0.01)$\uparrow$ & (0.05)$\uparrow$ \\
\midrule
\textcolor{gray}{VGGT} & $\xmark$ &
\textcolor{gray}{--} & \textcolor{gray}{1.144} & \textcolor{gray}{0.498} & \textcolor{gray}{61.85} & \textcolor{gray}{95.90} &
\textcolor{gray}{--} & \textcolor{gray}{1.396} & \textcolor{gray}{0.388} & \textcolor{gray}{65.98} & \textcolor{gray}{95.95} \\
\textcolor{gray}{MoGe} & $\xmark$ &
\textcolor{gray}{92.80} & \textcolor{gray}{1.284} & \textcolor{gray}{0.581} & \textcolor{gray}{58.54} & \textcolor{gray}{95.31} &
\textcolor{gray}{96.18} & \textcolor{gray}{1.743} & \textcolor{gray}{0.575} & \textcolor{gray}{58.99} & \textcolor{gray}{94.68} \\ 
OnePoseGen & $\cmark$ &
9.34 & 153.2 & 59.92 & 6.11 & 24.10 &
12.16 & 116.2 & 60.56 & 7.28 & 25.77 \\
LaRa & $\cmark$ &
68.11 & 32.15 & 16.59 & 18.57 & 57.67 &
70.63 & 34.23 & 19.36 & 13.48 & 49.95 \\
OpenLRM & $\cmark$ &
86.26 & 2.726 & 1.291 & 40.42 & 90.60 &
91.35 & 3.741 & 1.858 & 34.14 & 87.20 \\
Ours & $\cmark$ &
\textbf{92.43} & \textbf{2.534} & \textbf{0.236} & \textbf{69.82} & \textbf{97.76} &
\textbf{95.27} & \textbf{1.823} & \textbf{0.434} & \textbf{61.01} & \textbf{95.59} \\
\bottomrule
\end{tabular}
}
\end{table*}

\begin{table}[t]
\vspace{6pt}
\centering
\caption{\textbf{Comparison on full 3D quality. }We report CLIP image scores of novel views following \citep{gao2024cat3d}.}
\label{tab:clip}

\renewcommand{\arraystretch}{1.0}
\resizebox{1.0\linewidth}{!}{
\begin{tabular}{l|ccccc}
\toprule
 & OnePoseGen & LaRa & OpenLRM & TRELLIS & Cupid \\
\midrule
\textbf{ViT-B/16} & 0.7933 & 0.8334 & 0.8939 & 0.9465 & \textbf{0.9501} \\
\textbf{ViT-L/14} & 0.7193 & 0.7682 & 0.8410 & 0.9210 & \textbf{0.9291} \\
\bottomrule
\end{tabular}
}
\end{table}



\begin{figure}[t!]
\centering
\includegraphics[width=1.0\linewidth]{chapters/chapter8/figs/cupid/figures/nvs.pdf}
\caption{
\textbf{
    Qualitative comparison on input view consistency.
}We render the input view using its generated camera pose.
For view centric methods (LRM, LaRa), we use ground-truth intrinsic
for rendering as they do not model intrinsic.
Our method produces the highest-fidelity geometry and appearance; LRM
hallucinates incorrect details, LaRa is overly blurry due to 2D
diffusion inconsistencies, and 3D generation method OnePoseGen
frequently fails to register pose reliably.
}
\label{fig:view_consistency}
\end{figure}
\renewcommand{\arraystretch}{1.0}
\begin{table}[t]
\centering
\caption{
\textbf{Input-view consistency.} Cupid achieves superior input view consistency, producing accurate appearance alignment.
}
\label{tab:appearance}
\centering
\resizebox{\linewidth}{!}{
\begin{tabular}{lc|ccc|ccc}
\toprule
\textbf{Dataset} & &
\multicolumn{3}{c|}{\textbf{Toys4K}} &
\multicolumn{3}{c}{\textbf{GSO}} \\
\textbf{Method} & \textbf{Pose} & \textbf{PSNR}$\uparrow$ & \textbf{SSIM}$\uparrow$ & \textbf{LPIPS}$\downarrow$ & \textbf{PSNR}$\uparrow$ & \textbf{SSIM}$\uparrow$ & \textbf{LPIPS}$\downarrow$ \\
\midrule
LaRa & $\xmark$ & 22.00 & 93.42 & 0.0884 & 19.81 & 91.61 & 0.1119 \\
OpenLRM & $\xmark$ & 26.41 & 80.17 & 0.1156 & 25.79 & 78.80 & 0.1268 \\
OnePoseGen & $\cmark$ & 17.43 & 89.37 & 0.1174 & 14.87 & 86.46 & 0.1386 \\
Cupid & $\cmark$ & \textbf{30.05} & \textbf{96.81} & \textbf{0.0251} & \textbf{28.68} & \textbf{95.49} & \textbf{0.0354} \\
\bottomrule
\end{tabular}
}
\end{table}

\paragraph{Geometric Accuracy and Pipeline Simplification.}
In monocular geometry tasks, the proposed model demonstrated a clear
superiority over view-centric approaches (Table~\ref{tab:geometry}).
On the GSO dataset~\citep{downs2022google}, it reduced the average
Chamfer Distance (CD) by 50\% relative to LRM~\citep{hong2023lrm}.
This improvement is attributed to the model's dual capability: it
accurately models object intrinsics via generative priors while
simultaneously predicting precise extrinsic camera parameters.

Crucially, the model achieved geometric accuracy competitive with
specialized point-map regressors like MoGe~\citep{wang2025moge}. This
finding has significant implications for pipeline design: it suggests
that complex, multi-stage workflows—which traditionally use partial
geometry estimation as a prerequisite
step~\citep{yao2025castcomponentaligned3dscene}—can be streamlined.
Our approach achieves comparable accuracy in a single step, rendering
those intermediate representations unnecessary.

\paragraph{Visual Fidelity and Consistency.}
We evaluated input-view consistency to measure how well the 3D
reconstruction realigns with the source image
(Table~\ref{tab:appearance}). The results highlighted specific
weaknesses in the baselines: LaRa's texture quality degraded due to
inconsistent novel-view generation, while OnePoseGen suffered from
severe texture misalignment and color shifting during registration.

By contrast, our method maintained high pose accuracy without
sacrificing texture details (Figure~\ref{fig:view_consistency}). It
successfully reconstructed high-fidelity objects from both synthetic
datasets (Toys4K) and challenging ``in-the-wild'' images
(Figure~\ref{fig:cupid_results}), consistently outperforming
baselines in semantic alignment metrics such as CLIP
similarity~\citep{radford2021learning} (Table~\ref{tab:clip}).

\paragraph{Ablation Studies}
\label{sec:ablation}

\begin{figure}[t!]
\centering
\includegraphics[width=1\linewidth]{chapters/chapter8/figs/cupid/figures/ablation.pdf}
\caption{
\textbf{Qualitative comparison of various pose-aligned conditioning.}
Our method (e) achieves the best visual quality in terms of color
fidelity and detail.
}\label{fig:abl_appearance}
\end{figure}
% \begin{table}[t]
% \centering
% \renewcommand{\arraystretch}{1.0}
% \begin{minipage}[t]{\linewidth}
% \vspace{0pt} % Ensures baseline alignment
% \captionof{table}{Comparison on CLIP image scores of novel views.}
% \label{tab:clip}
% \centering
% \resizebox{0.55\linewidth}{!}{
% \begin{tabular}{l|cc}
% \toprule
% \textbf{Method} & \textbf{ViT-B/16} & \textbf{ViT-L/14} \\
% \midrule
% OnePoseGen   & 0.7933 & 0.7193 \\
% LaRa         & 0.8334 & 0.7682 \\
% OpenLRM      & 0.8939 & 0.8410 \\
% TRELLIS      & 0.9465 & 0.9210 \\
% Ours         & \textbf{0.9501} & \textbf{0.9291} \\
% \bottomrule
% \end{tabular}
% }
% \end{minipage}

% \hfill

% \renewcommand{\arraystretch}{1.065}

% \begin{minipage}[t]{\linewidth}
% \vspace{0pt} % Ensures baseline alignment
% \captionof{table}{Ablation studies of pose-aligned conditioning.}
% \label{tab:abl_appearance}
% \centering
% \resizebox{\linewidth}{!}{
% \begin{tabular}{l|ccc|ccc}
% \toprule
% \multirow{2}{*}{\textbf{Method}} & \multicolumn{3}{c|}{\textbf{GT Geo \& Pose}} & \multicolumn{3}{c}{\textbf{Sampled Geo \& Pose}} \\
%  & \textbf{PSNR} & \textbf{SSIM} & \textbf{LPIPS} & \textbf{PSNR} & \textbf{SSIM} & \textbf{LPIPS} \\
% \midrule
% (a) Baseline (TRELLIS)   & 31.84 & 97.50 & 0.0219 & 27.47 & 95.64 & 0.0327 \\
% (b) Position Embedding   & 32.07 & 97.58 & 0.0211 & 27.56 & 95.67 & 0.0323 \\
% (c) Latent (w/o Occ.)    & 32.37 & 97.72 & 0.0201 & 27.85 & 95.87 & 0.0309 \\
% (d) Latent (Occ.)        & 32.39 & 97.77 & 0.0199 & 27.74 & 95.80 & 0.0313 \\
% (e) Latent (Visual Feat.)& \textbf{34.86} & \textbf{98.24} & \textbf{0.0168} & \textbf{30.05} & \textbf{96.81} & \textbf{0.0251} \\
% \bottomrule
% \end{tabular}
% }
% \end{minipage}
% \end{table}


\begin{table}[t]
\centering
\caption{\textbf{Ablation studies of pose-aligned conditioning.}}
\label{tab:abl_appearance}
\renewcommand{\arraystretch}{1.065}
\resizebox{1.0\linewidth}{!}{
\begin{tabular}{l|ccc|ccc}
\toprule
\multirow{2}{*}{\textbf{Method}} & \multicolumn{3}{c|}{\textbf{GT Geo \& Pose}} & \multicolumn{3}{c}{\textbf{Sampled Geo \& Pose}} \\
 & \textbf{PSNR} & \textbf{SSIM} & \textbf{LPIPS} & \textbf{PSNR} & \textbf{SSIM} & \textbf{LPIPS} \\
\midrule
(a) Baseline (w/o PAC)   & 31.84 & 97.50 & 0.0219 & 27.47 & 95.64 & 0.0327 \\
(b) Position Embedding   & 32.07 & 97.58 & 0.0211 & 27.56 & 95.67 & 0.0323 \\
(c) Latent (w/o Occ.)    & 32.37 & 97.72 & 0.0201 & 27.85 & 95.87 & 0.0309 \\
(d) Latent (Occ.)        & 32.39 & 97.77 & 0.0199 & 27.74 & 95.80 & 0.0313 \\
(e) Latent (Visual Feat.)& \textbf{34.86} & \textbf{98.24} & \textbf{0.0168} & \textbf{30.05} & \textbf{96.81} & \textbf{0.0251} \\
\bottomrule
\end{tabular}
}
\end{table}

Ablation studies on the second-stage refinement confirmed the
critical role of Pose-Aligned Conditioning (PAC). We tested variants
ranging from standard latent baselines to those incorporating DINOv2
feature volumes.

As shown in Table~\ref{tab:abl_appearance}, the most effective
configuration proved to be a hybrid approach: concatenating DINOv2
features with low-level visual cues extracted via convolutional
layers. This combination yielded the highest visual quality,
effectively bridging semantic understanding with pixel-level detail
(Figure~\ref{fig:abl_appearance}). Furthermore, this conditioning
mechanism demonstrated remarkable robustness, maintaining high
performance even when the input geometry contained stochastic
perturbations from the first generation stage.

\paragraph{Application I: Compositional Scene Reconstruction.}

\begin{figure*}[t]
\centering
\includegraphics[width=\textwidth]{chapters/chapter8/figs/cupid_scene.pdf}
\caption{\small\textbf{Component-aligned scene reconstruction.} Each
object is reconstructed independently with \textbf{Cupid}; explicit
3D–2D correspondences enable precise placement into a shared frame.}
\label{fig:cupid_scene_reconstruction}
\end{figure*}

\label{sub:cupid_application}
To extend Cupid to scene-level reconstruction, we leverage foundation
models to decompose the scene into $K$ objects with masks $\{
\boldsymbol{M}^{(k)} \}_{k=1}^K$ and estimate a global metric
pointmap $\mathcal{P}$~\citep{wang2025moge}.

In the first stage, we reconstruct each object independently to
recover its canonical geometry, adapting the generator to handle
mutual occlusions via random-mask fine-tuning~\citep{wu2025amodal3r}.
For the $k$-th object, this yields a set of canonical 3D points
$\boldsymbol{x}_i^{(k)}$ corresponding to pixels $\boldsymbol{u}_i
\in \boldsymbol{M}^{(k)}$. To resolve the scale ambiguity inherent in
independent generation, we formulate the alignment as a
correspondence problem between the generated canonical space and the
estimated camera space:
\begin{equation}
\boldsymbol{p}_i = \mathcal{P}(\boldsymbol{u}_i), \quad \forall
\boldsymbol{u}_i \in \boldsymbol{M}^{(k)},
\end{equation}
where $\boldsymbol{p}_i$ represents the target metric coordinates in
the camera frame derived from the depth foundation model.

In the second stage, we compute a per-object similarity
transformation $\mathcal{T}_{k} = \{s_{k}, \boldsymbol{R}_{k},
\boldsymbol{t}_{k}\}$ to place all components into a unified
coordinate system. We solve for these parameters using the Umeyama
method~\citep{umeyama2002least} by minimizing the alignment error
over the visible pixels of each object:
\begin{equation}
\mathcal{T}_{k}^* = \argmin_{s_k, \boldsymbol{R}_k, \boldsymbol{t}_k}
\sum_{i} \left\| \boldsymbol{p}_i - \left(s_k \boldsymbol{R}_k
\boldsymbol{x}_i^{(k)} + \boldsymbol{t}_k\right) \right\|^2.
\end{equation}
Applying these optimal transformations results in a metric-consistent
scene composition where all generated components are correctly
positioned and scaled relative to one another. Please see
Figure~\ref{fig:cupid_scene} for more scene composition examples.

\begin{figure}[t!]
\centering
\includegraphics[width=0.95\linewidth]{chapters/chapter8/figs/cupid/figures/scene_result.pdf}
%
\caption{\textbf{Additional examples of component-aligned scene
reconstruction.} For each example shown, the panels display: (top
left) the input image, (top right or bottom left) the final rendered
output, and (bottom) the reconstructed individual components,
color-coded for clarity.
}
\label{fig:cupid_scene}
\vspace{-5mm}
\end{figure}

\paragraph{Application II: Multi-view Consistent Reconstruction.}

\begin{figure}[t!]
\centering
\includegraphics[width=\textwidth]{chapters/chapter8/figs/cupid/figures/mv.pdf}
\caption{
\textbf{Multi-view conditioning.} Our decoupled joint modeling
naturally supports multi-view conditioning. With multiple input views
available, we fuse the shared \textit{view-agnostic} object latent
across flow paths (like
MultiDiffusion~\citep{bar2023multidiffusion}), enabling object and
cameras refinement across all views. Top: inputs; Middle:
reconstructed 3D object and camera poses; Bottom: rendered images and geometry.}
\label{fig:mv_images}
\end{figure}

\label{sub:cupid_application_mv}
% To implement the multi-view capability described in
% Section~\ref{sec:flow_model}, we generalize the sampling process to
% accommodate $K$ input views $\{ \image^{(k)} \}_{k=1}^K$.
Thanks to our Bayesian formulation, our model generalizes to
multi-view scenarios even when trained only on single images. To
implement the multi-view capability described in
Section~\ref{sec:flow_model}, we extend the sampling process to
accommodate $K$ input views ${ \image^{(k)} }_{k=1}^K$.

In the first stage, we aim to recover a single canonical geometry
while simultaneously estimating $K$ distinct camera poses. We
initialize a shared geometry latent $\boldsymbol{S}_{o}$ and $K$
independent pose latents $\{ \boldsymbol{S}_{uv}^{(k)} \}_{k=1}^K$.
At each denoising step $t$, we predict the flow fields for both
geometry and pose in parallel for each view. Let the output of the
flow network for the $k$-th view be denoted as a concatenated vector
of geometry and pose updates:
\begin{equation}
[d_{o, t}^{(k)}, d_{uv, t}^{(k)}] =
\sflow(\text{concat}(\boldsymbol{S}_{o}^t, \boldsymbol{S}_{uv}^{(k),
t}), t, \image^{(k)}).
\end{equation}
To ensure geometric consistency across all observations, we aggregate
the geometry updates while maintaining view-specific pose
trajectories. The update rule at step $t$ becomes:
\begin{align}
\boldsymbol{S}_{o}^{t-1} &\leftarrow
\text{Step}\left(\boldsymbol{S}_{o}^{t}, \frac{1}{K} \sum_{k=1}^K
d_{o, t}^{(k)}\right), \\
\boldsymbol{S}_{uv}^{(k), t-1} &\leftarrow
\text{Step}\left(\boldsymbol{S}_{uv}^{(k), t}, d_{uv, t}^{(k)}\right),
\end{align}
where $\text{Step}(\cdot)$ denotes the update rule of the flow
sampler. The resulted feature grids will be decoded into a shared
coarse occupancy grid $\boldsymbol{G}_o$ and poses $\{\pose^{k}\}_{k=1}^{K}$.

In the second stage, with the geometry fixed, we apply a
Multi-Diffusion strategy~\citep{bar2023multidiffusion} to the
appearance latent $\boldsymbol{F}_{o}$. We maintain a single shared
latent and update it using the averaged flow predictions from all $K$
views, conditioned on the coarse geometry $\boldsymbol{G}_o$ and
poses $\{\pose^{k}\}_{k=1}^{K}$ estimated in the first stage:
\begin{equation}
d_{\text{feat}, t}^{(k)} = \lflow(\boldsymbol{F}_{o}^t,
\boldsymbol{G}_o, \pose^{k}, t, \image^{(k)}),
\end{equation}
\begin{equation}
\boldsymbol{F}_{o}^{t-1} \leftarrow
\text{Step}\left(\boldsymbol{F}_{o}^{t}, \frac{1}{K} \sum_{k=1}^K
d_{\text{feat}, t}^{(k)}\right).
\end{equation}
This averaging operation ensures that the generated 3D representation
satisfies the visual constraints imposed by all input views
simultaneously, yielding the SfM-like consistency shown in
Figure~\ref{fig:mv_images}.

\subsection{Conclusion}
This section presented a unified framework for 3D generation and
reconstruction, operationalizing the biological distinction between
object identity (``what'') and spatial context (``where''). By
jointly modeling the canonical 3D object $\object$ and the camera
pose $\pose$, our approach effectively disentangles intrinsic shape
from extrinsic viewpoint.

Leveraging pose priors and aligned conditioning, the model performs
Bayesian-like inference to resolve full 3D geometry from limited
single-view inputs. This ensures geometric consistency and robustness
to viewpoint variations. Ultimately, this structured representation
advances the capabilities of Embodied AI, equipping agents with the
necessary spatial understanding to navigate and interact with 3D
environments effectively.

\section{Conditioned Human Body Motion Generation}
\label{sec:egoallo}
% \yima{EgoAllo (Brent, Annie and Chen), move the stuff from Chapter
% 6 to here and elaborate on the methods and implementation more.}
% \subsection{Body Pose Generation Conditioned on Head and
% Hands}\label{sub:ego-allo}

%\input{chapters/chapter8/applications_egoallo}
% This type of
% conditional estimation or generation problem arises rather naturally in many
% practical applications.
% A typical problem of this kind is how to estimate and generate body
% pose and hand gesture conditioned on a given head pose and
% egocentric images, as illustrated in Figure \ref{fig:pose-teaser}.
% This is often the problem we need to solve when one is wearing a
% head-mounted device such as the Vision Pro from Apple or the
% Project Aria from Meta. The pose of the whole body and the gesture
% of the hands need to be inferred so that we can use the information
% to control virtual objects that the person interacts with.
% \begin{figure*}[t]
%   \centering
%   \includegraphics[width=\linewidth]{\toplevelprefix/chapters/chapter7/figs/pose-teaser.pdf}
%     \caption{
% A system that estimates human body height, pose, and hand
% parameters (middle), conditioned on egocentric SLAM poses and
% images (left). Outputs capture the wearer's actions in the
% allocentric reference frame of the scene, which we visualize here
% with 3D reconstructions (right).
%   }
%   \label{fig:pose-teaser}

% \end{figure*}

% Notice that in this case, one only has the head pose provided by
% the device and a very limited field of view for part of one's hands
% and upper limbs. The pose of the rest of the body needs to be
% ``inferred'' or ``completed'' based on such partial information.
% The only way one can estimate the body pose over time is by
% learning the joint distribution of the head and body pose sequences
% in advance and then sampling this prior distribution conditioned on
% the real-time partial inputs. Figure \ref{fig:pose-method} outlines
% a system called EgoAllo \cite{yi2024egoallo} to solve this problem
% based on a learned conditional diffusion-denoising model.
% \begin{figure*}[t]
%   \centering
%   \includegraphics[width=\linewidth]{\toplevelprefix/chapters/chapter7/figs/pose-method.pdf}
% \caption{
%     \textbf{Overview of technical components of EgoAllo \cite{yi2024egoallo}.}
%     A diffusion model is pretrained that can generate body pose
% sequence based on local body parameters (middle).
%     An invariant parameterization $g(\cdot)$ of SLAM poses (left)
% is used to condition the diffusion model. These can be placed into
% the global coordinate frame via global alignment to input poses.
%     When available, egocentric video is used for hand detection
% (left) via HaMeR~\cite{pavlakos2023reconstructing}, which can be
% incorporated into samples via guidance by the generated gesture.
%   }
%   \label{fig:pose-method}
% \end{figure*}

% Figure \ref{fig:comparison} compares some ground truth motion
% sequences with sampled results generated by the EgoAllo. Although
% the figure shows one result for each input head pose sequence,
% different runs can generate different body pose sequences that are
% consistent with the given head pose, all drawing from the
% distribution of natural full-body motion sequences.
% \begin{figure}[t]
% \centering
% \begin{subfigure}{0.45\textwidth}
%   \centering
%   \includegraphics[width=\linewidth]{\toplevelprefix/chapters/chapter7/figs/qual0_gt.png}
%   \caption{\centering Ground-truth}
% \end{subfigure}
% \hfill
% \begin{subfigure}{0.45\textwidth}
%   \centering
%   \includegraphics[width=\linewidth]{\toplevelprefix/chapters/chapter7/figs/qual0_ours.png}
%   \caption{EgoAllo}
% \end{subfigure}

% \begin{subfigure}{0.45\textwidth}
%   \centering
%   \includegraphics[width=\linewidth]{\toplevelprefix/chapters/chapter7/figs/qual1_gt.png}
%   \caption{\centering Ground-truth}
% \end{subfigure}
% \hfill
% \begin{subfigure}{0.45\textwidth}
%   \centering
%   \includegraphics[width=\linewidth]{\toplevelprefix/chapters/chapter7/figs/qual1_ours.png}
%   \caption{EgoAllo}
% \end{subfigure}

% \caption{
% \textbf{Egocentric human motion estimation for a running (top) and
% squatting (bottom) sequence.}
% The ground-truth motion is compared with one output from EgoAllo
% that is consistent with the given head pose sequence.
% }
% \label{fig:comparison}
% \end{figure}

% Strictly speaking, the solution proposed in EgoAllo \cite{yi2024egoallo} does
% not enforce measurement matching using the techniques introduced
% above. Instead
% it heuristically enforces the condition by utilizing  the cross-attention
% mechanism in a transformer architecture. As we will describe with
% more precision
% in the paired data setting in \Cref{sub:text-cond}, there is reason
% to believe that the
% cross-attention mechanism is in a way approximately realizing the conditional
% sampling of the denoising a posteriori. We believe the more principled
% techniques introduced here, if properly implemented, can lead to
% better methods
% that further improve the body pose and hand gesture estimation.

%EgoAllo: Egocentric Human Motion Estimation
% Section for Chapter 7 - Learning Representations for Real-World Data
%
% Following the structure of the Cupid (3D Object Generation) section

% \section{Egocentric Human Motion Estimation}
% One must know that the primary reason why our brain senses the
% external 3D world to develop a visual memory is for us to negotiate
% with the physical environment with our body and interact with
% objects in it with our hand. Hence, our body movement and hand
% motion are arguable another most important data that we would like
% to learn about. In fact, our brain dedicates a significant portion
% of it, the motion cortex, to record and coordinate our own body
% motion. In this section, we show how to use the methods introduced
% in this book to learn the distribution of human motions in advance
% and how it can be used to facilitate the task of sensing and
% estimating our own motion in real time.

So far, we have shown, in preceding sections, how to learn the
distributions of visual data including 2D images and 3D object
shapes. Note that the primary goal for us to sense the 3D world and
develop a model of it is for us to navigate, interact, and manipulate
objects within. Based on information perceived about the environment,
our brain makes plans and takes actions based on tasks at hand. The
level of dexterity achievable by ``eye-hand coordination'' is
arguably the highest form of physical intelligence in nature. Hence
accurately estimating and controlling our body and hand movement is
crucial for us to interact with objects within the environment. Our
brain dedicates significant resources to this: the {\em motor
cortex}, located in the brain's frontal lobe, records and coordinates
our body movements, while proprioceptive signals from our body's
nerve system provide continuous feedback about our limb positions.
Therefore, if we want robots, especially humanoid-like robots, to
emulate human actions, they need to learn (the distribution of) our
body movements.

In this section, we show how the conditional inference framework from
Chapter~\ref{ch:conditional-inference} can be applied to learn the
distribution of human motion and then exploited it for the body pose
estimation task. To this end, we consider a simple but natural setup:
a person wears a head-mounted device such as an AR headset or smart
glasses that captures forward-facing video and  tracks its 3D pose
via Structure from Motion (SFM) or Simultaneous Localization and
Mapping (SLAM). Such methods estimate the device's position by
tracking and triangulating visual features in the scene to recover 3D
structure and pose.\footnote{Interested readers may refer to the
textbook \cite{MaY2003} for an introduction to this line of work.}
Our goal is to estimate the wearer's full-body pose, height, and hand
configuration from these egocentric observations, even though the
body is rarely visible in the captured video. The key observation is
that \emph{head motion encodes rich information about body motion}.
When we walk, our head bobs; when we reach, it tilts; when we sit, it
follows a distinctive arc. These correlations, learned from motion
capture data, form a conditional prior $p(\text{body} \mid
\text{head})$. The method we feature here in this section is based on
an open-sourced project, called \textbf{EgoAllo} (for
egocentric-to-allocentric estimation):\\
\centerline{\url{https://egoallo.github.io}},\\
which incorporates additional measurements like hand detections and
ground contact via guided sampling. Figure~\ref{fig:egoallo_teaser}
illustrates the overall setup and task.

\begin{figure}[t!]
\centering
\includegraphics[width=\textwidth]{\toplevelprefix/chapters/chapter8/figs/egoallo/teaser_oct9.pdf}
\caption{\textbf{Egocentric Human Motion Estimation.} Given head
poses from SLAM and egocentric images from a head-mounted device
(left), EgoAllo estimates full body pose, height, and hand
parameters in the world (allocentric) reference frame (right).
The estimated bodies are grounded in the scene, with feet
contacting the floor and hands at physically plausible positions.}
\label{fig:egoallo_teaser}
\end{figure}

% ==============================================================================
\subsection{Representing Human Bodies}
\label{sub:egoallo_body_models}
% ==============================================================================

Before formulating the estimation problem, we need a mathematical
representation for human bodies. The skeleton of our body is a
kinematic tree: rigid bones connected by joints, with the pelvis as
root and chains extending to the head, hands, and feet. Each joint
has rotational degrees of freedom: the shoulder rotates in three
axes, the elbow primarily in one. For a more formal and detailed
introduction to kinematic representation of linked rigid bodies like
the human body, interested readers may refer to the classic textbook
\cite{MurrayR1994}.

\paragraph{The SMPL body model.}
SMPL (Skinned Multi-Person Linear model) \citep{loper2015smpl}
provides a compact, differentiable representation of the human body,
for both its (kinematic) pose and shape. It separates a body
configuration into two components:
\begin{itemize}
\item \textbf{Pose parameters} $\theta$: Joint rotations
specifying the body's articulation. SMPL uses 23 joints (plus
the root), each parameterized by a 3D rotation. Together,
these determine whether a person is standing, sitting,
reaching, walking, etc.
\item \textbf{Shape parameters} $\beta$: A low-dimensional
vector, typically 10 coefficients, controlling body
proportions such as height, weight, limb lengths, and torso
width. These are derived from principal component analysis of
a large corpus of 3D body scans.
\end{itemize}

Given pose $\theta$ and shape $\beta$, SMPL produces a 3D mesh via
learned blend skinning. It also provides joint positions through
forward kinematics: starting from the root (pelvis) position and
applying each joint's rotation in sequence along the kinematic chain,
we can compute where each joint ends up in 3D space. For example, to
find the wrist position, we compose the rotations of the pelvis,
spine, shoulder, and elbow, accumulating the translations defined by
bone lengths.

\paragraph{Adding hands: SMPL-H and MANO.}
The original SMPL model represents hands as simple endpoints. For
applications requiring detailed hand poses like grasping, gesturing,
or typing, this is insufficient. SMPL-H \citep{romero2017embodied}
extends SMPL with articulated hands based on the MANO hand model,
adding 15 joints per hand (three per finger) for 52 total joints.

Hand rotations $\phi$ follow the same parameterization as body
joints. The complete state includes body rotations $\theta$, hand
rotations $\phi$, and shape parameters $\beta$.

\paragraph{Why body models matter.}
Why use SMPL-H rather than, say, directly regressing joint positions
or mesh vertices? Parametric models offer key advantages. Outputs are
always valid human bodies, with no anatomically impossible
configurations. The model is differentiable, enabling gradient-based
optimization. The parameterization is compact, requiring only a few
hundred numbers versus a full mesh. And we get semantic structure,
reasoning about ``left wrist'' rather than vertex indices.

% ==============================================================================
\subsection{Task and Setup}
\label{sub:egoallo_objective}
% ==============================================================================

\paragraph{Inputs.}
The inputs consist of two streams from a head-mounted device:
\begin{itemize}
\item \textbf{Head poses} $\boldsymbol{h}_{1:T} \in
\mathsf{SE}(3)^T$\footnote{Here $\mathsf{SE}(3)$ represents the
    special Euclidean group in $\mathbb{R}^3$, essentially the group of
transformations that consist of rotation and translation.}:
Position and orientation of the device in world coordinates
at each timestep, from visual-inertial SLAM. Here $\mathsf{SE}(3)$
denotes the group of rigid transformations in 3D, comprising
rotations and translations.\footnote{Interested readers may
    refer to \cite{MaY2003} for a formal introduction to rigid
body motion and its representation and properties.} These
poses have metric scale, meaning we know actual distances in
meters rather than just relative geometry, and are
gravity-aligned so we know which direction is up.
\item \textbf{Egocentric images} $I_{1:T}$: Forward-facing video.
The wearer's body is mostly not visible, but hands frequently
appear when reaching or manipulating objects.
\end{itemize}

From images, an off-the-shelf hand estimator like HaMeR
\citep{pavlakos2023hamer} produces per-frame measurements when hands
are visible:
\begin{equation}
\boldsymbol{y}^{\text{hand},t} = \big(\hat{\phi}^{\text{MANO},t},\,
\hat{\boldsymbol{q}}^{3D,t},\, \hat{\boldsymbol{q}}^{2D,t}\big),
\end{equation}
consisting of estimated hand rotations, 3D keypoints in the camera
frame, and 2D keypoints.

\paragraph{Outputs.}
EgoAllo estimates a motion sequence over $T$ timesteps:
\begin{equation}
\boldsymbol{x}_0 = \{\boldsymbol{x}_0^{1}, \dots, \boldsymbol{x}_0^{T}\},
\quad\text{where}\quad
\boldsymbol{x}_0^{t} = \big(\theta^t,\, \phi^t,\, \beta,\, s\big).
\end{equation}
Here $\theta^t$ are body joint rotations, $\phi^t$ hand rotations,
$\beta$ shape parameters shared across time, and $s$ a height
parameter that is also shared. The subscript $0$ denotes clean
motion, following diffusion notation where $\boldsymbol{x}_n$ is the noisy
version at step $n$.

EgoAllo estimates height in addition to pose, which is essential for
grounding. Without it, the body floats at an arbitrary vertical
position; with it, feet contact the floor and the head aligns with the camera.

\paragraph{Posterior formulation.}
Following the Bayesian framework from
Chapter~\ref{ch:conditional-inference}, estimation is formulated as
sampling from a posterior. Let $\boldsymbol{y}$ denote all observations
and $\boldsymbol{c} = g(\boldsymbol{h}_{1:T})$ be a processed representation of
head poses (to be explained soon in
Section~\ref{sub:egoallo_representation}). The posterior factorizes
as:\footnote{For this factorization to be valid, we essentially
assume that given $\boldsymbol{x}_0$, $\boldsymbol{y}^{\text{hand}},
\boldsymbol{y}^{\text{phys}}$ is conditionally independent of the head
pose $\boldsymbol{c}$.}
\begin{equation}
p(\boldsymbol{x}_0 \mid \boldsymbol{y}) \;\propto\;
p_\theta(\boldsymbol{x}_0 \mid
\boldsymbol{c}) \cdot p(\boldsymbol{y}^{\text{hand}},
\boldsymbol{y}^{\text{phys}} \mid \boldsymbol{x}_0),
\end{equation}
where:
\begin{itemize}
\item $p_\theta(\boldsymbol{x}_0 \mid \boldsymbol{c})$ is a learned conditional
prior, specifically a diffusion-denoising model that
generates plausible body motion given head motion.
\item $p(\boldsymbol{y}^{\text{hand}}, \boldsymbol{y}^{\text{phys}} \mid
\boldsymbol{x}_0)$ is the likelihood of hand observations and
physical constraints such as ground contact and no foot
skating, given the body motion.
\end{itemize}

The likelihood is expressed via a guidance energy $\mathcal{L}_{\text{guide}}$:
\begin{equation}
p(\boldsymbol{y}^{\text{hand}}, \boldsymbol{y}^{\text{phys}} \mid
\boldsymbol{x}_0) \propto
\exp\big(-\mathcal{L}_{\text{guide}}(\boldsymbol{x}_0)\big),
\end{equation}
so the posterior becomes:
\begin{equation}
\label{eq:egoallo_posterior}
p(\boldsymbol{x}_0 \mid \boldsymbol{y}) \;\propto\;
p_\theta(\boldsymbol{x}_0 \mid
\boldsymbol{c}) \exp\big(-\mathcal{L}_{\text{guide}}(\boldsymbol{x}_0)\big).
\end{equation}
This is the form from Chapter~\ref{ch:conditional-inference}: a
learned prior combined with measurement constraints via an
energy-based likelihood.

% ==============================================================================
\subsection{Designing the Conditioning Representation}
\label{sub:egoallo_representation}
% ==============================================================================

The conditional prior $p_\theta(\boldsymbol{x}_0 \mid \boldsymbol{c})$ generates
body motion given a representation $\boldsymbol{c}$ of head motion. What
should this representation look like? As emphasized throughout this
book, representation choice fundamentally affects learning and generalization.

A naive approach would condition directly on the raw head poses
$\boldsymbol{h}_{1:T} \in \mathsf{SE}(3)^T$, the absolute position
and orientation
at each timestep. This works poorly. To understand why, consider what
properties the representation should satisfy.

\paragraph{Spatial Invariance.}
The mapping from head motion to body motion should not depend on
where in the world the motion occurs. Walking in the kitchen should
produce the same body pose estimates as walking in the living room.
Only the local motion pattern matters, not absolute position.

Conditioning on absolute poses violates this. As
Figure~\ref{fig:egoallo_conditioning} illustrates, identical body
motion corresponds to different absolute head trajectories depending
on world location. A model conditioned on absolute poses must learn
separately that ``head at $(3,2,1)$ moving forward'' and ``head at
$(10,5,2)$ moving forward'' both correspond to walking---an
unnecessary burden that harms generalization.

\paragraph{Temporal Invariance.}
The mapping should also not depend on when in an observation window
the motion occurs. Processing frames $0$--$100$ should behave the
same as processing frames $50$--$150$ if they contain the same
underlying motion. This rules out cumulative quantities like total
displacement from sequence start, as well as absolute time indices.

One might try canonicalizing to the first timestep by placing the
first frame at the origin. But this violates temporal invariance:
different temporal slices of the same motion, canonicalized to their
respective first frames, produce different trajectories. The model
sees different inputs for the same underlying body motion.

\paragraph{The Solution: Per-timestep Canonicalization.}
To achieve both spatial and temporal invariance, EgoAllo defines a
canonical coordinate frame at \emph{every} timestep rather than just
at the start. The representation $\boldsymbol{c} = \{\boldsymbol{c}^{1}, \dots,
\boldsymbol{c}^{T}\}$ couples:
\begin{itemize}
\item Relative head motion $\Delta\boldsymbol{T}_{\text{head}}^{t}$
between consecutive frames, capturing how the head moved from
$t-1$ to $t$.
\item Per-timestep canonicalized pose
$\boldsymbol{T}_{\text{canonical}}^{t}$: the head's orientation
relative to gravity and height above floor, with floor-plane
position and yaw removed.
\end{itemize}

These quantities are invariant to where the motion occurs in the
world and to which temporal window we extract. Two sequences with
identical local motion produce identical conditioning tokens
$\boldsymbol{c}$, regardless of absolute world position or time offset.

\begin{figure}[t!]
\centering
\includegraphics[width=2in]{\toplevelprefix/chapters/chapter8/figs/egoallo/conditioning_oct9.pdf}
\caption{\textbf{Invariant Head Motion Conditioning.} The
conditioning representation must be invariant to both spatial
transformations (same motion at different world locations) and
temporal shifts (same motion at different absolute times).
Per-timestep canonicalization achieves both properties.}
\label{fig:egoallo_conditioning}
\end{figure}

% ==============================================================================
\subsection{The Conditional Prior}
\label{sub:egoallo_architecture}
% ==============================================================================

The conditional prior $p_\theta(\boldsymbol{x}_0 \mid \boldsymbol{c})$ is a
diffusion model that generates body motion sequences given head
motion conditioning. The architecture follows the transformer-based
design discussed in Section~\ref{sec:image_classification}, adapted
for sequential motion data.

\paragraph{Architecture.}
The model takes three inputs:
\begin{itemize}
\item The noisy motion sequence $\boldsymbol{x}_n$ at diffusion step
$n$, tokenized per-timestep.
\item The diffusion timestep embedding, indicating the noise level.
\item The invariant head conditioning tokens $\boldsymbol{c}$, one per timestep.
\end{itemize}

The head tokens are first processed by an encoder consisting of six
transformer blocks with self-attention to produce contextualized
representations. These condition a decoder, another six transformer
blocks, that denoises the motion sequence. Conditioning is injected
via cross-attention: in each decoder layer, the motion tokens form
queries that attend to the head tokens as keys and values, allowing
the model to dynamically select which aspects of head motion are
relevant for predicting each body joint at each timestep. The model
outputs a direct estimate of the clean motion $\boldsymbol{x}_0$, rather
than predicting noise, parameterized as joint rotations in a
head-relative coordinate frame.

\paragraph{Training Data.}
The model is trained on AMASS \citep{mahmood2019amass}, a large-scale
motion capture dataset aggregating multiple sources into unified
SMPL-H format. AMASS contains diverse human motions---walking,
running, dancing, sitting, reaching, exercising---but no head-mounted
device recordings. Head trajectories are simulated by extracting head
joint position and orientation from the motion capture data, as if a
device were rigidly attached to the head.

\paragraph{Training Objective.}
Training uses the denoising score-matching loss from \Cref{ch:denoising}:
\begin{equation}
\mathcal{L}_{\text{diff}} = \mathbb{E}_{\boldsymbol{x}_0,
\boldsymbol{n}} \left[ w_n
\left\| D_\theta(\boldsymbol{x}_n, \boldsymbol{n}, \boldsymbol{c}) -
\boldsymbol{x}_0 \right\|^2 \right],
\end{equation}
where $D_\theta$ is the denoiser network, $w_n$ are noise-dependent
weights, and $\boldsymbol{x}_n$ is obtained by adding noise to clean motion
$\boldsymbol{x}_0$ according to the forward diffusion process. The model
learns to predict clean motion from noisy motion, conditioned on head
trajectories.

% ==============================================================================
\subsection{Incorporating Measurements via Guidance}
\label{sub:egoallo_guidance}
% ==============================================================================

The conditional prior captures the relationship between head motion
and body motion, but does not use image observations or physical
constraints. To incorporate these, EgoAllo defines a guidance energy
$\mathcal{L}_{\text{guide}}$ that measures consistency between a
motion hypothesis and the observations.

\paragraph{Guidance Energy.}
The total guidance energy combines several terms:
\begin{equation}
\mathcal{L}_{\text{guide}}(\boldsymbol{x}_0) =
\lambda_{3D}\mathcal{L}_{3D} + \lambda_{2D}\mathcal{L}_{2D} +
\lambda_{\text{skate}}\mathcal{L}_{\text{skate}} +
\lambda_{\text{prior}}\mathcal{L}_{\text{prior}},
\end{equation}
where:
\begin{itemize}
\item $\mathcal{L}_{3D}$: Alignment between estimated hand
keypoints and 3D hand detections from HaMeR, in the world frame.
\item $\mathcal{L}_{2D}$: Reprojection error. Estimated hand
keypoints, projected into the image, should match 2D detections.
\item $\mathcal{L}_{\text{skate}}$: Foot skating penalty. Feet
should not slide when in contact with ground.
\item $\mathcal{L}_{\text{prior}}$: Weak regularization toward
the diffusion prior's prediction, preventing guidance from
pushing too far from plausible motions.
\end{itemize}

Following the energy-based interpretation from
Chapter~\ref{ch:conditional-inference}, this guidance energy
corresponds to a negative log-likelihood: $p(\boldsymbol{y} \mid
\boldsymbol{x}_0) \propto
\exp(-\mathcal{L}_{\text{guide}}(\boldsymbol{x}_0))$. Low
energy means consistency with observations; high energy means inconsistency.

\paragraph{Why Reprojection Matters.}
Single-frame hand estimators like HaMeR produce 3D hand poses from
monocular images, but suffer from scale and depth ambiguity. A small
hand close to the camera looks the same as a large hand far away. The
3D positions are therefore unreliable in absolute terms, even when
local hand pose such as finger articulation is accurate.

The 2D reprojection loss addresses this by supervising in image
space, where the hand was actually observed. Estimated 3D hand
positions are projected back into the image and compared against
detected 2D keypoints. This sidesteps depth ambiguity while still
providing useful signal about hand location.

% ==============================================================================
\subsection{Conditional Sampling Algorithm}
\label{sub:egoallo_sampling}
% ==============================================================================

The measurement matching framework from
Chapter~\ref{ch:conditional-inference} tells us how to sample from
the posterior: modify the denoiser to include a gradient correction
toward measurement consistency. Recall that the conditional denoiser
decomposes as:
\begin{equation}
\mathbb{E}[\boldsymbol{x}_0 \mid \boldsymbol{x}_n, \boldsymbol{y}]
=
\mathbb{E}[\boldsymbol{x}_0 \mid \boldsymbol{x}_n]
+
\sigma_n^2 \nabla_{\boldsymbol{x}_n} \log p(\boldsymbol{y} \mid
\boldsymbol{x}_n),
\end{equation}
where the first term is the prior denoiser and the second is the
measurement matching correction.

In our setting, the prior denoiser is conditioned on head motion:
$D_n^{\text{head}}(\cdot; \boldsymbol{c})$. The measurement correction uses
the guidance energy. The full conditional denoiser becomes:
\begin{equation}
D_n^{\text{guided}}(\boldsymbol{z})
=
D_n^{\text{head}}(\boldsymbol{z};\boldsymbol{c})
-
\sigma_n^2
\nabla_{\boldsymbol{z}}
\mathcal{L}_{\text{guide}}(\boldsymbol{x}_0(\boldsymbol{z})),
\end{equation}
using the relationship $\nabla \log p(\boldsymbol{y} \mid \boldsymbol{x}) =
-\nabla \mathcal{L}_{\text{guide}}(\boldsymbol{x})$.

\paragraph{Alternating Optimization.}
In practice, rather than computing gradients through the full
guidance energy at each step, EgoAllo alternates between two operations:

\begin{enumerate}
\item \textbf{Prior denoising step.} Apply a DDIM-style update
using the head-conditioned denoiser
$D_n^{\text{head}}(\cdot;\boldsymbol{c})$ to obtain a proposal
$\tilde{\boldsymbol{x}}_{n-1}$. DDIM (Denoising Diffusion Implicit
Models) provides a deterministic sampling procedure that maps
noisy samples toward cleaner ones; unlike stochastic
samplers, each denoising step is a deterministic function of
the current state, which allows for faster sampling with fewer steps.

\item \textbf{Measurement-matching step.} Refine the proposal by
solving a proximal optimization problem:
\begin{equation}
    \boldsymbol{x}_{n-1}
    =
    \operatorname*{arg\,min}_{\boldsymbol{x}}
    \left\{
        \frac{1}{2}\|\boldsymbol{x}-\tilde{\boldsymbol{x}}_{n-1}\|^2
        +
        \gamma_n\,\mathcal{L}_{\text{guide}}(\boldsymbol{x})
    \right\},
\end{equation}
where $\gamma_n$ is a step-size parameter. This finds motion
close to the prior's prediction but more consistent with observations.
\end{enumerate}

The proximal optimization is solved using Levenberg-Marquardt, a
classic algorithm for nonlinear least-squares problems. It
interpolates between gradient descent and Gauss-Newton optimization,
adapting its step size based on how well the linearization
approximates the true objective. This is well-suited to SMPL-H
parameter estimation, where the objective involves differentiable
forward kinematics. A few iterations suffice at each denoising step.

\begin{figure}[t!]
\centering
\includegraphics[width=\linewidth]{\toplevelprefix/chapters/chapter8/figs/egoallo/egoallo-overview.png}
\caption{\textbf{Overview of components of the EgoAllo
framework.} The diffusion model is restricted to local body
parameters. An invariant
parameterization g(·) of SLAM (head) poses is used to
condition a diffusion model. These can be placed into the
global coordinate frame
via global alignment to input poses. When available, egocentric
video is used for hand detection, say via HaMeR
\citep{pavlakos2023hamer}, which can be incorporated into samples
via guidance.}
\label{fig:egoallo_framework}
\end{figure}

\paragraph{Global Alignment.}
The diffusion model operates in a head-relative coordinate frame,
predicting body pose relative to the head rather than in absolute
world coordinates. After the final denoising step, EgoAllo places the
body in the world by composing with the observed SLAM head poses:
\begin{equation}
\boldsymbol{T}_{\text{world},\text{root}}^t
=
\boldsymbol{T}_{\text{world},\text{head}}^t
\cdot
\boldsymbol{T}_{\text{head},\text{root}}^{(\theta^t, \beta)},
\end{equation}
where $\boldsymbol{T}_{\text{head},\text{root}}$ is computed via forward
kinematics from the estimated pose and shape. This ensures the body
is correctly grounded: feet on the floor, head aligned with the camera.

\paragraph{Handling Long Sequences.}
The diffusion model is trained on fixed-length sequences (32--128
frames). For longer recordings, EgoAllo uses the MultiDiffusion
approach \citep{bar2023multidiffusion}: the sequence is divided into
overlapping windows and denoised in parallel. After each denoising
step, the results in overlapping regions are blended by averaging
before proceeding to the next step. This per-step blending ensures
that windows remain consistent with each other throughout the
diffusion process, producing temporally coherent output without
discontinuities at window boundaries.

% ==============================================================================
\subsection{Insights from Experiments}
\label{sub:egoallo_experiments}
% ==============================================================================

EgoAllo is evaluated on three datasets: AMASS
\citep{mahmood2019amass} (motion capture with synthetic head
trajectories), RICH \citep{huang2022rich} (motion capture with
challenging human-scene interactions), and Aria Digital Twin
\citep{pan2023aria} (real recordings from Project Aria glasses with
ground-truth body poses). For hand estimation, EgoExo4D
\citep{grauman2023egoexo4d} is additionally used.

\paragraph{Representation Matters.}
Table~\ref{tab:egoallo_cond_ablation} shows the effect of different
conditioning representations on body estimation accuracy, comparing
five parameterizations that vary in their invariance properties:
\begin{itemize}
\item \textbf{Per-timestep invariant}: The representation
described in Section~\ref{sub:egoallo_representation}, with
both spatial and temporal invariance.
\item \textbf{Absolute + relative}: Absolute poses augmented with
local-frame relative motion.
\item \textbf{Absolute + global deltas}: Absolute poses with
world-frame position deltas.
\item \textbf{Sequence canonicalization}: Alignment to first
frame, providing spatial but not temporal invariance.
\item \textbf{Absolute}: Raw $\mathsf{SE}(3)$ poses with no
invariance properties.
\end{itemize}

\begin{table}[t]
\centering
\caption{\textbf{Effect of conditioning representation on AMASS.}
Different parameterizations of head motion conditioning lead to
substantially different estimation accuracy. The per-timestep
invariant representation, which satisfies both spatial and
temporal invariance, performs best. MPJPE is mean per-joint
position error in millimeters.}
\label{tab:egoallo_cond_ablation}
\small
\begin{tabular}{lcccc}
\toprule
& \multicolumn{2}{c}{\textbf{Seq Length 32}} &
\multicolumn{2}{c}{\textbf{Seq Length 128}} \\
\cmidrule(lr){2-3} \cmidrule(lr){4-5}
\textbf{Conditioning} & MPJPE & Effect & MPJPE & Effect \\
\midrule
Per-timestep invariant & \textbf{129.8} & --- & \textbf{119.7} & --- \\
Absolute + relative & 133.0 & +2.4\% & 124.5 & +4.0\% \\
Absolute + global deltas & 136.2 & +4.9\% & 127.4 & +6.4\% \\
Sequence canonicalization & 153.1 & +17.9\% & 134.0 & +11.9\% \\
Absolute & 159.9 & +23.2\% & 148.3 & +23.9\% \\
\bottomrule
\end{tabular}
\end{table}

The results show that conditioning representation has a dramatic
impact on accuracy, with up to 24\% difference in MPJPE between the
best and worst choices. This reinforces a central theme of this book:
the choice of representation fundamentally affects learning and generalization.

\paragraph{Height Estimation Enables Grounding.}
Table~\ref{tab:egoallo_body_results} compares full estimation
(including height) against a variant without height estimation.
Estimating height improves MPJPE by 6--7\% and is essential for
grounding: without it, estimated bodies often ``float'' above the
ground because the system cannot determine the correct vertical
position. The grounding metric (GND) measures whether estimated feet
ever contact the floor, and height estimation substantially improves this.

\begin{table}[t]
\centering
\caption{\textbf{Body estimation results across datasets.} MPJPE
and PA-MPJPE in millimeters; GND is the fraction of sequences
where feet contact the ground. Height estimation is essential for
proper grounding.}
\label{tab:egoallo_body_results}
\small
\begin{tabular}{llccc}
\toprule
\textbf{Dataset} & \textbf{Variant} &
\textbf{MPJPE}$\downarrow$ & \textbf{PA-MPJPE}$\downarrow$ &
\textbf{GND}$\uparrow$ \\
\midrule
AMASS & Full (with height) & \textbf{119.7} & \textbf{101.1}
& \textbf{1.00} \\
& Without height & 128.1 & 110.3 & 0.98 \\
\midrule
RICH  & Full (with height) & \textbf{176.2} & \textbf{160.1}
& \textbf{0.96} \\
& Without height & 185.7 & 169.9 & 0.82 \\
\midrule
ADT   & Full (with height) & \textbf{155.1} & \textbf{129.3} & 0.94 \\
& Without height & 163.7 & 140.0 & \textbf{0.96} \\
\bottomrule
\end{tabular}
\end{table}

\paragraph{Body Context Improves Hand Estimation.}
A notable finding is the symbiotic relationship between body and hand
estimation. Single-frame monocular hand estimators like HaMeR produce
accurate local hand poses, including finger articulation, but cannot
determine the hand's absolute 3D position due to scale and depth
ambiguity. By estimating the full body first, we gain kinematic
constraints: the hand is attached to the wrist, which connects to the
forearm, elbow, upper arm, and torso. This chain provides context for
global hand localization. Figure \ref{fig:egoallo_hands} shows some
qualitative results.

Table~\ref{tab:egoallo_hands} quantifies this effect on EgoExo4D. Raw
HaMeR estimates have 237.9mm world-frame error. Adding body context
reduces this to 131.5mm, a 45\% improvement, even using only
monocular hand observations. When additional 3D wrist observations
from stereo cameras are available, errors drop further to 60.1mm.

\begin{table}[t]
\centering
\caption{\textbf{Hand estimation on EgoExo4D.} MPJPE measures
world-frame accuracy; PA-MPJPE measures local pose accuracy. Body
context dramatically improves world-frame localization while
preserving local pose quality.}
\label{tab:egoallo_hands}
\small
\begin{tabular}{lcc}
\toprule
\textbf{Method} & \textbf{MPJPE}$\downarrow$ &
\textbf{PA-MPJPE}$\downarrow$ \\
\midrule
HaMeR (monocular, no body) & 237.9 & \textbf{13.0} \\
\midrule
With body (2D reprojection) & 131.5 & 14.7 \\
With body (3D wrist from stereo) & \textbf{60.1} & 14.4 \\
\bottomrule
\end{tabular}
\end{table}

The relationship is bidirectional: hand observations also improve
body estimation. When hand guidance is incorporated during sampling,
body MPJPE on AMASS improves from 119.7mm to 78.8mm, a 34\%
reduction. Additional measurement constraints tighten the posterior
and produce more accurate estimates throughout the kinematic chain.

\begin{figure}[t!]
\centering
\includegraphics[width=\textwidth]{\toplevelprefix/chapters/chapter8/figs/egoallo/hands_oct1_v4.png}
\caption{\textbf{Body Context Improves Hand Estimation.} Blue:
monocular hand estimates from HaMeR, which have accurate local
pose but incorrect world-frame position due to scale/depth
ambiguity. Purple: estimates with body context, correctly
grounded in the world frame.}
\label{fig:egoallo_hands}
\end{figure}

\paragraph{Qualitative Results.}
Figure~\ref{fig:egoallo_qualitative} shows estimated body poses from
real-world recordings with Project Aria glasses. The estimated bodies
are correctly grounded: feet contact the floor appropriately, body
height matches the environment, and poses are consistent with
observed head motion across diverse activities including walking,
sitting, and object manipulation.

\begin{figure}[t!]
\centering
\includegraphics[width=\textwidth]{\toplevelprefix/chapters/chapter8/figs/egoallo/qual_grounded_outputs_v3.png}
\caption{\textbf{Qualitative Results.} Estimated body poses from
real-world egocentric recordings, visualized in 3D scene
reconstructions. The method produces physically plausible poses
with appropriate grounding across diverse activities.}
\label{fig:egoallo_qualitative}
\end{figure}

% ==============================================================================
\subsection{Discussion}
\label{sub:egoallo_discussion}
% ==============================================================================

This application demonstrates how the conditional inference framework
from Chapter~\ref{ch:conditional-inference} extends beyond images to
structured, articulated outputs like human motion.

\paragraph{Representation Design Is Crucial.}
A recurring theme throughout this book is that representation choices
fundamentally affect learning and generalization. This principle is
particularly evident here. By analyzing what invariances the
conditioning representation should satisfy, we derived a
parameterization that substantially outperforms alternatives. Spatial
invariance ensures independence from world location; temporal
invariance ensures independence from time offset. The lesson: before
training a model, carefully consider what structure the
representation should encode.

\paragraph{The Prior-plus-guidance Pattern Generalizes.}
The sampling procedure directly applies the measurement matching
framework: a learned prior, here the head-conditioned diffusion
model, is combined with measurement constraints like hand
observations and physics via guided sampling. This same pattern of
prior plus guidance appears repeatedly across domains, from image
generation to motion estimation. The components are modular: we can
swap in different priors or different guidance terms without changing
the overall framework.

\paragraph{Holistic Estimation Outperforms Modular Pipelines.}
The symbiotic relationship between body and hand estimation suggests
that joint inference over the full kinematic chain outperforms
estimating components independently. Head motion informs body pose;
body pose provides context for hand position; hand observations
refine both body and hand estimates. This mutual benefit emerges
naturally from the posterior formulation, where all variables are
coupled through the likelihood and prior.

\paragraph{Limitations.}
Several limitations remain. First, the method requires accurate SLAM
with metric scale and gravity alignment; errors in head pose
estimation propagate directly to body estimates. Second,
generalization is limited by training data diversity, so unusual
poses or activities not represented in AMASS may be estimated poorly.
Third, hands can only be refined when visible in the egocentric view;
occluded hands rely entirely on the body prior. Finally, the method
assumes flat floors; staircases or uneven terrain would require
additional modeling.

\paragraph{Open Problems.}
Promising directions include incorporating scene geometry for better
physical grounding, where hands should contact surfaces and bodies
should avoid obstacles. Other directions include enabling real-time
inference for live AR/VR applications and extending to multi-person
scenarios where the wearer interacts with others. The framework is
flexible enough to accommodate these extensions through additional
guidance terms or modified priors.

\section{Natural Language Representation and Generation}\label{sec:clm_text}
So far, we have shown how to learn deep structured representations
for the distributions of data associated with visual perception and
body motions. To a large extent, the so learned structured
representations can be viewed as memories stored in the visual cortex
and motor cortex of the brain. These memories together can be loosely
referred to as a ``model of the world'', or simply a ``world model''
-- enable us to both perceive the world and take actions. Notice that
both human and animal have the ability to learn and develop such
memories for the world and their body, which record what can be
predicted from the perceived as well as motor skills crucial for
their survival within the environment.

On top of such a world model for the physical environment and body
motion, human has developed (spoken and written) languages so that we
are able to communicate knowledge learned about the world to others
-- the Broca's area and Wernicke's area in our brain are known for
speech processing and language comprehension, respectively. To a
large extent, languages record a small, but important, fraction of
our world models that are worth sharing among ourselves. Hence,
natural languages, as a very special type of data, must have very
special structures that encode very rich information. If machines
were able to learn the structures of natural languages, they may be
able to communicate better with humans.

Therefore, in this section, we show how to apply methods introduced
in this book to the task of \textit{language modeling}. Or more
specifically, how to train large language models (LLMs) to learn the
distribution of natural languages and then generate them. The task
and its setup is very much the same as that used to train GPT-2 and
many other language models.

\subsection{Data} \label{sub:clm_text_data}

The data we will use to investigate the performance of CRATE for
language tasks will be OpenWebText (OWT)~\cite{Gokaslan2019OpenWeb},
an open-source reproduction of the unreleased WebText dataset used by
OpenAI to train GPT2. Each sample in OWT is a web document, typically
sourced from high-quality web pages, blogs, articles, or online
discussions, that is written in well-formed natural language. The
OpenWebText dataset contains around 8.01M documents of varying
lengths, totaling around 41.70GB of text. For evaluation, we will use
several datasets, such as
WikiText~\cite{merity2016pointer}\footnote{For WikiText2 and
WikiText103~\cite{merity2016pointer}, the test splits are the same,
so we merge them as a single dataset referred to as WikiText.},
LAMBADA~\cite{paperno2016lambadadatasetwordprediction}\footnote{To
obtain the accuracy on the LAMBADA dataset, we use greedy decoding.},
and PTB~\cite{marcus-etal-1993-building}. PTB and OWT are generally
easier compared to other datasets. PTB focuses on simpler
journalistic text, ideal for traditional language modeling, while OWT
is diverse and informal, covering various topics but with less
complexity in language structure or long-range dependencies.
WikiText, with its formal structure and domain-specific content,
requires a more complex understanding than OWT but remains
manageable. LAMBADA is the most challenging, as it involves
long-range dependencies, requiring the model to grasp broader
contextual information to complete sentences accurately.

On a more formal level, our data \(\vX\) will be text, or strings of
characters; we let \(\cT\) be the set of all strings.

\subsection{Task and Objective} \label{sub:clm_text_objective}

For causal language modeling pre-training, the idea is that we want
to \textit{train the model to output human-like text}. The most
popular way to do this by far is to use a two-stage training
process:\footnote{Modern language model training has several
additional training steps which demand different data distributions
and algorithm approaches. However, training a model to merely mimic
human writing only requires these few presented steps.}
\begin{itemize}
\item \textit{First}, we wish to \textit{learn} a way to optimally
encode documents as a sequence of basic (``building block'')
strings, called \textit{tokens}. This is called
\textit{tokenization}, and we build a \textit{tokenizer}.
\item \textit{Second}, we wish to \textit{learn} a way to
\textit{predict the distribution of a token given all previous
tokens}. This is called \textit{next-token prediction}, and we
build a \textit{language model}.
\end{itemize}
This procedure actually dates back to Markov, who first noticed that
natural language could be modeled by the eponymous Markov chain
structure \citep{markov2006example} given an appropriate
tokenization, and then to Shannon, who proposed doing this exact
language modeling setup with a character-level tokenizer (i.e., each
character is a token) and so-called ``\(n\)-gram'' (i.e., an explicit
look-up table, calculated from training data, for the distribution of
a token given the \(n\) previous tokens) in place of the language
model \citep{Shannon-1948}.\footnote{A recent study
\citep{liu2024infini} scaling up \(n\)-gram models has shown that
they are able to model text reasonably well for large \(n\), but of
course the memory required to store such a lookup table is of order
\(V^{n}\) and hence completely intractable.}

\subsubsection{Training a Tokenizer}

To build a tokenizer amounts to building a vocabulary \(\cV\), which
is a set of tokens and has some pre-specified size \(V\). There are
several methods to do this. One popular algorithm is known as Byte
Pair Encoding (BPE), which can be described as:
\begin{itemize}
\item Start with a list of all unique characters in your training
data, and their frequencies. Ensure that there are fewer than
\(V\) such characters, and add each character as a separate
string (``token'') to the vocabulary along with its frequency.
\item Until there are \(V\) tokens in the vocabulary:
\begin{itemize}
    \item Construct a token by taking the two most frequent
        existing tokens and merging them.
    \item Compute this token's frequency in the dataset.
    \item Add it to the vocabulary (along with its frequency).
\end{itemize}
\item At this point, the frequency information is no longer needed
and can be discarded.
\end{itemize}
The overall process of BPE is in \Cref{fig:BPE}. Note that this
procedure is a modification of a classical information-theoretic
compression procedure for \textit{learning a lossless encoding} of
bytestream data (such as text), and as such, one can interpret it as
finding an optimal lossless compression of the data. Notice that this
is possible because (unlike images), the data here are fundamentally
discrete and noise-free.
\begin{figure}[t!]
\centering
\includegraphics[width=0.45\textwidth]{\toplevelprefix/chapters/chapter8/figs/BPE1.png}\hspace{0.3in}
\includegraphics[width=0.45\textwidth]{\toplevelprefix/chapters/chapter8/figs/BPE2.png}
\caption{\small {\bf The process of tokenizing text data using BPE.}
(Image credit to
\url{https://huggingface.co/learn/nlp-course/chapter6/5}). (Left) We
begin by analyzing the given text corpus and constructing an initial
vocabulary that consists of individual characters (or bytes in the
case of byte-level BPE). Then, we compute the frequencies of adjacent
character pairs in the corpus. This involves scanning the entire text
and counting how often each two-character sequence (bigram) appears.
(Right) After computing the frequencies of adjacent character pairs,
we identify the most frequent pair in the corpus. This pair is then
merged into a new subword unit, which is added to the vocabulary as a
single token. This process is repeated iteratively until the
predefined vocabulary size is reached. }
\label{fig:BPE}
\end{figure}
After such a vocabulary is built, a tokenizer can break down a
document into tokens (i.e., ``tokenize'' it). BPE uses a similar
procedure to tokenize data as in training:
\begin{itemize}
\item Separate the document into a long list of one-character-long
tokens. That is, if the document is ``Hello'' then the initial
list is `H', `e', `l', `l', `o'.
\item While any two adjacent tokens can be concatenated and their
concatenation is another token, we do it, i.e., we replace this
pair of tokens with the merged token. Namely, if `He' is a token
in the vocabulary, `H', `e', `l', `l', `o' would become `He', `l', `l', `o'.
\item Repeat the above process until no more merges can be done. At
this point, the document is partitioned into the final list
(sequence) of tokens.
\end{itemize}

There are many practical and efficiency-based considerations to take
into account during tokenization. The above algorithm, as presented,
is \textit{very far from optimal} if naively implemented, for
instance. We do not cover this topic in great detail; there are many
resources online to learn more, such as
\href{https://huggingface.co/learn/nlp-course/en/chapter6/5}{HuggingFace
tutorials}.

For instance, each token has a corresponding \textit{index} which is
just its index in the vocabulary (which after all is just a list of
length \(V\)). Thus, the output of most tokenizers is a list of
indices, say an element of \([V]^{*}\). Keep in mind that they
correspond to substrings of the original document, as shown above.

Once a tokenizer is learned, it can be used as a black box by any
language model. For instance, many models have the same
(OpenAI-based) tokenizer based on the \texttt{tiktoken} library. In
the remainder of this section, we will use such a fixed and pre-built
tokenizer for everything, and thus identify each text document \(\vX
\in \cT\) with its tokenized version in \([V]^{*}\). Therefore, we
may as well consider the text space \(\cT\) as \textit{equal} to the
space of token sequences \([V]^{*}\) (and lose nothing essential).

\subsubsection{Training a Language Model}

Once we have each document as a sequence of tokens \(\vX \in [V]^{N}
\subseteq [V]^{*} = \cT\), we wish to perform next-token prediction.
That is, given a \textit{context} \(\vX_{:n} \in [V]^{n}\) (i.e., the
first \(n\) tokens \(\vx_{1}, \dots, \vx_{n} \in [V]\) in the
document)\footnote{Note the incongruity with Python notation: here
the notation \textit{includes} index \(n\).}, we wish to predict the
token \(\vx_{n + 1} \in [V]\) at position \(n + 1\). To do this, we
compute the aggregate feature of \(\vX_{:n}\) via
\(\vz_{\theta}(\vX_{:n}) \doteq (f_{\theta}^{\ext} \circ
f_{\theta})(\vX_{:n}) \in \R^{d}\), and use a classification head
\(h_{\theta} \colon \R^{d} \to \Delta_{V}\) (implemented as either a
linear layer, MLP, or something slightly more complicated) to project
this feature into the \(V\)-dimensional probability simplex
\(\Delta_{V}\). This projection \(\vp_{\theta}(\vX_{:n}) \doteq
h_{\theta}(\vz_{\theta}(\vX_{:n}))\) serves as an estimated
probability distribution of the next token. Then, using the notation
\(\vone(\vx_{n + 1}) \in \Delta_{V}\) to be \(1\) in the \(\vx_{n +
1}\)th component and \(0\) elsewhere, the causal language modeling loss is
\begin{equation}\label{eq:clm_loss}
\min_{\theta}\bc{\cL_{\mathrm{CLM}}(\theta) \doteq
\Ex_{\vX}\rs{\frac{1}{N - 1}\sum_{n = 1}^{N - 1}\CE(\vone(\vx_{n +
1}), \vp_{\theta}(\vX_{:n}))}}
\end{equation}
Note how similar this is to a classification loss (say for images);
one uses the cross-entropy and tries to align a predicted probability
vector with the ground truth. The major difference between these two
losses is that in this one, we compute the loss on a whole sequence,
where each prediction is correlated with the others (unlike in the
i.i.d.~classification case).

Optimizing this loss is usually called ``pre-training'' in the
language model community (contrast with ``post-training'' and, more
recently, ``mid-training'', which are methodologies to modify a
next-token-predictor for useful tasks).

\textit{Side note:} Why does the first term of \eqref{eq:clm_loss}
predict \(\vone(\vx_{2})\), and there is no term which measures the
loss to predict the first token? It's because if we wanted to predict
the first token, we would have the \textit{empty sequence} as
context, and therefore make this first token prediction using a
qualitatively different mechanism than that which applies to the
other tokens. So actually this model is not trained to predict the
\textit{very first} token of any document. The reason this is OK is
due to an implementation detail of the tokenizer: often, after
building the tokenizer, we insert a \textit{special} token into its
vocabulary, called the \textit{beginning-of-string (or document)}
token and labeled \texttt{<|bos|>}.\footnote{There are usually
several special tokens for different purposes. Existing text
containing the special tokens are specially processed.} Then, while
processing each document, we add the \texttt{<|bos|>} token at the
beginning of the document's token sequence, increasing the length of
the tokenized sequence by \(1\). Thus the above causal language
modeling objective has a term which involves trying to predict the
first token of the document given only the \texttt{<|bos|>} token as
context, and so it is a conceptually correct loss.

\subsection{Architecture: Causal CRATE}

For the architecture, we use a standard GPT-2-style transformer,
substituting CRATE layers for the transformer layers.\footnote{In
direct contravention of the conventions in this book and those of
many other communities, the NLP community calls such GPT-2-style
transformers (encompassing nearly all current LLMs) ``decoder-only''
transformers. ``Encoder-only'' transformers have a different
architecture, and ``encoder-decoder'' transformers concatenate an
``encoder-only'' transformer with a ``decoder-only'' transformer.
This despite the fact that ``decoder-only'' transformers
\textit{also} compute an encoding of the data!} For completeness, we
specify the architecture here.

\paragraph{Embedding.} We first embed the token sequence \(\vX \in
[V]^{N}\) to Euclidean space. This is often done by associating each
index in \([V]\) with a vector in \(\R^{d}\) using a
\textit{massive}\footnote{By ``massive'' we mean that such a
structure is often a large fraction of the language model's total
size.} array \(\vE \in \R^{V \times d}\), and directly forming the
sequence \([\vE_{\vx_{1}}, \dots, \vE_{\vx_{N}}] \in \R^{d \times
N}\). The full embedding map \(f_{\theta}^{\emb}\) also applies a
positional encoding \(\vE^{\pos} \in \R^{d  \times N_{\max}}\) where
\(N_{\max}\) is the maximum number of tokens which are possible to
process,\footnote{Modern positional encoding methods have since taken
care of this issue and allowed for (in theory) infinite
extrapolation, but such methods are more complex to develop, and for
simplicity we only introduce the absolute additive positional
encoding here.} which yields the embedding map
\begin{equation}
f_{\theta}^{\emb}(\vX) \doteq [\vE_{\vx_{1}}, \dots, \vE_{\vx_{N}}] +
\vE_{:N}^{\pos}
\end{equation}
The parameters \(\vE\) and \(\vE^{\pos}\) are directly trainable.
Since \(\vE\) is so large (and the gradient update is very sparse
w.r.t.~it since only a small fraction of the vocabulary is used in
each sample), specialized software is used to make sure the memory
updates are not too onerous. Notice also that we do not use a class
token like in the other sections; more on this later.

\paragraph{Backbone.} We process the embeddings using a CRATE-like
backbone which uses causal masking. To motivate causal masking,
consider the causal language modeling loss \(\cL_{\mathrm{CLM}}\)
defined in \eqref{eq:clm_loss}. The most naive implementation would
require us to compute the forward pass \(N\) times in order to
backpropagate once. Obviously this is extremely inefficient, since
\(N\) can often be in the thousands. In order to scale training with
this loss efficiently, we impose a \textit{causal} constraint, i.e.,
\begin{equation}\label{eq:causal_backbone_def}
\vZ_{\theta}(\vX_{:n}) = \vZ_{\theta}(\vX)_{:n}
\end{equation}
i.e., the \(n\) columns of the token features
\(\vZ_{\theta}(\vX_{:n}) \in \R^{d \times n}\) should be the same as
the first \(n\) columns of the token features \(\vZ_{\theta}(\vX) \in
\R^{d \times N}\) regardless of the positive values of \(n\) and
\(N\) such that \(N \geq n\). In effect, this means we can apply the
backbone \textit{once} to the whole sequence and compute
\(\vZ_{\theta}(\vX)\), then apply \(f_{\theta}^{\ext}\) to each
increasing subset \(\vZ_{\theta}(\vX_{:n}) = \vZ_{\theta}(\vX)_{:n}\)
as \(n\) grows to the sequence length \(N\). Then we can use all of
those to compute the loss.

So now that we want a causal architecture for the backbone, how can
we get it? Since the MLP and layer normalizations inside each
transformer layer affect each token individually, the only thing that
matters for causality is the attention block (or \(\MSSA\) in the
case of CRATE). In order to make \(\MSSA\) causal, we define the
\(\mathrm{CMSSA}\) block as
\begin{align}
&\operatorname{CMSSA}_{\theta}^{\ell}(\vZ) \doteq
\vU_{\out}^{\ell}\mat{\operatorname{CSA}([\vU^{1,
    \ell}]^{\top}\vZ, [\vU^{1, \ell}]^{\top}\vZ, [\vU^{1,
\ell}]^{\top}\vZ) \\ \vdots \\ \operatorname{CSA}([\vU^{K,
    \ell}]^{\top}\vZ, [\vU^{K, \ell}]^{\top}\vZ, [\vU^{1,
\ell}]^{\top}\vZ)} + \vb_{\out}^{\ell}\vone_{N}^{\top} \\
&\text{where} \quad  \operatorname{CausalSA}(\vQ, \vK, \vV) \doteq
\vV\softmax\rp{\frac{\operatorname{CausalMask}(\vK^{\top}\vQ)}{\sqrt{p}}}, \\
&\text{where} \quad  \operatorname{CausalMask}(\vM)_{ij} =
\casework{M_{ij}, & \text{if}\ i \geq j, \\ -\infty, & \text{if}\ i < j}.
\end{align}
Here, practitioners say that the causal mask \textit{allows future
tokens \(i\) to attend to past tokens \(j\) but not vice versa}. To
see why, let us write out the expression for the \(t\)th column of
\(\operatorname{CSA}(\vQ, \vK, \vV)\):
\begin{equation}
\operatorname{CSA}(\vQ, \vK, \vV)_{t} = \sum_{i =
1}^{t}\vV_{i}\softmax\rp{[\vK_{:t}]^{\top}\vQ_{t}}_{i}
\end{equation}
(where here the non-colon subscript denotes the column). This
expression for the \(t\)th token uses no information about any token
beyond index \(t\). Therefore \(\operatorname{CSA}\), hence
\(\operatorname{CMSSA}\), hence the whole causal CRATE backbone
is causal in terms of the definition in
\eqref{eq:causal_backbone_def}, and we unlock the considerable
efficiency gains that we were promised.

\paragraph{Feature extractor.} We use a post-processing step
\(f_{\theta}^{\ext}\) which extracts the \textit{feature vector of
the last known token} so as to predict the next token. In theory,
this means that each token \(\vZ_{\theta}(\vX)_{n}\) should contain
rich information about all tokens that come before or at index \(n\),
i.e., \(\vx_{1}, \dots, \vx_{n}\), as all of this information should
be available for predicting the next token at index \(n + 1\). In
practice, only a few of these tokens are really needed for each
prediction task. Anyways, the equation for \(f_{\theta}^{\ext}\) is
\begin{equation}
f_{\theta}^{\ext}(\vZ_{\theta}(\vX_{:n})) \doteq (\vZ_{\theta}(\vX))_{n}
\end{equation}
where (again) the non-colon subscript is the column. In this case, as
promised, we just directly extract the feature vector of the last
token in the sequence.

\paragraph{Task-specific head.} For our classification head
\(h_{\theta}\), the GPT-2 architecture uses a simple linear layer and
a softmax to get the desired probability vectors:
\begin{equation}
h_{\theta}(\vz) \doteq \softmax(\vW^{\out}\vz + \vb^{\out}),
\end{equation}
where \(\vW^{\out} \in \R^{V \times d}, \vb^{\out} \in \R^{V}\). Some
other more modern architectures use small MLPs and layer
normalizations, but the idea is very much the same. Note that this
linear layers also have large memory usage (because \(V\) is very
large), and form a bottleneck in training; there has been significant
effort attempting to circumvent it.

All these architectural choices mean that causal training is
extremely efficient relative to non-causal training:
\begin{itemize}
\item We only need \textit{one forward pass} through the backbone to
compute the loss for the whole sequence.
\item The feature extraction is basically \textit{free}.
\item All tokens can be pushed through the task-specific head
\textit{in parallel}.
\end{itemize}

\subsection{Optimization Strategy}

We train our language model using end-to-end stochastic optimization.
One remaining issue is that, in practice, different documents in a
batch have different lengths (in terms of the number of tokens
required for each sequence), but as of the time of writing this book,
the main deep learning frameworks by and large allow only
``rectangular'' tensors, which do not accommodate this behavior. To
try to resolve this issue, we just insert a special padding token
\texttt{<|pad|>} for all shorter samples in the batch, so that we can
batch-process everything using rectangular tensors. At each timestep \(k\), we:
\begin{itemize}
\item Subsample \(B\) different tokenized documents
\(\{\vX_{b}^{(k)}\}_{b = 1}^{B} \subseteq \cT = [V]^{*}\), each
with length \(N_{b}^{(k)}\).
\item Compute \(N_{\max}^{(k)} \doteq \max_{b \in [B]}N_{b}^{(k)}\)
and pad each \(\vX_{b}^{(k)}\) to length \(N_{\max}^{(k)}\) using
a special padding token.
\item Compute the features \(\vZ_{\theta}(\vX_{b}^{(k)})\).
\item Compute the predicted distributions \(\vp_{\theta}(\vX_{b,
    :n}^{(k)}) \doteq (h_{\theta} \circ
f_{\theta}^{\ext})(\vZ_{\theta}(\vX_{b}^{(k)})_{:n})\).
\item Form the surrogate stochastic loss
\begin{equation}
    \hat{\cL}_{\mathrm{CLM}}^{(k)}(\theta) \doteq
    \frac{1}{B(N_{\max}^{(k)} - 1)}\sum_{b = 1}^{B}\sum_{n =
    1}^{N_{\max}^{(k)}-1}\CE(\vone(\vx_{b, n + 1}^{(k)}),
\vp_{\theta}(\vX_{b, :n}^{(k)}))).
\end{equation}
\item Compute one step of an optimization algorithm on \(\theta\),
giving the following iteration:
\begin{equation}
\theta^{(k + 1)} \doteq \textsc{OptUpdate}^{(k)}(\theta^{(k)};
\nabla_{\theta}\hat{\cL}_{\mathrm{CLM}}^{(k)}).
\end{equation}
\end{itemize}

\subsection{Evaluation Methodology} \label{sub:clm_text_evals}

There are several ways to evaluate a trained transformer language model.
\begin{itemize}
\item On a holdout dataset of arbitrary text, we can evaluate
\(\cL_{\mathrm{CLM}}\) on it; lower losses are better since they mean
that the model has learned that distribution well.
\item On a multiple choice question dataset, for each question we can
put it as the context and check the estimated probability of the
correct answer being generated.
\item We can also test the \textit{text generation} capabilities.
Namely, we can repeatedly sample from the model's probability
distribution over the next token given the context. Each time we
sample we generate a new token, which we print and add to the
context. This allows us to sample from the LLM, and judge the
generated samples however we please.\footnote{Having to re-run the
model on each token every time can become prohibitively expensive.
Clever storages of different internal features of the language model
(such as the so-called \textit{\(K\)-\(V\) cache}), along with the
causality of the architecture, can dramatically reduce the cost of sampling.}
\end{itemize}
In this section, we perform the first kind of evaluation exclusively.

\subsection{Experimental Setup and Results}

Since our causal CRATE architecture is directly built upon GPT-2, we
compare the optimal settings for GPT-2 as given by the NanoGPT
repository \citep{nanogpt} with the same settings applied to CRATE
for a fair comparison.

\paragraph{Model architecture.} We use the GPT-2 tokenizer, which has
vocabulary size \(V = 50257\), including a special token for
\texttt{<|pad|>}.\footnote{The \texttt{<|bos|>} token is not included
in this setup, although it is very common in modern language models.}
The context length is \(N_{\max} = 1024\). The backbone model follows
the GPT2-Base architecture \citep{radford2019language} with the
appropriate alterations to have causal CRATE layers, and we compare
against GPT2-Small and GPT2-Base.

\paragraph{Datasets and optimization.} For training causal CRATE, we
follow the implementations in the NanoGPT repository \citep{nanogpt}.
Specifically, we use a batch size of 384 and train for 600,000 steps
with the Adam optimizer~\citep{kingma2014adam}. For the Adam
optimizer, we use $(\beta_1, \beta_2)=(0.9, 0.95)$ and weight decay
of $0.1$. For the learning rate schedule, we apply a linear warm-up
and cosine decay, with a peak value of $\eta=6\times 10^{-4}$ at the
$2,000$th iteration, and minimum value $6\times 10^{-5}$. The
training and validation losses over iterations are shown in
\Cref{fig:crate-text-evals}. The training/validation loss converges
around $3.37$ after training with a batch size of $384$ and $600,000$
iterations. In comparison, the open GPT-2 implementation is
pre-trained on OpenWebText with a batch size of $512$ and $600,000$
steps and converges to a validation loss of $2.85$ \citep{nanogpt}.

\begin{figure}[t!]
\centering
\includegraphics[width=0.5\textwidth]{\toplevelprefix/chapters/chapter8/figs/gpt-loss.pdf}
\caption{\bf The loss curve of CRATE-GPT-Base trained on the
OpenWebText dataset.}
\label{fig:crate-text-evals}
\end{figure}

\paragraph{Experiment results.}

\Cref{tab:gpt-eval} demonstrates that CRATE models achieve reasonable
performance on the causal language modeling loss across a variety of
datasets compared to GPT-2 models with similar parameter counts and
similar architectures.

\begin{table}
\def\arraystretch{1.1}
\small
\caption{\small Zero-shot cross-entropy loss of the CRATE-GPT2-Base
model and GPT2-Small, GPT2-Base model evaluated on the test split of
the datasets ($\downarrow$ lower is better).
}
\centering
\resizebox{1\columnwidth}{!}
{
\begin{tabular}{ccccccc}
\hline
& \#parameters & \textbf{OWT} & \textbf{LAMBADA} &
\textbf{WikiText} & \textbf{PTB} & \textbf{Avg} \\
\hline
GPT2-Base  & {124M} & 2.85$\downarrow$ & 4.12$\downarrow$ &
3.89$\downarrow$ & 4.63$\downarrow$ & 3.87$\downarrow$ \\
{GPT2-Small } &  {64M} & {3.04} & {4.49} & {4.31} & {5.15} & {4.25} \\
Causal-CRATE-Base & {60M} & 3.37 & 4.91 & 4.61 & 5.53 & 4.61 \\
\hline
\end{tabular}
}
\label{tab:gpt-eval}
\end{table}

\section{Scaling and Improving White-Box Transformers
}\label{sec:scalable}\label{sec:improvements}

In this last section, we will discuss several ways in which various
parts of CRATE-type models can be scaled up or made more efficient
for certain special tasks while still remaining fully interpretable
white-box. These developments mix both conceptual and empirical
insights, and can be viewed as case studies about how to use
white-box understanding to improve deep learning models in practice.
The tasks that we use to evaluate the methods will be image
classification and next-token-prediction, the data will be ImageNet
and OpenWebText respectively, the optimization procedure will be the
same backpropagation, and the only thing that changes is the architecture.

\subsection{Increasing Network Width:
CRATE-\texorpdfstring{\(\alpha\)}{alpha}}\label{sub:crate_alpha_experiments}

\begin{figure}[t!]
\centering
\includegraphics[width=\textwidth]{\toplevelprefix/chapters/chapter8/figs/crate_alpha_backbone.pdf}
\caption{\small\textbf{One layer of the CRATE-\(\alpha\) backbone.}
The difference from CRATE is that the \(\ISTA_{\theta}^{\ell}\) block
is replaced by the \(\operatorname{ODL}_{\theta}^{\ell}\) block,
which performs several \(\ISTA\) steps with an overcomplete dictionary.}
\label{fig:crate_alpha_backbone}
\end{figure}

One design decision enforced by the CRATE framework is the
\textit{width} of the nonlinearity in the network. In a regular
transformer, the width is usually set to \(4\), \(8\), or
\(\frac{11}{3}\) times the feature dimension. However, CRATE enforces
that the width is exactly equal to the feature dimension, i.e., the
dictionaries \(\vD^{\ell}\) are square, which could lead to reduced
performance. The fundamental reason that the CRATE framework
constrains us to this choice is as follows:
\begin{itemize}
\item The ISTA block takes a \textit{single} step of optimization for
dictionary learning.
\item Usually one step of any iterative optimization algorithm cannot
effectively optimize the objective. So then why does this work?
\item Optimization algorithms usually converge very quickly if and
only if they have good initializations, or \textit{warm starts}. The
ISTA block has a warm start --- it treats the input features as an
initialization to the resulting sparse codes.
\item This enforces that the input features and sparse codes have the
same dimension. Namely, ISTA learns a complete sparsifying dictionary
(cf \Cref{ch:classic-models}).
\end{itemize}
Thus if we want to use a wide dictionary, we need ISTA to perform
\textit{overcomplete} dictionary learning. This means we cannot have
the same warm start (as our sparse codes have a larger dimension than
our features), and need more iterations to converge to a sparse code.
Hence the step from features \(\vZ_{\theta}^{\ell + 1/2}\) to sparse
codes \(\vZ_{\theta}^{\ell + 1}\) would no longer be
\begin{equation}
\vZ_{\theta}^{\ell + 1} = \ISTA_{\theta}^{\ell}(\vZ_{\theta}^{\ell +
1/2} \mid \vZ_{\theta}^{\ell + 1/2})
\end{equation}
where the \(\ISTA_{\theta}^{\ell}\) function is defined as (by an
abuse of notation from earlier sections)
\begin{equation}
\ISTA_{\theta}^{\ell}(\vZ \mid \vY) \doteq \ReLU(\vZ - \beta
(\vD^{\ell})^{\top}(\vD^{\ell}\vZ - \vY) + \beta \lambda
\vone_{s}\vone_{n}^{\top})
\end{equation}
but rather the following iteration:
\begin{equation}
\vZ_{\theta}^{\ell + 1} = \vA_{\theta}^{\ell, T}; \qquad
\vA_{\theta}^{\ell, t + 1} =
\ISTA_{\theta}^{\ell}(\vA_{\theta}^{\ell, t} \mid \vZ_{\theta}^{\ell
+ 1/2}) \quad \forall 0 \leq t < T; \qquad \vA_{\theta}^{\ell, 0} =
\vzero_{s \times n},
\end{equation}
i.e., running proximal gradient on the LASSO objective for \(T \geq
1\) steps in the forward pass at each layer, initialized at
\(\vzero_{s \times n}\). In this circumstance, the dictionary can be
as wide as needed, i.e., \(\vD^{\ell} \in \R^{s \times d}\) where \(s
\geq d\) (usually one takes \(s = 4d\) in practice).

\begin{table}[b]
\caption{\small \textbf{Object detection and fine-grained
segmentation via MaskCut on COCO val2017~\citep{lin2014microsoft}}.
Here all models are trained with patch size \(8\) instead of \(16\).
Compared with existing models such as CRATE and ViT, the
CRATE-\(\alpha\) model family noticeably has improved performance as
well as scalability.}
\centering
\begin{tabular}{@{}lcccccccc@{}}
\toprule
&  & \multicolumn{3}{c}{Detection} &  \multicolumn{3}{c}{Segmentation} \\
Model & AP$_{50} \uparrow $ & AP$_{75} \uparrow $ & AP $\uparrow$ &
AP$_{50} \uparrow$ & AP$_{75} \uparrow $ & AP $\uparrow$ \\
\midrule
\midrule
CRATE-\(\alpha\)-B/8 & 3.5 & 1.1 & 1.5 & 2.2 & 1.0 & 1.1 \\
CRATE-\(\alpha\)-L/8 & \textbf{4.0} & \textbf{1.7} & \textbf{2.0} &
\textbf{2.7} & \textbf{1.1} & \textbf{1.4} \\
\midrule
\color{gray}CRATE-B/8 & \color{gray}2.9 & \color{gray}1.0 &
\color{gray}1.3 & \color{gray}2.2 & \color{gray}0.7 & \color{gray}1.0 \\
\color{gray}ViT-B/8 & \color{gray}0.8 & \color{gray}0.2 &
\color{gray}0.4 & \color{gray}0.7 & \color{gray}0.5 & \color{gray}0.4 \\
\bottomrule
\end{tabular}
\label{tab:crate_alpha_detection_segmentation}
\end{table}

However, this presents an empirical problem. Using the above
configuration, if \(\vZ^{\ell + 1/2} \in \R^{d \times n}\), then
\(\vZ^{\ell + 1} \in \R^{s \times n}\), which can have an arbitrarily
large feature dimension. In practice, we want the feature dimension
at each layer to be the same. So this sets up a practical trichotomy
for designing the network, namely, we \textit{cannot} have
\textit{all} of the following desiderata:
\begin{enumerate}
\item The feature dimension at each layer is the same.
\item The dictionary is wide, i.e., overcomplete.
\item The output of the nonlinearity is the sparse codes of the input
with respect to the dictionary.
\end{enumerate}
In practice, giving up (1) is less tractable for efficiency reasons.
Giving up (2) leads to the usual CRATE framework. Giving up (3) leads
to a wide version of CRATE, i.e., CRATE-\(\alpha\), which has the
following nonlinearity to get from \(\vZ^{\ell + 1/2}\) to \(\vZ^{\ell + 1}\):
\begin{equation}
\vZ_{\theta}^{\ell + 1} = \vD^{\ell}\vA_{\theta}^{\ell, T}; \qquad
\vA_{\theta}^{\ell, t + 1} =
\ISTA_{\theta}^{\ell}(\vA_{\theta}^{\ell, t} \mid \vZ_{\theta}^{\ell
+ 1/2}); \qquad \vA_{\theta}^{\ell, 0} = \vzero,
\end{equation}
i.e., it takes the sparse codes obtained via proximal gradient
descent and multiplies by the dictionary to get the denoised version
of the input. Thus CRATE-\(\alpha\)'s nonlinearity computes a
denoised version of the input which is amenable to sparse coding, not
the actual sparse codes themselves. The map from \(\vZ_{\theta}^{\ell
+ 1/2}\) to \(\vZ_{\theta}^{\ell + 1}\) here is called the
Overcomplete Dictionary Learning (ODL) block and denoted
\(\operatorname{ODL}_{\theta}^{\ell}\), i.e.,
\begin{equation}
\vZ_{\theta}^{\ell + 1}(\vX) \doteq
\operatorname{ODL}_{\theta}^{\ell}(\vZ_{\theta}^{\ell + 1/2}(\vX)).
\end{equation}

\begin{figure}[t!]
\centering
\includegraphics[width=\textwidth]{\toplevelprefix/chapters/chapter8/figs/crate_alpha_semantic_heads.pdf}
\caption{\small\textbf{Saliency maps from CRATE-\(\alpha\) with patch
size \(8\).} Each row is a different image and each column
corresponds to a different attention head in the last layer. We
observe that the saliency maps strongly correspond to the objects in
the input image.}
\label{fig:crate_alpha_saliency_maps}
\end{figure}

\begin{table}
\caption{\small\textbf{Validation loss in language modeling.} Here
all models are pre-trained on most of OpenWebText, and the validation
cross-entropy loss is measured on a hold-out subset of OpenWebText.
CRATE-\(\alpha\) shows significant improvement over the CRATE design,
though there still exists a gap with traditional transformers like GPT-2.}
\centering
\begin{tabular}{@{}lcccc@{}}
\toprule
Model & GPT-2-B(ase) & CRATE-B & CRATE-\(\alpha\)-S(mall) &
CRATE-\(\alpha\)-B \\
\midrule
\midrule
\# parameters & 124M & 60M & 57M & 120M \\
OWT val.~loss & 2.85 & 3.37 & 3.28 & 3.14 \\
\bottomrule
\end{tabular}
\label{tab:crate_alpha_lm}
\end{table}

The CRATE-\(\alpha\) layer is shown in
\Cref{fig:crate_alpha_backbone}. In practice this modification of
CRATE performs very well at larger scales. For example, when we
pre-train CRATE-\(\alpha\) models on ImageNet-21K, unsupervised tasks
like segmentation (see \Cref{fig:crate_alpha_saliency_maps} and
\Cref{tab:crate_alpha_detection_segmentation}) generally have
significantly improved performance compared to CRATE. Similar trends
are present in language model training using causal self-attention
(see \Cref{tab:crate_alpha_lm}). Overall, it is a promising avenue to
scaling up the performance to match black-box models such as
transformers.\footnote{Note that the experimental results in this
section use a slightly different model architecture, which add very
slight empirical gains. The changes are: (1) an additional residual
connection on the ODL block, (2) modifying \(\ISTA\) to use two
separate dictionaries instead of \(\vD^{\ell}\) and \((\vD^{\ell})^{\top}\).}

\subsection{Linear Time Complexity Transformers}\label{sub:tost_experiments}

In practice, deep learning models suffer from bottlenecks in space
and time complexity, representing problem sizes beyond which they
cannot scale given fixed resources. One such bottleneck, particularly
meaningful when dealing with data where each sample is itself
high-dimensional and rich (such as long streams of text or videos),
is the \textit{time complexity} of processing long sequences of data.
In order to alleviate the time complexity of processing data using
transformers, in \Cref{sub:tost} we proposed a \textit{token
statistics self-attention} operator \(\TSSA_{\theta}^{\ell}\). We now
build a \textit{token statistics transformer}, called ToST:\\
\centerline{\url{https://robinwu218.github.io/ToST/},} around it,
which we can use for long-context tasks. In particular, we can
use the following layer (depicted in \Cref{fig:tost_backbone}) as a
drop-in replacement for a backbone layer in CRATE:
\begin{align}
\vZ_{\theta}^{\ell + 1/2}(\vX)
&= \vZ_{\theta}^{\ell}(\vX) + \TSSA_{\theta}^{\ell}(\LN_{\theta}^{1,
\ell}(\vZ_{\theta}^{\ell}(\vX))) \\
\vZ_{\theta}^{\ell + 1}(\vX)
&= \vZ_{\theta}^{\ell + 1/2}(\vX) +
\MLP_{\theta}^{\ell}(\LN_{\theta}^{2, \ell}(\vZ_{\theta}^{\ell + 1/2}(\vX)))
\end{align}
where the \(\TSSA\) block is defined as in \Cref{sub:tost}. Notice
that this is exactly the same as the vision transformer architecture
discussed in \Cref{sub:contrastive_learning_architecture}, except
that \(\TSSA\) replaces the conventional multi-head self-attention
block \(\MHSA\). Regardless, the computational complexity of the
forward pass of this layer is linear in all problem variables ---
sequence length, feature dimension, number of heads, and head dimension.

\begin{figure}[t!]
\centering
\includegraphics[width=\textwidth]{\toplevelprefix/chapters/chapter8/figs/tost_backbone.pdf}
\caption{\small\textbf{One layer of the ToST backbone}. Token
representations go through layer-norms, the token statistics
self-attention (TSSA) operator, and an MLP, in order to form the
layer's output.}
\label{fig:tost_backbone}
\end{figure}

\begin{table}[t]
\centering
\resizebox{\textwidth}{!}{
\begin{tabular}{@{}lccc|cc|cc@{}}
\toprule
Datasets & ToST-T(iny) & ToST-S(mall) & ToST-M(edium) &  {
\color{gray} XCiT-S} &  { \color{gray}XCiT-M} &
{\color{gray}ViT-S} & {\color{gray}ViT-B(ase)}\\
\toprule
\# parameters & 5.8M & 22.6M & 68.1M &  { \color{gray}24.9M }& {
\color{gray}80.2M } & { \color{gray} 22.1M } & { \color{gray}86.6 M } \\
\midrule
ImageNet & 67.3 & 77.9 & 80.3 &  { \color{gray} 80.5} & {
\color{gray} 81.5 } & { \color{gray} 79.8} & { \color{gray} 81.8} \\
ImageNet ReaL & 72.2  & 84.1 & 85.6 &  { \color{gray} 85.6 } & {
\color{gray} 85.9} & { \color{gray} 85.6 } & { \color{gray} 86.7 }\\
\midrule
CIFAR10 & 95.5 & 96.5 & 97.5 &  { \color{gray} 98.1} & {
\color{gray} 98.3} & { \color{gray} 98.6} & { \color{gray} 98.8}\\
CIFAR100 & 78.3 & 82.7 & 84.5 &  { \color{gray} 86.1} & {
\color{gray} 87.6} & { \color{gray} 88.8} & { \color{gray} 89.3}\\
Oxford Flowers-102 & 88.6 & 92.8 & 94.2 &  { \color{gray} 93.9} &
{ \color{gray} 94.0} & { \color{gray} 94.0}& { \color{gray} 95.7}\\
Oxford-IIIT-Pets & 85.6 & 91.1 & 92.8 &  { \color{gray} 92.9} & {
\color{gray} 94.0} & { \color{gray} 92.8} & { \color{gray} 94.1} \\
\bottomrule
\end{tabular}%
}
\caption{\small \textbf{Linear probing classification accuracy of
ToST} on various datasets with different model sizes when the
backbone is pre-trained for ImageNet-1K classification. We observe
that, compared to the XCiT (a empirically-designed transformer-like
architecture specialized for efficient processing of long sequences)
and the ViT, ToST maintains relatively similar performance, even
while enjoying benefits like faster runtime and white-box design.}
\label{tab:tost_linear_probing}
\end{table}

\begin{table}[!htbp]
\centering
\begin{tabular}{@{}lcccccc@{}}
\toprule
Model & \# params & OWT & Lambada & Wikitext & PTB & Avg $\downarrow$
\\ \midrule
GPT-2-Base & 124M & 2.84 & 4.32 & 4.13 & 5.75 & 4.26 \\
ToST-Base & 110M & 3.20 & 4.98 & 4.77 & 6.39 & 4.84 \\
ToST-Medium & 304M & 2.88 & 4.45 & 4.30 & 5.64 & 4.32 \\
ToST-Large & 655M & 2.72 & 4.32 & 3.99 & 5.03 & 4.02 \\ \bottomrule
\end{tabular}%
\caption{\small\textbf{Language modeling validation loss} computed on
(holdout sets of) a variety of natural language datasets, after
pre-training the model on that dataset. We observe that ToST scales
well, so that ToST-Large surpasses the baseline GPT-2-Base in causal
language modeling, while enjoying superior efficiency in long contexts.}
\label{tab:tost_lm}
\end{table}

\begin{table}
\centering

\begin{tabular}{@{}lccccccc@{}}
\toprule
Model        & ListOps  & Text     & Retrieval & Image    &
Pathfinder & Avg      \\
\midrule
\midrule
Reformer              & \textbf{37.27} & 56.10             & 53.40
& 38.07             & 68.50               & 50.56             \\
BigBird               & 36.05             & 64.02             & 59.29
& 40.83             & 74.87                & 54.17             \\
LinFormer         & 16.13             & \underline{65.90} & 53.09
& 42.34             & \underline{75.30}               & 50.46             \\
Performer             & 18.01             & 65.40             & 53.82
& 42.77             & \textbf{77.05}                & 51.18             \\
Transformer           & 37.11             & 65.21             &
\underline{79.14}              & \underline{42.94}             &
71.83              & \underline{59.24}            \\
ToST & \underline{37.25}    & \textbf{66.75}    & \textbf{79.46}
& \textbf{46.62}    &    69.41      &     \textbf{59.90}\\

\bottomrule
\end{tabular}%
\caption{\small \textbf{Long-Range Arena (LRA) performance comparison
of ToST(-B) versus the top transformer variants optimized for
long-context.} Long-Range Arena is a family of benchmarks that test
the long sequence modeling capability of algorithms and
architectures, by fixing the dataset and evaluation mechanism. ToST
scores at the top of the leaderboard compared to all known
transformer variants, including XCiT and the regular (ViT)
transformer (cf \Cref{tab:tost_linear_probing}). Moreover, ToST has
the lowest time- and space-complexity inference. (In this table, the
best score for a particular benchmark is bolded, and the second-best
score is underlined.)}
\label{tab:tost_lra_results}
\end{table}

Moreover, the proposed architecture, named ToST (for ``Token
Statistics Transformer'') performs well at vision tasks (i.e.,
\Cref{tab:tost_linear_probing}) and language tasks (i.e.,
\Cref{tab:tost_lm}). This is especially true for long-sequence-length
tasks (cf \Cref{tab:tost_lra_results}), where it is both more
performant and much more efficient than conventional transformers and
all other transformer-like architectures.

\subsection{Attention-Only Transformers} \label{sub:aot_experiments}

Another bottleneck to remove from deep learning models, specifically
transformer-like architectures, is the memory bottleneck that comes
from massive matrix multiplications in MLPs, where the internal
dimension is far greater than the feature dimension \(d\). It thus is
an interesting and important question to ask: do we \textit{really}
need the MLP inside a transformer, and how good can the performance
get without it? To explore this question, we use the
attention-only-transformer (AoT) architecture (see \Cref{sub:aot}),
depicted in \Cref{fig:aot_backbone}. Namely, each layer is simply of the form
\begin{equation}
\vZ_{\theta}^{\ell + 1}(\vX) = \vZ_{\theta}^{\ell}(\vX) +
\MSSA_{\theta}^{\ell}(\LN_{\theta}^{\ell}(\vZ_{\theta}^{\ell}(\vX))).
\end{equation}
In our implementation, we also experimented with using multi-head
self-attention (MHSA) in place of MSSA. It turns out that this
architecture is \textit{also} viable, though the depth of the network
needs to be much greater in order to achieve equivalent performance
to the usual CRATE or transformer architecture. %

\begin{figure}[t!]
\centering
\includegraphics[width=0.7\textwidth]{\toplevelprefix/chapters/chapter8/figs/aot_backbone.pdf}
\caption{\small\textbf{One layer of the AoT backbone.} Token
representations merely go through a layer-norm and the multi-head
(subspace) self-attention operator to form the layer's output. Notice
that there is no token-wise nonlinearity such as MLP or ISTA or ODL.}
\label{fig:aot_backbone}
\end{figure}

We conduct experiments using the proposed AoT architecture and
demonstrate its potential. We pre-train the AoT-MSSA and AoT-MHSA
models of different sizes, along with GPT-2, on
OpenWebText~\citep{Gokaslan2019OpenWeb}. We plot the training loss
and validation loss against the number of training iterations in
\Cref{fig:loss}(a) and (b), respectively. It is observed that medium-
and large-sized AoT-based models achieve training and validation
losses comparable to those of the GPT-2 base model. In addition,
compared to the GPT-2 base model, the AoT-MHSA model is identical to
the GPT-2 base model, except for the absence of MLP layers in the
architecture. As shown in \Cref{fig:loss}, incorporating MLP layers
can accelerate the training process. Using the above pre-trained
models, we compute the cross-entropy validation loss without training
on different datasets in \Cref{tab:zeroshot}. It is observed that the
AoT models with medium and large parameter sizes can achieve
comparable performance to the GPT-2 base model. Moreover, we found
that adding MLP layers to AoT does not improve the zero-shot
performance. These results highlight the potential of attention-only
models to achieve competitive results while maintaining interpretability.

\begin{table*}[!htbp]
\caption{Zero-shot results on several language benchmark datasets and
tasks: Evaluation of different sizes of AoT with the MSSA and MHSA
operators and comparison to the GPT2 model.}\vskip 0.1in
\centering
\resizebox{1\columnwidth}{!}
{
\begin{tabular}{l|cccccc}
\toprule
Models  & {\bf LAMBADA} & {\bf PTB} & {\bf WikiText} & {\bf
LAMBADA} & {\bf CBT CN} & {\bf CBT NE} \\
\# of parameters  & (val loss) $\downarrow$ &  (val loss)
$\downarrow$ &(val loss) $\downarrow$ & (acc) $\uparrow$ &(acc)
$\uparrow$ &(acc) $\uparrow$ \\
\midrule
AoT-MSSA Base (102M) & 4.70 & 6.03 & 4.65 & 0.25 & 0.80 & 0.74\\
AoT-MSSA Medium (182M) & 4.47 & 5.08 & 4.22 & 0.29 & 0.84 & 0.77 \\
AoT-MHSA Base (122M) & 4.42 & 5.52 & 4.19 & 0.38 & 0.86 & 0.82\\
GPT-2 Base (124M) & 4.32 & 5.75 & 4.13 &  0.40 &  0.87 &  0.84 \\
\bottomrule
\end{tabular}
}
\label{tab:zeroshot}
\end{table*}
\vspace{-0.05in}

\begin{figure*}[t]
\begin{center}
\includegraphics[width=0.4\linewidth]{\toplevelprefix/chapters/chapter8/figs/training_loss.png}
\hspace{0.4in}
\includegraphics[width =
0.4\linewidth]{\toplevelprefix/chapters/chapter8/figs/val_loss.png}
\vspace{-0.15in}
\caption{\centering \textbf{Evaluating models on language tasks.} We
plot the training loss (left) and validation loss (right) of the AoT
and GPT-2 models pretrained on OpenWebText.}  \label{fig:loss}
\end{center}
\vspace{-0.15in}
\end{figure*}

\section{Summary and Notes}

All work in this chapter is downstream of the Transformer
architecture, which was introduced by \citet{vaswani2017attention}.
The Transformer architecture is formally described in
\Cref{sec:contrastive_learning}. A main empirical innovation in
recent years, spurred by the prevalence and performance of the
Transformer architecture, is to formulate a given learning problem as
a sequence-to-sequence problem and apply the Transformer
architecture. This has enabled the Transformer architecture to be
essentially ubiquitous in (almost) all deep learning applications. As
such, direct improvements to the Transformer can propagate to become
solutions to many problems and have considerable impact; similarly,
we can apply our white-box understanding of transformer-like
architectures to many modern problems. The material covered in this
chapter is merely a subset of the work that has already been done;
other work includes masked completion for text data (i.e., BERT)
\citep{devlin2019bert,yu2024white}, (mechanistic) interpretability of
language and vision models \citep{bai2024improving}, and error
correcting coding \citep{zheng2025white}. There is much more to do.

There is also much more theory specifically about the practice of
scaling neural networks, which is enormously practically viable, and
we at least remark on it here. This line of work was popularized by
the ``Tensor Programs'' line of work \citep{yang2022tensor}. The
basic prescription is that we want the initial gradient updates in a
transformer to be constant size, and by working through the
backpropagation equations (\Cref{app:optimization}) carefully, we can
determine the scale of the initialization and learning rates (chosen
layer-wise) that are required to achieve this. In practice, such
prescriptions greatly increase the stability and convergence of
training at large scales; they also prescribe a way to find the
``optimal''\footnote{The word ``optimal'' is used in quotes because
the work on this merely uses some desiderata about the weight size,
feature size, and gradient size at initialization to determine
``optimality'', as opposed to, say, the test loss at convergence.}
hyperparameters for large-scale training using only small-scale
training. Follow-ups to this work attempt to accommodate the feature
geometry \citep{bernstein2024oldoptimizernewnorm}, which could be
informed by the work in this book about representation learning.
Other follow-ups incorporate this weight-wise information into the
optimizer itself to obtain these scaling benefits automatically,
obtaining optimizers such as Muon \citep{jordan6muon}, which have
recently been used for training trillion-parameter models very stably
\citep{moonshot2025kimi}. Overall, the two approaches to deep
learning theory are orthogonal or complementary.

\section{Exercises and Extensions}

\begin{exercise}
Read the DINO paper \cite{caron2021emerging}.
\end{exercise}

\begin{exercise}
DINO v2 \cite{oquab2023dinov2} uses everything from DINO v1 but also,
during the data augmentation phase, randomly masks out patches within
each view. This kind of augmentation should enforce that the features
of images with similar local information are similar. Formulate an
optimization problem that promotes this in the encoder, and implement it.
\end{exercise}

\begin{exercise}
This exercise considers the implementation of stochastic optimization
algorithms to minimize losses involving expectations.
\begin{enumerate}[(a)]
\item Propose an alternative to the term involving \(R_{\eps}\) in
\eqref{eq:simdino_loss_teacherstudent_empirical} for
approximating the covariance regularization term in
\eqref{eq:simdino_loss_teacherstudent}. Evaluate the time
complexity required to compute your proposed term and its
gradient. Include an analysis for computing it on a single
compute node vs.~multiple nodes.
\item Evaluate the time complexity required to compute the existing
term in \eqref{eq:simdino_loss_teacherstudent_empirical} and its gradient.
\end{enumerate}
\end{exercise}

\begin{exercise}
Prove that \eqref{eq:linear_probing} and
\eqref{eq:linear_probing_empirical} are convex optimization problems.
\end{exercise}

\begin{exercise}
\phantom{}
\begin{enumerate}[(a)]
\item Implement the CRATE and CRATE-\(\alpha\) models.
\item Compare their performance and efficiency on the CIFAR-10 dataset.
\item Compare their interpretability in two ways:
\begin{itemize}
    \item The sparsity \(\norm{\vZ}_{0}\) of the representation \(\vZ\)
    \item The attention maps \(\va_{\theta}^{k, \ell}\)
\end{itemize}
\end{enumerate}
\end{exercise}

\end{document}
